% !TEX program = xelatex
% 第三章完善版；主文件、原图与补充图目录须一起保存。
\documentclass[UTF8,zihao=-4,fontset=fandol,a4paper]{ctexart}
\usepackage[margin=24mm,headheight=16pt,headsep=8mm,footskip=12mm]{geometry}
\usepackage{amsmath,amssymb,bm,mathtools}
\usepackage{xcolor,graphicx,tikz,booktabs,array,tabularx,longtable}
\usepackage{enumitem,fancyhdr,caption,hyperref}
\usetikzlibrary{arrows.meta,calc,positioning,angles,quotes,patterns,decorations.pathmorphing}
\definecolor{ink}{HTML}{233E50}
\definecolor{soft}{HTML}{F0F4F6}
\definecolor{muted}{HTML}{586774}
\hypersetup{hypertexnames=false,colorlinks=true,linkcolor=ink,urlcolor=ink,pdftitle={理论力学讲义：第三章 刚体力学（完善版）},pdfauthor={}}
\ctexset{section={format=\Large\bfseries\color{ink}},subsection={format=\large\bfseries},subsubsection={format=\normalsize\bfseries}}
\renewcommand{\thesection}{3.\arabic{section}}
\renewcommand{\theequation}{3.\arabic{equation}}
\renewcommand{\thefigure}{3.\arabic{figure}}
\setcounter{tocdepth}{2}
\linespread{1.18}
\tikzset{every picture/.style={>=Stealth},every node/.style={font=\small},force/.style={->,thick,ink},guide/.style={dashed,muted}}
\setlength{\parskip}{3pt}
\setlength{\emergencystretch}{2em}
\clubpenalty=10000 \widowpenalty=10000
\setlist{nosep,topsep=5pt,leftmargin=2em}
\setlength{\textfloatsep}{13pt}
\setlength{\intextsep}{12pt}
\captionsetup{font=small,labelfont=bf,skip=6pt,hypcap=false}
\pagestyle{fancy}
\fancyhf{}
\fancyhead[L]{\small\color{muted}理论力学讲义}
\fancyhead[R]{\small\color{muted}第三章\quad 刚体力学}
\fancyfoot[C]{\small\thepage}
\renewcommand{\headrulewidth}{0pt}
\newcommand{\vect}[1]{\bm{#1}}
\newcommand{\dd}{\mathrm d}
\newcommand{\er}{\vect e_r}
\newcommand{\et}{\vect e_\theta}
\newcommand{\vt}{\vect e_t}
\newcommand{\vn}{\vect e_n}
\newcommand{\important}[1]{\par #1\par}
\newcounter{example}
\newcommand{\example}[1]{\Needspace{5\baselineskip}\refstepcounter{example}\subsubsection*{例\theexample\quad #1}}
\newcommand{\solution}{\setcounter{solutionstep}{0}\par\nopagebreak[4]\smallskip\nopagebreak[4]\noindent\textbf{解答}\quad}


\newcommand{\Mtot}{M}
\newcommand{\CM}{\mathrm C}
\newcommand{\figstart}{\par\addvspace{8pt}\noindent\begin{minipage}{\linewidth}\centering}
\newcommand{\figend}[1]{\captionof{figure}{#1}\end{minipage}\par\addvspace{8pt}}
\newcounter{exercise}
\newcommand{\exercise}[1]{\Needspace{5\baselineskip}\refstepcounter{exercise}\subsubsection*{训练题\theexercise\quad #1}}


\makeatletter
\newcommand{\Needspace}[1]{\par\if@nobreak\else\penalty0\begingroup\dimen0=#1\relax\dimen2=\pagegoal\advance\dimen2 by-\pagetotal\ifdim\dimen0>\dimen2\newpage\fi\endgroup\fi}
\makeatother
\let\SavedSection\section
\renewcommand{\section}{\Needspace{8\baselineskip}\SavedSection}
\let\SavedSubsection\subsection
\renewcommand{\subsection}{\Needspace{6\baselineskip}\SavedSubsection}
\newsavebox{\TeachingBox}
\newcommand{\TeachingFigure}[2]{%
\par\addvspace{8pt}\noindent\begin{minipage}{\linewidth}\centering
\sbox{\TeachingBox}{\input{补充图/#1.tex}}%
\ifdim\wd\TeachingBox>.98\linewidth\resizebox{.98\linewidth}{!}{\usebox{\TeachingBox}}\else\usebox{\TeachingBox}\fi
\captionof{figure}{#2}\label{fig:lesson-#1}\end{minipage}\par\addvspace{8pt}}
\newcommand{\OriginalFigure}[3][6.2cm]{%
\par\addvspace{8pt}\noindent\begin{minipage}{\linewidth}\centering
\includegraphics[width=.94\linewidth,height=#1,keepaspectratio]{原图/#2}%
\captionof{figure}{#3}\label{fig:orig-#2}\end{minipage}\par\addvspace{8pt}}


\newsavebox{\QuestionDrawing}
\newcommand{\QuestionStatement}[3]{%
\par\noindent\begin{minipage}{\linewidth}%
\begin{minipage}[t]{.63\linewidth}\vspace{0pt}\raggedright\setlength{\parindent}{0pt}#1\end{minipage}\hfill%
\begin{minipage}[t]{.34\linewidth}\vspace{2pt}\centering%
\sbox{\QuestionDrawing}{\input{题图/#2.tex}}%
\resizebox{.98\linewidth}{!}{\usebox{\QuestionDrawing}}\par\vspace{3pt}{\small #3}%
\end{minipage}\end{minipage}\par\smallskip}
\newcommand{\LongQuestionStatement}[3]{\par #1\par\smallskip\noindent\begin{minipage}{\linewidth}\centering\sbox{\QuestionDrawing}{\input{题图/#2.tex}}\resizebox{.42\linewidth}{!}{\usebox{\QuestionDrawing}}\par{\small #3}\end{minipage}\par\smallskip}
\newcounter{solutionstep}
\newcommand{\Step}[1]{\ifnum\value{solutionstep}>0\par\smallskip\noindent\fi\stepcounter{solutionstep}\textbf{（\arabic{solutionstep}）#1}\quad}


\newsavebox{\SolutionDrawing}
\newsavebox{\SolutionScaled}
\newcommand{\SolutionFigure}[2]{%
\par\addvspace{6pt}\noindent\begin{minipage}{\linewidth}\centering%
\sbox{\SolutionDrawing}{\input{解答图/#1.tex}}%
\sbox{\SolutionScaled}{\resizebox{.88\linewidth}{!}{\usebox{\SolutionDrawing}}}%
\ifdim\ht\SolutionScaled>5.6cm\resizebox{!}{5.6cm}{\usebox{\SolutionDrawing}}\else\usebox{\SolutionScaled}\fi%
\captionof{figure}{#2}\label{fig:solution-#1}\end{minipage}\par\addvspace{6pt}}

\begin{document}
\begin{titlepage}
\centering\vspace*{21mm}
{\small\color{muted}理论力学讲义\par}
\vspace{13mm}{\Huge\bfseries 第三章\par}
\vspace{7mm}{\LARGE\bfseries\color{ink}刚体力学\par}
\vspace{13mm}{\large 从空间姿态到转动动力学\par}
\vspace{17mm}
\begin{tikzpicture}
\draw[thick,ink,fill=soft,rotate=-18](0,0)ellipse(2.1 and .65);
\fill[ink](0,0)circle(2pt)node[below]{$O$};
\draw[force](0,0)--(0,2.25)node[above]{$\vect L$};
\draw[force](0,0)--(-1.2,1.9)node[above left]{$\vect\omega$};
\node[align=center]at(0,-1.65){$\vect L=\mathsf I_O\vect\omega$\qquad
$\vect M_O=\dfrac{\dd\vect L}{\dd t}$};
\node[font=\small,muted]at(0,-2.35){定点 $O$ 的转动示意};
\end{tikzpicture}
\vfill{\normalsize 几何描述\quad 定理推导\quad 完整例题\quad 练习与解答\par}
\vspace{12mm}
\end{titlepage}
\pagenumbering{roman}
\begingroup\small\linespread{1.05}\selectfont
\makeatletter
\renewcommand*{\@pnumwidth}{2.2em}
\renewcommand*{\@tocrmarg}{3.2em}
\renewcommand*\l@section[2]{\addvspace{4pt}\@dottedtocline{1}{0em}{2.7em}{\bfseries #1}{\bfseries #2}}
\renewcommand*\l@subsection{\@dottedtocline{2}{1.5em}{3.2em}}
\makeatother
\tableofcontents\endgroup\clearpage
\pagenumbering{arabic}

\section{刚体模型与运动分解}\label{sec:c3-1}
\subsection{刚性约束、自由度与适用范围}\label{sub:c3-1.1}
刚体是任意两质元间距离都不随时间变化的理想质点组。若质元的位置为
$\vect r_i(t)$，其约束可写为
\begin{equation}
|\vect r_i-\vect r_j|^2=d_{ij}^2,\qquad d_{ij}=\text{常数}.
\end{equation}
这一模型忽略弹性变形。它保留物体的尺寸与质量分布，因此能够描述姿态、转动惯量、
角动量和支承反力；把物体当成质点则会丢掉这些信息。

一个一般的三维刚体有六个自由度。先确定一个体内点 $O$ 的三个坐标；
再确定一个固定在物体上的右手正交三轴组的方向，需要三个独立参数。
也可用三个不共线点说明：第一个点有三个位置参数，第二个点与它距离固定，
还需两个方向参数，第三个点与前两点的距离固定，还需一个参数。
其余各点由刚性约束确定。这里所说的六个自由度针对没有外加约束的非退化刚体；
定点、定轴和保持平面运动都会减少自由度。

\begin{center}
\begin{tabularx}{\linewidth}{l c X}
\toprule
运动条件&自由度&一组适用变量\\\midrule
自由空间运动&6&$\vect r_O$ 与三个姿态角\\
定点转动&3&三个姿态角\\
平面运动&3&$x_C,y_C,\varphi$\\
定轴转动&1&绕固定轴的有向转角 $\varphi$\\
方向保持不变的平动&3&任一体内点的三个坐标\\\bottomrule
\end{tabularx}
\end{center}
内部作用力在刚体整体方程中如何消去，要分别考察合力、合力矩与功。
成对满足牛顿第三定律的内力合力为零；若成对内力沿两质元连线，
内力总矩也为零。理想刚性约束的内力总功为零。实际轴承、柔性连接、马达和耗散
应按选定系统边界计入，不能把一切内部机制都无条件删去。


\subsubsection*{从距离约束到六个独立参数}
设刚体含 $N$ 个可辨认质元，初始位置一共需要 $3N$ 个坐标。
不能把所有 $N(N-1)/2$ 条距离条件直接相减来计算自由度，
因为这些条件在 $N$ 较大时彼此相关。
取三个不共线的标记点 $A,B,D$，先给出 $A$ 的位置，消耗三个参数。
点 $B$ 必须位于以 $A$ 为中心、$AB$ 为半径的球面上，球面上的位置仍有两个参数。
点 $D$ 必须同时满足 $AD$ 与 $BD$ 两个定长条件，
两个球面的交集为一条圆，其位置还有一个参数。
于是
\[
3+2+1=6.
\]
三个标记点确定一个随体的有向三角形。
在已选择的构形分支内，其余质元相对于该三角形的位置不再独立。
这种数法同时说明了定点约束为何减少三个自由度，
固定轴约束为何最终只剩一个转角。

若全部质元严格共线，绕这条直线旋转不能改变任何质元的位置，
理想无厚度线段的可辨认空间构形只有五个参数。
若还要描述真实杆的截面标记方向，就应恢复其有限截面，
而不能一面忽略截面、一面要求无厚度模型分辨轴向姿态。
质量对称也不等于姿态不存在：球体上的标记点仍可随转动改变位置。

\subsubsection*{内力为什么在整体方程中消去}
记 $\vect F_{ij}$ 为 $j$ 对 $i$ 的作用力，
$\vect F_{ji}=-\vect F_{ij}$。
把质元方程逐项相加，每一对内力给出
\[
\vect F_{ij}+\vect F_{ji}=0.
\]
但角动量方程和能量方程还分别需要检查矩与功率。
对于沿两质元连线的中心内力，
\[
\vect r_i\times\vect F_{ij}
+\vect r_j\times\vect F_{ji}
=(\vect r_i-\vect r_j)\times\vect F_{ij}=0.
\]
功率则为
\[
\mathcal P_{ij}
=\vect F_{ij}\cdot(\vect v_i-\vect v_j).
\]
令 $\vect r_{ij}=\vect r_i-\vect r_j$，
对刚性约束 $r_{ij}^2=d_{ij}^2$ 求导，得到
\[
2\vect r_{ij}\cdot(\vect v_i-\vect v_j)=0.
\]
由于 $\vect F_{ij}$ 平行于 $\vect r_{ij}$，所以 $\mathcal P_{ij}=0$。
总功率为所有内力对功率之和，故理想刚性约束内力不改变整体动能。
这不是说单个内力对单个质元不做功，而是成对功率相消。

在系统内装有电机、化学能源或可变形机构时，
其储能与耗散不能由以上理想中心约束推论自动排除。
应先确定把哪些部件和能源计入系统，
再决定哪一项属于外功、哪一项属于系统内部能量转化。


\subsection{空间坐标系与本体坐标系}\label{sub:c3-1.2}
取惯性坐标系 $S$，在刚体内选定基点 $O$，并附着正交坐标系
$B=(\vect e_x,\vect e_y,\vect e_z)$。体内质元 $A$ 在 $B$ 中的坐标
$\vect\rho_A$ 为常数。令 $\mathsf R(t)$ 的三列为本体三轴在空间坐标系中的分量，
则
\begin{equation}\label{eq:position}
\vect r_A(t)=\vect r_O(t)+\mathsf R(t)\vect\rho_A,\qquad
\mathsf R^{\mathsf T}\mathsf R=\mathsf 1,\quad \det\mathsf R=1.
\end{equation}
$\mathsf R$ 把本体分量映到空间分量，故空间分量变回本体分量用
$\mathsf R^{\mathsf T}$。正交性保证长度和夹角不变，行列式为 $+1$
排除镜像反射。只写“坐标变换矩阵”而不说明变换方向，容易把所有正弦项的符号写反。

\TeachingFigure{c3-diagram-01}{位置由基点平移与本体姿态共同确定}
% BEGIN ORIGINAL FIGURES
\OriginalFigure[6.2cm]{c3-1-p004-body-frame}{附着在刚体质心处的坐标轴}
% END ORIGINAL FIGURES

若选质心 $C$ 作基点，$\int\vect\rho\,\dd m=0$。
质心系便于分解动能和角动量；固定铰点、轴上点或接触点则可能使几何约束更简单。
选择基点是计算手段，不会改变物体的真实运动。
% BEGIN ORIGINAL FIGURES
\OriginalFigure[6.2cm]{c3-1-p005-base-frame}{改选体内基点后的刚体姿态描述}
% END ORIGINAL FIGURES



\subsubsection*{方向余弦与坐标变换的逐项含义}
空间系单位轴记为 $\vect E_1,\vect E_2,\vect E_3$，
本体单位轴记为 $\vect e_1,\vect e_2,\vect e_3$。
矩阵元
\[
R_{ij}=\vect E_i\cdot\vect e_j
\]
是第 $j$ 根本体轴对第 $i$ 根空间轴的方向余弦。
若某矢量在本体系分量为 $a_j$，则
\[
\vect a=\sum_j a_j\vect e_j,\qquad
a_i^{(S)}=\vect E_i\cdot\vect a=\sum_jR_{ij}a_j^{(B)}.
\]
因此 $\vect a_S=\mathsf R\vect a_B$。
取点积，
\[
\vect a_S^{\mathsf T}\vect b_S
=\vect a_B^{\mathsf T}\mathsf R^{\mathsf T}\mathsf R\vect b_B
=\vect a_B^{\mathsf T}\vect b_B,
\]
说明变换同时保持长度与夹角。
逆变换是 $\vect a_B=\mathsf R^{\mathsf T}\vect a_S$。
位移矢量和速度矢量的分量可以这样变换，
但若比较两个不同原点下的绝对位置坐标，还须计入原点平移。

九个矩阵元受到六条独立正交条件约束：
三条列向量单位长度条件与三条列向量互相垂直条件。
所以姿态矩阵仍只有三个独立连续参数。
$\det\mathsf R=+1$ 选定保持右手性的连通分支；
它不增加一个连续自由度。

\subsubsection*{更换基点时哪些量改变}
以 $O$ 和 $C$ 分别描述同一质元时，
\[
\vect\rho_{CA}=\vect\rho_{OA}-\vect\rho_{OC},\qquad
\vect r_C=\vect r_O+\mathsf R\vect\rho_{OC}.
\]
把两式代入，恰好得到相同的 $\vect r_A$。
若本体坐标轴方向保持相同，姿态矩阵无须改变。
若同时另选本体轴组，令常正交矩阵 $\mathsf Q$ 把新本体分量映到旧本体分量，
则新姿态矩阵为 $\mathsf R_{\text{新}}=\mathsf R\mathsf Q$。
基点平移和轴组重选是两种不同操作，不宜混成一项“坐标变化”。
% BEGIN ORIGINAL FIGURES
\OriginalFigure[6.2cm]{c3-1-p037-axis-polar}{用极角和方位角指定一根空间轴}
% END ORIGINAL FIGURES



\subsection{平动、转动与一般运动}\label{sub:c3-1.3}
若 $\mathsf R$ 不变，任意两点连线在空间中始终平行于自身，
物体做平动。各点位移相同、速度相同、加速度也相同。
平动可以沿直线，也可以沿曲线：座舱始终保持竖直而随摩天轮运动，
就是曲线平动。是否转动取决于姿态是否改变，不能只看质心轨迹是否弯曲。

若某体内点 $O$ 在惯性系中固定，物体做定点转动。一般运动可分为
基点的平动和相对于基点的转动。分解取决于基点；角速度则是整个刚体共有的量。

\TeachingFigure{c3-diagram-02}{相同的中心轨迹可以对应平动或转动}
% BEGIN ORIGINAL FIGURES
\OriginalFigure[6.2cm]{c3-1-p006-ferris-wheel}{摩天轮座舱的曲线平动}
% END ORIGINAL FIGURES

\example{曲线平动的速度与加速度}\label{ex:c3-1}

\QuestionStatement{一个矩形平台的中心沿半径 $R$ 的圆周以恒定角速率 $\Omega$ 运动，
但平台相对于地面保持水平且各边方向不变。求平台上任意点的速度和加速度。
取圆心为固定原点，$t=0$ 时中心位于 $+x$ 轴。写出任意材料点的速度、加速度，并判断运动类型。}{ex01}{例1题图}
% BEGIN DETAILED ANSWER ex01
\solution \Step{确定运动类型与基点}
以图~\ref{fig:solution-ex01}~中中心 $C$ 的圆周轨迹与固定的 $\overrightarrow{CA}$为依据，先写出图右侧第一行的几何或动力学关系，再代入其下方的计算式。

平台的各边方向保持不变，因此姿态矩阵为常矩阵、刚体角速度为零。取中心 $C$ 为基点，对任意材料点 $A$，有
\[
\vect r_A(t)=\vect r_C(t)+\vect\rho_A,\qquad
\vect r_C=R(\cos\Omega t\,\vect i+\sin\Omega t\,\vect j),
\]
其中 $\vect\rho_A=\vect r_A-\vect r_C$ 在空间中为常矢量。
\SolutionFigure{ex01}{例1的计算示意}
\Step{分别求导，不引入绕中心的转动}
固定空间基矢的导数为零，故
\[
\vect v_A=\dot{\vect r}_C
=R\Omega(-\sin\Omega t\,\vect i+\cos\Omega t\,\vect j),
\]
\[
\vect a_A=\ddot{\vect r}_C
=-R\Omega^2(\cos\Omega t\,\vect i+\sin\Omega t\,\vect j).
\]
因 $\dot{\vect\rho}_A=0$，所有点在同一时刻的速度和加速度均相同。
\Step{轨迹与参数极限}
$|\vect v_A|=R|\Omega|$、$|\vect a_A|=R\Omega^2$，且两者正交。
$t=0$ 时 $\vect v_A=R\Omega\vect j$、$\vect a_A=-R\Omega^2\vect i$，与切向速度、法向加速度一致。
每点的圆形轨迹圆心相对原点平移了 $\vect\rho_A$，所以这是曲线平动。
$\Omega\to0$ 时两者同时趋零。这里的向心加速度指向各点自身轨迹的圆心；它不意味着平台绕 $C$ 转动。

% END DETAILED ANSWER ex01
\subsubsection*{曲线平动例的两次求导}
平台任一点相对于中心的固定偏移记为
$\vect\rho_A=x_A\vect i+y_A\vect j+z_A\vect k$。
因为平台的三根边方向在空间中不变，$x_A,y_A,z_A$ 及三个单位轴都不随时间变化。
所以
\begin{align*}
\vect r_A(t)&=(R\cos\Omega t+x_A)\vect i
 +(R\sin\Omega t+y_A)\vect j+z_A\vect k,\\
\dot{\vect r}_A(t)&=-R\Omega\sin\Omega t\,\vect i
 +R\Omega\cos\Omega t\,\vect j,\\
\ddot{\vect r}_A(t)&=-R\Omega^2\cos\Omega t\,\vect i
 -R\Omega^2\sin\Omega t\,\vect j.
\end{align*}
所有点走的是相互平移后的圆，而不是绕同一个空间中心的同一圆。
因此各点向心加速度相同，所指向的各自轨迹中心却不同。
这里 $\Omega$ 是中心轨迹的角速率，
平台自身的姿态角速率为零，二者属于不同几何对象。

若撤去保持方向的机构，把平台刚性固定在轮辐上，
则 $\vect\rho_A$ 的空间方向随时间变化，
应改为 $\vect r_A=\vect r_C+\mathsf R(t)\vect\rho_A$。
新增的 $\dot{\mathsf R}\vect\rho_A$ 正是平台相对中心的转动速度。


\subsection{定点转动的欧拉定理}\label{sub:c3-1.4}
% BEGIN ORIGINAL FIGURES
\OriginalFigure[6.2cm]{c3-1-p007-euler}{欧拉与刚体定点转动定理}
% END ORIGINAL FIGURES

定点刚体从一个姿态变到另一个姿态，总可以等效为绕过定点的一根轴转过一个角度。
这是有限转动的几何结论，等效轴不一定是实际运动过程中的固定轴。

证明如下。姿态之间的相对变换是三维实正交矩阵 $\mathsf Q$，
且 $\det\mathsf Q=1$。奇数阶实矩阵至少有一个实特征值；
正交矩阵的特征值模为一，因此实特征值只能为 $\pm1$。
复特征值成共轭对、乘积为一，再由行列式为一可知一定存在特征值 $+1$。
相应单位特征矢量 $\vect n$ 满足 $\mathsf Q\vect n=\vect n$，给出不动的轴。
垂直于它的二维平面保持不变，$\mathsf Q$ 在此平面上的限制是二维旋转，
于是所有点都等效绕 $\vect n$ 转动。转角满足
\[
\cos\gamma=\frac{\operatorname{tr}\mathsf Q-1}{2}.
\]
若 $\mathsf Q=\mathsf1$，所有轴都可选；转角为 $\pi$ 时，轴方向的正负无法由转角区分。

有限姿态变换的等效轴、某一时刻的瞬时转轴、空间中始终固定的轴是三个概念。
一般定点运动只保证每一瞬间有角速度方向，并不保证这根方向不变。


\subsubsection*{定点欧拉定理的几何证明}
除矩阵证明外，还可以从球面上两标记点的位移看等效轴。
在固定点 $O$ 为球心的球面上选两个不共线方向 $OA,OB$，
变换后为 $OA',OB'$。
若候选不动方向为 $ON$，转动必须保持夹角，
故
\[
\vect n\cdot(\vect a'-\vect a)=0,\qquad
\vect n\cdot(\vect b'-\vect b)=0.
\]
这两式要求 $\vect n$ 位于两条弦的垂直平分大圆对应的平面中。
一般情况下，两平面的交线给出两个相反的轴向候选。
在垂直轴的截面中，一个方向的转角确定另一个方向的转角；
由于两构形之间保持有向夹角，两者相同。
这样得到的绕轴旋转与原姿态变换在三个独立方向上作用相同，
因而对刚体全部点作用相同。
若弦重合、某点本来不动或姿态变换为恒等变换，
应单独选其他标记方向或按退化情形处理。

\subsubsection*{等效转角与轴方向的实际提取}
当 $0<\gamma<\pi$ 时，
\[
\vect n=
\frac1{2\sin\gamma}
\begin{pmatrix}
Q_{32}-Q_{23}\\Q_{13}-Q_{31}\\Q_{21}-Q_{12}
\end{pmatrix}.
\]
求解次序是：先用迹求 $\gamma$，再用反对称部分求有向轴，
最后代入 $\mathsf Q\vect n=\vect n$ 检查。
$\gamma=\pi$ 时 $\sin\gamma=0$，上式不能使用；
此时应直接解 $(\mathsf Q-\mathsf1)\vect n=0$，
或利用 $\mathsf Q+\mathsf1=2\vect n\vect n^{\mathsf T}$ 提取轴。
$\gamma=0$ 时任何轴都可以作为等效轴。
这些退化处理属于反解方法的必要部分。


\subsection{运动描述的计算次序}\label{sub:c3-1.5}
实际问题宜先定系统与约束，再选空间系和本体系，明确基点及姿态角；
由几何约束求位置，从位置或角速度求速度与加速度，最后进入动力学。
遇到滚动，要区分接触处的体内质点与每一时刻重新选出的几何接触位置。
遇到多个转动叠加，要说明每个角速率相对于哪个参考系测量。

\example{更换基点后的位置分解}\label{ex:c3-2}

\QuestionStatement{刚体的质心为 $C$，某铰接点为 $O$。已知本体常矢量
$\vect\rho_{OC}$ 与 $\vect\rho_{OA}$，写出以 $C$ 和 $O$ 为基点的等价描述。
题给 $\vect\rho$ 均为本体坐标中的常矢量，以 $\mathsf R$ 将其转换为空间分量。分别以 $C$、$O$ 为基点写出 $A$ 的位置与速度，并说明换基点是否改变姿态矩阵。}{ex02}{例2题图}
% BEGIN DETAILED ANSWER ex02
\solution \Step{统一位置矢量的表达基底}
以图~\ref{fig:solution-ex02}~中三角形 $OCA$ 的有向位置分解为依据，先写出图右侧第一行的几何或动力学关系，再代入其下方的计算式。

题给 $\vect\rho_{OC}$、$\vect\rho_{OA}$ 是本体坐标中的常列矢量。设 $\mathsf R(t)$ 将本体分量变为空间分量，则
\[
\vect r_C=\vect r_O+\mathsf R\vect\rho_{OC},\quad
\vect\rho_{CA}=\vect\rho_{OA}-\vect\rho_{OC}.
\]
$\mathsf R$ 只描述三根体轴的取向，改变基点不会改变它。
\SolutionFigure{ex02}{例2的计算示意}
\Step{两种等价的位置分解}
以 $O$ 为基点：$\vect r_A=\vect r_O+\mathsf R\vect\rho_{OA}$；以 $C$ 为基点：
\[
\vect r_A=\vect r_C+\mathsf R\vect\rho_{CA}
=\vect r_O+\mathsf R\vect\rho_{OC}
+\mathsf R(\vect\rho_{OA}-\vect\rho_{OC}).
\]
中间两项相消，恢复第一种表达。
\Step{求导得到同一速度场}
由 $\dot{\mathsf R}\vect\rho=\vect\omega\times(\mathsf R\vect\rho)$，
\[
\vect v_C=\vect v_O+\vect\omega\times\vect r_{OC},\qquad
\vect v_A=\vect v_O+\vect\omega\times\vect r_{OA}
=\vect v_C+\vect\omega\times\vect r_{CA}.
\]
若 $O$ 是地面固定铰点，令 $\vect v_O=0$；若铰点随机构运动，则保留该项。
变化的是平移部分和相对位置，角速度及物体姿态完全相同。

% END DETAILED ANSWER ex02
\subsubsection*{自由度减少与运动分解的适用边界}
一般空间运动的六参数并不意味着任何刚体问题都要解六个独立微分方程。
如果定轴约束已经把姿态限制为 $\mathsf R_z(q)$，
便只需要一个角 $q(t)$；其余方程用于求约束反力。
对于平面运动，基点法向坐标固定，且本体的法向轴始终平行于空间法向轴，
位置可写成
\[
\vect r_A=x_C\vect i+y_C\vect j+
\mathsf R_z(q)\vect\rho_{CA}.
\]
独立运动变量为 $x_C,y_C,q$。
对于定点运动，$\vect r_O$ 为常量，但三姿态参数通常仍全部需要。
这三种限制必须由真实铰链、轴承或接触条件支持，
不能因为某个初始瞬间的速度恰好符合限制，就认为此后永远保持该运动类型。

选基点的常用次序是：
先看有无固定点，再看质心是否便于分解动力学，
最后看接触点是否能简化当前瞬间的运动学。
接触点可以随时间更换；它通常不是从开始到结束都固定的同一材料点。
后续使用速度、加速度或取矩公式时，都应保留这一区别。


\clearpage
\section{角速度与刚体速度场}\label{sec:c3-2}
\subsection{无穷小转动与角速度矢量}\label{sub:c3-2.1}
% BEGIN ORIGINAL FIGURES
\OriginalFigure[6.2cm]{c3-1-p030-infinitesimal}{无穷小转角与质点位移：图中 $\Delta n$ 表示无穷小转角矢量}
% END ORIGINAL FIGURES

绕单位轴 $\vect n$ 转过无穷小有向角 $\dd\gamma$，
体内点相对于基点的位置变化为
\[
\dd\vect r=\dd\vect\vartheta\times\vect r,\qquad
\dd\vect\vartheta=\vect n\,\dd\gamma.
\]
若先后施加两个无穷小转动，第二次转动作用于第一次转动后的点，
多出的项为二阶小量：
\[
\dd\vect r_{\text{合}}
=\dd\vect\vartheta_1\times\vect r+
\dd\vect\vartheta_2\times\vect r+O(\dd\gamma^2).
\]
所以在一阶精度下，两次转动交换次序不改变结果，且
$\dd\vect\vartheta=\dd\vect\vartheta_1+\dd\vect\vartheta_2$。
由此定义
\begin{equation}
\vect\omega=\lim_{\Delta t\to0}\frac{\Delta\vect\vartheta}{\Delta t}.
\end{equation}
方向由右手定则规定，量纲为 $\mathrm{s}^{-1}$。
有限转角也可带轴方向，但不满足普通矢量的加法规则；
角速度的矢量性质来自瞬时极限。


\subsubsection*{无穷小转动的二阶差异}
令两次小角位移为 $\epsilon_1\vect n_1$ 和 $\epsilon_2\vect n_2$。
第一种次序的质元位置为
\[
\vect r'=\vect r+\epsilon_1\vect n_1\times\vect r
+\epsilon_2\vect n_2\times\vect r
+\epsilon_1\epsilon_2\vect n_2\times(\vect n_1\times\vect r)
+\text{各次转动自身的二阶项}.
\]
交换次序时，混合项变成
$\epsilon_1\epsilon_2\vect n_1\times(\vect n_2\times\vect r)$。
二者相减，
\[
\Delta\vect r=
\epsilon_1\epsilon_2(\vect n_2\times\vect n_1)\times\vect r
+O(\epsilon^3).
\]
所以非交换性没有消失，而是推迟到了二阶。
取 $\epsilon_i=\omega_i\Delta t$ 后，混合项除以 $\Delta t$ 仍趋于零，
于是瞬时角速度严格满足普通矢量合成。
把很多不同轴的小转动积累成有限时间运动时，
二阶效应的总累积不能任意删掉。

角速度是轴矢量。镜像反射下，长度与时间不变，
但右手定则方向改变；这也是处理空间对称时需区分普通位置矢量与角速度的原因。
在本讲义所用的右手正交坐标系之间，
角速度分量按保持定向的转动矩阵变换。


\subsection{任意两点的速度与加速度}\label{sub:c3-2.2}
对式\eqref{eq:position}求导。$\dot{\mathsf R}\mathsf R^{\mathsf T}$
由正交关系可知为反对称矩阵，故存在唯一矢量 $\vect\omega$ 使
$\dot{\mathsf R}\mathsf R^{\mathsf T}\vect a=\vect\omega\times\vect a$。
于是
\begin{equation}\label{eq:velocity}
\vect v_A=\vect v_O+\vect\omega\times\vect r_{OA}.
\end{equation}
进一步求导，得到
\begin{equation}\label{eq:acceleration}
\vect a_A=\vect a_O+\vect\alpha\times\vect r_{OA}
+\vect\omega\times(\vect\omega\times\vect r_{OA}),\qquad
\vect\alpha=\frac{\dd\vect\omega}{\dd t}.
\end{equation}
最后一项可写为
$\vect\omega(\vect\omega\cdot\vect r_{OA})-\omega^2\vect r_{OA}$，
即指向瞬时转轴的法向加速度。它不总是指向基点，
因为平行于转轴的位置分量不贡献这项。


\subsubsection*{速度公式从姿态矩阵逐步得到}
对 $\vect r_A=\vect r_O+\mathsf R\vect\rho_A$ 求导，
注意本体坐标 $\vect\rho_A$ 恒定：
\[
\vect v_A=\vect v_O+\dot{\mathsf R}\vect\rho_A
=\vect v_O+\dot{\mathsf R}\mathsf R^{\mathsf T}\vect r_{OA}.
\]
对 $\mathsf R\mathsf R^{\mathsf T}=\mathsf1$ 求导，
\[
\dot{\mathsf R}\mathsf R^{\mathsf T}
+\mathsf R\dot{\mathsf R}^{\mathsf T}=0.
\]
因此 $\mathsf W=\dot{\mathsf R}\mathsf R^{\mathsf T}$ 为反对称矩阵。
其三个独立矩阵元恰构成 $\vect\omega$，
而 $\mathsf W\vect r=\vect\omega\times\vect r$。
这一步把保持长度的几何性质转换成唯一的刚体速度场。

再对速度式求导：
\begin{align*}
\vect a_A-\vect a_O
&=\frac{\dd}{\dd t}(\vect\omega\times\vect r_{OA})\\
&=\dot{\vect\omega}\times\vect r_{OA}
+\vect\omega\times\dot{\vect r}_{OA}\\
&=\vect\alpha\times\vect r_{OA}
+\vect\omega\times(\vect v_A-\vect v_O)\\
&=\vect\alpha\times\vect r_{OA}
+\vect\omega\times(\vect\omega\times\vect r_{OA}).
\end{align*}
其中 $\dot{\vect r}_{OA}$ 是连线矢量的空间导数，
不能因为其长度固定就把它写成零。

\subsubsection*{切向项与法向项怎样判断方向}
令 $\vect n=\vect\omega/\omega$，
分解 $\vect r_{OA}=r_\parallel\vect n+\vect r_\perp$。
由三重积公式，
\[
\vect\omega\times(\vect\omega\times\vect r_{OA})
=-\omega^2\vect r_\perp.
\]
所以其大小为 $\omega^2r_\perp$，方向从质元指向过基点的瞬时轴。
若固定轴方向恒定，$\vect\alpha=\dot\omega\vect n$，
另一项大小为 $|\dot\omega|r_\perp$，沿圆轨迹的切线。
若瞬时轴本身转向，$\vect\alpha$ 不一定平行 $\vect\omega$；
此时仍可使用矢量式，但不应把全部 $\vect\alpha\times\vect r$ 都解释成同一圆上的切向加速度。

\subsubsection*{相对速度公式的三个推论}
由 $\vect v_A-\vect v_B=\vect\omega\times\vect r_{BA}$ 得
\[
(\vect v_A-\vect v_B)\cdot\vect r_{BA}=0.
\]
因此任意两点的速度沿两点连线投影相等。
若相对位置平行于 $\vect\omega$，两点的速度相等；
若相对位置垂直于 $\vect\omega$，
相对速度大小为 $\omega|\vect r_{BA}|$。
反过来，一条连线上投影相等只是刚性约束的必要条件，
还不足以用两个质元的速度唯一确定全部空间角速度。
确定角速度一般还要另一条不平行的连线。


\TeachingFigure{c3-diagram-03}{相对速度垂直于转轴和位置矢量，法向加速度指向轴}
% BEGIN ORIGINAL FIGURES
\OriginalFigure[6.2cm]{c3-1-p008-velocity}{定点转动中位置、角速度与速度的相互方向}
% END ORIGINAL FIGURES

以另一点 $B$ 为基点，同样有
$\vect v_A=\vect v_B+\vect\omega\times\vect r_{BA}$。
因此角速度与基点选择无关。若两个候选角速度都能描述全部点的相对速度，
它们的差与两根不平行的体内连线叉积都为零，差只能为零。

\example{定点转动中一点的速度}\label{ex:c3-3}

\QuestionStatement{刚体绕固定点 $O$ 运动，某瞬间
$\vect\omega=(2,1,-1)\,\mathrm{s}^{-1}$，
$\vect r_{OA}=(1,-2,3)\,\mathrm m$。求点 $A$ 的速度。
同时计算速率，并检验速度与角速度、位置矢量的正交关系。}{ex03}{例3题图}
% BEGIN DETAILED ANSWER ex03
\solution \Step{明确模型与计算依据}
以图~\ref{fig:solution-ex03}~中从固定点 $O$ 指向 $A$ 的位置矢量为依据，先写出图右侧第一行的几何或动力学关系，再代入其下方的计算式。
 \[
\vect v_A=\begin{vmatrix}\vect i&\vect j&\vect k\\2&1&-1\\1&-2&3\end{vmatrix}
=(1,-7,-5)\,\mathrm{m\,s^{-1}}.
\]
检查 $\vect v_A\cdot\vect r_{OA}=1+14-15=0$，
以及 $\vect v_A\cdot\vect\omega=2-7+5=0$。
相对速度同时垂直于两个参与叉积的矢量。


把行列式完全展开可看出每个分量的来源：
\begin{align*}
v_x&=\omega_yz-\omega_zy=3-2=1,\\
v_y&=\omega_zx-\omega_xz=-1-6=-7,\\
v_z&=\omega_xy-\omega_yx=-4-1=-5.
\end{align*}

\SolutionFigure{ex03}{例3的计算示意}
\Step{逐项计算与求解}
再用标量恒等式作第三个独立核查：
\[
v^2=\omega^2r^2-(\vect\omega\cdot\vect r)^2.
\]

\Step{结果与适用条件}
本题 $\omega^2=6$、$r^2=14$、$\vect\omega\cdot\vect r=-3$，
故 $v^2=84-9=75$，与 $1^2+(-7)^2+(-5)^2=75$ 一致。
若还给定 $\vect\alpha$，才可完整确定加速度；
单给某瞬间 $\vect\omega$ 不能自行假定 $\vect\alpha=0$。
\Step{给出速率} 由 $v^2=75$，速率为 $5\sqrt3\simeq8.66025\,\mathrm{m\,s^{-1}}$。三个分量的符号描述运动方向，取模时不能把分量代数相加。

% END DETAILED ANSWER ex03
\subsection{转动参考系的导数公式}\label{sub:c3-2.3}
对任意矢量 $\vect A=A_x\vect e_x+A_y\vect e_y+A_z\vect e_z$，
本体轴满足 $\dot{\vect e}_i=\vect\omega\times\vect e_i$，所以
\begin{equation}\label{eq:transport}
\left(\frac{\dd\vect A}{\dd t}\right)_S
=\left(\frac{\dd\vect A}{\dd t}\right)_B+\vect\omega\times\vect A.
\end{equation}
左边是在空间系中观察到的真实导数；右边第一项仅对本体分量求导。
本体坐标系是随刚体转动的轴组；把原点移到体内点而不转动坐标轴的参考系，
与本体坐标系并不相同。

将 $\vect A=\vect\omega$ 代入时，叉积项为零，因此角加速度在本体轴上的分量
恰为 $\dot\omega_x,\dot\omega_y,\dot\omega_z$。
将 $\vect A=\vect L$ 代入则叉积项通常不为零，后面欧拉动力学方程正是由它产生。


\subsubsection*{运输公式的逐分量证明}
空间导数作用于分量和基矢两个部分：
\[
\left(\frac{\dd\vect A}{\dd t}\right)_S
=\sum_i\dot A_i\vect e_i+\sum_iA_i\dot{\vect e}_i.
\]
第一项定义为转动系分量导数，第二项由轴组转动引起。
利用 $\dot{\vect e}_i=\vect\omega\times\vect e_i$，
\[
\sum_iA_i\dot{\vect e}_i
=\vect\omega\times\sum_iA_i\vect e_i
=\vect\omega\times\vect A,
\]
于是得到式\eqref{eq:transport}。
“分量恒定”只保证第一项为零，不保证空间矢量恒定。
例如固定在转台上的径向单位轴，在转台内分量始终为 $(1,0,0)$，
却有空间导数 $\dot{\vect e}_r=\Omega\vect e_\varphi$。

若空间中有固定矢量 $\vect A_0$，
在转动系中恰满足
\[
\left(\frac{\dd\vect A_0}{\dd t}\right)_B
=-\vect\omega\times\vect A_0.
\]
这个负号表示坐标分量的变化抵消了轴组的真实转动。
它在固定空间角动量投影到陀螺本体系时尤其重要。

\subsubsection*{移动观察点的完整速度与加速度}
运输公式也可用于相对运动，但需把参考架角速度记成 $\vect\Omega$，
避免与运动物体自身的角速度混淆。
若参考原点为 $O'$，相对位置为 $\vect s$，相对速度为 $\vect v'$，则
\[
\vect v=\vect v_{O'}+\vect v'+\vect\Omega\times\vect s.
\]
再求一次空间导数，依次得到
\[
\vect a=\vect a_{O'}+\vect a'
+2\vect\Omega\times\vect v'
+\dot{\vect\Omega}\times\vect s
+\vect\Omega\times(\vect\Omega\times\vect s).
\]
两个 $\vect\Omega\times\vect v'$ 分别来自
相对速度的轴组导数与位置运输项的导数。
刚体内部质元相对于随体架静止时，$\vect v'=\vect a'=0$，
上式才退化为刚体加速度场。
若质元沿杆滑动，或观察螺旋桨相对于转弯机体的运动，
则不能删除相对运动项。


\subsection{角速度的合成及角加速度}\label{sub:c3-2.4}
设参考架相对空间以 $\vect\Omega$ 转动，
刚体相对架以 $\vect\omega_{\text{相}}$ 转动。
对每根本体轴应用导数公式可得
\begin{equation}
\vect\omega=\vect\Omega+\vect\omega_{\text{相}}.
\end{equation}
相加前须把两矢量投影到同一组轴上。有限转动矩阵一般不交换，
并不妨碍角速度在同一瞬间相加。


\subsubsection*{角速度合成的证明与有向角加速度}
取任一本体单位轴 $\vect e$。
若参考架的角速度为 $\vect\Omega$，
该轴相对于参考架的角速度为 $\vect\omega_{\text{相}}$，则
\[
\left(\frac{\dd\vect e}{\dd t}\right)_S
=\vect\omega_{\text{相}}\times\vect e+\vect\Omega\times\vect e
=(\vect\Omega+\vect\omega_{\text{相}})\times\vect e.
\]
同一个等式对三根正交轴都成立，
由刚体角速度的唯一性，合角速度就是两者之和。
这里三个矢量可以先用空间分量计算，也可以先用参考架分量计算，
但不能让一个矢量用空间分量、另一个用相对架分量后直接逐项相加。

把相对角速度的导数也分成分量变化和方向变化，
\[
\vect\alpha
=\dot{\vect\Omega}_S+
\left(\frac{\dd\vect\omega_{\text{相}}}{\dd t}\right)_{\text{架}}
+\vect\Omega\times\vect\omega_{\text{相}}.
\]
若两相对速率是常数，最后一项仍可能非零。
若有第三个串联转动，再对下一架应用同一公式，
角速度仍可逐架相加；角加速度则须逐架处理运输项。


\example{转台上的水平轴圆盘}\label{ex:c3-4}
% BEGIN ORIGINAL FIGURES
\OriginalFigure[6.2cm]{c3-1-p026-colored-turntable}{支架上的圆盘及两根转动轴}
% END ORIGINAL FIGURES
\QuestionStatement{转台绕竖直轴以常速 $\Omega$ 转动，盘轴始终水平且随转台转动，
圆盘相对于转台绕盘轴以常速 $\omega$ 转动。
令瞬时盘轴方向为 $\vect e_r$，竖直向上为 $\vect k$。
求圆盘的绝对角速度及其大小、绝对角加速度，并解释恒定分转速为何仍可有角加速度。}{ex04}{例4题图}
% BEGIN DETAILED ANSWER ex04
\solution \Step{明确模型与计算依据}
以图~\ref{fig:solution-ex04}~中竖直 $\vect k$ 轴与水平 $\vect e_r$ 轴为依据，先写出图右侧第一行的几何或动力学关系，再代入其下方的计算式。
 \[
\vect\omega_{\text{总}}=\omega\vect e_r+\Omega\vect k,\qquad
|\vect\omega_{\text{总}}|=\sqrt{\omega^2+\Omega^2}.
\]
虽两个速率都恒定，方向 $\vect e_r$ 仍在变化，故
\[
\vect\alpha=\omega\Omega\vect e_\varphi.
\]
这说明恒定的分转速不保证总角速度矢量恒定。


本例可以写成固定空间分量，直接检查上述角加速度。
取转台初始径向沿 $\vect i$：
\[
\vect e_r(t)=\cos\Omega t\,\vect i+\sin\Omega t\,\vect j,
\qquad
\vect e_\varphi(t)=-\sin\Omega t\,\vect i+\cos\Omega t\,\vect j.
\]

\SolutionFigure{ex04}{例4的计算示意}
\Step{逐项计算与求解}
因此
\[
\vect\omega_{\text{总}}
=\omega\cos\Omega t\,\vect i+\omega\sin\Omega t\,\vect j+\Omega\vect k,
\]
\[
\vect\alpha
=-\omega\Omega\sin\Omega t\,\vect i+\omega\Omega\cos\Omega t\,\vect j.
\]

\Step{结果与适用条件}
角速率大小恒定，意味着 $\vect\omega\cdot\vect\alpha=0$，
而不是意味着 $\vect\alpha=0$。
总角速度端点绕竖直轴画圆，其方向作圆锥运动。

若盘自转速率与转台速率都允许变化，则
\[
\vect\alpha
=\dot\omega\vect e_r+\omega\Omega\vect e_\varphi+\dot\Omega\vect k.
\]
三个项分别代表盘相对加速、盘轴随台转向、转台自身加速。

% END DETAILED ANSWER ex04
\TeachingFigure{c3-diagram-04}{转台转动和圆盘相对转动的矢量合成}
% BEGIN ORIGINAL FIGURES
\OriginalFigure[6.2cm]{c3-1-p010-turntable}{水平轴圆盘与转动框架：图中径向单位矢量 $\vect i$ 对应正文 $\vect e_r$}
% END ORIGINAL FIGURES


\example{倾斜轴陀螺的角速度}\label{ex:c3-5}
% BEGIN ORIGINAL FIGURES
\OriginalFigure[6.2cm]{c3-1-p011-inclined-top}{倾斜自转轴与竖直进动轴：$\omega_1$ 为自转角速率，$\omega_2$ 为进动角速率}
% END ORIGINAL FIGURES
\QuestionStatement{陀螺对称轴 $\vect n$ 与竖直轴夹角 $\theta$，
轴绕竖直方向以角速率 $\Omega$ 进动，物体相对于进动架绕 $\vect n$
以角速率 $\omega$ 自转，$\theta$ 恒定。求总角速度。
同时求总角速率大小与角加速度，说明进动对轴向角速度的贡献。}{ex05}{例5题图}
% BEGIN DETAILED ANSWER ex05
\solution \Step{明确模型与计算依据}
以图~\ref{fig:solution-ex05}~中本体轴 $\vect n$ 在竖直与径向的投影为依据，先写出图右侧第一行的几何或动力学关系，再代入其下方的计算式。
 令 $\vect n=\sin\theta\,\vect e_r+\cos\theta\,\vect k$，则
\[
\vect\omega_{\text{总}}
=\omega\sin\theta\,\vect e_r+(\Omega+\omega\cos\theta)\vect k,
\quad
\omega_{\text{总}}^2=\omega^2+\Omega^2+2\omega\Omega\cos\theta.
\]

\SolutionFigure{ex05}{例5的计算示意}
\Step{逐项计算与求解}
投影到对称轴得到
$\vect\omega_{\text{总}}\cdot\vect n=\omega+\Omega\cos\theta$。
相对自转速率与总角速度沿对称轴的分量不是同一个量。


把总角速度向对称轴方向和平面内垂直方向分解。
在 $\vect n$ 与 $\vect k$ 张成的平面中，
令 $\vect t$ 沿增大轴倾角的方向，满足
$\vect k=\cos\theta\,\vect n-\sin\theta\,\vect t$。
于是
\[
\vect\omega_{\text{总}}
=(\omega+\Omega\cos\theta)\vect n-\Omega\sin\theta\,\vect t.
\]
第一项使盘面内的材料方向绕本体轴转动，
第二项负责改变本体轴在空间中的方向。
把两项平方相加，
\begin{align*}
|\vect\omega_{\text{总}}|^2
&=(\omega+\Omega\cos\theta)^2+\Omega^2\sin^2\theta\\
&=\omega^2+\Omega^2+2\omega\Omega\cos\theta.
\end{align*}

\Step{结果与适用条件}
这也说明即使轴向投影在某一瞬间抵消，物体仍可能有非零横向角速度。
当 $\theta=0$，两转轴同向，合角速度为 $(\omega+\Omega)\vect k$；
当 $\theta=\pi/2$，两分角速度正交，大小为平方和的平方根。
有向速率允许 $\omega$、$\Omega$ 取负值，不能在叉积运算前都改成绝对值。

% END DETAILED ANSWER ex05
\example{悬挂杆端圆盘的两种约束}\label{ex:c3-6}
% BEGIN ORIGINAL FIGURES
\OriginalFigure[6.2cm]{c3-1-p012-hanging-disks}{悬挂杆端圆盘及两种轴线约束}
% END ORIGINAL FIGURES
\QuestionStatement{悬挂点附近的支架绕竖直轴以 $\Omega$ 转动，杆与竖直方向成固定夹角 $\beta$。
比较两种机构：第一种圆盘平面始终垂直于杆，并相对支架绕杆以 $\omega$ 转动；
第二种用万向机构保持圆盘水平，其相对支架的方位角以 $\omega$ 增长。
这里的两种 $\omega$ 均相对于支架测量。求各机构的绝对角速度及大小，比较 $\beta=0$ 时的结果。}{ex06}{例6题图}
% BEGIN DETAILED ANSWER ex06
\solution \Step{明确两个相对转轴}
以图~\ref{fig:solution-ex06}~中两种盘面法线对应的不同转轴为依据，先写出图右侧第一行的几何或动力学关系，再代入其下方的计算式。

角速度合成要把相对转动和支架转动写在同一空间中。取竖直向上为 $\vect k$；第一种盘法线 $\vect n$ 与它夹角 $\beta$，故
$\vect n\cdot\vect k=\cos\beta$。第二种盘法线始终为 $\vect k$。
\SolutionFigure{ex06}{例6的计算示意}
\Step{求垂直于杆的圆盘}
第一种机构有
\[
\vect\omega_1=\Omega\vect k+\omega\vect n,\qquad
|\vect\omega_1|=\sqrt{\Omega^2+\omega^2+2\Omega\omega\cos\beta}.
\]
若 $\Omega,\omega$ 为常数，$\dot{\vect n}=\Omega\vect k\times\vect n$，所以
$\vect\alpha_1=\Omega\omega\vect k\times\vect n$。
其非零值来自杆方向的改变；盘面垂直于杆，不能把盘法线当作竖直方向。
\Step{求保持水平的圆盘并比较}
第二种机构的两次姿态转动都绕竖直方向，因此
\[
\vect\omega_2=(\Omega+\omega)\vect k,\qquad
|\vect\omega_2|=|\Omega+\omega|,\qquad \vect\alpha_2=0.
\]
杆的倾斜改变圆盘中心的平移轨迹，却不改变已经由万向约束确定的盘面方向。
当 $\beta=0$，两式一致；当 $\omega=-\Omega$，第二种圆盘姿态在空间固定，第一种一般仍有转动。
若 $\omega$ 改为相对地面测得的绝对盘转速，第二式应直接写 $\omega\vect k$，不可再加支架转速。

% END DETAILED ANSWER ex06
\subsubsection*{两种悬挂机构中轴向约束的差别}
第一种机构中圆盘法线为杆轴 $\vect n$，
盘面倾斜并随支架进动，故
\[
|\vect\omega_1|^2=\Omega^2+\omega^2+2\Omega\omega\cos\beta,
\quad
(\vect\omega_1\cdot\vect n)=\omega+\Omega\cos\beta.
\]
第二种机构中万向连接允许杆与盘相对改变夹角，
盘面法线仍竖直，故只有同向的两个转动：
\[
|\vect\omega_2|=|\Omega+\omega|,
\qquad
\vect\alpha_2=(\dot\Omega+\dot\omega)\vect k.
\]
圆盘中心的位置可以随倾斜杆绕竖直轴走圆，
但这不会给水平圆盘自动增加一项沿杆轴的角速度。
圆心平移由 $\vect r_A(t)$ 决定，圆盘姿态由盘上的两个材料方向决定；
必须分别写出，不能只从支杆形状猜测圆盘自转轴。


\subsection{摆杆转盘的完整速度计算}\label{sub:c3-2.5}
\example{铰接摆杆端的旋转圆盘}\label{ex:c3-7}

\QuestionStatement{长 $l$ 的杆 $OA$ 绕固定水平铰轴以有向角速率 $\omega_1$ 转动。
杆端装半径 $a$ 的圆盘，盘面始终垂直杆，相对杆绕杆以 $\omega_2$ 转动。
$t=0$ 杆竖直向下，圆盘上半径 $AC$ 沿水平铰轴正向，
半径 $AB$ 垂直铰轴并在水平面内，初始沿 $\vect e_2=\vect e_3\times\vect e_1$，其中 $\vect e_1$ 沿杆向下、$\vect e_3$ 沿铰轴正向。求 $B,C$ 点随时间的速度。
两转速均取恒定有向值。给出任意时刻两点的速度，并写出 $t=0$ 时的结果以检验符号。}{ex07}{例7题图}
% BEGIN DETAILED ANSWER ex07
\solution \Step{明确模型与计算依据}
以图~\ref{fig:solution-ex07}~中圆心 $A$ 及材料点 $B,C$ 的有向半径为依据，先写出图右侧第一行的几何或动力学关系，再代入其下方的计算式。
 取右手随杆轴组：$\vect e_1$ 沿 $OA$，$\vect e_3$ 沿水平铰轴，
$\vect e_2=\vect e_3\times\vect e_1$。
杆倾角 $q=\omega_1t$，圆盘相对转角 $s=\omega_2t$，则
\[
\vect r_{AB}=a(\cos s\,\vect e_2+\sin s\,\vect e_3),\qquad
\vect r_{AC}=a(-\sin s\,\vect e_2+\cos s\,\vect e_3).
\]

\SolutionFigure{ex07}{例7的计算示意}
\Step{逐项计算与求解}
总角速度为 $\vect\omega=\omega_1\vect e_3+\omega_2\vect e_1$，
圆心速度 $\vect v_A=l\omega_1\vect e_2$。利用右手叉积，
\begin{align*}
\vect v_B&=-a\omega_1\cos s\,\vect e_1
+(l\omega_1-a\omega_2\sin s)\vect e_2
+a\omega_2\cos s\,\vect e_3,\\
\vect v_C&=a\omega_1\sin s\,\vect e_1
+(l\omega_1-a\omega_2\cos s)\vect e_2
-a\omega_2\sin s\,\vect e_3.
\end{align*}

\Step{结果与适用条件}
在最低位置且 $s=0$ 时，
$\vect v_B=-a\omega_1\vect e_1+l\omega_1\vect e_2+a\omega_2\vect e_3$，
$\vect v_C=(l\omega_1-a\omega_2)\vect e_2$。
两运动的速度可能在 $C$ 点抵消。
\Step{初始位置与单一转动极限} 题设 $AB$ 还须取向 $+\vect e_2$，才与上述有向相位一致。若 $\omega_2=0$，两点只随杆绕铰轴转动；若 $\omega_1=0$，圆心静止，两点速度大小均为 $a|\omega_2|$。两项的方向与初值均由同一个右手轴组确定。

% END DETAILED ANSWER ex07
\subsubsection*{摆杆转盘例的叉积展开}
随杆轴组满足
$\vect e_1\times\vect e_2=\vect e_3$、
$\vect e_2\times\vect e_3=\vect e_1$、
$\vect e_3\times\vect e_1=\vect e_2$。
对 $B$，分开计算摆杆与自转的相对速度：
\begin{align*}
\omega_1\vect e_3\times\vect r_{AB}
&=-a\omega_1\cos s\,\vect e_1,\\
\omega_2\vect e_1\times\vect r_{AB}
&=-a\omega_2\sin s\,\vect e_2+a\omega_2\cos s\,\vect e_3.
\end{align*}
然后加上 $\vect v_A=l\omega_1\vect e_2$，
得到 $\vect v_B$。
对 $C$，
\begin{align*}
\omega_1\vect e_3\times\vect r_{AC}
&=a\omega_1\sin s\,\vect e_1,\\
\omega_2\vect e_1\times\vect r_{AC}
&=-a\omega_2\cos s\,\vect e_2-a\omega_2\sin s\,\vect e_3.
\end{align*}
这两步显示了各分量的符号，
也说明为什么 $B,C$ 的速度不能简单用同一个速率乘圆盘半径。

\subsubsection*{把随杆表达变回空间表达}
令空间 $\vect i$ 沿铰轴，$\vect k$ 竖直向上，$\vect j$ 补成右手系。
则
\[
\vect e_1=\sin q\,\vect j-\cos q\,\vect k,\qquad
\vect e_2=\cos q\,\vect j+\sin q\,\vect k,\qquad
\vect e_3=\vect i.
\]
记 $V_B=l\omega_1-a\omega_2\sin s$、
$V_C=l\omega_1-a\omega_2\cos s$，有
\begin{align*}
\vect v_B={}&a\omega_2\cos s\,\vect i
 +(V_B\cos q-a\omega_1\cos s\sin q)\vect j\\
&+(V_B\sin q+a\omega_1\cos s\cos q)\vect k,\\
\vect v_C={}&-a\omega_2\sin s\,\vect i
 +(V_C\cos q+a\omega_1\sin s\sin q)\vect j\\
&+(V_C\sin q-a\omega_1\sin s\cos q)\vect k.
\end{align*}
这是完整的时间函数，$q=\omega_1t,s=\omega_2t$。
随体系给出紧凑的几何解释，空间系给出直接可观察的速度分量；
二者经上述轴变换严格一致。

若两速率恒定，轴组还满足
$\dot{\vect e}_1=\omega_1\vect e_2$、
$\dot{\vect e}_2=-\omega_1\vect e_1$。
对速度再次求导，得到可用于核查的加速度：
\begin{align*}
\vect a_B={}&(2a\omega_1\omega_2\sin s-l\omega_1^2)\vect e_1
-a(\omega_1^2+\omega_2^2)\cos s\,\vect e_2
-a\omega_2^2\sin s\,\vect e_3,\\
\vect a_C={}&(2a\omega_1\omega_2\cos s-l\omega_1^2)\vect e_1
+a(\omega_1^2+\omega_2^2)\sin s\,\vect e_2
-a\omega_2^2\cos s\,\vect e_3.
\end{align*}
交叉项 $2a\omega_1\omega_2$ 来自参考架与相对转动的耦合，
不是两个独立向心项能够代替的。
令 $\omega_2=0$，结果退化为整个机构随杆转动；
令 $\omega_1=0$，结果退化为圆心固定的圆盘自转。

在 $q=s=0$ 且 $l\omega_1=a\omega_2$ 时，$C$ 的瞬时速度为零，
但加速度仍由上式确定，一般并不为零。
这一具体例子再次区分了“材料点某瞬间静止”与“材料点始终固定”。


\TeachingFigure{c3-diagram-05}{最低位置的摆杆、圆盘及计算轴组}
% BEGIN ORIGINAL FIGURES
\OriginalFigure[6.2cm]{c3-1-p024-swing-disk}{铰接摆杆端的旋转圆盘与盘上标记点}
% END ORIGINAL FIGURES


\clearpage
\section{欧拉角、姿态矩阵与转动几何}\label{sec:c3-3}
\subsection{欧拉角的几何定义与取向约定}\label{sub:c3-3.1}
设固定空间轴为 $(\xi,\eta,\zeta)$，本体轴为 $(x,y,z)$。
当本体 $z$ 轴不与 $\zeta$ 轴平行时，定义正向节点轴
\[
\vect e_N=\frac{\vect e_\zeta\times\vect e_z}
{|\vect e_\zeta\times\vect e_z|}.
\]
采用 $Z\!-\!X\!-\!Z$ 型欧拉角：
$\varphi$ 为绕空间 $\zeta$ 轴从 $\xi$ 轴到节点轴的角；
$\theta$ 为绕节点轴从 $\zeta$ 轴到本体 $z$ 轴的角；
$\psi$ 为绕本体 $z$ 轴从节点轴到本体 $x$ 轴的角。
角度均由各正轴的右手定则规定。依次完成进动、倾斜、自转，
就得到本体姿态。

\TeachingFigure{c3-diagram-06}{欧拉角的三个定义平面；从各正转轴端看向原点，按右手定则取正}
% BEGIN ORIGINAL FIGURES
\OriginalFigure[6.2cm]{c3-1-p014-euler-angles}{空间轴、本体轴与三个欧拉角：角弧用于识别夹角，转动正向按正文右手定则}
% END ORIGINAL FIGURES

这些欧拉角描述姿态，不要求物体真实运动恰好依次完成三次转动。
$\theta$ 又称章动角，$\varphi$ 为进动角，$\psi$ 为自转角。
“章动”在动力学中还常指 $\theta$ 的周期变化；固定的章动角意味着没有此种摆动。

本体 $z$ 轴在空间中若有极角 $\theta_z$ 和方位角 $\varphi_z$，则本约定下
\[
\theta=\theta_z,\qquad \varphi=\varphi_z+\frac{\pi}{2}\pmod{2\pi}.
\]
节点轴比 $z$ 轴的水平投影超前 $\pi/2$。
最后再在本体横截面内测量 $\psi$，才能确定完整姿态；
单知道对称轴方向还不能确定标记点在盘上的位置。


\subsubsection*{用节点轴逐步建立正交轴组}
先在空间水平面内写节点方向及与它垂直的方向：
\[
\vect e_N=\cos\varphi\,\vect e_\xi+\sin\varphi\,\vect e_\eta,
\qquad
\vect e_H=-\sin\varphi\,\vect e_\xi+\cos\varphi\,\vect e_\eta.
\]
两者满足 $\vect e_N\times\vect e_H=\vect e_\zeta$。
第二次转动绕 $\vect e_N$ 进行，
将 $\vect e_H,\vect e_\zeta$ 分别变成
\[
\vect e_{H'}=\cos\theta\,\vect e_H+\sin\theta\,\vect e_\zeta,
\qquad
\vect e_z=-\sin\theta\,\vect e_H+\cos\theta\,\vect e_\zeta.
\]
注意倾斜后本体 $z$ 轴的水平投影沿 $-\vect e_H$，
其方位角为 $\varphi-\pi/2$。
最后绕 $\vect e_z$ 转角 $\psi$，得到
\[
\vect e_x=\cos\psi\,\vect e_N+\sin\psi\,\vect e_{H'},\qquad
\vect e_y=-\sin\psi\,\vect e_N+\cos\psi\,\vect e_{H'}.
\]
这一组表达可以逐项验算
$\vect e_x\cdot\vect e_y=0$、
$\vect e_x\cdot\vect e_z=0$、
$\vect e_x\times\vect e_y=\vect e_z$。
由第三次转动的逆变换，
$\vect e_N=\cos\psi\vect e_x-\sin\psi\vect e_y$，
后续角速度投影中的节点分量便有了明确来源。

\TeachingFigure{c3-diagram-07}{欧拉姿态的三步构造；虚线保留初始空间轴供比较}
% BEGIN ORIGINAL FIGURES
\OriginalFigure[6.2cm]{c3-1-p015-euler-steps}{进动、章动与自转的三步姿态构造；转动正向按正文右手定则}
% END ORIGINAL FIGURES


\subsubsection*{怎样从可见标记确定三个角}
首先观察本体 $z$ 轴，测定其与空间竖直方向的夹角 $\theta$；
再由 $\vect e_\zeta\times\vect e_z$ 求正向节点。
用节点相对于 $\xi$ 的方位确定 $\varphi$，
最后观察本体 $x$ 轴在垂直于 $z$ 的平面中的方向，
用从节点到 $x$ 的有向角确定 $\psi$。
圆盘若没有任何盘面标记，外形只能提供法线方向，
无法从照片独立测出 $\psi$；这不是数学自由度消失，而是观测量不足。

$\varphi$ 与 $\psi$ 通常可按连续运动取超出 $2\pi$ 的展开角，
以保留累计转动次数。
若只研究瞬时姿态，则可把它们取模 $2\pi$。
两种记录方式的矩阵姿态相同，却适用于不同的时间分析。


\subsection{欧拉运动学方程及其反解}\label{sub:c3-3.2}
三个瞬时分转动的角速度相加：
\[
\vect\omega=\dot\varphi\,\vect e_\zeta
+\dot\theta\,\vect e_N+\dot\psi\,\vect e_z.
\]
节点轴和空间竖直轴在本体轴上的展开为
\[
\vect e_N=\cos\psi\,\vect e_x-\sin\psi\,\vect e_y,\qquad
\vect e_\zeta=\sin\theta\sin\psi\,\vect e_x
+\sin\theta\cos\psi\,\vect e_y+\cos\theta\,\vect e_z.
\]
由此得到
\begin{equation}\label{eq:eulerkin}
\begin{aligned}
\omega_x&=\dot\varphi\sin\theta\sin\psi+\dot\theta\cos\psi,\\
\omega_y&=\dot\varphi\sin\theta\cos\psi-\dot\theta\sin\psi,\\
\omega_z&=\dot\varphi\cos\theta+\dot\psi.
\end{aligned}
\end{equation}
已知角度函数便能求本体分量。反过来，当 $\sin\theta\ne0$ 时，
\begin{equation}\label{eq:eulerinverse}
\begin{aligned}
\dot\theta&=\omega_x\cos\psi-\omega_y\sin\psi,\\
\dot\varphi&=\frac{\omega_x\sin\psi+\omega_y\cos\psi}{\sin\theta},\\
\dot\psi&=\omega_z-\dot\varphi\cos\theta.
\end{aligned}
\end{equation}
同时，角速度大小满足
\[
\omega^2=\dot\theta^2+\dot\varphi^2+\dot\psi^2
+2\dot\varphi\dot\psi\cos\theta.
\]
$\dot\psi$ 是第三次相对转动的速率，而 $\omega_z$ 是总角速度在本体 $z$ 轴上的投影。

当 $\theta=0$ 时节点轴没有唯一方向，只能确定 $\varphi+\psi$；
当 $\theta=\pi$ 时只能确定相应的角度差。
式\eqref{eq:eulerinverse}的发散属于坐标参数的奇异性，
并不说明刚体的真实角速度发散。姿态矩阵仍然可以保持光滑。
欧拉角还有 $2\pi$ 周期的多值性，故不能在全部姿态上声称三个角唯一。


\subsubsection*{欧拉运动学三式的逐项投影}
绕 $\zeta$ 的进动速度在本体轴上的投影分别是
\[
\dot\varphi\sin\theta\sin\psi,\qquad
\dot\varphi\sin\theta\cos\psi,\qquad
\dot\varphi\cos\theta.
\]
绕节点轴的章动速度投影为
$\dot\theta\cos\psi,-\dot\theta\sin\psi,0$；
绕本体轴的自转速度投影为 $0,0,\dot\psi$。
把三行贡献逐分量相加，就得到式\eqref{eq:eulerkin}。
因此本体 $z$ 分量中的 $\dot\varphi\cos\theta$ 是进动投影，
不是附加的一次独立转动。

反解时不必求三阶矩阵的通用逆。
把前两式分别乘 $\cos\psi,-\sin\psi$ 后相加：
\[
\omega_x\cos\psi-\omega_y\sin\psi
=\dot\theta(\cos^2\psi+\sin^2\psi)=\dot\theta.
\]
再分别乘 $\sin\psi,\cos\psi$ 后相加：
\[
\omega_x\sin\psi+\omega_y\cos\psi
=\dot\varphi\sin\theta.
\]
最后由第三式扣除进动投影，得到 $\dot\psi$。
每一步都可看出消去的是哪个分转动。

用同一角速度矢量在空间轴上直接展开，还有
\[
\begin{aligned}
\omega_\xi&=\dot\theta\cos\varphi+\dot\psi\sin\theta\sin\varphi,\\
\omega_\eta&=\dot\theta\sin\varphi-\dot\psi\sin\theta\cos\varphi,\\
\omega_\zeta&=\dot\varphi+\dot\psi\cos\theta.
\end{aligned}
\]
本体与空间的第三分量不同，
因为投影到的是两根不同方向的轴。

\subsubsection*{角速度平方与欧拉角速度的耦合}
前两个本体分量的平方和为
\begin{align*}
\omega_x^2+\omega_y^2
&=\dot\varphi^2\sin^2\theta(\sin^2\psi+\cos^2\psi)
+\dot\theta^2(\cos^2\psi+\sin^2\psi)\\
&\quad+2\dot\varphi\dot\theta\sin\theta
(\sin\psi\cos\psi-\cos\psi\sin\psi)\\
&=\dot\varphi^2\sin^2\theta+\dot\theta^2.
\end{align*}
再加 $\omega_z^2=(\dot\varphi\cos\theta+\dot\psi)^2$，
就得到角速度大小关系。
交叉项 $2\dot\varphi\dot\psi\cos\theta$ 表示进动轴和自转轴夹角为 $\theta$；
只有章动轴与它们都垂直。
不能把三种欧拉角速率平方直接相加来代替 $\omega^2$。

\subsubsection*{奇异性从哪里产生}
将欧拉运动学写成
\[
\begin{pmatrix}\omega_x\\\omega_y\\\omega_z\end{pmatrix}
=
\begin{pmatrix}
\sin\theta\sin\psi&\cos\psi&0\\
\sin\theta\cos\psi&-\sin\psi&0\\
\cos\theta&0&1
\end{pmatrix}
\begin{pmatrix}\dot\varphi\\\dot\theta\\\dot\psi\end{pmatrix}.
\]
此矩阵的行列式为 $-\sin\theta$。
当 $\theta=0$，进动轴和自转轴重合，
$\dot\varphi$ 与 $\dot\psi$ 两列对同一角速度产生相同贡献；
当 $\theta=\pi$，两根轴反向，只能从姿态中区分角度差。
这说明奇异性来自三个参数的瞬时独立性丢失，
而不是来自刚体物理运动不可继续。

数值计算中若已知角速度且运动将穿过极点，
可以直接积分
$\dot{\mathsf R}=\mathsf W\mathsf R$，
并以正交矩阵记录姿态。
只有需要输出欧拉角时，才进行带分支管理的反解。


\example{匀速进动陀螺的本体分量}\label{ex:c3-8}
% BEGIN ORIGINAL FIGURES
\OriginalFigure[6.2cm]{c3-1-p019-top-initial}{节点线与本体轴重合的一种陀螺初态}
% END ORIGINAL FIGURES
\QuestionStatement{取 $\theta=\theta_0$，
$\varphi=\pi/2+\Omega t$，$\psi=-\pi/2+\omega t$。
求角速度在本体坐标系的分量及初始空间分量。
说明 $\omega_z$ 与欧拉角自转率 $\dot\psi$ 是否相同。}{ex08}{例8题图}
% BEGIN DETAILED ANSWER ex08
\solution \Step{明确模型与计算依据}
以图~\ref{fig:solution-ex08}~中节点轴和本体 $z$ 轴的投影关系为依据，先写出图右侧第一行的几何或动力学关系，再代入其下方的计算式。
 \[
\omega_x=-\Omega\sin\theta_0\cos\omega t,\quad
\omega_y=\Omega\sin\theta_0\sin\omega t,\quad
\omega_z=\Omega\cos\theta_0+\omega.
\]
初始时 $z$ 轴位于 $\xi\zeta$ 平面且向 $\xi$ 倾斜，
故空间分量为
\[
\vect\omega(0)=\omega\sin\theta_0\,\vect e_\xi+
(\Omega+\omega\cos\theta_0)\vect e_\zeta.
\]

\SolutionFigure{ex08}{例8的计算示意}
\Step{逐项计算与求解}
两种表达使用不同轴组，不能直接逐项比较数值；乘姿态矩阵后应一致。


本例的三个速率为 $\dot\theta=0,\dot\varphi=\Omega,\dot\psi=\omega$。
代入 $\sin(-\pi/2+\omega t)=-\cos\omega t$、
$\cos(-\pi/2+\omega t)=\sin\omega t$，
即可得到三个本体分量。
横向分量在本体系中按相对自转速率 $\omega$ 变化，
但空间中的本体轴按进动速率 $\Omega$ 绕竖直转动。
它们描述不同参考系中的方向变化，因此没有矛盾。

按当前本体轴求角加速度分量时，
$\vect\omega\times\vect\omega=0$，所以只需对本体分量求导：
\[
\alpha_x=\Omega\omega\sin\theta_0\sin\omega t,\qquad
\alpha_y=\Omega\omega\sin\theta_0\cos\omega t,\qquad
\alpha_z=0.
\]

\Step{结果与适用条件}
空间表达则为
\[
\vect\alpha
=\Omega\omega\sin\theta_0
(-\sin\Omega t\,\vect e_\xi+\cos\Omega t\,\vect e_\eta).
\]
两者大小均为 $|\Omega\omega\sin\theta_0|$。
若误把三个本体分量都视为空间固定轴上的分量，
就会把真实轴转向与坐标分量变化混淆。
\Step{明确轴向分量的含义} 本体轴向分量是 $\Omega\cos\theta_0+\omega$，欧拉角自转率为 $\omega$，差值来自进动在本体轴上的投影。只有 $\Omega\cos\theta_0=0$ 时两者才相同。

% END DETAILED ANSWER ex08
\example{理想地球的姿态描述}\label{ex:c3-9}

\QuestionStatement{忽略自转轴的长期改变，以轨道法线为固定 $\zeta$ 轴，
把地球自转轴作为本体 $z$ 轴。设轴倾角为常数 $\varepsilon$，
固定其方位使节点角为 $\varphi_0$，地球日转角速率为 $\omega$。
采用本章 $z$--$x$--$z$ 欧拉角约定，设初始自转角为 $\psi_0$。求三个欧拉角和总角速度，说明年周公转是否须计入姿态自转。}{ex09}{例9题图}
% BEGIN DETAILED ANSWER ex09
\solution \Step{分开质心轨道与本体姿态}
以图~\ref{fig:solution-ex09}~中固定轴倾角 $\varepsilon$ 与地球本体轴 $z$为依据，先写出图右侧第一行的几何或动力学关系，再代入其下方的计算式。

地球年周运动确定质心位置 $\vect r_C(t)$；刚体姿态由自转轴和地表标记相对于空间轴的方向确定。题设自转轴不进动，节点方向固定，因而
$\theta(t)=\varepsilon$、$\varphi(t)=\varphi_0$。
\SolutionFigure{ex09}{例9的计算示意}
\Step{利用初值积分自转角}
由 $\dot\psi=\omega$、$\psi(0)=\psi_0$，
\[
\psi(t)=\psi_0+\omega t,\qquad
\mathsf R(t)=R_z(\varphi_0)R_x(\varepsilon)R_z(\psi_0+\omega t).
\]
它把固定轴倾角、固定方位和地表标记的日转动分别表示出来。
\Step{由欧拉运动学求总角速度}
一般合成为
$\vect\omega=\dot\varphi\vect e_\zeta+\dot\theta\vect e_N+\dot\psi\vect e_z$，故本题
\[
\vect\omega=\omega\vect e_z,\qquad
(\omega_x,\omega_y,\omega_z)=(0,0,\omega).
\]
惯性系中 $\vect e_z$ 固定，因此恒定 $\omega$ 还意味着 $\vect\alpha=0$。
若今后加入慢进动 $\dot\varphi\ne0$，则轴向分量改为
$\omega_z=\dot\psi+\dot\varphi\cos\varepsilon$。
公转不会自动产生额外的本体自转项；只有另用随公转的坐标架描述时，才须作参考系角速度换算。

% END DETAILED ANSWER ex09
\TeachingFigure{c3-diagram-08}{理想地球模型中的轴倾角与质心轨道}
% BEGIN ORIGINAL FIGURES
\OriginalFigure[6.2cm]{c3-1-p018-earth-orbit}{地球公转轨道、季节位置与黄赤交角}
% END ORIGINAL FIGURES

这个近似模型的轴倾角约为 $23^\circ26'$。
太阳和月球对地球赤道隆起的引力矩还使地轴发生缓慢岁差，
绕黄道面法线一周约需 $2.6$ 万年，
按 $360^\circ$ 除以该周期计算，对应方位角每年变化约 $50''$。
这是进动方位角的变化率，不是轴倾角每年的变化量；
轴端点球面速度还含 $\sin\theta$ 因子。
真实地球还叠加章动及长期轴倾角变化。
\footnote{近似轴倾角与轴向岁差参见
\href{https://nssdc.gsfc.nasa.gov/planetary/factsheet/earthfact.html}{NASA Earth Fact Sheet}及
\href{https://cdaweb.gsfc.nasa.gov/pub/documents/archived_websites/pwg.gsfc.nasa.gov/stargaze/Sprecess.htm}{NASA Precession}。}
东晋虞喜在中国天文学中识别并讨论了岁差，
南朝祖冲之将岁差纳入《大明历》的计算。
历史上的岁差认识与观测历法精度密切相关。


\subsubsection*{地球模型中的三种周期各自决定什么}
记自转周期为 $T_s$、质心绕太阳的公转周期为 $T_o$、
地轴方位进动周期为 $T_p$。
在理想圆轨道、固定倾角的近似中，
\[
\omega_s=\frac{2\pi}{T_s},\qquad
\Omega_o=\frac{2\pi}{T_o},\qquad
|\Omega_p|=\frac{2\pi}{T_p}.
\]
$\Omega_o$ 进入质心位置 $\vect r_C(t)$，
$\omega_s$ 与有向 $\Omega_p$ 进入姿态角速度。
令 $\dot\varphi=\Omega_p,\dot\theta=0,\dot\psi=\omega_s$，
则
\[
\omega_z=\omega_s+\Omega_p\cos\varepsilon,\qquad
\omega_x^2+\omega_y^2=\Omega_p^2\sin^2\varepsilon.
\]
若题目给定的自转速率已经定义为总轴向分量 $\omega_z$，
则应令 $\dot\psi=\omega_z-\Omega_p\cos\varepsilon$，
避免重复计入岁差投影。

按约 $2.6$ 万年一周的近似，方位进动速率大小约为
\[
\frac{360^\circ}{26000\,\mathrm{yr}}
\simeq0.01385^\circ/\mathrm{yr}
\simeq49.85''/\mathrm{yr}.
\]
若单位轴端点为 $\vect n(t)$，其单位球面路程速率是
$|\dot{\vect n}|=|\Omega_p|\sin\varepsilon$。
所以方位角变化率与轴端点的球面角距离变化率还差一个 $\sin\varepsilon$。
这里的模型分别控制位置与姿态，只用于理解运动学；
真实地球的长期倾角变化和章动须作为额外角度函数另行加入。

\footnote{人物与历法史参见
\href{https://mathshistory.st-andrews.ac.uk/BEA/zu_chongzhi_bea.pdf}{Biographical Encyclopedia of Astronomers：Zu Chongzhi}。}


\example{同步自转卫星}\label{ex:c3-10}

\QuestionStatement{卫星在同一平面内做匀速圆周运动，角速率为 $\Omega$，
始终以同一面朝向中心天体。取本体 $z$ 轴与轨道法线同向。
设轨道半径为 $R$，$t=0$ 时卫星位于 $+X$ 轴，标记体轴 $\vect e_1$ 指向中心天体。求惯性系中的姿态、角速度及周期，并说明“始终同面”与自转的关系。}{ex10}{例10题图}
% BEGIN DETAILED ANSWER ex10
\solution \Step{用朝向约束直接建立体轴}
以图~\ref{fig:solution-ex10}~中由卫星指向 $O$ 的标记方向为依据，先写出图右侧第一行的几何或动力学关系，再代入其下方的计算式。

令卫星轨道位置为
$\vect r_C=R(\cos\Omega t,\sin\Omega t,0)$。
朝向中心的体轴为
\[
\vect e_1=(-\cos\Omega t,-\sin\Omega t,0),\quad
\vect e_3=\vect k,\quad
\vect e_2=\vect e_3\times\vect e_1=(\sin\Omega t,-\cos\Omega t,0).
\]
这三根轴保持右手正交，并满足题设的朝向条件。
\SolutionFigure{ex10}{例10的计算示意}
\Step{写出姿态矩阵与角速度}
以三根体轴为空间分量列向量，得到
\[
\mathsf R(t)=R_z(\pi+\Omega t),\qquad
\vect\omega=\Omega\vect k.
\]
在欧拉角奇异位置，可取 $\theta=0,\varphi=0,\psi=\pi+\Omega t$。
此时单独的 $\varphi,\psi$ 不唯一，能够确定的是它们的和。
\Step{两种观察系与运动周期}
惯性系中自转周期为 $2\pi/|\Omega|$，与公转周期相同；每绕行一周，朝内的标记方向也转一周。
随轨道架以 $\Omega\vect k$ 转动时，卫星相对于它的角速度为零。
例如轨道转过 $\pi/2$，$\vect e_1$ 从 $-\vect i$ 变到 $-\vect j$，其空间方向已改变。
因此“同一面朝向中心”规定的是相对轨道架的固定姿态，同时要求惯性系中的同步自转。

% END DETAILED ANSWER ex10
\subsubsection*{同步卫星的标记点运动}
取卫星质心圆轨道为
$\vect r_C=R\vect e_r(t)$，
$\dot{\vect e}_r=\Omega\vect e_t$。
朝向中心天体的材料点 $A$ 相对于卫星质心位置为
$\vect r_{CA}=-b\vect e_r(t)$，故
\[
\vect r_A=(R-b)\vect e_r,\qquad
\vect v_A=(R-b)\Omega\vect e_t.
\]
也可从刚体速度场计算：
\[
\vect v_A=\vect v_C+\Omega\vect k\times(-b\vect e_r)
=R\Omega\vect e_t-b\Omega\vect e_t.
\]
结果一致。
若卫星完全不自转，朝向中心的材料方向会在一周期内改变，
而不是始终由同一个标记面朝向中心。

理想同步关系 $\omega_z=\Omega$ 对应圆轨道、共面的姿态模型。
椭圆轨道的轨道角速率随时间变化，
恒定自转速率与瞬时轨道速率不再逐时相等，
此时应区分平均同步与瞬时朝向变化。
这一限定来自轨道和姿态两个角度函数的区别。


\subsection{滚动圆锥与曲柄连杆}\label{sub:c3-3.3}
\example{固定顶点圆锥的无滑动滚动}\label{ex:c3-11}

\QuestionStatement{直圆锥顶点固定在水平面上，轴向高度 $h$，半顶角 $0<\alpha<\pi/2$。
圆锥沿一条母线接触地面并无滑动滚动。
以顶点为原点，竖直向上为 $\zeta$，
$t=0$ 时接触母线沿 $+\eta$ 方向。
圆锥相对于进动架绕对称轴的有向自转速率为常数 $\omega$。
求欧拉角、角速度本体分量及底面圆周最高点的瞬时速度。
采用本章的 $z$--$x$--$z$ 约定，初始本体 $x$ 轴取竖直方向在垂直于对称轴平面内的投影。}{ex11}{例11题图}
% BEGIN DETAILED ANSWER ex11
\solution \Step{明确模型与计算依据}
以图~\ref{fig:solution-ex11}~中接触母线 $\vect u$ 与最高圆周位置 $A$为依据，先写出图右侧第一行的几何或动力学关系，再代入其下方的计算式。
 对称轴与水平面的夹角为 $\alpha$，
$\theta=\pi/2-\alpha$。
初始方位为 $\pi/2$，故 $\varphi(0)=\pi$。
选初始本体 $x$ 轴沿竖直方向在垂直于 $z$ 轴平面内的投影，
则 $\psi(0)=-\pi/2$。这个投影与 $z$ 轴正交，才能构成正交本体系。
滚动瞬时转轴是接触母线，必水平，故
\[
\dot\varphi+\omega\cos\theta=0,\qquad
\dot\varphi=-\omega\sin\alpha.
\]
欧拉角及本体分量为
\[
\theta=\frac\pi2-\alpha,\quad
\varphi=\pi-\omega\sin\alpha\,t,\quad
\psi=-\frac\pi2+\omega t,
\]

\SolutionFigure{ex11}{例11的计算示意}
\Step{逐项计算与求解}
\[
\omega_x=-\omega\sin\alpha\cos\alpha\sin\psi,\quad
\omega_y=-\omega\sin\alpha\cos\alpha\cos\psi,\quad
\omega_z=\omega\cos^2\alpha.
\]
设接触母线水平单位方向为 $\vect u$，则
$\vect e_z=\cos\alpha\,\vect u+\sin\alpha\,\vect e_\zeta$，
总角速度为 $\vect\omega=\omega\cos\alpha\,\vect u$。

底面半径 $b=h\tan\alpha$。最高圆周位置相对于顶点为
\[
\vect r_A=\frac{h\cos2\alpha}{\cos\alpha}\vect u
+2h\sin\alpha\,\vect e_\zeta.
\]

\Step{结果与适用条件}
于是
\[
\vect v_A=\vect\omega\times\vect r_A
=2h\omega\sin\alpha\cos\alpha\,(\vect u\times\vect e_\zeta),
\qquad |v_A|=h|\omega|\sin2\alpha.
\]
这指的是恰在最高位置的体内质点的瞬时速度，
不是对每一时刻几何最高点的位置函数直接求导。
\Step{滚动轴与速度量纲} 接触母线上任一点的位置平行于 $\vect\omega$，因此其速度严格为零；最高圆周材料点的速度则垂直接触母线。$h|\omega|\sin2\alpha$ 的量纲为长度除时间，与速度一致。

% END DETAILED ANSWER ex11
\subsubsection*{圆锥接触条件为什么决定进动率}
固定顶点 $O$ 和接触母线上的另一个材料点在接触瞬间都静止。
非零角速度必须沿这两点连线，因此瞬时转轴沿母线。
令
\[
\vect e_z=\cos\alpha\,\vect u+\sin\alpha\,\vect e_\zeta,
\]
合角速度由进动与相对自转相加：
\[
\vect\omega_{\text{总}}
=\dot\varphi\vect e_\zeta+\omega\vect e_z
=\omega\cos\alpha\,\vect u+
(\dot\varphi+\omega\sin\alpha)\vect e_\zeta.
\]
要让它完全水平，括号内必须为零。
这一步给出 $\dot\varphi=-\omega\sin\alpha$，
并留下 $|\vect\omega_{\text{总}}|=|\omega|\cos\alpha$。
无滑动约束决定的是两个有向速率之间的关系，
不是把相对自转速率直接当作总角速率。

母线长为 $g=h/\cos\alpha$。
把圆锥侧面沿一条母线展开成扇形，
扇形半径为 $g$，弧长等于底面周长 $2\pi b$。
所以扇形中心角为
\[
\frac{2\pi b}{g}=2\pi\sin\alpha.
\]
这一几何关系与进动、自转的速率比例相容，
为滚动方向和大小提供独立核查。

\subsubsection*{最高圆周点的位置逐步求出}
轴截面内、垂直于对称轴且向上的单位方向为
\[
\vect m=-\sin\alpha\,\vect u+\cos\alpha\,\vect e_\zeta,
\qquad
\vect m\cdot\vect e_z=0.
\]
底面中心位置为 $h\vect e_z$。
最高点与最低点相对于底面中心分别为 $+b\vect m$ 和 $-b\vect m$。
因此
\begin{align*}
\vect r_A
&=h\vect e_z+b\vect m\\
&=(h\cos\alpha-b\sin\alpha)\vect u
 +(h\sin\alpha+b\cos\alpha)\vect e_\zeta\\
&=\frac{h\cos2\alpha}{\cos\alpha}\vect u
 +2h\sin\alpha\,\vect e_\zeta.
\end{align*}
最低圆周点则为
\[
\vect r_P=h\vect e_z-b\vect m
=\frac h{\cos\alpha}\vect u.
\]
它恰好位于水平接触母线上，且
$\vect\omega_{\text{总}}\times\vect r_P=0$，
从几何与速度两方面满足无滑动。

最高点的轴平行坐标不贡献叉积，
故只有高度 $2h\sin\alpha$ 产生速度：
\[
|v_A|=|\omega|\cos\alpha\,(2h\sin\alpha)
=h|\omega|\sin2\alpha.
\]
当 $\alpha=\pi/4$，几何最高位置始终在固定顶点的正上方，
其位置矢量为 $\sqrt2h\,\vect e_\zeta$。
然而经过该位置的材料点仍有 $|v_A|=h|\omega|$。
这是“移动材料点”与“按位置重新选择的最高点”不能混同的直接例子。


\TeachingFigure{c3-diagram-09}{圆锥滚动的轴截面与最高圆周位置}
% BEGIN ORIGINAL FIGURES
\OriginalFigure[6.2cm]{c3-1-p025-rolling-cone}{圆锥在水平面上绕固定顶点滚动}
% END ORIGINAL FIGURES


\example{匀速曲柄带动的连杆}\label{ex:c3-12}

\QuestionStatement{曲柄 $OA$ 长 $a$，以角速率 $\Omega$ 在平面内匀速转动，
滑块 $B$ 沿过 $O$ 的竖直直线运动，连杆 $AB$ 长 $L>a$。
令曲柄转角 $q=\Omega t$，使 $A$ 的水平坐标为 $a\sin q$。
取本体 $x$ 轴从 $B$ 指向 $A$，初始沿向下的固定 $\xi$ 轴，
$\eta$ 轴向右，$\zeta$ 轴垂直运动平面向外。求连杆的姿态及角速度。
取 $B$ 位于 $A$ 上方的装配支路，初始 $BA$ 向下，连杆转角 $\psi\in(-\pi/2,\pi/2)$。}{ex12}{例12题图}
% BEGIN DETAILED ANSWER ex12
\solution \Step{从装配支路建立几何方程}
以图~\ref{fig:solution-ex12}~中连杆 $BA$ 向水平轴的投影为依据，先写出图右侧第一行的几何或动力学关系，再代入其下方的计算式。

以向下的 $\xi$ 轴和向右的 $\eta$ 轴为固定正交轴，$B$ 位于 $A$ 上方，且 $BA$ 初始沿 $+\xi$。设连杆相对 $+\xi$ 的有向角为 $\psi$，由水平方向投影
\[
L\sin\psi=a\sin q,\qquad q=\Omega t.
\]
题设 $L>a$ 保证 $|\sin\psi|<1$；取连续支路 $\cos\psi>0$，有
$\psi(t)=\arcsin[a\sin(\Omega t)/L]$，并满足 $\psi(0)=0$。
\SolutionFigure{ex12}{例12的计算示意}
\Step{对约束求导，而非令连杆匀速}
\[
L\cos\psi\,\dot\psi=a\Omega\cos\Omega t,\qquad
L\cos\psi=\sqrt{L^2-a^2\sin^2\Omega t}.
\]
所以
\[
\vect\omega_{AB}=\frac{a\Omega\cos\Omega t}{\sqrt{L^2-a^2\sin^2\Omega t}}\vect e_\zeta.
\]
连杆姿态可用奇异位置的一种欧拉角选择表示：$\theta=0,\varphi=0$，第三角为上述 $\psi$；实际可测量的平面姿态由总平面转角给出。
\Step{特殊位置与变化率}
$t=0$ 时 $\dot\psi=a\Omega/L$；$q=\pi/2$ 时 $\dot\psi=0$，连杆瞬时平动而非全体静止。
再求导得
\[
\ddot\psi=-\frac{a\Omega^2(L^2-a^2)\sin\Omega t}{(L^2-a^2\sin^2\Omega t)^{3/2}}.
\]
因此曲柄匀速不保证连杆匀速。当 $a/L\ll1$，$\psi\simeq(a/L)\sin\Omega t$，恢复小摆角近似。

% END DETAILED ANSWER ex12
\subsubsection*{曲柄机构的长度约束如何求导}
记 $d(q)=\sqrt{L^2-a^2\sin^2q}$。
连杆从 $B$ 到 $A$ 的水平投影为 $a\sin q$，沿 $\xi$ 的投影为 $d$，
故
\[
\sin\psi=\frac{a\sin q}{L},\qquad
\cos\psi=\frac{d}{L}>0.
\]
求导时保留 $\cos\psi$：
\[
\cos\psi\,\dot\psi=\frac aL\cos q\,\dot q,
\qquad
\dot\psi=\frac{a\Omega\cos q}{d}.
\]
对分母再求导，
$\dot d=-a^2\Omega\sin q\cos q/d$，于是
\begin{align*}
\ddot\psi
&=-\frac{a\Omega^2\sin q}{d}
+\frac{a^3\Omega^2\sin q\cos^2q}{d^3}\\
&=-\frac{a\Omega^2(L^2-a^2)\sin q}
{(L^2-a^2\sin^2q)^{3/2}}.
\end{align*}
曲柄虽匀速，连杆的角速率和角加速度均随相位改变。
在 $q=\pi/2$，$\dot\psi=0$，但
$\ddot\psi=-a\Omega^2/\sqrt{L^2-a^2}$ 非零，
说明连杆处于角运动的转向瞬间。

为了同时检查滑块运动，可以明确选取曲柄的相位原点，
令 $A_\xi=a\cos q,A_\eta=a\sin q$，并选 $B$ 位于 $A$ 的上方支路：
\[
B_\xi=a\cos q-d,\qquad B_\eta=0.
\]
这样 $B$ 到 $A$ 的投影与上述姿态约定一致。
其速度和加速度为
\[
v_{B\xi}=-a\Omega\sin q+\frac{a^2\Omega\sin q\cos q}{d},
\]
\[
a_{B\xi}=-a\Omega^2\cos q+
\frac{a^2\Omega^2\cos2q}{d}
+\frac{a^4\Omega^2\sin^2q\cos^2q}{d^3}.
\]
选择另一装配支路时根号项整体变号；
不能跨越死点时仍机械沿用同一反正弦分支。
$L>a$ 使当前装配支路在全周转动中保持非退化。
若 $L\gg a$，
$\psi\simeq(a/L)\sin q$，连杆角运动近似为小幅周期摆动。


\TeachingFigure{c3-diagram-10}{曲柄转角与连杆姿态之间由长度约束联系}
% BEGIN ORIGINAL FIGURES
\OriginalFigure[6.2cm]{c3-1-p026-crank}{曲柄与滑块：$\xi$ 向下、$\eta$ 向右，正文 $\psi$ 从 $\xi$ 量至本体 $x$ 轴}
% END ORIGINAL FIGURES


\subsection{基本转动矩阵与欧拉矩阵}\label{sub:c3-3.4}
% BEGIN ORIGINAL FIGURES
\OriginalFigure[6.2cm]{c3-1-p043-planar-rotation}{平面内矢量的主动转动：初始角 $\alpha$、转角 $\theta$ 对应正文 $\beta,\gamma$}
% END ORIGINAL FIGURES

本节统一采用主动转动：空间坐标轴固定，矢量被转动。
绕 $z$ 轴正转角 $\gamma$ 后，平面极角从 $\beta$ 变成 $\beta+\gamma$，
故 $x'=x\cos\gamma-y\sin\gamma$，
$y'=x\sin\gamma+y\cos\gamma$。三个基本矩阵为
\[
\mathsf R_x(\gamma)=
\begin{pmatrix}1&0&0\\0&\cos\gamma&-\sin\gamma\\0&\sin\gamma&\cos\gamma\end{pmatrix},
\quad
\mathsf R_y(\gamma)=
\begin{pmatrix}\cos\gamma&0&\sin\gamma\\0&1&0\\-\sin\gamma&0&\cos\gamma\end{pmatrix},
\]
\[
\mathsf R_z(\gamma)=
\begin{pmatrix}\cos\gamma&-\sin\gamma&0\\\sin\gamma&\cos\gamma&0\\0&0&1\end{pmatrix}.
\]
若转动的是坐标轴而矢量不动，分量变换使用逆矩阵
$\mathsf R_i(-\gamma)=\mathsf R_i^{\mathsf T}(\gamma)$。

欧拉构造中的节点轴和最终本体轴都在移动。
把转轴写回初始固定轴，可得到
\begin{equation}
\mathsf R(\varphi,\theta,\psi)
=\mathsf R_z(\varphi)\mathsf R_x(\theta)\mathsf R_z(\psi).
\end{equation}
右边矩阵从右向左作用于列矢量；
这与按移动轴依次实施 $\varphi,\theta,\psi$ 的几何构造等价。
用 $c_\varphi=\cos\varphi,s_\varphi=\sin\varphi$ 等记号，
\[
\mathsf R=
\begin{pmatrix}
c_\varphi c_\psi-s_\varphi c_\theta s_\psi&
-c_\varphi s_\psi-s_\varphi c_\theta c_\psi&s_\varphi s_\theta\\
s_\varphi c_\psi+c_\varphi c_\theta s_\psi&
-s_\varphi s_\psi+c_\varphi c_\theta c_\psi&-c_\varphi s_\theta\\
s_\theta s_\psi&s_\theta c_\psi&c_\theta
\end{pmatrix}.
\]
第三列给出本体 $z$ 轴方向，再次验证其方位角为 $\varphi-\pi/2$。
第三行则给出空间 $\zeta$ 轴在本体轴上的分量，
与欧拉运动学方程中的投影一致。


\subsubsection*{动轴转动与矩阵乘法次序的完整对应}
第一次转动为 $\mathsf Q_1=\mathsf R_z(\varphi)$。
第二次转轴是已经随 $\mathsf Q_1$ 转过的节点轴，
它的空间转动矩阵为
\[
\mathsf A_2=\mathsf Q_1\mathsf R_x(\theta)\mathsf Q_1^{\mathsf T}.
\]
完成两次转动后，
\[
\mathsf Q_2=\mathsf A_2\mathsf Q_1
=\mathsf Q_1\mathsf R_x(\theta).
\]
第三次转轴是 $\mathsf Q_2$ 的本体 $z$ 轴，
相应空间矩阵为
$\mathsf A_3=\mathsf Q_2\mathsf R_z(\psi)\mathsf Q_2^{\mathsf T}$。
于是
\[
\mathsf Q_3=\mathsf A_3\mathsf Q_2
=\mathsf Q_2\mathsf R_z(\psi)
=\mathsf R_z(\varphi)\mathsf R_x(\theta)\mathsf R_z(\psi).
\]
这套构造按时间先后是 $\varphi,\theta,\psi$，
但它使用不断变化的转轴。
若把同一个矩阵乘积解释为固定轴主动转动，
右端才最先作用。二者不能在省略“轴是否移动”的条件下混为一种次序规则。

\subsubsection*{欧拉矩阵的乘法中间步骤}
先算右边两个矩阵：
\[
\mathsf R_x(\theta)\mathsf R_z(\psi)=
\begin{pmatrix}
c_\psi&-s_\psi&0\\
c_\theta s_\psi&c_\theta c_\psi&-s_\theta\\
s_\theta s_\psi&s_\theta c_\psi&c_\theta
\end{pmatrix}.
\]
再左乘 $\mathsf R_z(\varphi)$。
例如第一行第一列是
$c_\varphi c_\psi-s_\varphi c_\theta s_\psi$，
第一行第三列是 $s_\varphi s_\theta$，
第三行保留为
$(s_\theta s_\psi,s_\theta c_\psi,c_\theta)$。
其余元素按同样的“行乘列”得到。
计算后可逐列检查单位长度、两列正交，
并检查 $\theta=\psi=0$ 时是否还原为 $\mathsf R_z(\varphi)$。

对于 $\theta=0$，
\[
\mathsf R=\mathsf R_z(\varphi+\psi).
\]
对于 $\theta=\pi$，
\[
\mathsf R=\mathsf R_z(\varphi-\psi)\mathsf R_x(\pi).
\]
这两个矩阵恒等式比单说“欧拉角奇异”更具体：
它们直接指出了姿态能保留的角度组合。


\example{由姿态矩阵提取非奇异欧拉角}\label{ex:c3-13}

\QuestionStatement{已知本体到空间的旋转矩阵元素为 $R_{ij}$，
并约定 $0<\theta<\pi$。写出欧拉角的反解。
采用主动转动、列矢量和 $R=R_z(\varphi)R_x(\theta)R_z(\psi)$ 约定；解释反解的象限选择和奇异姿态。}{ex13}{例13题图}
% BEGIN DETAILED ANSWER ex13
\solution \Step{选取不含多角混合的矩阵元}
以图~\ref{fig:solution-ex13}~中倾角 $\theta$ 与节点方向的矩阵投影为依据，先写出图右侧第一行的几何或动力学关系，再代入其下方的计算式。

由约定 $R=R_z(\varphi)R_x(\theta)R_z(\psi)$ 的乘积，第三列与第三行给出
\[
R_{13}=\sin\varphi\sin\theta,\quad R_{23}=-\cos\varphi\sin\theta,
\quad R_{33}=\cos\theta,
\]
\[
R_{31}=\sin\theta\sin\psi,\qquad R_{32}=\sin\theta\cos\psi.
\]
这些元素分离了三个角的依赖。
\SolutionFigure{ex13}{例13的计算示意}
\Step{利用倾角范围确定分支}
$0<\theta<\pi$ 时 $\sin\theta>0$，因此
\[
\theta=\arccos R_{33},\quad
\varphi=\operatorname{atan2}(R_{13},-R_{23}),\quad
\psi=\operatorname{atan2}(R_{31},R_{32}).
\]
$\operatorname{atan2}(y,x)$ 同时使用正弦、余弦符号，可区分四个象限。各方位角在模 $2\pi$ 意义下确定；连续跟踪时还需展开跨越 $\pm\pi$ 的跳变。
\Step{结果与奇异位置}
将三角代回全部九个元素，应恢复给定的 $R$；数值输入还应满足 $R^{\mathsf T}R=\mathsf1,\det R=1$。
在 $\theta=0$ 时，$R=R_z(\varphi+\psi)$，只能求和角；在 $\theta=\pi$ 时只能确定对应的差角组合。
此时反正切的两个输入同时为零，不可强行给出唯一的两个方位角。
例如取 $\varphi=\pi/3,\theta=\pi/4,\psi=-\pi/6$，上述三组矩阵元保留了 $\psi<0$ 的象限信息，反解可恢复原角。

% END DETAILED ANSWER ex13
\subsubsection*{矩阵反解例的具体验算}
设真实角为
$\varphi=\pi/3,\theta=\pi/4,\psi=-\pi/6$。
只需从正向矩阵中取五个相关元素：
\[
R_{33}=\frac{\sqrt2}{2},\quad
R_{13}=\frac{\sqrt6}{4},\quad
R_{23}=-\frac{\sqrt2}{4},\quad
R_{31}=-\frac{\sqrt2}{4},\quad
R_{32}=\frac{\sqrt6}{4}.
\]
于是 $\arccos R_{33}=\pi/4$。
第二个双参数反正切的两个输入均为正，给出第一象限的 $\pi/3$；
第三个的纵输入负、横输入正，给出第四象限的 $-\pi/6$。
如果只写 $\arctan(R_{13}/(-R_{23}))$，
其商虽然在本例没有错误，却不能区分输入同为负时相差 $\pi$ 的另一姿态。
反解后再把三个角代回完整矩阵，
比只检查某一个比值更可靠。

靠近 $\theta=0$ 时，$R_{13},R_{23},R_{31},R_{32}$ 都很小，
测量误差会使独立的 $\varphi,\psi$ 发生大变化；
它们的可观测和角却可能仍很稳定。
这属于参数反解的病态，不应解释为物体转动剧烈。


\subsection{任意轴转动、生成元与不交换性}\label{sub:c3-3.5}
% BEGIN ORIGINAL FIGURES
\OriginalFigure[6.2cm]{c3-1-p047-arbitrary-axis}{任意轴转动的几何步骤：图中 $\theta,\varphi$ 对应正文 $\beta,\lambda$}
% END ORIGINAL FIGURES

设单位轴 $\vect n=(\sin\beta\cos\lambda,\sin\beta\sin\lambda,\cos\beta)$。
先将它转到 $z$ 轴，绕 $z$ 转动，再把轴转回，得到
\[
\mathsf R_{\vect n}(\gamma)=
\mathsf R_z(\lambda)\mathsf R_y(\beta)\mathsf R_z(\gamma)
\mathsf R_y(-\beta)\mathsf R_z(-\lambda).
\]
同样，节点轴转动可以用
$\mathsf R_N(\theta)=\mathsf R_z(\varphi)\mathsf R_x(\theta)
\mathsf R_z(-\varphi)$ 表示。
这说明“绕移动轴转动”能够全部改写成固定轴的矩阵乘积，
但改写时次序不可颠倒。


\subsubsection*{任意轴矩阵共轭的推导}
令 $\mathsf Q=\mathsf R_z(\lambda)\mathsf R_y(\beta)$，
则 $\mathsf Q\vect e_z=\vect n$。
为了使空间中的 $\vect n$ 轴转动成为计算轴上的 $z$ 轴转动，
先把被转动矢量的分量变成 $\mathsf Q^{\mathsf T}\vect a$，
再施加 $\mathsf R_z(\gamma)$，最后映回：
\[
\vect a'=\mathsf Q\mathsf R_z(\gamma)\mathsf Q^{\mathsf T}\vect a.
\]
因此
$\mathsf R_{\vect n}=\mathsf Q\mathsf R_z(\gamma)\mathsf Q^{\mathsf T}$。
该矩阵确实固定 $\vect n$：
\[
\mathsf R_{\vect n}\vect n
=\mathsf Q\mathsf R_z(\gamma)\vect e_z
=\mathsf Q\vect e_z=\vect n.
\]
它的行列式为1，且仍正交，
所以共轭没有改变转角，只是改变轴的空间方向。
当 $\beta=0$，$\lambda$ 不再影响轴方向，
共轭自然退化为绕 $z$ 的原转动。


把无穷小转动写成 $\mathsf R_i(\epsilon)=\mathsf1+\epsilon\mathsf T_i+O(\epsilon^2)$，
生成元为
\[
\mathsf T_x=\begin{pmatrix}0&0&0\\0&0&-1\\0&1&0\end{pmatrix},\quad
\mathsf T_y=\begin{pmatrix}0&0&1\\0&0&0\\-1&0&0\end{pmatrix},\quad
\mathsf T_z=\begin{pmatrix}0&-1&0\\1&0&0\\0&0&0\end{pmatrix}.
\]
令 $\mathsf N=n_x\mathsf T_x+n_y\mathsf T_y+n_z\mathsf T_z$，
它的作用为 $\mathsf N\vect a=\vect n\times\vect a$。
由 $\mathsf N^2=\vect n\vect n^{\mathsf T}-\mathsf1$，
$\mathsf N^3=-\mathsf N$，指数级数可整理为
\begin{equation}\label{eq:rodrigues}
\mathsf R_{\vect n}(\gamma)=e^{\gamma\mathsf N}
=\mathsf1+\sin\gamma\,\mathsf N+(1-\cos\gamma)\mathsf N^2.
\end{equation}
矩阵指数定义是
$e^{\mathsf A}=\sum_{k=0}^{\infty}\mathsf A^k/k!$，零阶项必须包括单位矩阵。
也可用极限
$\lim_{N\to\infty}(\mathsf1+\gamma\mathsf N/N)^N$
理解一个有限转动由很多同轴小转动组成。由式\eqref{eq:rodrigues}得到
\[
\vect a'=\vect a\cos\gamma+(\vect n\times\vect a)\sin\gamma
+\vect n(\vect n\cdot\vect a)(1-\cos\gamma).
\]

\example{两次直角转动为何不能当作矢量相加}\label{ex:c3-14}

\QuestionStatement{比较先绕 $y$ 后绕 $x$ 各转 $\pi/2$ 与反向次序的结果。
均指作用于列矢量的主动转动，$x,y,z$ 是固定空间右手轴。写出两个合成矩阵，跟踪初始 $\vect e_z$，并比较总转角与小转动极限。}{ex14}{例14题图}
% BEGIN DETAILED ANSWER ex14
\solution \Step{固定主动转动与乘法顺序}
以图~\ref{fig:solution-ex14}~中同一初始标记方向在两种转序下的像为依据，先写出图右侧第一行的几何或动力学关系，再代入其下方的计算式。

作用于列矢量时先发生的转动写在右侧。基本矩阵为
\[
R_x=\begin{pmatrix}1&0&0\\0&0&-1\\0&1&0\end{pmatrix},\qquad
R_y=\begin{pmatrix}0&0&1\\0&1&0\\-1&0&0\end{pmatrix}.
\]
因此“先 $y$ 后 $x$”是 $R_xR_y$，“先 $x$ 后 $y$”是 $R_yR_x$。
\SolutionFigure{ex14}{例14的计算示意}
\Step{逐项相乘与跟踪标记方向}
\[
R_xR_y=\begin{pmatrix}0&0&1\\1&0&0\\0&1&0\end{pmatrix},\qquad
R_yR_x=\begin{pmatrix}0&1&0\\0&0&-1\\-1&0&0\end{pmatrix}.
\]
各自第三列就是初始 $\vect e_z$ 的像：分别为 $\vect e_x$ 与 $-\vect e_y$，故最终姿态不同。
第一矩阵满足 $R(1,1,1)^{\mathsf T}=(1,1,1)^{\mathsf T}$，转轴为 $(1,1,1)/\sqrt3$。
由 $\operatorname{tr}R=1+2\cos\gamma=0$，合成转角 $\gamma=2\pi/3$。
它不同于把两个直角矢量相加所得到的模 $\pi/\sqrt2$。
\Step{说明角速度为何仍可相加}
对小参数 $\epsilon$，
\[
R_x(\epsilon)R_y(\epsilon)=\mathsf1+\epsilon(T_x+T_y)+O(\epsilon^2).
\]
交换次序仅影响二阶项 $\epsilon^2(T_xT_y-T_yT_x)$。
因此一阶无穷小转动可相加，并在除以时间后取极限得到角速度的矢量合成；有限姿态的组合仍须按顺序乘矩阵。

% END DETAILED ANSWER ex14
\subsubsection*{有限转动例的等效轴与角}
第一种合成矩阵 $\mathsf Q_1$ 满足
\[
\mathsf Q_1(1,1,1)^{\mathsf T}=(1,1,1)^{\mathsf T},
\]
故其轴为 $\vect n_1=(1,1,1)/\sqrt3$。
第二种合成矩阵 $\mathsf Q_2$ 的不动轴为
$\vect n_2=(1,1,-1)/\sqrt3$。
二者的迹都为零，所以
\[
\cos\gamma=\frac{0-1}{2}=-\frac12,\qquad \gamma=\frac{2\pi}{3}.
\]
虽然等效转角相同，等效轴不同，最终姿态仍不同。
对初始三个单位轴逐一作用，
第一种给出
$\vect e_x\mapsto\vect e_y$、
$\vect e_y\mapsto\vect e_z$、
$\vect e_z\mapsto\vect e_x$；
第二种给出
$\vect e_x\mapsto-\vect e_z$、
$\vect e_y\mapsto\vect e_x$、
$\vect e_z\mapsto-\vect e_y$。
三个标记方向完整记录了两种操作的差别。

三个生成元的交换子满足
\[
[\mathsf T_x,\mathsf T_y]=\mathsf T_z,\quad
[\mathsf T_y,\mathsf T_z]=\mathsf T_x,\quad
[\mathsf T_z,\mathsf T_x]=\mathsf T_y.
\]
故
\[
\mathsf R_x(\epsilon)\mathsf R_y(\epsilon)
-\mathsf R_y(\epsilon)\mathsf R_x(\epsilon)
=\epsilon^2\mathsf T_z+O(\epsilon^3).
\]
二阶差异本身又对应一个小转动。
这与前面从位置矢量展开得到的非交换性是同一个结果，
只是改用了矩阵语言。

“无穷小转动可交换”是保留一阶项的结论，不能理解为精确的矩阵恒等式。

\TeachingFigure{c3-diagram-11}{同样两根轴与两个转角，交换次序产生不同姿态}
% BEGIN ORIGINAL FIGURES
\OriginalFigure[6.2cm]{c3-1-p029-rotation-xy}{先绕 $x$ 再绕 $y$ 各转 $\pi/2$ 后的姿态}
\OriginalFigure[6.2cm]{c3-1-p029-rotation-yx}{先绕 $y$ 再绕 $x$ 各转 $\pi/2$ 后的姿态}
% END ORIGINAL FIGURES


最后从正交矩阵直接建立瞬时角速度。
若 $\mathsf R(t+\dd t)\mathsf R^{\mathsf T}(t)=\mathsf1+\mathsf W\,\dd t$，
正交条件要求 $\mathsf W+\mathsf W^{\mathsf T}=0$。任意三阶反对称矩阵都可写为
\[
\mathsf W=\begin{pmatrix}0&-\omega_z&\omega_y\\
\omega_z&0&-\omega_x\\-\omega_y&\omega_x&0\end{pmatrix}.
\]
因此所有微小姿态变化都等效绕 $\vect\omega$ 方向的瞬时转动。
$\mathsf W=\dot{\mathsf R}\mathsf R^{\mathsf T}$ 给出空间分量，
$\mathsf R^{\mathsf T}\dot{\mathsf R}$ 给出本体分量。
这将欧拉几何定理、无穷小角位移与速度场统一起来。


\subsubsection*{生成元指数如何整理为Rodrigues公式}
由 $\mathsf N\vect a=\vect n\times\vect a$，
对单位 $\vect n$ 两次作用，
\[
\mathsf N^2\vect a
=\vect n(\vect n\cdot\vect a)-\vect a.
\]
因此 $\mathsf N^2=\vect n\vect n^{\mathsf T}-\mathsf1$，
再左乘 $\mathsf N$，由于 $\mathsf N\vect n=0$，
得到 $\mathsf N^3=-\mathsf N$。
于是偶次幂与奇次幂分别交替：
\[
\mathsf N^{2k}=(-1)^{k-1}\mathsf N^2\ (k\ge1),\qquad
\mathsf N^{2k+1}=(-1)^k\mathsf N.
\]
把指数级数按两类幂收集：
\begin{align*}
e^{\gamma\mathsf N}
&=\mathsf1+
\left(\gamma-\frac{\gamma^3}{3!}+\frac{\gamma^5}{5!}-\cdots\right)\mathsf N\\
&\quad+
\left(\frac{\gamma^2}{2!}-\frac{\gamma^4}{4!}
+\frac{\gamma^6}{6!}-\cdots\right)\mathsf N^2\\
&=\mathsf1+\sin\gamma\,\mathsf N+(1-\cos\gamma)\mathsf N^2.
\end{align*}
零阶单位矩阵负责保持沿轴分量，
不可从指数定义中省略。
在轴平行和轴垂直两部分上分别看该结果：
\[
\vect a_\parallel'=\vect a_\parallel,\qquad
\vect a_\perp'=\cos\gamma\,\vect a_\perp
+\sin\gamma\,(\vect n\times\vect a_\perp).
\]
这恰是垂直轴平面中的普通二维旋转。

\subsubsection*{从姿态矩阵提取瞬时角速度}
若在相邻时刻比较同一材料方向，
\[
\mathsf R(t+\Delta t)\mathsf R^{\mathsf T}(t)
=\mathsf1+\Delta t\,\dot{\mathsf R}\mathsf R^{\mathsf T}+O(\Delta t^2).
\]
它是“当前姿态到下一姿态”的空间小转动。
从空间反对称矩阵 $\mathsf W$ 中读出
\[
\omega_\xi=W_{32},\qquad
\omega_\eta=W_{13},\qquad
\omega_\zeta=W_{21}.
\]
本体矩阵为
\[
\mathsf W_B=\mathsf R^{\mathsf T}\dot{\mathsf R}
=\mathsf R^{\mathsf T}\mathsf W\mathsf R,
\qquad
\vect\omega_B=\mathsf R^{\mathsf T}\vect\omega_S.
\]
这里两种矩阵表示同一个角速度，只是基底不同。
若 $\vect\omega_S$ 在空间中恒定，
可写
$\mathsf R(t)=e^{t\mathsf W}\mathsf R(0)$。
若方向随时间变化，不同时刻的 $\mathsf W(t)$ 一般不交换，
不能未经说明就把时间积分放进一个普通指数来表示全部运动。
这一边界正是有限转动非交换性在时间演化中的体现。


\clearpage
\section{刚体动力学基础、力系简化与平衡}\label{sec:c3-4}
\subsection{整体运动方程与功率}\label{sub:c3-4.1}
% BEGIN ORIGINAL FIGURES
\OriginalFigure[6.2cm]{c3-2-p002-force-system}{分布在刚体各点的外力}
% END ORIGINAL FIGURES

刚体的运动应同时回答两个问题：质心怎样移动，物体绕质心怎样改变姿态。
设总质量为 $m$，质心为 $C$，各外力为 $\vect F_i$，直接施加的外力偶为 $\vect C_j$。
先对全部质元应用动量定理。内力成对抵消，因此
\[
\frac{\dd}{\dd t}\sum_i m_i\vect v_i=\sum_i\vect F_i,
\qquad \sum_i m_i\vect v_i=m\vect v_C.
\]
外力主矢 $\vect F=\sum_i\vect F_i$ 决定质心加速度，得到
\begin{equation}\label{eq:newtoneuler}
m\vect a_C=\vect F,\qquad
\frac{\dd\vect L_C}{\dd t}=\vect M_C.
\end{equation}
第二式中的角动量是相对质心位置与相对质心速度定义的角动量，
$\vect M_C=\sum_i\vect r_{Ci}\times\vect F_i+\sum_j\vect C_j$。
每个矢量式有三个分量，共构成自由刚体的六条动力学方程。
但它们还需要速度场、姿态参数与角速度之间的运动学关系才能闭合。

\paragraph{先区分取矩点，再求导。}
若 $O$ 是惯性系中的固定点，用绝对速度定义
\[
\vect L_O=\sum_i\vect r_{Oi}\times m_i\vect v_i.
\]
求导时 $\dot{\vect r}_{Oi}=\vect v_i$，第一项
$\sum_i\vect v_i\times m_i\vect v_i=0$，故
\[
\frac{\dd\vect L_O}{\dd t}=\sum_i\vect r_{Oi}\times m_i\vect a_i
=\vect M_O.
\]
若取矩点 $A$ 在移动，仍用绝对速度定义 $\vect L_A$，则
$\dot{\vect r}_{Ai}=\vect v_i-\vect v_A$，于是
\[
\frac{\dd\vect L_A}{\dd t}
=\vect M_A-\vect v_A\times m\vect v_C.
\]
这不是直接套用固定点公式。若改用相对速度，定义
\[
\vect K_A=\sum_i\vect r_{Ai}\times m_i(\vect v_i-\vect v_A),
\]
则 $\vect K_A=\vect L_A-m\vect r_{AC}\times\vect v_A$，再求导可得
\[
\frac{\dd\vect K_A}{\dd t}
=\vect M_A-m\vect r_{AC}\times\vect a_A.
\]
取 $A=C$ 时最后一项消失，所以质心角动量方程成立，即使质心本身在加速。
取固定点时 $\vect v_A=\vect a_A=0$，两种角动量定义也一致。
这两种特例解释了为何质心与固定点最适合列转动方程。

在直角坐标中，质心方程为
$m\ddot x_C=F_x,m\ddot y_C=F_y,m\ddot z_C=F_z$。
在平面极坐标中则必须保留坐标基底的变化：
\[
m(\ddot r_C-r_C\dot\theta_C^2)=F_r,\qquad
m(r_C\ddot\theta_C+2\dot r_C\dot\theta_C)=F_\theta.
\]
径向加速度不是单独的 $\ddot r_C$。
角动量方程若用本体基底表示，也必须使用
\[
\left(\frac{\dd\vect L_C}{\dd t}\right)_{\!S}
=\left(\frac{\dd\vect L_C}{\dd t}\right)_{\!B}
+\vect\omega\times\vect L_C.
\]
本体分量的普通导数与惯性矢量导数有这一个明确差别。

\paragraph{由刚体速度场逐项求功率。}
取体内基点 $O$，各力作用点速度为
$\vect v_i=\vect v_O+\vect\omega\times\vect r_{Oi}$。
代入外力功率，利用标量三重积的循环性质，
\begin{align*}
\mathcal P
&=\sum_i\vect F_i\cdot\vect v_i+\sum_j\vect C_j\cdot\vect\omega\\
&=\left(\sum_i\vect F_i\right)\cdot\vect v_O
+\sum_i\vect F_i\cdot(\vect\omega\times\vect r_{Oi})
+\sum_j\vect C_j\cdot\vect\omega\\
&=\vect F\cdot\vect v_O+
\left(\sum_i\vect r_{Oi}\times\vect F_i+\sum_j\vect C_j\right)\cdot\vect\omega\\
&=\vect F\cdot\vect v_O+\vect M_O\cdot\vect\omega.
\end{align*}
第一项是主矢通过基点运动做功的速率，第二项是主矩通过转动做功的速率。
改变体内基点时这两项分别会变，但其和不变。例如移到体内点 $A$，
$\vect v_A=\vect v_O+\vect\omega\times\vect r_{OA}$，
$\vect M_A=\vect M_O-\vect r_{OA}\times\vect F$，两项变化正好抵消。
不能只移动主矩而把基点速度保持不变。

\paragraph{内力为什么不增加刚体总动能？}
对一对理想中心内力 $\vect f_{ij}=-\vect f_{ji}$，配对功率为
\[
\vect f_{ij}\cdot\vect v_i+\vect f_{ji}\cdot\vect v_j
=\vect f_{ij}\cdot(\vect v_i-\vect v_j).
\]
刚性约束保持 $|\vect r_i-\vect r_j|$ 不变，求导得
\[
(\vect r_i-\vect r_j)\cdot(\vect v_i-\vect v_j)=0.
\]
内力沿两点连线，因此配对功率为零。这里消失的是相对速度在连线方向上的分量，
相对速度本身通常不为零；转动刚体上的两点完全可以有不同速度。
对更一般的理想刚性约束，其约束内力的总虚功和实际功率均为零。
由全部质元动能定理相加，得到
\[
\frac{\dd T}{\dd t}=\mathcal P,
\qquad T(t)-T(t_0)=\int_{t_0}^{t}\mathcal P\,\dd t.
\]
若做功外力中的保守部分由势能 $V(\vect q,t)$ 导出，沿实际运动轨迹有
\[
\mathcal P_{\text{保守}}
=-\frac{\dd V}{\dd t}+\frac{\partial V}{\partial t},
\qquad
\frac{\dd (T+V)}{\dd t}
=\frac{\partial V}{\partial t}+\mathcal P_{\text{非保守}}.
\]
这里 $\mathcal P_{\text{非保守}}$ 包括其余做功外力的总功率。
当势能无显含时间、非保守作用的总功率为零时，才有 $T+V$ 守恒；
固定理想支承的反力虽可非零，其功率仍为零。
力为零功率并不等于力为零，例如固定光滑支承可以提供非零反力而不做功。

质心表达 $T=mv_C^2/2+T_C$ 把平动和相对转动分开；
若有惯性固定点，全部速度均可写成绕该点的转动，全部动能也可写成定点转动形式。
两种表达是同一总动能的不同分解，不能重复相加。

\subsection{力的可传递性与主矢、主矩}\label{sub:c3-4.2}
力对理想刚体有作用线上的可传递性。设原力 $\vect F$ 作用于 $A$，
改在 $B$ 施加相同矢量。关于任意点 $O$ 的矩差为
\[
\vect M_O^{(B)}-\vect M_O^{(A)}
=(\vect r_{OB}-\vect r_{OA})\times\vect F
=\vect r_{AB}\times\vect F.
\]
如果 $A,B$ 在同一作用线上，$\vect r_{AB}\parallel\vect F$，矩差为零，
主矢也没有改变，所以整体平动与转动效果相同。
力因此称为滑动矢量。沿作用线滑动与任意平行搬移有不同的补偿要求。
可传递性只判断刚体总体外部效果，局部应力和变形仍依赖实际施力处。
% BEGIN ORIGINAL FIGURES
\OriginalFigure[6.2cm]{c3-2-p003-sliding-force}{同一作用线上的力可在刚体上滑移}
% END ORIGINAL FIGURES


\paragraph{把偏心力搬至指定点。}
把作用于 $A$ 的力搬到 $O$ 时，在 $O$ 人为加入
$\vect F$ 与 $-\vect F$ 这一零力对。将 $O$ 的 $\vect F$ 留作新力，
将 $A$ 的原力与 $O$ 的 $-\vect F$ 合成力偶，得到
\[
\{\text{$A$处的力}\}
\ \sim\ \{\text{$O$处的同一力}\}
+\{\text{力偶 }\vect r_{OA}\times\vect F\}.
\]
辅助零力对作用在同一刚体上，不是分别作用于两物体的第三定律作用反作用对。
矩的正负由 $\vect r_{OA}\times\vect F$ 决定，不能凭弯箭头的视觉印象猜测。
例如 $A=(d,0,0)$、$\vect F=F\vect j$，附加矩为 $dF\vect k$；
若同一位置的力向下，附加矩随之反号。

对每个力重复此操作，将全部力集中在 $O$，所有自由力偶相加，得到
\begin{equation}
\vect F=\sum_i\vect F_i,\qquad
\vect M_O=\sum_i\vect r_{Oi}\times\vect F_i+\sum_j\vect C_j.
\end{equation}
$\vect F$ 是主矢，不依赖取矩中心；$\vect M_O$ 是关于中心 $O$ 的主矩。
二者一共具有六个标量分量，任何有限空间力系都可由它们描述。
已有自由力偶必须加入主矩；仅对作用力求叉积会漏掉直接施加的扭矩。
% BEGIN ORIGINAL FIGURES
\OriginalFigure[6.2cm]{c3-2-p009-spatial-reduction}{空间力系移到同一中心时附加相应力偶}
% END ORIGINAL FIGURES


将中心从 $O$ 改为 $A$，由 $\vect r_{Ai}=\vect r_{Oi}-\vect r_{OA}$，
\begin{align*}
\vect M_A
&=\sum_i(\vect r_{Oi}-\vect r_{OA})\times\vect F_i+\sum_j\vect C_j\\
&=\vect M_O-\vect r_{OA}\times\sum_i\vect F_i,
\end{align*}
故
\begin{equation}\label{eq:momenttranslation}
\vect M_A=\vect M_O-\vect r_{OA}\times\vect F.
\end{equation}
这也是判定两个力系是否等效的实用方法：检查主矢相同，并检查关于同一个点的主矩相同。
若这两个条件满足，关于所有其他点的主矩也相同。
若 $\vect F=0$，主矩对取矩点完全独立；若 $\vect F\ne0$，沿主矢方向移动中心不改变主矩。
单独保持主矢而不保持主矩，无法保证刚体转动效果相同。

\TeachingFigure{c3-diagram-12}{偏离作用线平移一个力时必须补上力偶}
% BEGIN ORIGINAL FIGURES
\OriginalFigure[6.2cm]{c3-2-p008-transfer-pair}{在新作用点加入零力对，实现力与力偶的等效转换}
% END ORIGINAL FIGURES


\subsection{共面力系、平行力与力偶}\label{sub:c3-4.3}
若两个共面非平行力的作用线交于 $P$，先利用可传递性把两力移到 $P$，
再作平行四边形。合力 $\vect F_{12}=\vect F_1+\vect F_2$ 也作用于 $P$，因此
\[
\vect r_{OP}\times\vect F_{12}
=\vect r_{OP}\times\vect F_1+\vect r_{OP}\times\vect F_2
=\vect M_{O1}+\vect M_{O2}.
\]
这一步同时保持主矢与主矩。第三个力可与 $\vect F_{12}$ 递次合成，
但仅当作用线确有可用交点时，这个作图操作才可继续。
中间主矢为零、剩余主矩不为零时已经形成力偶；
中间力与下一力平行时则要改用平行力的矩条件。
因此对任意力系，应先用主矢和主矩记录信息，再决定能否用交点作图。
% BEGIN ORIGINAL FIGURES
\OriginalFigure[6.2cm]{c3-2-p005-three-resultant}{三个共面非平行力的递次合成}
% END ORIGINAL FIGURES


\TeachingFigure{c3-diagram-13}{先沿原作用线移至交点，再在交点合成两个力}
% BEGIN ORIGINAL FIGURES
\OriginalFigure[6.2cm]{c3-2-p004-two-forces}{两个非平行力及其作用线}
\OriginalFigure[6.2cm]{c3-2-p004-two-resultant}{将两力移至作用线交点后合成}
% END ORIGINAL FIGURES


\paragraph{从矩方程求合力整条作用线。}
在平面内取 $z$ 垂直纸面并朝向观察者，主矢与主矩为
\[
F_x=\sum_i F_{ix},\quad F_y=\sum_i F_{iy},\quad
M_O=\sum_i(x_iF_{iy}-y_iF_{ix})+\sum_jC_j.
\]
假设 $\vect F\ne0$，若一个等效单力作用于 $A=(x,y)$，则必须有
\[
\vect r_{OA}\times\vect F=M_O\vect k,
\qquad xF_y-yF_x=M_O.
\]
这是一条直线，而不是唯一作用点。其法向可取 $(F_y,-F_x)$，
与合力 $(F_x,F_y)$ 垂直，所以直线本身平行合力。
距 $O$ 最近的一点是
\[
(x_0,y_0)=\frac{M_O}{F_x^2+F_y^2}(F_y,-F_x),
\]
一般作用点为 $(x,y)=(x_0,y_0)+\lambda(F_x,F_y)$。
这同时给出最短力臂 $|M_O|/|\vect F|$，
正负矩则确定合力作用线在中心的哪一侧。

三种平面简化结果应分开讨论：
\begin{enumerate}
\item $\vect F\ne0$：存在上述有限作用线，全部力系可等效为一个力。
\item $\vect F=0$、$M_O\ne0$：只剩力偶。任何单力若非零就有非零主矢，
因此没有单力可代替；取矩中心改变也不能消掉它。
\item $\vect F=0$、$M_O=0$：力系对刚体整体等效于零。
这并不说每个外力都消失，约束反力仍可能逐个非零。
\end{enumerate}
共面力系的单力简化总能成立于第一种分支，是因为其矩垂直于所在平面，
而主矢在平面内，两者内积必为零。空间力系没有这个保证。

\example{同向与反向平行力的合成}\label{ex:c3-15}
% BEGIN ORIGINAL FIGURES
\OriginalFigure[6.2cm]{c3-2-p006-parallel-opposite}{反向不等平行力的合力位于较大力一侧}
% END ORIGINAL FIGURES
\QuestionStatement{沿 $y$ 方向的两力 $F_1,F_2$ 作用于横坐标 $x_1,x_2$。
$F_i$ 按向上为正取代数值，求合力作用线。
讨论合力为零和两力反向的特殊情况，说明什么条件下可化为单一合力。}{ex15}{例15题图}
% BEGIN DETAILED ANSWER ex15
\solution \Step{明确模型与计算依据}
以图~\ref{fig:solution-ex15}~中两条平行作用线的横向坐标 $x_1,x_2$为依据，先写出图右侧第一行的几何或动力学关系，再代入其下方的计算式。
 先以 $x_1$ 为取矩点，两力总矩为 $(x_2-x_1)F_2$。
若合力为 $F=F_1+F_2$，作用线位于 $x_R$，其矩必须满足
\[
(x_R-x_1)(F_1+F_2)=(x_2-x_1)F_2.
\]
当 $F_1+F_2\ne0$ 时，解得
\[
F=F_1+F_2,\qquad
x_R=x_1+\frac{F_2}{F_1+F_2}(x_2-x_1)
=\frac{x_1F_1+x_2F_2}{F_1+F_2}.
\]

\SolutionFigure{ex15}{例15的计算示意}
\Step{逐项计算与求解}
令 $x_2>x_1$。两力同向时，$F_2/(F_1+F_2)$ 在0与1之间，
所以 $x_R$ 位于两力之间，并有
$F_1(x_R-x_1)=F_2(x_2-x_R)$。
取 $F_1=2F_0,F_2=3F_0$，合力就在从第一作用线向第二作用线走 $3/5$ 间距的位置。
合力靠近较大的力，而不是靠近较小的力。

反向时仍使用带符号代数值。若 $F_1=3F_0,F_2=-F_0$，
合力为向上的 $2F_0$，并有 $x_R=x_1-(x_2-x_1)/2$，
作用线位于第一力外侧。若反向较大的力在 $x_2$，则作用线在 $x_2$ 外侧。
无需另记一套比例，只要保持符号与矩方程即可。

若 $F_2=-F_1$，前面的除法失效，因为主矢为零。
此时关于原点的主矩为
\[
M_O=x_1F_1+x_2F_2=F_1(x_1-x_2).
\]
两作用线分离就形成非零力偶；两线重合时主矩也为零，力系完全相消。
当反向两力趋于等大而保持间距不变时，普通合力越来越小，
$|x_R|$ 越来越大，但乘积“合力乘力臂”保持有限。这说明等大反向极限必须用力偶描述。

% END DETAILED ANSWER ex15
\TeachingFigure{c3-diagram-14}{同向平行力的合力位置由总力矩确定}
% BEGIN ORIGINAL FIGURES
\OriginalFigure[6.2cm]{c3-2-p006-parallel-same}{同向平行力的合力位于两作用线之间}
% END ORIGINAL FIGURES

力偶由等大、反向、不同线的两力组成。设 $A$ 处作用 $\vect F$，
$B$ 处作用 $-\vect F$。主矢显然为零，关于任意点 $O$ 的矩为
\[
\vect M_O=\vect r_{OA}\times\vect F+\vect r_{OB}\times(-\vect F)
=(\vect r_{OA}-\vect r_{OB})\times\vect F
=\vect r_{BA}\times\vect F.
\]
结果只含两作用点差，与 $O$ 无关，因此力偶矩是自由矢量。
设两点连线与力方向夹角为 $\gamma$，矩大小为
$|\vect r_{BA}|F\sin\gamma=Fd$，$d$ 是两作用线的垂直距离。
若两线重合，$d=0$，虽有两反向力，却不构成非零力偶。

平面中规定 $\vect k$ 朝纸外，逆时针矩为正。
若上方力向右、下方力向左，合力为零但顺时针矩为负；
若箭头全部反向，力偶矩反号。判断应由叉积或旋转趋势统一完成。
多个力偶的矩可直接矢量相加；空间力偶不必位于同一个平面，
其结果仍是全部自由矩的矢量和。

\example{扳手上的力偶}\label{ex:c3-16}
% BEGIN ORIGINAL FIGURES
\OriginalFigure[6.2cm]{c3-2-p007-couple}{等大反向的非共线二力构成力偶；图中为顺时针力偶，例~\ref{ex:c3-16} 取两力反向后的逆时针情形}
% END ORIGINAL FIGURES
\QuestionStatement{在相距 $d$ 的两条平行作用线上施加大小相等的反向力 $F$，
方向使物体逆时针转动。再把整个力偶平移到其他位置，求主矢、主矩。
求主矢、关于任意点的主矩，以及移动取矩点后的变化。}{ex16}{例16题图}
% BEGIN DETAILED ANSWER ex16
\solution \Step{计算主矢与矩的方向}
以图~\ref{fig:solution-ex16}~中两条反向力的作用线间垂距 $d$为依据，先写出图右侧第一行的几何或动力学关系，再代入其下方的计算式。

可取左力作用于 $(0,0)$、方向向下，右力作用于 $(d,0)$、方向向上。两力分别为 $-F\vect j,F\vect j$，于是
\[
\vect F_{\text{主}}=-F\vect j+F\vect j=\vect0,\qquad
\vect M_O=(d\vect i)\times(F\vect j)=Fd\vect k.
\]
从 $+z$ 看为逆时针，符合题设；$d$ 必须是两条作用线的垂直距离。
\SolutionFigure{ex16}{例16的计算示意}
\Step{取矩点平移下的不变量}
对任意 $A=(x_A,y_A)$，两力矩的 $z$ 分量为
\[
(-x_A)(-F)+(d-x_A)F=Fd.
\]
也可用式\eqref{eq:momenttranslation}：
$\vect M_A=\vect M_O-\vect r_{OA}\times\vect F_{\text{主}}=Fd\vect k$。
把整个力偶平移任意矢量 $\vect b$，新增矩为 $\vect b\times(-F\vect j+F\vect j)=0$，结论相同。
\Step{解释等效条件与限制}
只要 $Fd$ 和矩方向不变，改变两力大小、间距或整体位置都保持刚体整体的力学效果。
若以一力替代，非零单力必有非零主矢，而本题主矢为零，故不能等效。
$F\to0$ 或 $d\to0$ 时力偶矩趋零；单独给出两个反向力而不说明作用线间距，不能确定其转动效果。
力偶可作为自由矩处理，但真实施力位置仍会影响物体内部应力。

% END DETAILED ANSWER ex16
\subsection{空间力系的中心轴}\label{sub:c3-4.4}
空间中一个力的矩永远垂直于它自身的方向。因此主矢与主矩都非零时，
不能直接断言存在一个单力同时代表二者。先将主矩分解为沿主矢与垂直主矢的部分：
\[
\vect M_O=\vect M_\parallel+\vect M_\perp,
\quad \vect M_\parallel=
\frac{\vect F\cdot\vect M_O}{F^2}\vect F,
\quad \vect F\cdot\vect M_\perp=0.
\]
移动简化中心只能减去 $\vect r_{OA}\times\vect F$，它也垂直主矢，
所以只能改变 $\vect M_\perp$，不能改变 $\vect M_\parallel$。
直接验证
\[
\vect F\cdot\vect M_A
=\vect F\cdot\vect M_O-
\vect F\cdot(\vect r_{OA}\times\vect F)
=\vect F\cdot\vect M_O.
\]
内积 $\vect F\cdot\vect M_O$ 是取矩点变换的不变量。

\paragraph{构造使主矩最小的简化中心。}
希望 $\vect r_{OA}\times\vect F=\vect M_\perp$，可选
\[
\vect r_{OA}=\frac{\vect F\times\vect M_O}{F^2}.
\]
利用 $(\vect a\times\vect b)\times\vect c
=\vect b(\vect a\cdot\vect c)-\vect a(\vect b\cdot\vect c)$，
\[
\vect r_{OA}\times\vect F
=\vect M_O-\frac{\vect F\cdot\vect M_O}{F^2}\vect F.
\]
代回主矩变换，得到
\[
\vect M_A=\frac{\vect F\cdot\vect M_O}{F^2}\vect F.
\]
沿 $\vect F$ 再移动任意距离也不改变这个矩，故所有这样的中心组成中心轴
\[
\vect r=\frac{\vect F\times\vect M_O}{F^2}+\lambda\vect F.
\]
写成单位轴方向 $\vect n=\vect F/|\vect F|$ 时，也可用
$\vect r=\vect r_0+s\vect n$，其中 $s$ 是实际长度。
$\vect r_0\perp\vect F$，所以它是中心轴上距原简化中心最近的一点。

若中心离该轴的垂距为 $d$，额外横向主矩的大小为 $|\vect F|d$，
故主矩大小平方为
\[
|\vect M|^2=|\vect M_\parallel|^2+F^2d^2.
\]
中心轴使主矩达到最小，最小值仍可能非零。
定义力系的螺距 $h_s=(\vect F\cdot\vect M_O)/F^2$，
它有长度量纲，中心轴上的主矩为 $h_s\vect F$。
轴向推力与轴向扭矩并存的螺旋机构，就是这一组合的直观例子。

\paragraph{所有分支的判定顺序。}
先算 $\vect F$，再算关于同一点的 $\vect M_O$：
\begin{enumerate}
\item $\vect F=0,\vect M_O=0$：整体等效零力系。
\item $\vect F=0,\vect M_O\ne0$：纯力偶，没有中心轴公式中的除数 $F^2$，
不可硬套中心轴；主矩对中心独立。
\item $\vect F\ne0,\vect F\cdot\vect M_O=0$：中心轴上主矩为零，
可化成沿中心轴的一个力。
\item $\vect F\ne0,\vect F\cdot\vect M_O\ne0$：中心轴上仍留平行力偶，
只能化成轴向力与轴向力偶的组合。
\end{enumerate}
若原中心的 $\vect M_O=0$，第三种分支的中心轴就通过原点。
若原中心主矩恰与主矢平行，第四种分支的中心轴也通过原点，
但力偶仍须保留。位置相同并不意味着两个分支的转动效果相同。

\example{不能合成单力的空间力系}\label{ex:c3-17}

\QuestionStatement{在原点作用 $\vect F=(0,0,F)$，同时施加力偶
$\vect M_O=(0,0,C)$，$F,C\ne0$。判断能否以一力替代。
从力矩与合力的内积证明结论，并给出力系简化后的最简表示。}{ex17}{例17题图}
% BEGIN DETAILED ANSWER ex17
\solution \Step{写出力系的不变量}
以图~\ref{fig:solution-ex17}~中同轴的主矢 $F\vect k$ 和主矩 $C\vect k$为依据，先写出图右侧第一行的几何或动力学关系，再代入其下方的计算式。

本题主矢 $\vect F=F\vect k$，主矩 $\vect M_O=C\vect k$，且
\[
\vect F\cdot\vect M_O=FC\ne0.
\]
将取矩点移至任意 $A$，主矩改为 $\vect M_A=\vect M_O-\vect r_{OA}\times\vect F$；与主矢点乘时，叉积项恒为零。因此上述内积与取矩点无关。
\SolutionFigure{ex17}{例17的计算示意}
\Step{证明不存在等效单力}
若一力 $\vect F$ 作用于某位置 $\vect r=(x,y,z)$，它关于 $O$ 的矩为
\[
\vect r\times\vect F=(Fy,-Fx,0).
\]
其轴向分量恒为零，无法等于 $(0,0,C)$。
等价地，任何单力都满足 $\vect F\cdot(\vect r\times\vect F)=0$，与本题的非零不变量矛盾。
\Step{确定最简形式与中心轴}
关于 $A=(x,y,z)$ 的主矩为 $(-Fy,Fx,C)$，其模平方为
$F^2(x^2+y^2)+C^2$。
最小值在 $x=y=0$，因此中心轴是整个 $z$ 轴。最简形式仍为沿该轴的力 $F\vect k$ 加同轴自由矩 $C\vect k$，可称为力螺旋，螺距 $C/F$ 的量纲为长度。
只有 $C=0$ 才能退化为本题原点处的单力；给定条件明确排除了这一情形。

% END DETAILED ANSWER ex17
\subsection{平衡条件与等价的平面表述}\label{sub:c3-4.5}
静力平衡要求刚体的位形保持不变。对于初始静止且约束持续成立的理想刚体，
平衡条件是主矢和主矩同时为零：
\begin{equation}
\vect F=0,\qquad\vect M_O=0.
\end{equation}
由于主矢为零，若关于一个点的主矩为零，则关于所有点都为零。
空间中可写成六个标量条件
\[
\sum F_x=\sum F_y=\sum F_z=0,
\qquad \sum M_x=\sum M_y=\sum M_z=0.
\]
平面内各力在 $xy$ 平面，矩只有 $z$ 分量，因此只需
\[
\sum F_x=0,\qquad\sum F_y=0,\qquad\sum M_O=0.
\]
若力系平面汇交，取交点矩自然为零，只剩两条力平衡式；
若各力均沿 $y$ 平行，只剩 $\sum F_y=0$ 和一条矩平衡式。
这些减少的方程来自力系的几何性质，不能任意删除仍独立的方程。

\paragraph{零主矢、零主矩与稳定性。}
若刚体一开始已有速度，零主矢可以对应匀速平移；零外矩可以对应自由转动。
此时恒定的是惯性系动量和角动量，姿态及角速度方向未必恒定。
所以静力题还隐含初始静止。稳定性更是另一个问题：
满足平衡方程的姿态受到小扰动后，可能回到邻近位置，也可能离开它。
悬挂物质心在悬点上方与下方都能有零重力矩，但前者不稳定、后者稳定。

\paragraph{两点矩加一投影为什么等价？}
设平面内 $A,B$ 不重合，已知 $M_A=M_B=0$。
由取矩中心变换得到
\[
0=M_B-M_A=-\vect r_{AB}\times\vect F.
\]
这说明可能剩下的合力只能沿 $AB$。再取单位投影方向 $\vect e$，
若 $\vect e\cdot\vect r_{AB}\ne0$，写 $\vect F=f\vect r_{AB}$ 后，
$\vect e\cdot\vect F=0$ 就迫使 $f=0$，于是主矢为零。
结合 $M_A=0$，得到全部平衡条件。

用坐标更容易看到非退化要求。取 $A=(0,0),B=(d,0)$，$d\ne0$。
$M_B-M_A=-dF_y=0$ 先给出 $F_y=0$；
第三条若选 $F_x=0$，平衡成立；若仍选 $F_y=0$，只是重复已有结论。
沿 $AB$ 的非零合力此时不会对 $A,B$ 产生力矩，也没有竖直投影，因而会漏掉。

\paragraph{三点矩法为什么要求不共线？}
若 $M_A=M_B=M_C=0$，两次相减得
\[
\vect r_{AB}\times\vect F=0,
\qquad \vect r_{AC}\times\vect F=0.
\]
若 $AB,AC$ 不平行，同一个非零矢量不可能同时平行两者，故 $\vect F=0$。
再保留一个矩方程，便得到完全平衡。
若三点共线，非零合力仍可沿这条公共直线，三点矩都会是零。
三点矩法在本节是平面力系的替代表述；不可把平面标量矩条件直接搬到空间而少列分量。

选方程的目的在于消去不需要的未知量。对铰支点取矩可以消去铰反力，
对两个支承点分别取矩可以直接求另一支反力。
但代数解还必须通过接触检验：负的法向力意味着所设接触不能维持，
负的绳张力意味着绳应松弛，所需摩擦超过上限意味着静止解不可实现。

\TeachingFigure{c3-diagram-15}{三点取矩法的几何非退化条件}
% BEGIN ORIGINAL FIGURES
\OriginalFigure[6.2cm]{c3-2-p026-three-moments}{三个不共线取矩点的位置关系}
% END ORIGINAL FIGURES


\subsection{三力定理与静力学建模}\label{sub:c3-4.6}
% BEGIN ORIGINAL FIGURES
\OriginalFigure[6.2cm]{c3-2-p018-skew-forces}{异面二力产生的主矩不能由单一反向力抵消}
% END ORIGINAL FIGURES

若刚体只受三个非零、非平行力而平衡，三力的实际作用线必共面且共点。
这一定理有两步：方向矢量共面不能直接代替作用线共面；
只有先排除空间异面，再取交点矩，才得到严格的作用线结论。

\paragraph{第一种证明：逐分量排除异面作用线。}
力平衡给出 $\vect F_3=-\vect F_1-\vect F_2$。
以第一力作用线上的一点为原点，取 $x$ 沿第一力，
再取 $xy$ 平面容纳第二力的方向。于是
\[
\vect F_1=(f_1,0,0),\quad
\vect F_2=(f_{2x},f_{2y},0),\quad
\vect F_3=(-f_1-f_{2x},-f_{2y},0),
\]
其中 $f_1f_{2y}\ne0$。
令第二、三力的作用点为 $(x_2,y_2,z_2)$ 与 $(x_3,y_3,z_3)$。
第一力的矩为零，剩余总矩的 $x$ 分量为
\[
M_x=-z_2f_{2y}+z_3f_{2y}=f_{2y}(z_3-z_2)=0,
\]
所以 $z_2=z_3$。
再计算 $y$ 分量，
\[
M_y=z_2f_{2x}+z_3(-f_1-f_{2x})=-f_1z_3=0,
\]
所以 $z_2=z_3=0$。三力的方向均在 $xy$ 平面，作用点也在其中，
因此实际作用线共面。前两力不平行，设作用线交于 $P$，对 $P$ 取矩，
$\vect M_{P1}=\vect M_{P2}=0$；平衡要求 $\vect M_{P3}=0$，
第三力非零，所以其作用线也必须经过 $P$。

\paragraph{第二种证明：利用主矢方向上的不变矩。}
取第一力线上一点为 $O$，第二力线上一点的位置为 $\vect r_2$。
先只合成前两力，主矢和主矩为
\[
\vect F_{12}=\vect F_1+\vect F_2,
\qquad\vect M_{12}=\vect r_2\times\vect F_2.
\]
它们的内积为
\[
\vect F_{12}\cdot\vect M_{12}
=\vect r_2\cdot(\vect F_2\times\vect F_1).
\]
非平行两线异面时，$\vect r_2$ 有沿两方向公共法线的非零分量，
故这个标量三重积非零。平衡要求第三力为 $-\vect F_{12}$，
但一个沿此方向的力，无论作用在哪里，其矩都垂直 $\vect F_{12}$，
永远消不掉前两力主矩沿 $\vect F_{12}$ 的非零分量。
这与总矩为零矛盾，因而前两作用线不能异面。
它们共面且相交后，仍通过对交点取矩得到第三作用线共点。
这一证明把三力定理与空间力系中心轴的不可消去轴向力偶联系起来。

三力共点只确定作用线的几何关系，不直接给出力的大小。
大小可再由力三角形、两个分量平衡式，或正弦定理求出。
若三力彼此平行，应使用平行力系条件；若有一个力为零，就退化为两力平衡，
两个非零力必须等大、反向、共线；若另有独立力偶，就不符合“只受三力”的前提。

\paragraph{隔离物体时怎样画支承作用？}
光滑面只提供垂直接触面的法向力；
光滑棱边接触细杆时，理想化反力沿垂直杆的方向，不能误画成竖墙面的水平法向。
粗糙面具有法向 $N$ 与切向 $f$，静止要求 $N\ge0,|f|\le\mu N$。
柔绳只有沿绳方向的非负张力；理想平面铰提供两个力分量，不提供阻转力偶；
固定端既可提供两个力分量，也可提供一个平面力偶。
对滚动物体，即使力矩对圆心由法向抵消，切向摩擦仍会提供力矩，必须保留。

两刚体相接触时，每个物体自己的隔离图上都要画接触反力；
只有把两者组成一个整体时，这对内作用才互相消失。
分析木堆的整个重量不能替代每根木材的平衡，分析整辆车也不能自动知道钢绳张力。
正确的计算顺序是：定几何与接触方向，选隔离体，列独立平衡式，解反力，
最后检查单向接触、摩擦上限与稳定分支。静摩擦的大小由平衡需求决定，
只有研究即将滑动时才把某一接触的 $|f|$ 设为 $\mu N$。

\subsection{平衡算例、几何方法与模型边界}\label{sub:c3-4.7}
\example{粗糙地面上靠光滑棱边的均匀杆}\label{ex:c3-18}

\QuestionStatement{均匀杆长 $2l$、重 $P$，下端 $A$ 放在粗糙水平面，
中途支在距地面高度 $h$ 的光滑棱边 $B$，杆与水平面成角 $\varphi$。
地面静摩擦系数为 $\mu$。求平衡反力以及临界滑动角满足的方程。
棱边是局部接触而不是整块竖直光滑墙，接触反力垂直于杆。
设 $0<\varphi<\pi/2$，所有接触仅能受压。除临界等式外，还须写出接触与静摩擦可行条件。}{ex18}{例18题图}
% BEGIN DETAILED ANSWER ex18
\solution \Step{明确模型与计算依据}
以图~\ref{fig:solution-ex18}~中杆上接触点 $B$ 的法线和 $A$ 处的摩擦分解为依据，先写出图右侧第一行的几何或动力学关系，再代入其下方的计算式。
 先确定接触方向，再列矩方程。取 $A$ 为原点，地面向右为 $x$、向上为 $y$，
杆的单位方向为 $\vect u=(\cos\varphi,\sin\varphi)$。
中点 $D$ 与棱边接点 $B$ 的位置分别为
\[
\vect r_{AD}=l\vect u,
\qquad \vect r_{AB}=\frac h{\sin\varphi}\vect u
=(h\cot\varphi,h).
\]
光滑棱边反力垂直杆且推向左上，故为
$\vect N_1=N_1(-\sin\varphi,\cos\varphi)$。
地面反力拆成 $(f,N_2)$，先把向右的 $f$ 作为正方向；
重力为 $(0,-P)$。三条平衡式逐一写成
\begin{align*}
\sum F_x=0&:\quad f-N_1\sin\varphi=0,\\
\sum F_y=0&:\quad N_2+N_1\cos\varphi-P=0,\\
\sum M_A=0&:\quad
(h\cot\varphi)N_1\cos\varphi+hN_1\sin\varphi-Pl\cos\varphi=0.
\end{align*}
最后一式中的前两项合为
$N_1h(\cos^2\varphi+\sin^2\varphi)/\sin\varphi$。
也可以直接看出力矩大小为 $N_1AB$，因为 $N_1$ 垂直杆。
因此
\[
N_1\frac h{\sin\varphi}=Pl\cos\varphi,
\quad N_1=\frac{Pl}{h}\sin\varphi\cos\varphi.
\]
依次回代水平、竖直平衡，得
\[
f=\frac{Pl}{h}\sin^2\varphi\cos\varphi,
\qquad N_2=P-\frac{Pl}{h}\sin\varphi\cos^2\varphi.
\]
在 $0<\varphi<\pi/2$ 内，$N_1>0,f>0$，所以假定的摩擦方向正确。
不能将地面反力预先设成重力 $P$：棱边也承担了一部分竖直载荷。

\SolutionFigure{ex18}{例18的计算示意}
\Step{反力解与可行静止解}
必须同时检验
\[
\frac h{\sin\varphi}\le2l,
\qquad h-l\sin\varphi\cos^2\varphi>0,
\qquad
\frac{l\sin^2\varphi\cos\varphi}
{h-l\sin\varphi\cos^2\varphi}\le\mu.
\]
第一式保证 $B$ 仍在杆上；第二式保证地面仍受压；第三式是摩擦需求不超过上限。
只有将要滑动时第三式取等号，临界角为
\begin{equation}
\mu=\frac{l\sin^2\varphi_0\cos\varphi_0}
{h-l\sin\varphi_0\cos^2\varphi_0}.
\end{equation}
若把 $h/l=\eta$、$u=\sin\varphi_0$ 作为无量纲量，待解条件为
\[
\mu=\frac{u^2\sqrt{1-u^2}}{\eta-u+u^3},
\qquad 0<u<1,\quad u\ge\eta/2,\quad\eta-u+u^3>0.
\]
数值求根时可直接使用这一式。若为消去根号而平方，会得到
$\mu^2(\eta-u+u^3)^2=u^4(1-u^2)$，
但还须重新检查原式及上述符号限制，不能把所有多项式根都视为临界姿态。
有多个临界根时，应沿实际改变角度的连续可行分支寻找首次失稳位置。

\Step{第二解：求三力的交点坐标}
将地面两分量合为 $\vect R_A=(f,N_2)$，
其与竖直的夹角 $\beta$ 满足 $\tan\beta=f/N_2$。
重力作用线是 $x=l\cos\varphi$；棱边法线可写成
\[
(x,y)=(h\cot\varphi,h)+t(-\sin\varphi,\cos\varphi).
\]
令其横坐标等于重力线横坐标，求得
\[
t=\frac{h\cot\varphi-l\cos\varphi}{\sin\varphi},
\quad
Q=\left(l\cos\varphi,
\frac{h-l\sin\varphi\cos^2\varphi}{\sin^2\varphi}\right).
\]
三力共点要求 $\vect R_A$ 沿 $AQ$。
取其水平距离与竖直高度之比，得到
\begin{align*}
\tan\beta
&=\frac{l\cos\varphi}{Q_y}
=\frac{l\sin^2\varphi\cos\varphi}{h-l\sin\varphi\cos^2\varphi}\\
&=\frac{l\cos\varphi}
{l\sin\varphi+(h/\sin\varphi-l)/\sin\varphi}.
\end{align*}
在临界角 $\varphi_0$，$\tan\beta=\mu$，因此与第一解完全一致。
$Q_y>0$ 与 $N_2>0$ 是同一条件：若交点落到地面下方，
要求的合地面反力就会包含向下的法向分量，单向接触不能实现。
三力法把摩擦锥条件变成了合反力方向的几何条件。

% END DETAILED ANSWER ex18
\TeachingFigure{c3-diagram-16}{光滑棱边的反力垂直于杆，地面摩擦须由平衡求出}
% BEGIN ORIGINAL FIGURES
\OriginalFigure[6.2cm]{c3-2-p016-rough-rod}{粗糙地面与光滑棱边支承的均匀杆：$P$ 为重量，$\varphi_0$ 为临界倾角}
\OriginalFigure[6.2cm]{c3-2-p019-concurrent-rod}{利用三力共点确定地面反力方向}
% END ORIGINAL FIGURES


\example{两根直角刚接杆的悬挂平衡}\label{ex:c3-19}

\QuestionStatement{两均匀细杆长 $a,b$、线密度相同，刚接成直角。
$a$ 杆的上端 $A$ 用理想铰悬挂。
令 $a$ 杆从 $A$ 向下并向右偏离竖直角 $\theta$，
$b$ 杆从接点向左下伸出。求平衡角。
同时给出铰反力，判断所得平衡姿态的稳定性。}{ex19}{例19题图}
% BEGIN DETAILED ANSWER ex19
\solution \Step{明确模型与计算依据}
以图~\ref{fig:solution-ex19}~中两杆中点重力的水平力臂为依据，先写出图右侧第一行的几何或动力学关系，再代入其下方的计算式。
 取 $A$ 为原点，向右为 $x$、向上为 $y$。
$a$ 杆接点为 $B=(a\sin\theta,-a\cos\theta)$，
$b$ 杆从 $B$ 沿 $(-\cos\theta,-\sin\theta)$ 伸出。
两个杆方向点积为零，确实构成直角。
各杆质心是对应线段中点，位置为
\[
\vect r_1=\left(\frac a2\sin\theta,-\frac a2\cos\theta\right),
\quad
\vect r_2=\left(a\sin\theta-\frac b2\cos\theta,
-a\cos\theta-\frac b2\sin\theta\right).
\]
理想铰反力通过 $A$，对 $A$ 不产生矩。两杆重力的矩分别为
$-m_1gx_1$ 与 $-m_2gx_2$，因此
\[
m_1\frac a2\sin\theta+
 m_2\left(a\sin\theta-\frac b2\cos\theta\right)=0.
\]
将正弦项合并，
\[
a(m_1+2m_2)\sin\theta=m_2b\cos\theta,
\qquad\tan\theta=\frac{m_2b}{a(m_1+2m_2)}.
\]
同线密度时 $m_1=\lambda a,m_2=\lambda b$，约去 $\lambda$ 得
\[
\tan\theta=\frac{b^2}{a^2+2ab}.
\]
不应把“均匀”自动理解为两杆线密度相同：各杆内部均匀仍允许两杆材料或截面不同，
所以题干中的相同线密度是获得这个简式的必要条件。

从总质心看，
$x_C=(m_1x_1+m_2x_2)/(m_1+m_2)$，矩平衡正是 $x_C=0$。
于是总重力作用线经过悬点，铰反力的水平分量为零、竖直分量为 $(m_1+m_2)g$。
用质量加权质心求角度，与把两重力分别取矩是同一个计算。

\SolutionFigure{ex19}{例19的计算示意}
\Step{哪一个零力矩姿态稳定？}
令 $A_0=a(m_1/2+m_2)$、$B_0=m_2b/2$，重力势能除去常数后为
\[
V(\theta)=-(A_0\cos\theta+B_0\sin\theta)g.
\]
求导得到 $V'=g(A_0\sin\theta-B_0\cos\theta)$，
驻点条件与上面的矩方程相同。取 $0<\theta<\pi/2$ 的分支时，
\[
V''=g(A_0\cos\theta+B_0\sin\theta)>0,
\]
质心在 $A$ 下方，势能为极小值，是稳定平衡。
反向姿态 $\theta+\pi$ 的 $V''$ 反号，质心在 $A$ 上方，是不稳定平衡。
当 $b\to0$，$\theta\to0$，单杆竖直悬挂；
当 $a=b$，$\tan\theta=1/3$，这给出一个可直接核验的几何特例。

% END DETAILED ANSWER ex19
\TeachingFigure{c3-diagram-17}{悬挂刚体平衡时质心重力作用线通过悬点}
% BEGIN ORIGINAL FIGURES
\OriginalFigure[6.2cm]{c3-2-p017-right-angle-rods}{直角刚接双杆的悬挂平衡}
% END ORIGINAL FIGURES


\example{同时接触竖墙和棱边的光滑杆}\label{ex:c3-20}

\QuestionStatement{均匀杆长 $2l$、重 $P$，一端 $O$ 接触光滑竖墙，
杆中部支于光滑棱边 $B$；棱边距竖墙水平距离为 $d$。
杆从 $O$ 向右上倾斜，与水平面夹角 $\theta$，
且 $B$ 在质心之前：$d\le l\cos\theta$。
求平衡角及反力。
求平衡倾角及墙面、棱边的反力，并筛选接触力非负的几何分支。}{ex20}{例20题图}
% BEGIN DETAILED ANSWER ex20
\solution \Step{明确模型与计算依据}
以图~\ref{fig:solution-ex20}~中墙面水平法线和棱边处的杆法线为依据，先写出图右侧第一行的几何或动力学关系，再代入其下方的计算式。
 取 $O=(0,0)$，杆方向为 $\vect u=(\cos\theta,\sin\theta)$。
均匀杆中点为 $D=(l\cos\theta,l\sin\theta)$，
棱边接点为 $B=(d,d\tan\theta)$，沿杆距离 $OB=d/\cos\theta$。
墙在 $O$ 左侧，其光滑反力只能水平向右，记大小 $R$；
棱边法向垂直杆向左上，记 $\vect N=N(-\sin\theta,\cos\theta)$。

水平、竖直平衡分别为
\[
R-N\sin\theta=0,\qquad N\cos\theta-P=0,
\]
因此
\[
N=\frac P{\cos\theta},\qquad R=P\tan\theta.
\]
取 $O$ 矩消掉 $R$，棱边力与杆垂直，力矩为 $NOB$；
重力矩为 $-Pl\cos\theta$，所以
\[
N\frac d{\cos\theta}=Pl\cos\theta.
\]
代入 $N$ 后逐步消元，
\[
\frac{Pd}{\cos^2\theta}=Pl\cos\theta,
\quad d=l\cos^3\theta,
\quad\theta=\arccos\left(\frac dl\right)^{1/3}.
\]
长度 $l$ 是半杆长，不能把它误用成全长 $2l$；重力臂来自中点位置。

\SolutionFigure{ex20}{例20的计算示意}
\Step{第二解：两个正交投影怎样产生三个余弦？}
墙反力作用线是过 $O$ 的水平线；重力作用线是过 $D$ 的竖直线。
两线交于 $Q=(l\cos\theta,0)$。三力共点要求 $BQ$ 为棱边法线，
即 $BQ\perp OB$。直角三角形 $OBQ$ 中，
\[
OQ=l\cos\theta,
\qquad OB=OQ\cos\theta=l\cos^2\theta.
\]
再把 $OB$ 投影到水平方向，
$d=OB\cos\theta=l\cos^3\theta$。
第一个余弦来自质心位置的水平投影，第二个来自直角三角形，
第三个来自支点距离的水平投影。三力共点说的是作用线相交，
不要求三个力箭头都朝向交点；本例的 $N$ 箭头沿 $BQ$ 的反向也符合定理。

在非退化双接触中 $0<d<l$，所以 $0<\theta<\pi/2$，两反力均为正。
又因 $d=l\cos^3\theta\le l\cos\theta$，支点确在质心之前。
当 $d=l$ 时 $\theta=0$，墙反力 $R=0$，成为退化极限；
当 $d\to0$ 时所需 $N,R$ 发散，理想式的极限不意味着有限支承可承受无限反力。
本题光滑墙的水平法向与光滑棱边的斜法向必须同时保留。

% END DETAILED ANSWER ex20
\TeachingFigure{c3-diagram-18}{两个光滑接触的法向方向不同}
% BEGIN ORIGINAL FIGURES
\OriginalFigure[6.2cm]{c3-2-p020-smooth-corner}{光滑墙面与棱边对杆的两种法向约束}
% END ORIGINAL FIGURES


\example{光滑半球碗中的均匀杆}\label{ex:c3-21}

\QuestionStatement{半径 $r$ 的光滑半球碗固定，碗口水平。
均匀细杆一端 $B$ 位于碗内，穿过碗缘 $C$ 伸向碗外，
碗内杆段长 $BC=c$，杆总长为 $l>c$。
忽略杆厚度，碗内端反力沿球面法线，碗缘反力垂直杆。
求这个平衡分支中 $l$ 与 $c$ 的关系。
同时求两处反力，列出满足杆有碗外部分且接触受压的长度范围。}{ex21}{例21题图}
% BEGIN DETAILED ANSWER ex21
\solution \Step{明确模型与计算依据}
以图~\ref{fig:solution-ex21}~中碗心 $O$、碗内端 $B$ 和碗缘 $C$ 的弦关系为依据，先写出图右侧第一行的几何或动力学关系，再代入其下方的计算式。
 在经过杆的竖直截面内，球心为 $O=(0,0)$，右碗缘为 $C=(r,0)$。
设杆从碗内端 $B$ 向右上伸出，单位方向为
$\vect u=(\cos\theta,\sin\theta)$，故 $B=C-c\vect u$。
由于 $B$ 在圆上，
\[
|B|^2=(r-c\cos\theta)^2+c^2\sin^2\theta=r^2.
\]
展开并约去 $r^2$，得到 $c^2-2rc\cos\theta=0$；$c>0$，所以
\[
c=2r\cos\theta.
\]
这就是弦长关系，不需要把 $c$ 误作半径。
代回 $B$，可写成 $B=(-r\cos2\theta,-r\sin2\theta)$。

碗内端支持沿半径 $BO$，朝向球心；碗缘反力垂直杆，向左上。
设两条法线交于 $Q$。碗缘法线写成
\[
Q=(r,0)+t(-\sin\theta,\cos\theta).
\]
另一法线是过原点且平行 $B$ 的直线。联立可得
\[
t=2r\sin\theta=c\tan\theta,
\qquad Q=(r\cos2\theta,r\sin2\theta).
\]
因此 $Q=-B$：这一平衡几何里两法线的交点正好是 $B$ 的径向对称点。
重力作用线必须经过 $Q$。杆的中点 $D=B+(l/2)\vect u$，
重力竖直，故 $x_D=x_Q$，得到
\[
-r\cos2\theta+\frac l2\cos\theta=r\cos2\theta,
\qquad \frac l2=\frac{2r\cos2\theta}{\cos\theta}
=c\frac{\cos2\theta}{\cos^2\theta}.
\]
利用 $\cos2\theta=2\cos^2\theta-1$ 以及 $c=2r\cos\theta$，
\[
l=2c(1-\tan^2\theta)
=2c\left(2-\frac1{\cos^2\theta}\right)
=4c-\frac{8r^2}{c}.
\]
几何上，中点 $D$ 在碗内段上，$CD=c-l/2=c\tan^2\theta$。
这提供另一条课堂作图路线：先求 $CQ=c\tan\theta$，
再由 $CQ$ 的水平投影与 $CD$ 的水平投影相等，得到 $CD=CQ\tan\theta$。

\SolutionFigure{ex21}{例21的计算示意}
\Step{力的大小与几何分支}
设两反力大小为 $N_B,N_C$。内端法向为
$(\cos2\theta,\sin2\theta)$，碗缘法向为 $(-\sin\theta,\cos\theta)$，故
\[
N_B\cos2\theta-N_C\sin\theta=0,
\quad N_B\sin2\theta+N_C\cos\theta=mg.
\]
消去 $N_C$，得
\[
N_B=mg\tan\theta,
\qquad N_C=mg\frac{\cos2\theta}{\cos\theta}.
\]
题设 $l>c$ 要求
$4c-8r^2/c>c$，即 $c>\sqrt{8/3}\,r$。
弦长给出 $c<2r$，所以非退化分支为
\[
\sqrt{8/3}\,r<c<2r,
\qquad 0<\theta<\arccos\sqrt{2/3}.
\]
在这个角度范围中，两反力均为正，几何与受力相容。
$c=2r$ 时杆沿水平直径，$N_B=0$，三力问题退化；
$c=\sqrt{8/3}\,r$ 时 $l=c$，恰无碗外杆段，不符合题设严格条件。
只检查 $l>0$ 会保留一些没有题设构型的解，因此长度、接触与题干要求须一起筛选。

% END DETAILED ANSWER ex21
\TeachingFigure{c3-diagram-19}{碗杆的共点几何：径向线BQ、碗缘法线CQ与中点竖直线DQ}
% BEGIN ORIGINAL FIGURES
\OriginalFigure[6.2cm]{c3-2-p024-bowl-rod}{半球碗内杆的三力共点：图中交点 $A$ 对应正文 $Q$，$D$ 为杆中点}
% END ORIGINAL FIGURES

\TeachingFigure{c3-diagram-20}{半球碗内的杆：端点反力沿半径，碗缘反力垂直杆}

\example{圆柱木材堆积与运输的建模研究}\label{ex:c3-22}
\TeachingFigure{c3-diagram-21}{木材堆积研究题的几何模型}
% BEGIN ORIGINAL FIGURES
\OriginalFigure[6.2cm]{c3-2-p025-stacked-logs}{受两侧夹杆约束的圆柱木材堆积}
% END ORIGINAL FIGURES
\LongQuestionStatement{相同圆柱木材长 $l$、半径 $a$、密度 $\rho$，
每根质量 $m=\rho\pi a^2l$，截面按等边三角形紧密堆积。
先考虑完整五层堆，底层五根、上层依次四、三、二、一根。研究：
\begin{enumerate}
\item 木材表面光滑，两侧各压一根挡杆，两杆上端连接，
在下端施加水平向内的力 $F$，求连接处作用力与 $F$。
\item 缓慢撤去挡杆后，若木材不能滚落，最多能堆几层？
\item 木材接触静摩擦系数为 $\mu$，在底层左右最外侧木材施加水平向内 $F$，
最多能保持几层？
\item 仅保留原五层堆中的最下三层，即 $5+4+3$ 根，
装在货车上并用允许最大拉力 $T_{\max}$ 的钢绳固定。
以速度 $v$ 转弯时允许的最小半径是多少？
以速度 $V$ 紧急换道时允许的最大方向盘角速度是多少？
\end{enumerate}}{ex22}{例22题图}
% BEGIN DETAILED ANSWER ex22
\solution \Step{明确模型与计算依据}
以图~\ref{fig:solution-ex22}~中相邻圆柱的法向接触与底层侧推位置为依据，先写出图右侧第一行的几何或动力学关系，再代入其下方的计算式。
 四问共有同一个截面几何，但外部约束分别改变。
先建立每根圆柱的接触模型，再分析挡杆、撤杆、摩擦与运输。
仅给质量、半径和木材间摩擦系数，不能唯一确定全部问项；
应明确指出哪个量缺少哪个条件，并在明示假设下求可确定的结果。

\SolutionFigure{ex22}{例22的计算示意}
\Step{共同几何与接触方程}
以地面为 $y=0$。完整 $n$ 层堆的圆心可编号为
\[
\vect r_{kj}=\bigl((2j+k-n+1)a,\ a+\sqrt3ka\bigr),
\quad k=0,\ldots,n-1,\quad j=0,\ldots,n-1-k.
\]
同层邻圆心水平相距 $2a$，上下相邻圆心差为 $(\pm a,\sqrt3a)$，
长度也为 $2a$。它们分别给出接触法向的水平角 $0^\circ$、$60^\circ$、$120^\circ$。
总根数是 $n(n+1)/2$，五层时为15。
接触方向由这些圆心坐标决定，不能将所有斜接触力假定竖直。

对圆柱 $i$ 与 $j$ 的接触，取指向 $i$ 的法向
$\vect n_{ij}=(\vect r_i-\vect r_j)/(2a)$，
切向 $\vect t_{ij}=\vect k\times\vect n_{ij}$。
作用在 $i$ 上的接触力写成
\[
\vect f_{ij}=N_{ij}\vect n_{ij}+\tau_{ij}\vect t_{ij},
\qquad N_{ij}\ge0,\quad |\tau_{ij}|\le\mu N_{ij}.
\]
对 $j$ 的作用力取相反矢量。接触点相对 $i$ 圆心的力臂为 $-a\vect n_{ij}$，
所以法向力通过圆心而没有矩，切向力的矩为
$(-a\vect n_{ij})\times(\tau_{ij}\vect t_{ij})=-a\tau_{ij}\vect k$。
地面接触的力臂为 $(0,-a)$，若地面摩擦向右为 $f_{ig}$，其圆心矩为 $af_{ig}\vect k$。
每根圆柱必须独立列两条力式与一条矩式：
\[
\sum_{j\in\mathcal C_i}\vect f_{ij}
+\vect f_{ig}+\vect f_{i,\text{侧}}-mg\vect j=0,
\qquad
\sum_{j\in\mathcal C_i}(-a\vect n_{ij})\times\vect f_{ij}
+(0,-a)\times\vect f_{ig}+M_{i,\text{侧}}\vect k=0.
\]
若侧推力通过圆心，则 $M_{i,\text{侧}}=0$；若施力线有高度偏差，必须加其矩。
这些矩条件解释了为什么仅列水平力平衡不能判断圆柱是否滚落。

% END DETAILED ANSWER ex22
\TeachingFigure{c3-diagram-22}{接触切向力与地面摩擦必须同时通过圆柱的力矩平衡检验}

\paragraph{第一问：把木堆与两根挡杆分别隔离。}
左右挡杆沿侧边圆柱的公共外切线，倾角为 $60^\circ$。
对左挡杆而言，木材的法向反作用沿左上方向，
可写为 $N_k(-\sqrt3/2,1/2)$；作用点坐标由圆心向外切线投影得到。
设杆脚为 $A$、顶连接点为 $P$，脚部反力为 $(S_x,S_y)$，
顶端连接力为 $(H_x,H_y)$，杆重为 $W_b$，重心为 $G_b$。
把题设向内水平力 $F$ 与脚部未知反力分开，杆平衡为
\begin{align*}
F+S_x+H_x-\frac{\sqrt3}{2}\sum_kN_k&=0,\\
S_y+H_y+\frac12\sum_kN_k-W_b&=0,\\
\vect r_{AP}\times(H_x,H_y)
+\sum_k\vect r_{Ak}\times N_k(-\sqrt3/2,1/2)
-\vect r_{AG_b}\times(0,W_b)&=0.
\end{align*}
若脚部的 $F$ 与支承力作用在同一点，对 $A$ 取矩正好消去它们。
右杆列同样的方程；理想铰接要求两杆顶端作用力等大反向，
对称受力时其竖直分量为零。若顶端用拉索，方向应由索线决定；
若是刚接，还要加入连接力偶。这些结构不是同一个边界条件。

一个可立即做的模型检查是：假定两挡杆无重、顶端对称铰接、脚仅放在光滑地面，
外施 $F$ 水平。左杆竖直方程便要求
$S_y=-\tfrac12\sum_kN_k<0$，与地面只能向上推矛盾。
所以这一组额外理想化并不能维持静态木堆；必须给出杆重、脚部向下锚固、
或其他提供竖直约束的机构，再计算连接力与 $F$。
这比只给一个形式上的 $F$ 数值更能检查结构是否实际可用。

在任一完整可行模型中，木材总重确定为 $15mg$，左右侧作用的水平分量对称抵消。
但侧杆给木堆的力还有向下分量，会改变地面总支持；
挡杆连接力还取决于每个 $N_k$ 的作用臂。故 $15mg$ 不能单独确定 $F$ 或顶端连接力。
求解时将全部圆柱和两杆方程联合起来，检查全部 $N_k\ge0$、脚部接触方向与连接条件。
若出现多组反力解，说明仅刚体静力模型存在超静定分配，
还需要接触刚度或其他变形关系来确定每个接触的数值。

\paragraph{第二问：先看最小的两层结构。}
在圆柱间及地面均光滑、没有端部限位的模型下，
取一个上圆柱压在两个下圆柱上。对称接触法向大小均为 $N$，
$2N\sin60^\circ=mg$，因此
\[
N=\frac{mg}{\sqrt3},\qquad N\cos60^\circ=\frac{mg}{2\sqrt3}.
\]
上圆柱对每个下圆柱都有这个向外水平分量。
最外下圆柱从内邻圆柱得到的水平法向也只可能向外，
而地面支持竖直，因此水平式不可能为零。两层已经不成立，故最多一层。
这一最外圆柱论证也适用于更多层，不必逐层试算所有法向力。

若只将地面改为粗糙，光滑的木材接触力仍通过各圆心，重力也通过圆心。
为了抵消水平外推，地面摩擦须非零；它在圆心下方产生非零矩，
其他现有力却不能抵消该矩。所以单凭地面摩擦大，并不能使光滑自由圆柱静止。
若增加制动块、端部限位或木材间切向摩擦，才构成另一种可行模型。

\paragraph{第三问：摩擦锥可行性，而不是一个总摩擦力。}
对候选层数 $n$ 建立全部接触，收集未知数
$\vect q=(N_e,\tau_e,N_{ig},f_{ig},F,\ldots)^{\mathsf T}$。
每根圆柱三条平衡式可整理成
\[
\mathsf A\vect q=\vect b,
\quad N_e\ge0,\quad -\mu N_e\le\tau_e\le\mu N_e,
\quad -\mu_gN_{ig}\le f_{ig}\le\mu_gN_{ig}.
\]
$\mathsf A$ 由法向、切向和作用臂组成，$\vect b$ 由每根的 $mg$ 与规定外载组成。
$\mu_g$ 是木材对地面的摩擦系数，应与木材彼此的 $\mu$ 区分。
还要写明 $F$ 是已知输入、允许任意调整，还是满足 $0\le F\le F_{\max}$；
限制侧力会直接改变可保持的层数。

一个实用算法是从 $n=1$ 开始增加层数，逐个判定是否存在满足所有等式与不等式的反力组。
若要找所需最小侧力，可在该可行域上最小化 $F$。
若有接触的 $N_e=0$，它可以脱离，应重新检查剩余接触网络；
若有 $|\tau_e|=\mu N_e$，它在摩擦边界上，但并不保证所有其他接触也临界。
不能一次把所有接触都设成 $\mu N_e$，否则往往人为制造无解或错误上界。
本问缺少地面条件、侧力上限与作用线，因此没有唯一的最大整数层数。
完整解答应给出明确模型和对应可行域，而不是把未知参数隐藏到一个猜测数值中。

\paragraph{第四问：先算真实装载质量与质心。}
保留底三层是5、4、3根，共12根，并非重新排成3、2、1根。
各层高度为 $a,a+\sqrt3a,a+2\sqrt3a$，故
\begin{align*}
M&=12m,\\
H&=\frac{5ma+4m(a+\sqrt3a)+3m(a+2\sqrt3a)}{12m}
=a+\frac{5\sqrt3}{6}a.
\end{align*}
每根均匀圆柱的质心就是圆心，这个求 $H$ 的步骤不需要把圆柱自身惯量忽略。
若木材靠固定件形成车厢上的等效刚体载荷，
设有效半支承宽度为 $b$，横向总抗滑能力为 $Q_{\max}$，
固定件能提供附加抗倾覆矩 $C_{\max}$。
在车厢随动系中，恒定转弯时载荷有水平惯性力 $Ma_y$，作用于高度 $H$。
抗滑与抗倾覆分别要求
\[
Ma_y\le Q_{\max},\qquad Ma_yH\le Mgb+C_{\max}.
\]
第二式也可由内侧地面支持不变为负推导；等号时内侧支持消失，开始绕外侧支承边翻倒。
因此
\[
a_{\max}=\min\left(\frac{Q_{\max}}M,
\frac{Mgb+C_{\max}}{MH}\right),
\qquad R\ge R_{\min}=\frac{v^2}{a_{\max}}.
\]
若绳索 $\alpha$ 的单位拉力方向为 $\vect u_\alpha$、作用点为 $P_\alpha$，
其水平贡献是 $T_\alpha u_{\alpha x}$，关于倾覆边的贡献是
$\vect r_{\text{边},P_\alpha}\times T_\alpha\vect u_\alpha$。
须在 $0\le T_\alpha\le T_{\max,\alpha}$ 和地面支持非负条件下同时求这些量。
同一绳索的侧向力与抗倾覆矩有关联，分别取最大后再拼起来可能不可实现。
钢绳路径、角度、锚点与预紧力未给时，不能令 $Q_{\max}=T_{\max}$ 来得到唯一半径。
木材若还会彼此滑动或滚动，更需把每根的横向惯性力 $m a_y$ 加入前述接触可行性方程。

\paragraph{转角限制与转角变化率限制。}
采用轴距 $L_w$ 的低速无侧滑自行车模型，前轮角 $\delta$ 的曲率为
$\kappa=\tan\delta/L_w$。恒定车速 $V$ 时
\[
a_y=V^2\kappa=\frac{V^2}{L_w}\tan\delta,
\qquad\delta_{\text{方向盘}}=i_s\delta.
\]
$a_y\le a_{\max}$ 只给出
$|\delta|\le\arctan(L_wa_{\max}/V^2)$，不会给出角度变化率上限。
若另外规定横向加加速度 $|\dot a_y|\le j_{\max}$，求导得
\[
\dot a_y=\frac{V^2}{L_w}\sec^2\delta\,\dot\delta,
\qquad
|\dot\delta_{\text{方向盘}}|
\le i_s\frac{L_wj_{\max}}{V^2}\cos^2\delta.
\]
若车速也在变化，则另有 $2V\dot V\tan\delta/L_w$，不能继续使用恒速导数式。
实际紧急换道还涉及轮胎侧偏、车身滚转、悬架响应与运动持续时间。
这些逐项说明使四问变为可明确输入、可核查输出的建模问题；
在输入未闭合前，保留其条件不足的结论，同时给出以上可确定的几何与动力学链条。

\addtocontents{toc}{\protect\newpage}
\clearpage
\section{惯量张量、主轴与转动动能}\label{sec:c3-5}
\subsection{从质元求角动量和动能}\label{sub:c3-5.1}
先考虑绕惯性固定点 $O$ 的转动。质元相对位置为 $\vect r=(x,y,z)$，
角速度为 $\vect\omega=(\omega_x,\omega_y,\omega_z)$。速度场逐分量为
\[
\vect v=\vect\omega\times\vect r,\quad
v_x=\omega_yz-\omega_zy,\quad v_y=\omega_zx-\omega_xz,\quad v_z=\omega_xy-\omega_yx.
\]
全部质元共享同一角速度，线速度却随位置改变。
% BEGIN ORIGINAL FIGURES
\OriginalFigure[6.2cm]{c3-3-p005-mass-element}{定点转动中质元的位置矢量、动量与角速度}
% END ORIGINAL FIGURES


\paragraph{从角动量三个分量识别矩阵元。}
质元角动量为 $\dd\vect L_O=\vect r\times\vect v\dd m$。第一分量逐步展开：
\begin{align*}
\dd L_x&=(yv_z-zv_y)\dd m\\
&=[y(\omega_xy-\omega_yx)-z(\omega_zx-\omega_xz)]\dd m\\
&=[(y^2+z^2)\omega_x-xy\omega_y-xz\omega_z]\dd m.
\end{align*}
其余两分量为
\begin{align*}
\dd L_y&=[-xy\omega_x+(z^2+x^2)\omega_y-yz\omega_z]\dd m,\\
\dd L_z&=[-xz\omega_x-yz\omega_y+(x^2+y^2)\omega_z]\dd m.
\end{align*}
角速度分量与质元位置无关，可提到积分外。依次收集系数得到
\[
\mathsf I_O=\begin{pmatrix}
\int(y^2+z^2)\dd m&-\int xy\dd m&-\int xz\dd m\\
-\int xy\dd m&\int(z^2+x^2)\dd m&-\int yz\dd m\\
-\int xz\dd m&-\int yz\dd m&\int(x^2+y^2)\dd m
\end{pmatrix}.
\]
$y^2+z^2$ 是到 $x$ 轴的垂距平方，故对角元是坐标轴惯量。
若定义 $D_{xy}=\int xy\dd m$，实际非对角元为 $I_{xy}=-D_{xy}$。
本章 $I_{ij}$ 均指实际矩阵元，后续交叉项不另换符号定义。
$xy=yx$，所以矩阵实对称，九个位置只有六个独立量。

也可用三重积统一三个分量：
\[
\vect r\times(\vect\omega\times\vect r)
=r^2\vect\omega-\vect r(\vect r\cdot\vect\omega),
\quad L_i=\sum_j\left[\int(r^2\delta_{ij}-r_ir_j)\dd m\right]\omega_j.
\]
于是
\begin{equation}\label{eq:tensor}
\vect L_O=\mathsf I_O\vect\omega,\qquad
(I_O)_{ij}=\int(r^2\delta_{ij}-r_ir_j)\dd m.
\end{equation}
$i=j$ 时从 $r^2$ 减掉该坐标平方；$i\ne j$ 时只保留负的惯性积。
这个逐项检查解释了通式中不能遗漏 $\delta_{ij}$。

\paragraph{独立从速度平方求动能。}
$|\vect\omega\times\vect r|^2=\omega^2r^2-(\vect\omega\cdot\vect r)^2$，展开并积分得
\begin{align*}
2T_O={}&I_{xx}\omega_x^2+I_{yy}\omega_y^2+I_{zz}\omega_z^2\\
&+2I_{xy}\omega_x\omega_y+2I_{xz}\omega_x\omega_z+2I_{yz}\omega_y\omega_z\\
={}&I_{xx}\omega_x^2+I_{yy}\omega_y^2+I_{zz}\omega_z^2
-2D_{xy}\omega_x\omega_y-2D_{xz}\omega_x\omega_z-2D_{yz}\omega_y\omega_z.
\end{align*}
交叉项的2来自对称位置 $ij,ji$ 各贡献一次。因此
\[
T_O=\frac12\int|\vect\omega\times\vect r|^2\dd m
=\frac12\vect\omega^{\mathsf T}\mathsf I_O\vect\omega
=\frac12\vect\omega\cdot\vect L_O.
\]
最后一个等式也可由
$\vect\omega\cdot(\vect r\times\vect v)
=(\vect\omega\times\vect r)\cdot\vect v=v^2$ 直接得到。
动能虽是标量，非对角元仍参与其数值。

\paragraph{任意轴惯量与角动量方向。}
取单位方向 $\vect n=(n_x,n_y,n_z)$，令 $\vect\omega=\omega\vect n$。
质元到轴垂距为 $|\vect n\times\vect r|$，所以
\begin{equation}\label{eq:directioninertia}
I_{\vect n}=\vect n^{\mathsf T}\mathsf I_O\vect n
=\int|\vect n\times\vect r|^2\dd m.
\end{equation}
展开为
\[
I_{\vect n}=I_{xx}n_x^2+I_{yy}n_y^2+I_{zz}n_z^2
+2I_{xy}n_xn_y+2I_{xz}n_xn_z+2I_{yz}n_yn_z.
\]
$T_O=I_{\vect n}\omega^2/2$，角动量轴向投影为
$\vect L_O\cdot\vect n=I_{\vect n}\omega$。
整个角动量 $\omega\mathsf I_O\vect n$ 未必沿轴，只有本征方向才能写成 $I_{\vect n}\omega\vect n$。
$I_{\vect n}\ge0$；等于零要求全部质量落在该轴上。
三维实体所有轴惯量均正，理想细杆有一个零惯量。
量纲为 $\mathrm{kg\,m^2}$；质量不变而全部长度放大 $k$，惯量乘 $k^2$。

\subsection{连续、离散描述与惯量的变换}\label{sub:c3-5.2}
离散与连续是同一质量二阶矩的不同计算表示。对位置 $\vect r_\alpha$ 的质点，
\[
(I_O)_{ij}=\sum_\alpha m_\alpha
(r_\alpha^2\delta_{ij}-r_{\alpha i}r_{\alpha j}).
\]
位置 $(a,b,c)$、质量 $m_\alpha$ 的质点贡献
\[
m_\alpha\begin{pmatrix}b^2+c^2&-ab&-ac\\-ab&a^2+c^2&-bc\\-ac&-bc&a^2+b^2\end{pmatrix}.
\]
原点质点贡献零；逐个求和时质量各乘一次，不能再额外乘总质量。
连续体用 $\dd m=\rho\dd V$，薄片用 $\sigma\dd S$，细杆用 $\lambda\dd s$。
柱坐标和球坐标体元分别带 $s$ 与 $r^2\sin\vartheta$ 的雅可比，必须保留。

\paragraph{分块极限与狄拉克表示。}
第 $\alpha$ 块质心为 $\vect a_\alpha$、自身张量为 $\mathsf I_{\alpha,C}$，其贡献为
\[
\mathsf I_{\alpha,O}=\mathsf I_{\alpha,C}
+m_\alpha(a_\alpha^2\mathsf1-\vect a_\alpha\vect a_\alpha^{\mathsf T}).
\]
把有限块当质点会丢掉第一项。若分块最大尺寸为 $\varepsilon$，
自身惯量总误差为量级 $m\varepsilon^2$，尺寸趋零时才消失。
有限部件组装须保留自身惯量。

反向可定义
\[
\rho(\vect r)=\sum_\alpha m_\alpha\delta^{(3)}(\vect r-\vect r_\alpha),
\quad \delta^{(3)}=\delta(x-x_\alpha)\delta(y-y_\alpha)\delta(z-z_\alpha).
\]
抽样性质为 $\int g(\vect r)\delta^{(3)}(\vect r-\vect r_\alpha)\dd V=g(\vect r_\alpha)$。
代入张量积分便逐项得到离散和。
上标3指三维分布，不是函数三次方；其量纲为体积倒数，保证 $\rho$ 的密度量纲。

\paragraph{转动基底：明确映射方向。}
$\mathsf R$ 的列是本体单位轴在空间基底中的坐标，因而
\[
\vect r_S=\mathsf R\vect r_B,\quad
\vect\omega_S=\mathsf R\vect\omega_B,\quad
\vect L_S=\mathsf R\vect L_B,\quad \mathsf R^{\mathsf T}\mathsf R=\mathsf1.
\]
由 $\vect L_B=\mathsf I_B\vect\omega_B$ 推得
\[
\vect L_S=\mathsf R\mathsf I_B\mathsf R^{\mathsf T}\vect\omega_S,
\qquad\mathsf I_S=\mathsf R\mathsf I_B\mathsf R^{\mathsf T}.
\]
若 $\mathsf Q$ 的列是新轴在旧基底中的表示，新坐标为 $\mathsf Q^{\mathsf T}\vect r$，
故 $\mathsf I'=\mathsf Q^{\mathsf T}\mathsf I\mathsf Q$。
矩阵顺序由映射方向决定。动能满足
\[
\vect\omega_S^{\mathsf T}\mathsf I_S\vect\omega_S
=\vect\omega_B^{\mathsf T}\mathsf I_B\vect\omega_B.
\]
迹、行列式、本征值也不变。本体矩阵为常数，空间矩阵随姿态改变。
% BEGIN ORIGINAL FIGURES
\OriginalFigure[6.2cm]{c3-3-p015-body-bases}{同一刚体上选择两组本体坐标轴}
% END ORIGINAL FIGURES


\paragraph{平行惯性轴与本体轴的速度计算。}
定点 $O$ 处取三根与实验室轴平行的固定基矢，质元写为
$\vect r=\xi\vect e_\xi+\eta\vect e_\eta+\zeta\vect e_\zeta$，则
$\vect v=\dot\xi\vect e_\xi+\dot\eta\vect e_\eta+\dot\zeta\vect e_\zeta$。
质元坐标随时间变化，基矢不变。若改用固连刚体的基矢，
$\vect r=x\vect e_1(t)+y\vect e_2(t)+z\vect e_3(t)$，
$x,y,z$ 恒定而 $\dot{\vect e}_i=\vect\omega\times\vect e_i$，因此
\[
\vect v=x\dot{\vect e}_1+y\dot{\vect e}_2+z\dot{\vect e}_3
=\vect\omega\times\vect r.
\]
两种计算是同一实际速度。若 $O$ 随体移动，两式都需另外加 $\vect v_O$。
本体轴使惯量积分中的材料坐标不变，却不意味着它是惯性坐标系。

\paragraph{平移原点：展开并消去一阶矩。}
令 $C$ 为质心，$\vect a=\vect r_{OC}$，$\vect r=\vect a+\vect s$，
$\int\vect s\dd m=0$。逐分量展开
\begin{align*}
r^2\delta_{ij}-r_ir_j
={}&s^2\delta_{ij}-s_is_j+a^2\delta_{ij}-a_ia_j\\
&+2(\vect a\cdot\vect s)\delta_{ij}-a_is_j-s_ia_j.
\end{align*}
第一组积分是 $\mathsf I_C$，第二组乘总质量，最后三项由质心一阶矩全消，得
\begin{equation}\label{eq:paraxis}
\mathsf I_O=\mathsf I_C+m(a^2\mathsf1-\vect a\vect a^{\mathsf T}).
\end{equation}
显式增量为
\[
\Delta\mathsf I=m\begin{pmatrix}
a_y^2+a_z^2&-a_xa_y&-a_xa_z\\
-a_xa_y&a_z^2+a_x^2&-a_ya_z\\
-a_xa_z&-a_ya_z&a_x^2+a_y^2\end{pmatrix}.
\]
沿 $z$ 平移只增 $I_x,I_y$，$I_z$ 不变；一般斜平移还会生成惯性积。
从偏心点回质心必须减去增量。两个原点都不是质心时，交叉项不自动消失，
应分别联系到质心或保留原一阶矩计算。

沿同一单位方向左右乘，得到
\[
I_{O,\vect n}=I_{C,\vect n}
+m[a^2-(\vect a\cdot\vect n)^2]
=I_{C,\vect n}+m|\vect a\times\vect n|^2.
\]
叉积模是两轴间距 $d_\perp$，即 $I_O=I_C+md_\perp^2$。
轴向平移不改转轴；偏离质心的平行轴惯量只增不减。
% BEGIN ORIGINAL FIGURES
\OriginalFigure[6.2cm]{c3-3-p013-parallel-axes}{两条平行转轴及其垂直间距}
% END ORIGINAL FIGURES


薄片全部在 $z=0$ 时，$I_z=\int(x^2+y^2)\dd m=I_x+I_y$。
一般体则 $I_x+I_y-I_z=2\int z^2\dd m\ge0$，不能照搬薄片的等号。
% BEGIN ORIGINAL FIGURES
\OriginalFigure[6.2cm]{c3-3-p013-perpendicular-axes}{薄片面内两轴与垂直于片面的转轴}
% END ORIGINAL FIGURES


\example{小块组装后的总惯量}\label{ex:c3-23}

\QuestionStatement{两个质量均为 $m$、自身中心惯量均为 $\mathsf I_0$ 的刚体，
其质心分别位于总质心两侧 $\pm a\vect e_x$，且轴方向相同。
两部件惯量张量 $\mathsf I_0$ 均在同向正交轴组中表达。求组合体关于总质心的完整惯量张量，解释不同方向的平行轴修正。}{ex23}{例23题图}
% BEGIN DETAILED ANSWER ex23
\solution \Step{确认总质心与公共基底}
以图~\ref{fig:solution-ex23}~中两个部件中心到横向轴的距离 $a$为依据，先写出图右侧第一行的几何或动力学关系，再代入其下方的计算式。

两质量相等且中心位于 $\pm a\vect e_x$，总质心确在原点 $C$。
部件张量都已沿同一组平行正交轴表达，因此不必先作姿态变换；总质量为 $2m$。
\SolutionFigure{ex23}{例23的计算示意}
\Step{逐个应用张量平行轴定理}
对于任一部件，位移矢量 $\vect a_\pm=\pm a\vect e_x$，有
\[
\vect a_\pm\vect a_\pm^{\mathsf T}=a^2\operatorname{diag}(1,0,0),\qquad
m(a^2\mathsf1-\vect a_\pm\vect a_\pm^{\mathsf T})=ma^2\operatorname{diag}(0,1,1).
\]
两部分直接相加，得到
\[
\mathsf I_C=2\mathsf I_0+2ma^2\begin{pmatrix}0&0&0\\0&1&0\\0&0&1\end{pmatrix}.
\]
交叉元仍为原部件对应交叉元的两倍，并没有因为组合对称就自动消失。
\Step{从轴距理解结果并推广}
$x$ 轴经过两个部件质心，惯量没有平移增量；$y,z$ 轴各与部件中心轴相距 $a$，每块都增加 $ma^2$。
若部件间距从 $2a$ 变为 $4a$，平移部分乘四，自身部分不变。
当两部件姿态不同，正确通式是
\[
\mathsf I_{\text{总},O}=\sum_\alpha\left[R_\alpha\mathsf I_{\alpha,C}R_\alpha^{\mathsf T}
+m_\alpha(a_\alpha^2\mathsf1-\vect a_\alpha\vect a_\alpha^{\mathsf T})\right].
\]
所有项都须用同一个原点和同一公共基底；$a=0$ 时恢复 $2\mathsf I_0$。

% END DETAILED ANSWER ex23
\subsection{一般运动的角动量与动能分解}\label{sub:c3-5.3}
以 $\vect s$ 为相对质心位置，速度场为
$\vect v=\vect v_C+\vect\omega\times\vect s$，完整动能展开为
\begin{align*}
T&=\frac12\int v_C^2\dd m+\int\vect v_C\cdot(\vect\omega\times\vect s)\dd m
+\frac12\int|\vect\omega\times\vect s|^2\dd m\\
&=\frac12mv_C^2+\vect v_C\cdot\left(\vect\omega\times\int\vect s\dd m\right)
+\frac12\vect\omega^{\mathsf T}\mathsf I_C\vect\omega.
\end{align*}
质心一阶矩为零，混合项消失。首项是全部质量集中在质心的平动能，末项是相对转动能。

取矩位置为 $O$，$\vect a=\vect r_{OC}$，用绝对速度定义角动量：
\begin{align*}
\vect L_O&=\int(\vect a+\vect s)\times(\vect v_C+\vect\omega\times\vect s)\dd m\\
&=m\vect a\times\vect v_C+
\left(\int\vect s\dd m\right)\times\vect v_C
+\vect a\times\left(\vect\omega\times\int\vect s\dd m\right)
+\mathsf I_C\vect\omega.
\end{align*}
中间两项同样消失，因此
\begin{equation}\label{eq:generalenergy}
T=\frac12mv_C^2+\frac12\vect\omega^{\mathsf T}\mathsf I_C\vect\omega,
\qquad\vect L_O=\vect r_{OC}\times m\vect v_C+\mathsf I_C\vect\omega.
\end{equation}
轨道角动量与相对转动角动量是矢量相加，其方向可以不同；取 $O=C$ 时前项为零。

\paragraph{非质心移动基点的混合项。}
取体内点 $A$，$\vect b=\vect r_{AC}$，$\int\vect r_{Ai}\dd m=m\vect b$。
以 $\vect v_A$ 为基点展开同一动能得到
\[
T=\frac12mv_A^2+m\vect v_A\cdot(\vect\omega\times\vect b)
+\frac12\vect\omega^{\mathsf T}\mathsf I_A\vect\omega.
\]
混合项一般不为零。把 $\vect v_C=\vect v_A+\vect\omega\times\vect b$
及平行轴定理代入质心式，也会得到这一结果。
若 $A$ 固定或此刻速度为零，就可只写关于 $A$ 的转动动能；
这是充分条件。一般移动基点需保留平动项与混合项，
特殊运动中这两项也可能恰好相消，其准确判据为
\[
\frac12v_A^2+\vect v_A\cdot(\vect\omega\times\vect b)=0.
\]
瞬间速度为零不保证加速度为零，不能自动把速度瞬心当永久固定的动力学取矩点。
动能由当前速度决定，角动量方程还涉及取矩点的时间变化，使用条件不同。

\example{质心平动能与定点转动能的一致性}\label{ex:c3-24}

\QuestionStatement{均匀细杆长 $L$、质量 $m$，绕端点以角速率 $\omega$ 在平面内转动。
杆在垂直于端点固定转轴的平面内转动。分别用定点积分和质心分解求动能，验证二者一致。}{ex24}{例24题图}
% BEGIN DETAILED ANSWER ex24
\solution \Step{直接用各质元速度积分}
以图~\ref{fig:solution-ex24}~中固定端 $O$、杆中点 $C$ 与全长 $L$为依据，先写出图右侧第一行的几何或动力学关系，再代入其下方的计算式。

从固定端点 $O$ 沿杆取 $s\in[0,L]$，$\dd m=(m/L)\dd s$。
杆与转轴垂直，质元速度为 $|\omega|s$，所以
\[
T=\frac12\frac mL\int_0^L\omega^2s^2\dd s
=\frac12\frac mL\omega^2\frac{L^3}{3}
=\frac{mL^2\omega^2}{6},\qquad I_O=\frac{mL^2}{3}.
\]
\SolutionFigure{ex24}{例24的计算示意}
\Step{改以质心分解平动和转动}
质心在 $L/2$ 处，故 $v_C=|\omega|L/2$。用质心坐标 $u\in[-L/2,L/2]$，
\[
I_C=\frac mL\int_{-L/2}^{L/2}u^2\dd u=\frac{mL^2}{12}.
\]
由于 $\int u\dd m=0$，速度平方展开的平动和转动混合项积分为零。于是
\[
T=\frac12mv_C^2+\frac12I_C\omega^2
=\frac{mL^2\omega^2}{8}+\frac{mL^2\omega^2}{24}
=\frac{mL^2\omega^2}{6}.
\]
\Step{平行轴定理与能量比例}
$I_O=I_C+m(L/2)^2=mL^2/3$，故 $\frac12I_O\omega^2$ 与前两条路线一致。
质心平动能与相对质心转动能之比是 $3:1$，即总能量的四分之三来自质心平动。
若把固定端点的转动能之外再加一次质心平动能，会重复计数；若只保留 $\frac12I_C\omega^2$，则遗漏平动能。
两种分解均描述同一个物体、同一真实运动。

% END DETAILED ANSWER ex24
\subsection{惯量主轴与惯量椭球}\label{sub:c3-5.4}
实对称惯量矩阵存在正交单位本征向量组。通过选定原点沿这些方向的轴称为主轴，
本征值为主惯量：
\[
\mathsf I\vect n_i=I_i\vect n_i,\qquad\vect n_i\cdot\vect n_j=\delta_{ij}.
\]
将三向量作列组成 $\mathsf Q$，有
$\mathsf Q^{\mathsf T}\mathsf I\mathsf Q=\operatorname{diag}(I_1,I_2,I_3)$。
取第三方向为前两方向的叉积，便构成右手轴组；向量反号不改变轴线。

\paragraph{不同本征方向为何正交？}
对称性保证 $\vect n_j^{\mathsf T}\mathsf I\vect n_i$
与 $\vect n_i^{\mathsf T}\mathsf I\vect n_j$ 相等，代入本征方程，
\[
(I_i-I_j)(\vect n_i\cdot\vect n_j)=0.
\]
不同本征值的方向正交；重根则须在完整子空间中选正交基。
已有单位向量 $\vect n_1$ 及另一个独立向量 $\vect v_2$ 时，可取
$\vect w_2=\vect v_2-(\vect n_1\cdot\vect v_2)\vect n_1$，再归一化。
同一重根空间的线性组合仍是本征向量，因此这一操作保留主惯量。
二重根平面内任意正交组可用；三重根时任意正交轴组可用。
不能用一组非正交主方向直接使用对角公式。

在主轴中 $\vect\omega=\sum_i\omega_i\vect n_i$，
\[
\vect L=\sum_iI_i\omega_i\vect n_i,
\qquad T=\frac12\sum_iI_i\omega_i^2.
\]
只有沿单轴或位于同一重根空间的角速度，才必与角动量平行。
同时具有不同主惯量方向的非零分量时，方向一般会被惯量重新加权。

\paragraph{方向惯量的极值法。}
在单位约束下极值化 $\vect n^{\mathsf T}\mathsf I\vect n$，
取 $G=\vect n^{\mathsf T}\mathsf I\vect n-\lambda(\vect n^{\mathsf T}\vect n-1)$，
$\nabla G=0$ 给出 $\mathsf I\vect n=\lambda\vect n$。
所以极值方向为主方向，极值为主惯量。
写 $\vect n=\sum_i c_i\vect n_i$、$\sum_i c_i^2=1$，
$I_{\vect n}=\sum_iI_ic_i^2$，因而
\[
I_{\min}\le I_{\vect n}\le I_{\max}.
\]
斜轴结果应落在这个范围内。

\paragraph{惯量椭球的辅助坐标与量纲。}
选参考惯量 $I_*>0$，用无量纲辅助位置描述二次型：
\[
\vect q^{\mathsf T}\mathsf I\vect q=I_*.
\]
这些坐标不是真实质元位置。令 $\vect q=q\vect n$，
有 $q^2I_{\vect n}=I_*$，截距为 $\sqrt{I_*/I_{\vect n}}$。
主轴中为 $I_1q_1^2+I_2q_2^2+I_3q_3^2=I_*$，半轴为 $\sqrt{I_*/I_i}$。
小惯量对应长半轴、大惯量对应短半轴，惯量椭球一般不等于物体外形。
若右端写数值1而不引入参考惯量，辅助坐标应有惯量负二分之一次幂量纲。

\paragraph{法向几何法。}
$f(\vect q)=\vect q^{\mathsf T}\mathsf I\vect q-I_*$ 的法向为
$\nabla f=2\mathsf I\vect q$。主轴端点的切平面垂直径向，因而
\[
\mathsf I\vect q=\lambda\vect q,
\qquad\text{或}\qquad\vect q\times\mathsf I\vect q=0.
\]
反过来满足该式的曲面点，径向就是主方向，故几何法与本征法等价。
叉积式无需除以坐标分量，不会漏掉含零分量的轴向解。
零主惯量对应方向没有有限截距，应按退化二次曲面理解。

\TeachingFigure{c3-diagram-23}{惯量椭球的主截面与方向惯量}

\paragraph{本征求解的具体顺序。}
先提取有量纲系数 $\mathsf I=I_0\mathsf A$，再求无量纲根
$\det(\mathsf A-\lambda\mathsf1)=0$。三阶多项式由不变量给出
\[
\lambda^3-(\operatorname{tr}\mathsf A)\lambda^2
+\frac12[(\operatorname{tr}\mathsf A)^2-\operatorname{tr}(\mathsf A^2)]\lambda
-\det\mathsf A=0.
\]
每个根代回 $(\mathsf A-\lambda_i\mathsf1)\vect v_i=0$，
解非零向量并归一化。物理主惯量为 $I_0\lambda_i$，不能遗漏量纲系数。
重根时求完整零空间，再选正交基。
检查根的和、积、两两点积以及本征残差，数值截断只应引入小误差。

真实惯量还满足三角不等式。以主轴中的实际质元坐标 $s_i$ 计算，
\[
I_1+I_2-I_3
=\int[(s_2^2+s_3^2)+(s_3^2+s_1^2)-(s_1^2+s_2^2)]\dd m
=2\int s_3^2\dd m\ge0.
\]
循环交换指标得另外两式。质量全在 $s_3=0$ 平面内才取等号，
三维实体通常严格为正。任意正定矩阵未必满足这些额外约束，
因此正定性、三角不等式应一起用于物理核验。

\subsection{质量对称性与常见实体的积分}\label{sub:c3-5.5}
对称性可在积分前消去若干矩阵元，但前提是质量分布关于选定原点真正对称。
相同外形的两半若材料密度不同，就不能按镜面对称消去惯性积。
原点移开后也不能继续沿用原来的对称判断。

\paragraph{镜面对称怎样消去积分？}
若 $xy$ 平面是质量镜面，每个 $(x,y,z)$ 质元都有等质量的 $(x,y,-z)$ 质元。
两者的 $xz$ 与 $yz$ 贡献互相抵消，所以
\[
\int xz\dd m=\int yz\dd m=0,
\qquad\mathsf I=\begin{pmatrix}I_{xx}&I_{xy}&0\\I_{xy}&I_{yy}&0\\0&0&I_{zz}\end{pmatrix}.
\]
$z$ 因此是主轴，但 $xy$ 块仍可能含非零 $I_{xy}$，片面镜面对称不自动选出另外两轴。
若还有一个与它垂直的质量镜面，两个法线与两面交线构成三根正交主轴。
例如关于 $xz,yz$ 都对称时，$xy,xz,yz$ 各至少对某一反射奇变，全部惯性积为零。

若关于 $z$ 轴转 $180^\circ$ 不改变质量分布，点 $(x,y,z)$ 配对为 $(-x,-y,z)$。
$xz,yz$ 变号而 $xy$ 不变，所以同样只保证 $z$ 是主轴，不能据此令 $I_x=I_y$。
长方体绕一棱平行中心轴有二重旋转对称，两个横向惯量仍可不同。
% BEGIN ORIGINAL FIGURES
\OriginalFigure[6.2cm]{c3-3-p019-paired-symmetry}{绕 $z$ 轴半周旋转对应的成对质元：$(x,y,z)$ 与 $(-x,-y,z)$}
% END ORIGINAL FIGURES


\paragraph{四重旋转为什么强制横向惯量相等？}
写一般对称矩阵
\[
\mathsf I=\begin{pmatrix}A&D&E\\D&B&G\\E&G&C\end{pmatrix},
\qquad
\mathsf R_z(\pi/2)=\begin{pmatrix}0&-1&0\\1&0&0\\0&0&1\end{pmatrix}.
\]
质量分布转 $90^\circ$ 不变，张量也应不变。直接相乘，
\[
\mathsf R_z\mathsf I\mathsf R_z^{\mathsf T}
=\begin{pmatrix}B&-D&-G\\-D&A&E\\-G&E&C\end{pmatrix}=\mathsf I.
\]
比较对应位置得 $A=B,D=0,E=-G,G=E$，所以 $E=G=0$。
因此 $\mathsf I=\operatorname{diag}(A,A,C)$，过原点的任意横向方向都是主方向。
四个等质、等形侧部的塔体或理想四扇旋转门可满足这一条件，
即使绕其他角度转过之后外形不同。

更一般，$n\ge3$ 重旋转角不是0或 $\pi$。
一个具有两个不同本征值的二维二次型，其不变轴只允许半周旋转，
不能在这样的角度下不变，故横向块必须为常数乘二维单位矩阵。
同时横向的 $(I_{xz},I_{yz})$ 若在非零旋转下不变，只能是零矢量。
连续轴对称自然满足这些要求，但不是二阶惯量横向相等的必要条件。

人体左右外形、镜像和瓶形可帮助识别候选轴，实际使用时仍要确认密度、姿态与原点。
一个人与镜像不是两个同时存在的质量体；瓶塞、手柄、附着部件也会破坏理想对称。
以下积分先用对称性选主轴，再计算数值，两步不可互相代替。

\TeachingFigure{c3-diagram-24}{离散旋转对称也能确定一组主轴与横向主惯量的相等性}
% BEGIN ORIGINAL FIGURES
\OriginalFigure[6.2cm]{c3-3-p022-symmetry-examples}{旋转对称、镜面对称与四重旋转对称的例子；推断主轴还需相应的质量分布对称性}
% END ORIGINAL FIGURES



\example{均匀矩形薄片与椭圆薄片}\label{ex:c3-25}

\QuestionStatement{质量为 $m$ 的均匀矩形薄片边长 $a,b$，质心为原点，$x,y$ 沿边，$z$ 垂直片面。
再考虑质量同为 $m$、半轴为 $a,b$ 的均匀椭圆薄片，质心与轴向约定相同。分别从积分求两种薄片的完整惯量张量，并检验垂直轴定理。}{ex25}{例25题图}
% BEGIN DETAILED ANSWER ex25
\solution \Step{明确模型与计算依据}
以图~\ref{fig:solution-ex25}~中矩形全边长与椭圆半轴的区别为依据，先写出图右侧第一行的几何或动力学关系，再代入其下方的计算式。
 矩形区域为 $-a/2\le x\le a/2$、$-b/2\le y\le b/2$，$z=0$，
质量面密度 $\sigma=m/(ab)$。先计算两个平方二阶矩：
\begin{align*}
\int x^2\dd m
&=\sigma\int_{-a/2}^{a/2}x^2\dd x\int_{-b/2}^{b/2}\dd y
=\sigma\frac{a^3}{12}b=\frac{ma^2}{12},\\
\int y^2\dd m
&=\sigma\int_{-a/2}^{a/2}\dd x\int_{-b/2}^{b/2}y^2\dd y
=\sigma a\frac{b^3}{12}=\frac{mb^2}{12}.
\end{align*}
$\int xy\dd m=\sigma(\int x\dd x)(\int y\dd y)=0$，
另外两个惯性积因为 $z=0$ 而为零。
因此
\[
I_x=\frac{mb^2}{12},\qquad I_y=\frac{ma^2}{12},
\qquad I_z=\frac{m(a^2+b^2)}{12}.
\]
$I_x$ 含的是垂直 $x$ 轴的边长 $b$，不是沿 $x$ 的边长 $a$。
垂直轴结果也由 $I_z=I_x+I_y$ 校验。

\SolutionFigure{ex25}{例25的计算示意}
\Step{椭圆薄片的变量代换}
另取半轴 $a,b$ 的椭圆，区域 $x^2/a^2+y^2/b^2\le1$，面密度为 $m/(\pi ab)$。
令 $x=au\cos\chi,y=bu\sin\chi$，$0\le u\le1,0\le\chi<2\pi$。
雅可比为
\[
\left|\frac{\partial(x,y)}{\partial(u,\chi)}\right|
=\left|\begin{matrix}a\cos\chi&-au\sin\chi\\b\sin\chi&bu\cos\chi\end{matrix}\right|
=ab\,u,
\qquad\dd S=ab\,u\dd u\dd\chi.
\]
检验质量为 $\sigma ab(\int_0^1u\dd u)(\int_0^{2\pi}\dd\chi)=m$。
对 $x$ 轴，质元垂距平方是 $y^2=b^2u^2\sin^2\chi$，故
\[
I_x=\sigma ab\,b^2\int_0^1u^3\dd u\int_0^{2\pi}\sin^2\chi\dd\chi
=\frac{m}{\pi}b^2\frac14\pi=\frac{mb^2}{4}.
\]
同理 $I_y=ma^2/4$，$
\int xy\dd m=\sigma a^2b^2\int_0^1u^3\dd u\int_0^{2\pi}\sin\chi\cos\chi\dd\chi=0$，
所以
\[
I_z=I_x+I_y=\frac{m(a^2+b^2)}4.
\]
$a=b$ 时得到均匀圆盘 $I_x=I_y=ma^2/4,I_z=ma^2/2$。
矩形里的 $a,b$ 是全边长，椭圆里的 $a,b$ 是半轴；
若比较相同外接尺寸，应将矩形边长取为 $2a,2b$，再比较系数。
二者几何主轴方向相同，质量分布与数值惯量不同。

% END DETAILED ANSWER ex25

% BEGIN ORIGINAL FIGURES
\OriginalFigure[6.2cm]{c3-3-p021-lamina-axes}{矩形薄片与椭圆薄片的几何主轴}
% END ORIGINAL FIGURES


\example{均匀圆柱与细杆}\label{ex:c3-26}

\QuestionStatement{圆柱质量 $m$、半径 $a$、高 $L$，质心 $z$ 轴沿圆柱轴。
求完整质心惯量张量，并取半径趋零且质量、长度保持不变的细杆极限，求过质心、与杆夹角为 $\alpha$ 的轴惯量。}{ex26}{例26题图}
% BEGIN DETAILED ANSWER ex26
\solution \Step{明确模型与计算依据}
以图~\ref{fig:solution-ex26}~中圆柱中心轴 $z$ 与横向距离为依据，先写出图右侧第一行的几何或动力学关系，再代入其下方的计算式。
 密度为 $\rho=m/(\pi a^2L)$。柱坐标取
$x=s\cos\chi,y=s\sin\chi$，区域为
\[
0\le s\le a,\quad0\le\chi<2\pi,\quad-L/2\le z\le L/2,
\quad\dd V=s\dd s\dd\chi\dd z.
\]
体积积分先给 $\pi a^2L$，可检查积分上下限和密度。
轴向垂距平方为 $s^2$，故
\[
I_z=\rho\int_{-L/2}^{L/2}\dd z
\int_0^{2\pi}\dd\chi\int_0^a s^3\dd s
=\rho L(2\pi)\frac{a^4}{4}=\frac12ma^2.
\]
横向惯量 $I_x=\int(y^2+z^2)\dd m$ 必须分别计算两项。
\begin{align*}
\int y^2\dd m
&=\rho\int_{-L/2}^{L/2}\dd z
\int_0^{2\pi}\sin^2\chi\dd\chi\int_0^as^3\dd s
=\rho L\pi\frac{a^4}{4}=\frac{ma^2}{4},\\
\int z^2\dd m
&=\rho\int_{-L/2}^{L/2}z^2\dd z
\int_0^{2\pi}\dd\chi\int_0^as\dd s
=\rho\frac{L^3}{12}(2\pi)\frac{a^2}{2}=\frac{mL^2}{12}.
\end{align*}
于是
\[
I_x=I_y=\frac{m(3a^2+L^2)}{12},\qquad I_z=\frac{ma^2}{2}.
\]
$xy$ 的角积分为零；$xz,yz$ 至少对 $z$ 为奇函数，故交叉元全零。

\SolutionFigure{ex26}{例26的计算示意}
\Step{圆盘、圆环与细杆的极限}
$L\to0$ 且质量保持不变，变成薄圆盘，横向系数为1/4、轴向为1/2。
质量若集中在半径 $a$ 的圆周上，轴向垂距处处为 $a$，则 $I_z=ma^2$；
由平面与轴对称得到两个横向惯量为 $ma^2/2$。
盘和环差别来自径向质量分布，不能只看外半径相同。

$a\to0$、总质量不变时，圆柱趋于细杆，
$\mathsf I_C=(mL^2/12)\operatorname{diag}(1,1,0)$，这里杆沿 $z$。
也可直接用线积分 $\frac mL\int_{-L/2}^{L/2}z^2\dd z=mL^2/12$。
过质心的轴与杆夹角为 $\alpha$ 时，
\[
\vect n=(\sin\alpha\cos\gamma,\sin\alpha\sin\gamma,\cos\alpha),
\qquad I_{\vect n}=\frac{mL^2}{12}\sin^2\alpha.
\]
沿杆为零，垂直杆为 $mL^2/12$，与方位角 $\gamma$ 无关。
忽略截面厚度使张量半正定，在求反矩阵时必须处理这个退化。

% END DETAILED ANSWER ex26

% BEGIN ORIGINAL FIGURES
\OriginalFigure[6.2cm]{c3-3-p014-rod-direction}{细杆的方向惯量：图中杆沿 $x$、夹角为 $\theta$，对应正文杆沿 $z$、夹角为 $\alpha$ 的同一关系}
% END ORIGINAL FIGURES


\TeachingFigure{c3-diagram-25}{质量对称性使惯性积消失，但主惯量仍须计算}

\example{均匀实心球的惯量}\label{ex:c3-27}
% BEGIN ORIGINAL FIGURES
\OriginalFigure[6.2cm]{c3-3-p038-solid-sphere}{均匀实心球的球心主轴}
% END ORIGINAL FIGURES
\QuestionStatement{均匀球半径 $a$、质量 $m$，球心为原点。
求球心惯量张量和任意过球心轴的惯量，用质量积分与对称性相互核验。}{ex27}{例27题图}
% BEGIN DETAILED ANSWER ex27
\solution \Step{明确模型与计算依据}
以图~\ref{fig:solution-ex27}~中过球心的任意轴与半径 $a$为依据，先写出图右侧第一行的几何或动力学关系，再代入其下方的计算式。
 密度为 $\rho=3m/(4\pi a^3)$，质量区域是整个实心球。
球坐标取
\[
x=r\sin\vartheta\cos\chi,\quad y=r\sin\vartheta\sin\chi,
\quad z=r\cos\vartheta,
\]
\[
0\le r\le a,\quad0\le\vartheta\le\pi,\quad0\le\chi<2\pi,
\quad\dd V=r^2\sin\vartheta\dd r\dd\vartheta\dd\chi.
\]
对 $z$ 轴，$x^2+y^2=r^2\sin^2\vartheta$，所以
\[
I_z=\rho\int_0^ar^4\dd r
\int_0^\pi\sin^3\vartheta\dd\vartheta\int_0^{2\pi}\dd\chi.
\]
角积分用 $u=\cos\vartheta$，$\dd u=-\sin\vartheta\dd\vartheta$，得到
\[
\int_0^\pi\sin^3\vartheta\dd\vartheta
=\int_{-1}^{1}(1-u^2)\dd u
=\left[u-\frac{u^3}{3}\right]_{-1}^{1}=\frac43.
\]
径向为 $a^5/5$，方位角为 $2\pi$，故
\[
I_z=\rho\frac{a^5}{5}\frac43(2\pi)
=\frac{8\pi\rho a^5}{15}=\frac25ma^2.
\]
坐标反射使三个惯性积为零，球旋转对称使三个对角元相等，因而
$\mathsf I_C=(2ma^2/5)\mathsf1$。

\SolutionFigure{ex27}{例27的计算示意}
\Step{第二路线：先求迹}
逐对角元相加，
\[
\operatorname{tr}\mathsf I=2\int r^2\dd m
=2\rho\int_0^ar^4\dd r\int_0^\pi\sin\vartheta\dd\vartheta\int_0^{2\pi}\dd\chi
=\frac{8\pi\rho a^5}{5}.
\]
由三对角相等，各惯量等于迹的1/3，再得 $2ma^2/5$。
还可计算单个二阶矩，
\[
\int z^2\dd m=\rho\frac{a^5}{5}
\int_{-1}^{1}u^2\dd u(2\pi)=\frac{ma^2}{5},
\]
$I_x$ 是 $y^2,z^2$ 两项之和。这三种积分相互校验，避免把角度雅可比漏掉。
任意过球心轴都有 $I_{\vect n}=2ma^2/5$，角动量与角速度同向。

薄球壳质量集中在 $r=a$，面密度 $\sigma=m/(4\pi a^2)$，面积元为
$a^2\sin\vartheta\dd\vartheta\dd\chi$，故
\[
I_z=\sigma a^4\int_0^\pi\sin^3\vartheta\dd\vartheta\int_0^{2\pi}\dd\chi
=\frac23ma^2.
\]
薄壳没有实心球的径向平均，不能沿用2/5。偏心平行轴则在相应中心结果上加 $md^2$。

% END DETAILED ANSWER ex27
\example{均匀实心椭球的惯量}\label{ex:c3-28}

\QuestionStatement{椭球边界为 $x^2/a^2+y^2/b^2+z^2/c^2=1$，质量 $m$。
求质心惯量张量和三根主轴，列出变量代换、体积元及积分限，并检验球形极限。}{ex28}{例28题图}
% BEGIN DETAILED ANSWER ex28
\solution \Step{明确模型与计算依据}
以图~\ref{fig:solution-ex28}~中椭球的三条几何半轴为依据，先写出图右侧第一行的几何或动力学关系，再代入其下方的计算式。
 椭球体积为 $4\pi abc/3$，均匀密度为 $\rho=3m/(4\pi abc)$。
直接球坐标的径向上限依角度变化：
\[
r_0(\vartheta,\chi)=\left(
\frac{\sin^2\vartheta\cos^2\chi}{a^2}
+\frac{\sin^2\vartheta\sin^2\chi}{b^2}
+\frac{\cos^2\vartheta}{c^2}\right)^{-1/2}.
\]
虽然可以按此积分，仿射代换更能分离三个半轴的影响。
令 $x=au,y=bv,z=cw$，区域化为 $u^2+v^2+w^2\le1$，
雅可比为 $abc$，质量元变成 $\rho abc\dd u\dd v\dd w$。
变换后密度 $\rho'=\rho abc=3m/(4\pi)$，单位球总质量仍为 $m$。

对单位球，再令 $u=s\sin\vartheta\cos\chi$，
$v=s\sin\vartheta\sin\chi,w=s\cos\vartheta$，
其中 $0\le s\le1,0\le\vartheta\le\pi,0\le\chi<2\pi$。
计算一个平方矩即可：
\[
\int_{\text{单位球}}w^2\dd V'
=\int_0^1s^4\dd s\int_0^\pi\cos^2\vartheta\sin\vartheta\dd\vartheta
\int_0^{2\pi}\dd\chi
=\frac15\frac23(2\pi)=\frac{4\pi}{15}.
\]

\SolutionFigure{ex28}{例28的计算示意}
\Step{逐项计算与求解}
球对称给 $u^2,v^2$ 同样积分，交叉矩全零。
因此
\begin{align*}
\int x^2\dd m&=a^2\rho abc\frac{4\pi}{15}=\frac{ma^2}{5},\\
\int y^2\dd m&=\frac{mb^2}{5},\qquad
\int z^2\dd m=\frac{mc^2}{5}.
\end{align*}

\Step{结果与适用条件}
将每两个平方矩相加，
\[
\mathsf I_C=\frac m5\operatorname{diag}(b^2+c^2,c^2+a^2,a^2+b^2).
\]
$xz,xy,yz$ 在单位球各至少对一个坐标为奇函数，所以无交叉元。
$a=b=c$ 退化为球 $2ma^2/5$；$c\to0$ 且质量保持不变时所得极薄实心椭球
不是均匀椭圆薄片，因为投影后的面密度随位置变化，不能改用薄片1/4系数。
这里是积分变量代换，并非声称材料真实变形后惯量不变。
长方体全边长为 $a,b,c$ 时二阶矩系数为1/12，
与椭球半轴符号比较时应先统一实际外接尺寸。

% END DETAILED ANSWER ex28
\example{均匀实心半球：球心与质心原点}\label{ex:c3-29}

\QuestionStatement{半球为 $x^2+y^2+z^2\le a^2,z\ge0$，半球质量 $m$。
求球心 $O$ 与质心 $C$ 的惯量张量。
同时求质心位置，说明移点后三个惯量中哪一个保持不变。}{ex29}{例29题图}
% BEGIN DETAILED ANSWER ex29
\solution \Step{明确模型与计算依据}
以图~\ref{fig:solution-ex29}~中球心 $O$ 到质心 $C$ 的轴向位移为依据，先写出图右侧第一行的几何或动力学关系，再代入其下方的计算式。
 半球体积为 $2\pi a^3/3$，所以 $\rho=3m/(2\pi a^3)$。
球心 $O$ 是底面圆心，不是质量中心。横向坐标反射使 $xy,xz,yz$ 三个惯性积都为零。
对平方矩，半球是同密度完整球的一半，而完整球质量为 $2m$，故
\[
\mathsf I_O=\frac12\left[\frac25(2m)a^2\mathsf1\right]
=\frac25ma^2\mathsf1.
\]
此处各向同性是球心二阶矩的性质，不表示半球形状或质心张量各向同性。

\SolutionFigure{ex29}{例29的计算示意}
\Step{水平圆截面与平方矩}
高度 $z\in[0,a]$ 的截面半径为 $R(z)=\sqrt{a^2-z^2}$，面积为 $\pi(a^2-z^2)$。
于是
\[
\int z^2\dd m=\rho\pi\int_0^az^2(a^2-z^2)\dd z
=\rho\pi a^5\left(\frac13-\frac15\right)
=\frac{2\rho\pi a^5}{15}=\frac{ma^2}{5}.
\]
截面圆盘的 $x^2$ 面积分为 $\pi R^4/4$，由圆极坐标
$\int_0^{R}s^3\dd s\int_0^{2\pi}\cos^2\chi\dd\chi$ 可直接得到。
再沿 $z$ 积分，
\[
\int x^2\dd m=\frac{\rho\pi}{4}\int_0^a(a^2-z^2)^2\dd z
=\frac{\rho\pi a^5}{4}\left(1-\frac23+\frac15\right)
=\frac{ma^2}{5}.
\]
$y^2$ 同理，任意对角惯量是两项之和，再次得到 $2ma^2/5$。

\Step{一阶矩与质心平移}
$x_C=y_C=0$；质量积分为
$m=\rho\pi\int_0^a(a^2-z^2)\dd z=2\rho\pi a^3/3$。
质心高度为
\[
z_C=\frac{\rho\pi}{m}\int_0^az(a^2-z^2)\dd z
=\frac{\rho\pi}{m}\left[\frac{a^2z^2}{2}-\frac{z^4}{4}\right]_0^a
=\frac{\rho\pi a^4}{4m}=\frac{3a}{8}.
\]
也可用半球球坐标 $0\le\vartheta\le\pi/2$ 得
$\int z\dd m=\rho(\int_0^ar^3\dd r)(\int_0^{\pi/2}\cos\vartheta\sin\vartheta\dd\vartheta)(2\pi)
=\rho\pi a^4/4$。

$\vect r_{OC}=(3a/8)\vect k$，平移增量为
$ma^2\operatorname{diag}(9/64,9/64,0)$。
从 $O$ 回到 $C$ 减去它：
\[
\mathsf I_C=ma^2\operatorname{diag}
\left(\frac25-\frac9{64},\frac25-\frac9{64},\frac25\right)
=ma^2\operatorname{diag}\left(\frac{83}{320},\frac{83}{320},\frac25\right).
\]
两个横向轴变近质心，惯量减少；$z$ 轴前后是同一条轴，惯量不变。
穿质心并与 $z$ 成角 $\alpha$ 的方向有
$I_{\vect n}=ma^2[(83/320)\sin^2\alpha+(2/5)\cos^2\alpha]$，
落在两个主惯量之间，与极值界限一致。
从球心移到质心的符号、改变的对角位置与系数都须一起检查。

% END DETAILED ANSWER ex29
\TeachingFigure{c3-diagram-26}{半球的球心不在质心处；平移只改变横向轴惯量}
% BEGIN ORIGINAL FIGURES
\OriginalFigure[6.2cm]{c3-3-p042-hemisphere-pair}{由两个同密度半球组成完整球：两者共用球心原点}
% END ORIGINAL FIGURES


\subsection{正方体、长方体与主轴的具体求解}\label{sub:c3-5.6}
\example{同一个正方体在两个原点下的惯量}\label{ex:c3-30}
% BEGIN ORIGINAL FIGURES
\OriginalFigure[6.2cm]{c3-3-p009-cube-center}{以正方体质心为原点的本体坐标轴}
% END ORIGINAL FIGURES
\QuestionStatement{均匀实心正方体边长 $a$、质量 $m$，坐标轴平行棱。
分别取质心 $C$ 与一个顶点 $O$ 为原点。
分别求 $C$、$O$ 为原点时的完整惯量张量，以直接积分和平行轴定理交叉核验。}{ex30}{例30题图}
% BEGIN DETAILED ANSWER ex30
\solution \Step{明确模型与计算依据}
以图~\ref{fig:solution-ex30}~中正方体顶点 $O$ 与质心 $C$为依据，先写出图右侧第一行的几何或动力学关系，再代入其下方的计算式。
 密度为 $\rho=m/a^3$。先取质心原点，区域为 $[-a/2,a/2]^3$。
对一个平方二阶矩，
\[
\int x^2\dd m=\rho\int_{-a/2}^{a/2}x^2\dd x
\int_{-a/2}^{a/2}\dd y\int_{-a/2}^{a/2}\dd z
=\rho\frac{a^3}{12}a^2=\frac{ma^2}{12}.
\]
另两平方矩相同；$xy$ 积分可分成两个奇函数积分而为零，$xz,yz$ 同理。
对角惯量是两平方矩之和，所以
\[
\mathsf I_C=\frac{ma^2}{6}\mathsf1,
\qquad\vect L_C=\frac{ma^2}{6}\vect\omega,
\qquad T_C=\frac{ma^2\omega^2}{12}.
\]
任何过质心轴惯量都是 $ma^2/6$。
这不是因为正方体是球，而是它的质量二阶矩恰好各向同性，
外形或更高阶形状矩仍有方向差异。

\SolutionFigure{ex30}{例30的计算示意}
\Step{顶点原点的直接积分}
区域改为 $0\le x,y,z\le a$，密度不变。
\begin{align*}
\int x^2\dd m&=\rho\left(\int_0^ax^2\dd x\right)a^2
=\rho\frac{a^5}{3}=\frac{ma^2}{3},\\
\int xy\dd m&=\rho\left(\int_0^ax\dd x\right)
\left(\int_0^ay\dd y\right)\left(\int_0^a\dd z\right)
=\rho\frac{a^2}{2}\frac{a^2}{2}a=\frac{ma^2}{4}.
\end{align*}
循环交换变量，全部平方矩与全部交叉矩分别相同，故
\[
\mathsf I_O=ma^2\begin{pmatrix}
2/3&-1/4&-1/4\\-1/4&2/3&-1/4\\-1/4&-1/4&2/3
\end{pmatrix}.
\]
从顶点看，三个坐标都非负，交叉矩不再因奇对称抵消。

% END DETAILED ANSWER ex30

% BEGIN ORIGINAL FIGURES
\OriginalFigure[6.2cm]{c3-3-p011-cube-vertex}{以正方体顶点为原点、沿棱建立的本体坐标轴}
% END ORIGINAL FIGURES


\paragraph{张量平行轴的逐项校验。}
$\vect r_{OC}=\tfrac a2(\vect i+\vect j+\vect k)$，
$|OC|^2=3a^2/4$，外积的每个位置都是 $a^2/4$。
因此
\[
\Delta\mathsf I=ma^2\begin{pmatrix}
1/2&-1/4&-1/4\\-1/4&1/2&-1/4\\-1/4&-1/4&1/2
\end{pmatrix}.
\]
加上中心对角元1/6，得到顶点对角元 $1/6+1/2=2/3$，
交叉元保持 $-1/4$，与直接积分一致。

若绕顶点 $z$ 棱以 $\vect\omega=\omega\vect k$ 转动，
矩阵乘法选出第三列，
\[
\vect L_O=ma^2\omega(-\vect i/4-\vect j/4+2\vect k/3),
\qquad T=\frac12I_{zz}\omega^2=\frac13ma^2\omega^2.
\]
角动量有两个横向分量，说明棱方向不是主轴。
若固定这条轴运动，相应支承必须容许需要的约束作用，不能假定角动量必沿轴。

若绕体对角线，$\vect n=(\vect i+\vect j+\vect k)/\sqrt3$，
矩阵每行对 $(1,1,1)$ 的作用为 $2/3-1/4-1/4=1/6$，故
\[
\vect L_O=\frac{ma^2\omega}{6}\vect n,
\qquad I_{\vect n}=\frac{ma^2}{6},\qquad T=\frac{ma^2\omega^2}{12}.
\]
这条直线经过质心，顶点和质心给的是同一条转轴，所以惯量不增加。
两次转动比较显示：原点相同，方向不同，角动量是否平行与动能数值都可以不同。

\TeachingFigure{c3-diagram-27}{顶点原点下的正方体：体对角线为主轴，三条棱并非主轴}

\example{顶点正方体的三根正交主轴}\label{ex:c3-31}

\QuestionStatement{均匀正方体质量 $m$、边长 $a$，取一顶点 $O$ 为原点，三条正棱分别沿 $x,y,z$。求关于 $O$ 的全部主惯量和一组单位正交主轴，并说明重根主轴是否唯一。}{ex31}{例31题图}
% BEGIN DETAILED ANSWER ex31
\solution \Step{明确模型与计算依据}
以图~\ref{fig:solution-ex31}~中从顶点指向对顶点的体对角线为依据，先写出图右侧第一行的几何或动力学关系，再代入其下方的计算式。
 本例用三种互相独立的路线寻找全部主方向。
设 $\mathsf I_O=(ma^2/12)\mathsf A$，
\[
\mathsf A=\begin{pmatrix}8&-3&-3\\-3&8&-3\\-3&-3&8\end{pmatrix}.
\]

\SolutionFigure{ex31}{例31的计算示意}
\Step{第一法：代数本征求解}
令 $\vect u=(1,1,1)^{\mathsf T}$，有 $\mathsf A=11\mathsf1-3\vect u\vect u^{\mathsf T}$。
由于 $\vect u^{\mathsf T}\vect u=3$，
$\mathsf A\vect u=(11-9)\vect u=2\vect u$；
若 $\vect u^{\mathsf T}\vect v=0$，则 $\mathsf A\vect v=11\vect v$。
这已经确定一个一维本征空间与一个二维本征空间。

若按常规行列式求解，利用上述分解或从行列式两行相减，可以因式分解为
\[
\det(\mathsf A-\lambda\mathsf1)
=(2-\lambda)(11-\lambda)^2
=-\lambda^3+24\lambda^2-165\lambda+242.
\]
根为 $2,11,11$；换回有量纲主惯量得
\[
I_1=\frac{ma^2}{6},\qquad I_2=I_3=\frac{11ma^2}{12}.
\]
对 $\lambda=2$，齐次方程是
\[
\begin{pmatrix}6&-3&-3\\-3&6&-3\\-3&-3&6\end{pmatrix}
\begin{pmatrix}v_x\\v_y\\v_z\end{pmatrix}=0.
\]
前两式相减给 $v_x=v_y$，再代回得 $v_x=v_y=v_z$。
对 $\lambda=11$，三行均为 $-3(v_x+v_y+v_z)=0$，
所以完整解集为平面 $v_x+v_y+v_z=0$。
选择其内两个正交方向，得到
\[
\vect n_1=\frac{(1,1,1)}{\sqrt3},\quad
\vect n_2=\frac{(1,-1,0)}{\sqrt2},\quad
\vect n_3=\frac{(1,1,-2)}{\sqrt6}.
\]
逐一检查点积均为零、长度均为1，且 $\vect n_1\times\vect n_2=\vect n_3$。
平面内任意方向都是一个主方向，但必须再选正交的两根，才形成三维对角轴组。

\Step{第二法：从质心各向同性和平行轴几何求解}
质心张量为 $ma^2\mathsf1/6$，所有过质心方向惯量相同。
令 $\vect R=\vect r_{OC}$，$R^2=3a^2/4$。
与 $OC$ 夹角为 $\gamma$ 的顶点轴距质心的垂距为 $R\sin\gamma$，因此
\[
I_{\vect n}=\frac{ma^2}{6}+mR^2\sin^2\gamma
=ma^2\left(\frac16+\frac34\sin^2\gamma\right).
\]
$\gamma=0$ 时惯量最小，为1/6；$\gamma=\pi/2$ 时最大，为11/12。
沿 $OC$ 的方向与它的整个垂直平面因此分别是两个本征空间。
此法直接显示平移为何将中心的三重退化变为一维加二维，不需展开三次式。

\Step{第三法：惯量椭球的径向与法向}
设非零辅助径向为 $\vect q=(x,y,z)$，$\Sigma=x+y+z$。
$\mathsf A\vect q=11\vect q-3\Sigma(1,1,1)$，所以
\[
\vect q\times\mathsf A\vect q=-3\Sigma\vect q\times(1,1,1).
\]
径向平行法向要求叉积为零，得到
\[
(x-y)\Sigma=0,\qquad(y-z)\Sigma=0,\qquad(z-x)\Sigma=0.
\]
若 $\Sigma\ne0$，三个坐标相同，得到体对角线。
若 $\Sigma=0$，整个垂直体对角线的平面都是解，不能再把其中某几条特殊直线误列为独立分支。
代回张量仍得1/6与11/12。对应椭球沿体对角线的半轴较长，
横向两半轴相同；正方体外形并不因此变成旋转体。

三种方法分别利用谱分解、实际轴距与曲面法向，
都给出相同本征空间与物理主惯量，是比只比较三个数字更完整的互相核验。

% END DETAILED ANSWER ex31
\example{顶点长方体的惯量与数值主轴}\label{ex:c3-32}

\QuestionStatement{均匀长方体边长 $a,b,c$，取顶点原点，棱沿 $x,y,z$。
求惯量张量，并在 $a=3c,b=2c$ 时计算主轴。
数值结果用 $mc^2$ 作惯量单位，给出特征方程、本征值和一组单位正交本征矢量。}{ex32}{例32题图}
% BEGIN DETAILED ANSWER ex32
\solution \Step{明确模型与计算依据}
以图~\ref{fig:solution-ex32}~中顶点三条棱与质心位置为依据，先写出图右侧第一行的几何或动力学关系，再代入其下方的计算式。
 均匀密度为 $\rho=m/(abc)$，顶点区域为
$0\le x\le a,0\le y\le b,0\le z\le c$。
独立平方与交叉矩分别为
\[
\int x^2\dd m=\rho\frac{a^3}{3}bc=\frac{ma^2}{3},
\quad\int xy\dd m=\rho\frac{a^2}{2}\frac{b^2}{2}c=\frac{mab}{4},
\]
循环交换给另外四个量。因此
\[
\mathsf I_O=\frac m{12}\begin{pmatrix}
4(b^2+c^2)&-3ab&-3ac\\-3ab&4(c^2+a^2)&-3bc\\-3ac&-3bc&4(a^2+b^2)
\end{pmatrix}.
\]
也可用中心张量
$\mathsf I_C=(m/12)\operatorname{diag}(b^2+c^2,c^2+a^2,a^2+b^2)$，
质心位移 $(a/2,b/2,c/2)$，平行轴逐项给同一式。
取 $a=3c,b=2c$，公共系数为 $mc^2/12$，
\[
\mathsf I_O=\frac{mc^2}{12}\mathsf A,
\qquad\mathsf A=\begin{pmatrix}20&-18&-9\\-18&40&-6\\-9&-6&52\end{pmatrix}.
\]

\SolutionFigure{ex32}{例32的计算示意}
\Step{特征多项式的系数怎样得到？}
$\operatorname{tr}\mathsf A=20+40+52=112$。
三个二阶主子式为
\[
20\cdot40-18^2=476,\quad20\cdot52-9^2=959,
\quad40\cdot52-6^2=2044,
\]
和为3479。沿第一行展开行列式，
\[
\det\mathsf A=20(40\cdot52-6^2)+18(-18\cdot52-54)-9(108+360)=18848.
\]
因此特征方程为
\[
f(\lambda)=\lambda^3-112\lambda^2+3479\lambda-18848=0.
\]
本例三个主惯量互异，没有正方体的退化平面，三轴各自确定到正负号。

\Step{数值根的定位与迭代}
可先检查符号：$f(6)=-1790,f(7)=360$，
$f(50)=102,f(51)=-80$，$f(54)=-110,f(55)=72$。
因此三个根分别位于 $(6,7),(50,51),(54,55)$。
用二分法在各区间缩小范围，或用牛顿法
\[
\lambda^{(k+1)}=\lambda^{(k)}-
\frac{f(\lambda^{(k)})}{3(\lambda^{(k)})^2-224\lambda^{(k)}+3479}.
\]
例如从7出发，第一步为 $7-360/2058\approx6.825073$，继续得到 $6.826408$。
另外两根为50.506181和54.667411。对称矩阵保证所有根实数，
本例积分正定又保证三个物理主惯量正值。

\Step{从根求方向，不能只列最终数字}
$\mathsf A-\lambda\mathsf1$ 的前两行为
$\vect r_1=(20-\lambda,-18,-9)$、$\vect r_2=(-18,40-\lambda,-6)$。
根处矩阵秩为2，一个与这两行都垂直的向量就是零空间方向：
\[
\vect v(\lambda)=\vect r_1\times\vect r_2
=(468-9\lambda,\ 282-6\lambda,\ \lambda^2-60\lambda+476).
\]
代入各根后，用 $\vect n=\vect v/|\vect v|$，并按需要整体反号，得到
\begin{center}
\begin{tabular}{c r r r r}
\toprule
&$\lambda$&$n_x$&$n_y$&$n_z$\\\midrule
1&6.826408&.836600&.496001&.232556\\
2&50.506181&$-.533299$&.834479&.138697\\
3&54.667411&$-.125270$&$-.240056$&.962642\\\bottomrule
\end{tabular}
\end{center}
第一根例如未归一化方向约为 $(406.562,241.042,113.015)$，
归一化产生表中第一行。另两方向用同一叉积公式即可复现。

物理主惯量为 $mc^2\lambda_i/12$，约为
$0.568867mc^2,4.208848mc^2,4.555618mc^2$。
根之和为112，根两两乘积之和为3479，根积为18848。
还须检查 $|\vect n_i|=1$、$\vect n_i\cdot\vect n_j=0$、
$\mathsf A\vect n_i-\lambda_i\vect n_i\approx0$。
表中截断会带来末位残差，改变某轴整体符号不改变物理意义；
三根取成右手组后，即可用于后续角速度本体分量。

% END DETAILED ANSWER ex32
\example{长方体绕体对角线的方向惯量}\label{ex:c3-33}

\QuestionStatement{均匀长方体质量 $m$，边长 $a,b,c$，顶点 $O$ 为原点，三条棱沿坐标正轴。转轴为连接 $O$ 与对顶点的整条体对角线。
求该轴的方向惯量，说明它是否必为顶点惯量张量的主轴，并核验正方体特例。}{ex33}{例33题图}
% BEGIN DETAILED ANSWER ex33
\solution \Step{明确模型与计算依据}
以图~\ref{fig:solution-ex33}~中通过顶点和质心的同一体对角线为依据，先写出图右侧第一行的几何或动力学关系，再代入其下方的计算式。
 体对角线方向为 $\vect n=(a,b,c)/\sqrt{a^2+b^2+c^2}$，且经过质心。
先用中心原点计算二阶矩，
\[
\int x^2\dd m=\frac{m}{abc}\frac{a^3}{12}bc=\frac{ma^2}{12},
\quad\int y^2\dd m=\frac{mb^2}{12},\quad\int z^2\dd m=\frac{mc^2}{12}.
\]

\SolutionFigure{ex33}{例33的计算示意}
\Step{逐项计算与求解}
坐标上下限分别为 $[-a/2,a/2],[-b/2,b/2],[-c/2,c/2]$，
交叉积分均因奇对称为零。代入中心对角张量得
\begin{align*}
I_{\vect n}
&=\frac{m}{12}
\frac{a^2(b^2+c^2)+b^2(c^2+a^2)+c^2(a^2+b^2)}{a^2+b^2+c^2}\\
&=\frac m6\frac{a^2b^2+b^2c^2+c^2a^2}{a^2+b^2+c^2}.
\end{align*}


\Step{结果与适用条件}
再用顶点张量独立检验。令分母为 $S=a^2+b^2+c^2$，
三个对角贡献为 $4[a^2(b^2+c^2)+b^2(c^2+a^2)+c^2(a^2+b^2)]$，
交叉项贡献为 $-6(a^2b^2+a^2c^2+b^2c^2)$。
前者是 $8$ 倍同一个两两乘积和，减去后剩2倍，乘 $m/(12S)$ 仍得上式。
顶点计算中交叉项不可略去；只有中心对角张量才没有这些项。

体对角线过质心，轴距为零，所以两原点得到同一轴惯量。
但一般 $a,b,c$ 不相等时，中心三个主惯量不同，
$\mathsf I_C(a,b,c)$ 的三个分量比例也不同，角动量不会沿体对角线。
若三条边严格正值，体对角线要成为主轴需三个中心主惯量相等，也就是 $a=b=c$。
这一特例给 $ma^2/6$。方向惯量可计算与方向是主轴，是两个不同结论。

% END DETAILED ANSWER ex33
\example{三个质点组成的等边三角刚体}\label{ex:c3-34}

\QuestionStatement{三个质点各质量 $m$，固定在边长 $a$ 的等边三角形顶点。
取一个顶点为原点，另两点坐标为
$(a,0,0)$ 和 $(a/2,\sqrt3a/2,0)$。
求关于 $O$ 的完整惯量张量、三个主惯量及一组正交主轴，以质心和平行轴定理校验。}{ex34}{例34题图}
% BEGIN DETAILED ANSWER ex34
\solution \Step{明确模型与计算依据}
以图~\ref{fig:solution-ex34}~中三个离散质量到各坐标轴的距离为依据，先写出图右侧第一行的几何或动力学关系，再代入其下方的计算式。
 原点质点对 $O$ 的惯量贡献为零。
位于 $(a,0,0)$ 的质点贡献为 $ma^2\operatorname{diag}(0,1,1)$，
位于 $(a/2,\sqrt3a/2,0)$ 的质点贡献为
\[
ma^2\begin{pmatrix}3/4&-\sqrt3/4&0\\-\sqrt3/4&1/4&0\\0&0&1\end{pmatrix}.
\]
相加得到
\[
\mathsf I_O=ma^2\begin{pmatrix}3/4&-\sqrt3/4&0\\-\sqrt3/4&5/4&0\\0&0&2\end{pmatrix}.
\]
这里的 $m$ 是每个质点质量，总质量是 $3m$；已逐点相加后不能再乘3。
$z$ 方向直接是主轴，主惯量 $2ma^2$，并满足平面垂直轴定理。

\SolutionFigure{ex34}{例34的计算示意}
\Step{平面二阶本征方程}
令平面无量纲本征值为 $\lambda$，
\[
\left(\frac34-\lambda\right)\left(\frac54-\lambda\right)
-\frac3{16}=\lambda^2-2\lambda+\frac34=0.
\]
得到 $\lambda=1/2,3/2$。
对1/2根，第一行是 $v_x-\sqrt3v_y=0$，归一化得
$\vect n_1=(\sqrt3/2,1/2,0)$；另一正交方向可取
$\vect n_2=(-1/2,\sqrt3/2,0)$，并取 $\vect n_3=(0,0,1)$。
所以
\[
I_1=\frac{ma^2}{2},\qquad I_2=\frac{3ma^2}{2},\qquad I_3=2ma^2.
\]
$\vect n_1$ 沿原点至对边中点，角为 $30^\circ$，符合图中的几何主轴。

\Step{总质心与平行轴定理}
三点平均位置为 $C=(a/2,\sqrt3a/6,0)$，$|OC|^2=a^2/3$。
总质量 $3m$，平移增量为
\[
3m(R^2\mathsf1-\vect R\vect R^{\mathsf T})
=ma^2\begin{pmatrix}1/4&-\sqrt3/4&0\\-\sqrt3/4&3/4&0\\0&0&1\end{pmatrix}.
\]
从 $\mathsf I_O$ 减去它，得到
$\mathsf I_C=ma^2\operatorname{diag}(1/2,1/2,1)$。
质心平面张量具有二重退化；移到顶点后，沿 $OC$ 的轴距仍为零，
惯量保持 $ma^2/2$，另一个横向轴距为 $|OC|$，增加 $3m|OC|^2=ma^2$，
成为 $3ma^2/2$；垂直平面的轴也加 $ma^2$，成为 $2ma^2$。
这把离散求和、二阶特征方程与几何轴距连在一起。

% END DETAILED ANSWER ex34
\TeachingFigure{c3-diagram-28}{离散惯量求和与等边三角形的顶点主轴}
% BEGIN ORIGINAL FIGURES
\OriginalFigure[6.2cm]{c3-3-p047-triangle-masses}{三个等质量质点构成的等边三角刚体；原点取左下顶点}
% END ORIGINAL FIGURES


\subsection{八分之一球的完整张量计算}\label{sub:c3-5.7}
\example{第一卦限球体的惯量、角动量与主轴}\label{ex:c3-35}

\QuestionStatement{均匀实心球半径 $a$，只保留 $x,y,z\ge0$ 的部分，
剩余部分质量为 $m$。以球心 $O$ 为固定点，
求惯量张量；当
$\vect\omega=\omega(\vect i+\vect j)/\sqrt2$ 时求角动量和动能；
求主轴，并把原点移到质心后重做这些计算。
求角动量和动能时，角速度分量均指同一瞬时的本体坐标分量。}{ex35}{例35题图}
% BEGIN DETAILED ANSWER ex35
\solution \Step{明确模型与计算依据}
以图~\ref{fig:solution-ex35}~中第一卦限球体的三个坐标边界为依据，先写出图右侧第一行的几何或动力学关系，再代入其下方的计算式。
 以下按四个问题顺序计算，始终把 $m$ 作为剩余八分之一球的质量。
球坐标约定为 $x=r\sin\vartheta\cos\chi,y=r\sin\vartheta\sin\chi,z=r\cos\vartheta$。
第一卦限的积分区域与体元为
\[
0\le r\le a,\quad0\le\vartheta\le\pi/2,
\quad0\le\chi\le\pi/2,
\quad\dd V=r^2\sin\vartheta\dd r\dd\vartheta\dd\chi.
\]
体积为 $(a^3/3)\cdot1\cdot(\pi/2)=\pi a^3/6$，密度是
$\rho=6m/(\pi a^3)$。同密度完整球质量为 $8m$，不能用剩余 $m$ 当整球质量。

\SolutionFigure{ex35}{例35的计算示意}
\Step{第一问：完整计算六个独立二阶矩}
先算平方矩。$x^2=r^2\sin^2\vartheta\cos^2\chi$，因此
\[
\int x^2\dd m=\rho\int_0^ar^4\dd r
\int_0^{\pi/2}\sin^3\vartheta\dd\vartheta
\int_0^{\pi/2}\cos^2\chi\dd\chi.
\]
三个因子依次为 $a^5/5,2/3,\pi/4$，乘积为 $\rho\pi a^5/30=ma^2/5$。
$y^2$ 的方位角积分改为 $\sin^2\chi$，仍为 $\pi/4$。
对 $z^2$，极角积分变为
$\int_0^{\pi/2}\cos^2\vartheta\sin\vartheta\dd\vartheta=1/3$，
方位角为 $\pi/2$，所以也是 $ma^2/5$。
于是每个对角惯量都是 $2ma^2/5$。

但 $xy,xz,yz$ 在此区域内均非负，没有跨坐标平面的质元与它们配对抵消。
以 $xy$ 为例，
\begin{align*}
\int xy\dd m
&=\rho\int_0^ar^4\dd r
\int_0^{\pi/2}\sin^3\vartheta\dd\vartheta
\int_0^{\pi/2}\sin\chi\cos\chi\dd\chi\\
&=\rho\frac{a^5}{5}\frac23\frac12
=\frac{\rho a^5}{15}=\frac{2ma^2}{5\pi}.
\end{align*}
对 $xz$，角因子为
$\int_0^{\pi/2}\sin^2\vartheta\cos\vartheta\dd\vartheta=1/3$
与 $\int_0^{\pi/2}\cos\chi\dd\chi=1$，仍得 $\rho a^5/15$；
$yz$ 的方位角积分为1，也相同。
这些相等性也可由 $x,y,z$ 的排列对称直接判断，
但排列对称不是 $x\mapsto-x$ 一类镜面对称，不能把交叉矩设零。

令 $d=2/5,p=2/(5\pi)$，完整张量为
\[
\mathsf I_O=ma^2\begin{pmatrix}d&-p&-p\\-p&d&-p\\-p&-p&d\end{pmatrix}.
\]
迹为 $6ma^2/5$；非对角元负号来自张量定义，$p$ 本身为正惯性积系数。

\Step{第二问：逐行乘角速度并求动能}
$\vect\omega=(\omega/\sqrt2)(1,1,0)$。第一行给 $(d-p)\omega/\sqrt2$，
第二行相同，第三行给 $-2p\omega/\sqrt2$，故
\[
\vect L_O=\frac{ma^2\omega}{\sqrt2}
[(d-p)\vect i+(d-p)\vect j-2p\vect k].
\]
即使 $\omega_z=0$，角动量也有非零 $z$ 分量。
点乘时第三分量不参与，因为角速度第三分量为零；
\[
T_O=\frac12\frac{ma^2\omega^2}{2}[2(d-p)]
=\frac12ma^2\omega^2(d-p)
=\frac{ma^2\omega^2}{5}\left(1-\frac1\pi\right).
\]
用速度积分再次检查，$\vect v=(\omega/\sqrt2)(z,-z,y-x)$，
\[
T_O=\frac{\omega^2}{4}\int[2z^2+(y-x)^2]\dd m
=\frac{\omega^2}{4}
\left[\frac{4ma^2}{5}-\frac{4ma^2}{5\pi}\right],
\]
与张量结果一致。这里的 $(y-x)^2$ 特别显示交叉矩不能略去。

\Step{第三问：主轴空间与重根}
写 $\vect u=(1,1,1)^{\mathsf T}$，则
$\mathsf I_O/(ma^2)=(d+p)\mathsf1-p\vect u\vect u^{\mathsf T}$。
沿 $\vect u$，本征系数为 $d+p-3p=d-2p$；
在 $\vect u^\perp$ 平面，外积项为零，本征系数为 $d+p$。
特征行列式因此分解为
\[
\left[(d-2p)-\lambda\right]\left[(d+p)-\lambda\right]^2=0.
\]
可以选择正交主轴
\[
\vect n_1=\frac{(1,1,1)}{\sqrt3},\quad
\vect n_2=\frac{(1,-1,0)}{\sqrt2},\quad
\vect n_3=\frac{(1,1,-2)}{\sqrt6},
\]
主惯量为
\[
I_{O1}=\frac{2ma^2}{5}\left(1-\frac2\pi\right),
\qquad I_{O2}=I_{O3}=\frac{2ma^2}{5}\left(1+\frac1\pi\right).
\]
给定角速度在此轴组中的分量为
$\omega_1=\sqrt{2/3}\,\omega,\omega_2=0,\omega_3=\omega/\sqrt3$。
于是主轴动能为
\[
\frac12\left(\frac23I_{O1}+\frac13I_{O3}\right)\omega^2
=\frac12ma^2\omega^2(d-p),
\]
再次与坐标矩阵及速度积分相符。

\Step{第四问之一：先求质心一阶矩}
排列对称给 $x_C=y_C=z_C$，仍需积分求其数值。
\[
\int x\dd m=\rho\int_0^ar^3\dd r
\int_0^{\pi/2}\sin^2\vartheta\dd\vartheta
\int_0^{\pi/2}\cos\chi\dd\chi
=\rho\frac{a^4}{4}\frac\pi4\cdot1=\frac{3ma}{8}.
\]
对 $z$，极角因子为1/2，方位角为 $\pi/2$，得到同一值。
所以
\[
\vect R=\vect r_{OC}=\frac{3a}{8}(\vect i+\vect j+\vect k),
\qquad R^2=\frac{27a^2}{64}.
\]
也可固定 $x$，剩下 $yz$ 截面是半径 $\sqrt{a^2-x^2}$ 的四分之一圆，
$\int x\dd m=(\rho\pi/4)\int_0^ax(a^2-x^2)\dd x=\rho\pi a^4/16$，独立验证一阶矩。

\Step{第四问之二：逐矩阵元移至质心}
平行轴的增量为
\[
m(R^2\mathsf1-\vect R\vect R^{\mathsf T})
=\frac{ma^2}{64}\begin{pmatrix}18&-9&-9\\-9&18&-9\\-9&-9&18\end{pmatrix}.
\]
从 $\mathsf I_O$ 减去它，对角为 $d-18/64=19/160$，
交叉元为 $-p+9/64$。令
\[
D=\frac{19}{160},\qquad B=\frac9{64}-\frac{2}{5\pi},
\qquad
\mathsf I_C=ma^2\begin{pmatrix}D&B&B\\B&D&B\\B&B&D\end{pmatrix}.
\]
$B>0$，意味着实际中心矩阵的交叉元此时为正，不要仍机械添加负号。
迹减少 $2mR^2=27ma^2/32$，等于三对角减少量，另有一个整体核验。

与相同角速度相乘，得到质心相对角动量与相对转动能
\[
\vect L_C=\frac{ma^2\omega}{\sqrt2}
[(D+B)\vect i+(D+B)\vect j+2B\vect k],
\]
\[
T_C=\frac12ma^2\omega^2(D+B)
=ma^2\omega^2\left(\frac{83}{640}-\frac1{5\pi}\right).
\]
此处中心角动量 $z$ 分量已变为正，与 $O$ 处角动量方向不同。
原点改变同时改变轨道项；这并非同一个角动量矢量被随意改了坐标。

中心矩阵可写成 $(D-B)\mathsf1+B\vect u\vect u^{\mathsf T}$，主方向空间不变，
\[
I_{C1}=ma^2(D+2B)=ma^2\left(\frac25-\frac4{5\pi}\right),
\]
\[
I_{C2}=I_{C3}=ma^2(D-B)
=ma^2\left(\frac25+\frac2{5\pi}-\frac{27}{64}\right).
\]
第一根轴沿 $OC$，轴距为零，所以 $I_{C1}=I_{O1}$；
横向两轴移到质心各减 $mR^2=27ma^2/64$。
数值上第一项约 $0.145352ma^2$、横向项约 $0.105449ma^2$，都为正。
主方向空间未变，但最小最大顺序已改变：在 $O$ 处第一轴最小，在 $C$ 处第一轴最大。
这能提醒我们“主方向不变”不等于“惯量大小排序不变”。

\Step{同一真实运动的动能与角动量如何恢复？}
题设物体真实绕固定 $O$ 转动，质心速度不能消失：
\[
\vect v_C=\vect\omega\times\vect R
=\frac{3a\omega}{8\sqrt2}(\vect i-\vect j),
\qquad v_C^2=\frac{9a^2\omega^2}{64}.
\]
因此质心平动能为 $9ma^2\omega^2/128$，
\[
T_C+\frac12mv_C^2
=ma^2\omega^2\left(\frac{83}{640}+\frac9{128}-\frac1{5\pi}\right)
=\frac{ma^2\omega^2}{5}\left(1-\frac1\pi\right)=T_O.
\]
角动量的轨道项为
\[
\vect R\times m\vect v_C
=\frac{9ma^2\omega}{64\sqrt2}(\vect i+\vect j-2\vect k).
\]
把它加到 $\vect L_C$ 上，前两分量由 $D+B+9/64=d-p$，
第三分量由 $2B-18/64=-2p$，逐分量恰还原 $\vect L_O$。
所以将参考原点移至质心后，必须同时重建轨道平动部分，
不能只更换惯量矩阵就把中心相对量当作原定点运动的总量。

% END DETAILED ANSWER ex35
\TeachingFigure{c3-diagram-29}{第一卦限球体有排列对称性，却没有三个坐标镜面对称面}
% BEGIN ORIGINAL FIGURES
\OriginalFigure[6.2cm]{c3-3-p024-octant-sphere}{第一卦限球体的边界与球心坐标轴}
% END ORIGINAL FIGURES


\subsection{复杂转动中的动能分解}\label{sub:c3-5.8}
\example{顶点固定正方体的斜轴转动}\label{ex:c3-36}

\QuestionStatement{正方体边长 $a$、质量 $m$，顶点 $O$ 固定，
$\vect\omega=\omega(\vect i+\vect j)/\sqrt2$。
比较定点动能与质心分解。
同时求质心速度，并说明角速度方向不是主轴时，动能仍应如何计算。}{ex36}{例36题图}
% BEGIN DETAILED ANSWER ex36
\solution \Step{明确模型与计算依据}
以图~\ref{fig:solution-ex36}~中顶点固定时的质心位置矢量为依据，先写出图右侧第一行的几何或动力学关系，再代入其下方的计算式。
 首先直接用完整顶点张量作二次型。
$\omega_x=\omega_y=\omega/\sqrt2,\omega_z=0$，两个对角贡献相加为
$(2/3)ma^2\omega^2$，交叉贡献 $2I_{xy}\omega_x\omega_y$ 为
$-(1/4)ma^2\omega^2$，故
\[
T_O=\frac12ma^2\omega^2\left(\frac23-\frac14\right)
=\frac5{24}ma^2\omega^2.
\]
角动量也可逐行乘得
\[
\vect L_O=\frac{ma^2\omega}{\sqrt2}
\left(\frac5{12}\vect i+\frac5{12}\vect j-\frac12\vect k\right),
\]
与角速度点乘的一半仍给 $5ma^2\omega^2/24$。

\SolutionFigure{ex36}{例36的计算示意}
\Step{由质元速度积分计算动能}
$\vect v=(\omega/\sqrt2)(z,-z,y-x)$，
顶点区域 $[0,a]^3$ 上有
$\int x^2\dd m=\int y^2\dd m=\int z^2\dd m=ma^2/3$，
$\int xy\dd m=ma^2/4$。于是
\begin{align*}
T_O&=\frac{\omega^2}{4}\int[2z^2+(y-x)^2]\dd m\\
&=\frac{ma^2\omega^2}{4}
\left(\frac23+\frac13+\frac13-\frac12\right)
=\frac5{24}ma^2\omega^2.
\end{align*}
这条计算不预先使用张量动能式，可以独立检验交叉项的系数和符号。

\Step{质心平动与相对转动}
中心张量各向同性，$T_C=ma^2\omega^2/12$。
$\vect R=(a/2)(1,1,1)$，
\[
\vect v_C=\vect\omega\times\vect R
=\frac{a\omega}{2\sqrt2}(\vect i-\vect j),
\quad v_C^2=\frac{a^2\omega^2}{4}.
\]
所以
\[
T_C+\frac12mv_C^2
=\frac{ma^2\omega^2}{12}+\frac{ma^2\omega^2}{8}
=\frac5{24}ma^2\omega^2.
\]
相同物体、相同瞬间的三条路线一致。
若只保留矩阵两个对角元，就会漏掉 $xy$ 交叉项；
若在定点动能上再加一次质心平动，又会重复计数。
这两类错误都可通过上述独立检查直接识别。

% END DETAILED ANSWER ex36
\example{固定点上方的圆盘陀螺动能}\label{ex:c3-37}
% BEGIN ORIGINAL FIGURES
\OriginalFigure[6.2cm]{c3-3-p049-gyroscope-axes}{圆盘姿态与空间轴 $(\xi,\eta,\zeta)$、本体轴 $(x,y,z)$；图中 $O$ 为盘心，正文悬点距盘心为 $h$}
% END ORIGINAL FIGURES
\QuestionStatement{无质量杆一端固定于 $O$，另一端为均匀薄圆盘中心，
中心距 $O$ 为 $h$，盘质量 $m$、半径 $a$，
杆沿盘对称轴。求关于固定点的惯量及欧拉角动能。
采用本章欧拉角约定，保留所有进动、章动与自转交叉项，并用质心平动加相对转动核验。}{ex37}{例37题图}
% BEGIN DETAILED ANSWER ex37
\solution \Step{明确模型与计算依据}
以图~\ref{fig:solution-ex37}~中固定支点 $O$ 至盘心的距离 $h$为依据，先写出图右侧第一行的几何或动力学关系，再代入其下方的计算式。
 以杆方向为本体第三轴，先求圆盘质心惯量。
均匀薄盘由前面积分有
\[
\mathsf I_C=\operatorname{diag}(ma^2/4,ma^2/4,ma^2/2).
\]
位移 $\vect r_{OC}=h\vect e_3$，平行轴增量为 $mh^2\operatorname{diag}(1,1,0)$，所以
\[
I_1=I_2=m(h^2+a^2/4),\qquad I_3=ma^2/2.
\]
轴向转轴是同一条直线，不因 $h$ 增加而增加轴向惯量；
横向转动使盘心绕 $O$ 运动，因此增加 $mh^2$。

本体欧拉角速度分量为
\[
\omega_1=\dot\theta\cos\psi+\dot\varphi\sin\theta\sin\psi,
\quad\omega_2=-\dot\theta\sin\psi+\dot\varphi\sin\theta\cos\psi,
\quad\omega_3=\dot\psi+\dot\varphi\cos\theta.
\]
前两式平方相加，混合项正负抵消，得到
\[
\omega_1^2+\omega_2^2=\dot\theta^2+\dot\varphi^2\sin^2\theta.
\]
由于横向两惯量相等，动能只需这一个和，故
\[
T=\frac12I_1(\dot\theta^2+\dot\varphi^2\sin^2\theta)
+\frac12I_3(\dot\psi+\dot\varphi\cos\theta)^2.
\]
轴向角速度是两项之和，不能把自转、进动能量分别平方后漏掉交叉项。
$\psi$ 本身不出现，反映均匀盘绕自身轴转后的质量分布不变。

\SolutionFigure{ex37}{例37的计算示意}
\Step{用质心分解解释各项}
盘心位置为 $h\vect e_3$，
$\vect v_C=h\vect\omega\times\vect e_3$，故
\[
v_C^2=h^2(\omega_1^2+\omega_2^2)
=h^2(\dot\theta^2+\dot\varphi^2\sin^2\theta).
\]
质心相对转动动能为
\[
T_C=\frac{ma^2}{8}(\omega_1^2+\omega_2^2)+\frac{ma^2}{4}\omega_3^2.
\]
与 $mv_C^2/2$ 相加后，横向系数成为
$m(h^2+a^2/4)/2$，恰为定点张量表达。
因此 $mh^2$ 项已经包含盘心平动，不能在定点动能后再加一次。

匀速进动时 $\theta=\theta_0,\dot\theta=0,\dot\varphi=\Omega,\dot\psi=\omega$，
\[
T=\frac12m(h^2+a^2/4)\Omega^2\sin^2\theta_0
+\frac14ma^2(\omega+\Omega\cos\theta_0)^2.
\]
这里 $\omega$ 表示 $\dot\psi$，真实轴向角速度为 $\omega+\Omega\cos\theta_0$。
从空间矢量合成看，总角速度为 $\vect\omega_{\text{总}}=\Omega\vect e_Z+\omega\vect e_3$，
其中 $\vect e_Z$ 竖直，$\vect e_Z\cdot\vect e_3=\cos\theta_0$，故
\[
|\vect\omega_{\text{总}}|^2=\Omega^2+\omega^2+2\Omega\omega\cos\theta_0.
\]
在某瞬间把第一本体轴取向竖直轴在横向平面投影的反方向，可以写成
\[
\vect\omega_{\text{总}}=-\Omega\sin\theta_0\vect e_1
+(\omega+\Omega\cos\theta_0)\vect e_3.
\]
横向正负取决于本体第一轴取向，其平方和与动能不受这一选择影响。

$h=0$ 时恢复绕盘中心的公式；$\Omega=0$ 且章动为零时，只剩纯轴向自转能。
$\theta_0\to0$ 时欧拉角发生几何退化，但能量只依赖有限组合 $\dot\psi+\dot\varphi$，
不能把两个退化角速率理解为两项额外独立的物理转动。

% END DETAILED ANSWER ex37
\clearpage
\section{定轴转动、物理摆与轴承压力}\label{sec:c3-6}
\subsection{绕固定轴的动力学与能量}\label{sub:c3-6.1}
设固定转轴为 $z$，有向转角为 $q$，则
$\vect\omega=\dot q\,\vect e_z$、$\vect\alpha=\ddot q\,\vect e_z$。
关于轴的角动量分量为 $L_z=I_z\dot q$，
其中 $I_z$ 为物体关于实际转轴的惯量，恒定。
% BEGIN ORIGINAL FIGURES
\OriginalFigure[6.2cm]{c3-4-p002-fixed-axis-cases}{定轴转动中的转轴位置：穿过实体内部、边缘或实体之外}
% END ORIGINAL FIGURES

角动量定理投影到轴上给出
\begin{equation}\label{eq:fixedaxis}
I_z\ddot q=M_z.
\end{equation}
理想轴承的约束力对转轴的轴向矩为零，
但轴承摩擦偶矩、马达驱动矩和主动外力矩均须计入 $M_z$。

物体的角动量矢量不一定沿固定转轴。
定轴运动只要求角速度沿轴，并不要求该轴是主轴。
% BEGIN ORIGINAL FIGURES
\OriginalFigure[6.2cm]{c3-4-p002-axial-translation}{绕轴旋转并沿轴平移的两个时刻：轴向移动使它不属于定轴转动}
% END ORIGINAL FIGURES

若转轴不是主轴，横向角动量分量一般非零，轴承必须提供相应力偶，
以维持转轴不变。仅用一个标量转动方程能求 $q(t)$，
却不能求全部支承反力。

动能与功的轴向形式为
\[
T=\frac12I_z\dot q^2,\qquad
\dd T=M_z\,\dd q,\qquad \mathcal P=M_z\dot q.
\]
若 $M_z$ 只依赖 $q$，积分得
\[
\frac12I_z(\dot q^2-\dot q_0^2)=\int_{q_0}^qM_z(s)\,\dd s.
\]
若驱动矩随时间变化，用时间积分转动方程；
若摩擦力矩随速度方向变号，必须分段判断而不能在过零后仍沿用同一符号。

\example{恒定力矩驱动圆盘}\label{ex:c3-38}

\QuestionStatement{均匀圆盘质量 $m$、半径 $a$，绕中心轴转动，
恒定驱动矩为 $M$，恒定阻矩大小为 $M_f<M$，初始静止。
取驱动方向为正、初始转角为 $q_0$；假设装置能克服启动静摩擦。求角加速度、$\omega(t)$、$q(t)$ 及转过 $\Delta q>0$ 后的转速，并作功和能量核验。}{ex38}{例38题图}
% BEGIN DETAILED ANSWER ex38
\solution \Step{确定惯量与净驱动力矩}
以图~\ref{fig:solution-ex38}~中两个相反方向的力矩为依据，先写出图右侧第一行的几何或动力学关系，再代入其下方的计算式。

均匀盘绕中心对称轴的惯量为 $I_z=ma^2/2$。按驱动方向取正，盘转动时驱动力矩为 $+M$，阻矩为 $-M_f$；故净矩为 $M-M_f>0$。
\SolutionFigure{ex38}{例38的计算示意}
\Step{列转动方程并积分初值}
\[
I_z\dot\omega=M-M_f,\qquad
\alpha=\frac{2(M-M_f)}{ma^2}=\text{常数}.
\]
先用 $\omega(0)=0$ 积分，再用 $q(0)=q_0$ 积分，得到
\[
\omega(t)=\frac{2(M-M_f)}{ma^2}t,\qquad
q(t)=q_0+\frac{M-M_f}{ma^2}t^2.
\]
惯量把力矩的单位 $\mathrm{kg\,m^2s^{-2}}$ 转为角加速度单位 $\mathrm{s^{-2}}$。
\Step{时间消元与能量关系}
由 $\omega\dd\omega=\alpha\dd q$，
\[
\omega^2=2\alpha\Delta q=\frac{4(M-M_f)\Delta q}{ma^2}.
\]
净功 $(M-M_f)\Delta q$ 恰为 $\frac12I_z\omega^2$；驱动输入功 $M\Delta q$ 分成动能与耗散 $M_f\Delta q$。
在 $M_f=0$ 时恢复理想恒矩加速；$M\downarrow M_f$ 时加速度趋零。
本题已明示可启动，若真实启动静摩擦上限大于 $M$，则初始静止不能直接套用此滑动阻矩模型。

% END DETAILED ANSWER ex38
\subsubsection*{由各点圆周运动到六条整体方程}
% BEGIN ORIGINAL FIGURES
\OriginalFigure[6.2cm]{c3-4-p003-circular-path}{定轴转动中质元的圆周轨迹；括号内为与本体轴瞬时重合的空间轴标记}
% END ORIGINAL FIGURES

固定转轴不是一条随时间寻找的瞬时轴，而是一条在空间中始终不动的直线。
在轴上取固定点 $O$，令本体轴 $x,y$ 随刚体转动，$z$ 沿固定轴。
这里及以下所有 $I_{ij}$ 均表示关于指定原点的实际惯量矩阵元。
对任意质元，写 $\vect r=x\vect e_x+y\vect e_y+z\vect e_z$，则
\[
\vect v=\omega\vect e_z\times\vect r
=-\omega y\vect e_x+\omega x\vect e_y.
\]
材料点的 $x,y,z$ 在本体系中是常数，但本体基矢在惯性空间中变化：
\[
\dot{\vect e}_x=\omega\vect e_y,\qquad
\dot{\vect e}_y=-\omega\vect e_x,\qquad
\dot{\vect e}_z=0.
\]
所以速度的空间导数为
\[
\vect a=(-\alpha y-\omega^2x)\vect e_x
+(\alpha x-\omega^2y)\vect e_y,
\qquad \alpha=\dot\omega.
\]
其中 $\alpha$ 项是沿圆周切线的加速度，$\omega^2$ 项指向转轴。
当 $\omega=0$ 而 $\alpha\ne0$ 时仍有切向加速度；
当 $\alpha=0$ 而 $\omega\ne0$ 时仍有向心加速度。
因而“匀速转动”不能当成合力为零。
将上式按质量平均，直接得到质心的加速度及合力三式。

对轴上固定点 $O$，角动量是
\[
\vect L_O=\omega(I_{xz}\vect e_x+I_{yz}\vect e_y+I_{zz}\vect e_z).
\]
用上述基矢导数逐项求导，而不是只对 $\omega$ 求导，得到
\[
\dot{\vect L}_O=
(I_{xz}\alpha-I_{yz}\omega^2)\vect e_x
+(I_{yz}\alpha+I_{xz}\omega^2)\vect e_y
+I_{zz}\alpha\vect e_z.
\]
这里 $I_{xz}=-\int xz\,\dd m$、$I_{yz}=-\int yz\,\dd m$；
若改用不带这个负号的积分，就必须同时改写横向角动量式的号数。
这一步既说明横向约束矩来自哪里，也说明不能把固定轴默认成惯量主轴。

设主动合力为 $\vect F$，主动主矩关于 $O$ 为 $\vect M_O$，
两个轴承位于 $z=z_A,z_B$，对刚体的力为 $\vect R_A,\vect R_B$。
轴承力矩由叉积直接给出
\[
z_A\vect e_z\times\vect R_A
=-z_AR_{Ay}\vect e_x+z_AR_{Ax}\vect e_y,
\]
$B$ 点同理。因此完整的六条关系为
\begin{align*}
m(-\alpha y_C-\omega^2x_C)&=F_x+R_{Ax}+R_{Bx},\\
m(\alpha x_C-\omega^2y_C)&=F_y+R_{Ay}+R_{By},\\
0&=F_z+R_{Az}+R_{Bz},\\
I_{xz}\alpha-I_{yz}\omega^2&=M_{Ox}-z_AR_{Ay}-z_BR_{By},\\
I_{yz}\alpha+I_{xz}\omega^2&=M_{Oy}+z_AR_{Ax}+z_BR_{Bx},\\
I_{zz}\alpha&=M_{Oz}.
\end{align*}
最后一式以一个有向转角决定运动，前五式体现保持该运动所需的约束。
最后一式没有理想轴承力，是因为轴承力作用点在轴上，其轴向力矩为零；
如果轴承本身施加摩擦力偶，该力偶的轴向分量仍须计入 $M_{Oz}$。
这些惯量全部关于 $O$。将其中的横向惯量换成质心惯量而仍保留对 $O$ 的外矩，
会漏掉质心绕轴运动的轨道角动量。

\subsubsection*{反力联立的计算次序与可确定性}
给定初值 $q(0)=q_i,\omega(0)=\omega_i$，先求
$I_{zz}\ddot q=M_{Oz}$，再代入其余关系。
为便于实际计算，定义
\begin{align*}
S_x&=m(a_C)_x-F_x,&S_y&=m(a_C)_y-F_y,\\
K_x&=I_{xz}\alpha-I_{yz}\omega^2-M_{Ox},&
K_y&=I_{yz}\alpha+I_{xz}\omega^2-M_{Oy}.
\end{align*}
横向四式化为
\[
R_{Ax}+R_{Bx}=S_x,\quad z_AR_{Ax}+z_BR_{Bx}=K_y,
\]
\[
R_{Ay}+R_{By}=S_y,\quad z_AR_{Ay}+z_BR_{By}=-K_x.
\]
令 $d=z_B-z_A\ne0$，消元得到
\begin{align*}
R_{Ax}&=\frac{z_BS_x-K_y}{d},&
R_{Bx}&=\frac{K_y-z_AS_x}{d},\\
R_{Ay}&=\frac{z_BS_y+K_x}{d},&
R_{By}&=\frac{-K_x-z_AS_y}{d}.
\end{align*}
可见合力条件决定两个反力的和，力矩条件决定分配。
若两个支点轴向位置重合，四式不能独立确定两处横向力；
这正是 $d=0$ 的几何退化，而不是代数计算出了偶然的零分母。
轴向仅有 $R_{Az}+R_{Bz}=-F_z$。
例如 $A$ 是止推轴承而 $B$ 轴向自由时，才可加 $R_{Bz}=0$，求出 $R_{Az}$。
支承对刚体的力与刚体对支承的压力方向相反，二者大小相同。

\subsubsection*{功率式如何推出能量式}
把轴向转动方程乘以 $\omega$，利用 $I_{zz}$ 恒定，得
\[
\frac{\dd}{\dd t}\left(\frac12I_{zz}\omega^2\right)
=M_{Oz}\omega.
\]
这一推导在释放瞬间 $\omega=0$ 仍成立，并没有用 $\omega$ 作除数。
若某主动矩来自势能 $V(q)$，则 $M_{Oz}=-\dd V/\dd q$，从而
\[
\frac{\dd}{\dd t}(T+V)=0.
\]
若同时有驱动和耗散，则右边保留其非保守功率。
将 $\dd T=M_{Oz}\dd q$ 对转角积分，得到速度平方与转角的关系；
将 $I_{zz}\ddot q=M_{Oz}(t)$ 对时间积分，得到角速度与时间的关系。
两种积分选择取决于外矩的依赖变量。

对前面的圆盘，令净运动力矩 $M-M_f>0$。
先从动力学得到 $\omega=2(M-M_f)t/(ma^2)$，再积一次得到转角。
也可从做功关系求
\[
\frac14ma^2\omega^2=(M-M_f)\Delta q.
\]
于是马达输入能量 $M\Delta q$ 分成两部分：
圆盘增加 $\tfrac14ma^2\omega^2$ 的动能，接触摩擦消耗 $M_f\Delta q$。
若转速以后过零，运动阻矩方向必须随转向重定；
在静止瞬间，实际静摩擦矩在其允许区间内调整，启动条件由可用驱动矩决定。


\subsection{物理摆与等效单摆长度}\label{sub:c3-6.2}
质量 $m$ 的刚体绕水平固定轴 $O$ 摆动。
质心到轴的垂距为 $h$，令 $q$ 为质心连线相对于竖直向下的夹角。
重力矩为 $-mgh\sin q$，因此
\begin{equation}\label{eq:physicalpend}
I_O\ddot q+mgh\sin q=0.
\end{equation}
能量为
\[
E=\frac12I_O\dot q^2+mgh(1-\cos q).
\]
下垂位置 $q=0$ 是稳定平衡；直立位置 $q=\pi$ 是不稳定平衡。
小振幅近似 $\sin q\simeq q$ 得
\[
\omega_0=\sqrt{\frac{mgh}{I_O}},\qquad
\tau_0=2\pi\sqrt{\frac{I_O}{mgh}},\qquad
L_{\text{等效}}=\frac{I_O}{mh}.
\]
对于小振动初值 $q(0)=q_i,\dot q(0)=u_i$，显式解为
\[
q(t)=q_i\cos\omega_0t+\frac{u_i}{\omega_0}\sin\omega_0t.
\]
若从最大角 $q_0$ 静止释放，能量给出
$\dot q^2=(2mgh/I_O)(\cos q-\cos q_0)$。
在往返振动支路 $0<|q_0|<\pi$ 上，有限振幅的周期为
\[
\tau=4\sqrt{\frac{I_O}{mgh}}\,
K\left(\sin\frac{q_0}{2}\right),
\quad
K(k)=\int_0^{\pi/2}\frac{\dd u}{\sqrt{1-k^2\sin^2u}},
\]
并有 $\tau=\tau_0(1+q_0^2/16+\cdots)$。当 $|q_0|\to\pi$，周期发散；直立且精确静止的初态停留在不稳定平衡，不对应有限周期振动。
小振幅公式忽略的是非线性振幅效应，不是转动惯量。

\TeachingFigure{c3-diagram-30}{物理摆的重力矩由质心到轴的距离决定}
% BEGIN ORIGINAL FIGURES
\OriginalFigure[6.2cm]{c3-4-p008-physical-pendulum}{复摆的悬点、质心位置与重力；图中 $l,\theta$ 分别对应正文的 $h,q$}
% END ORIGINAL FIGURES


\example{均匀杆端点物理摆}\label{ex:c3-39}

\QuestionStatement{均匀细杆长 $L$，绕端点理想水平轴在竖直平面摆动，初始偏角 $q_0$、初始静止。
取偏离竖直向下的有向角为 $q$，杆质量为 $m$，$0<|q_0|<\pi$，忽略轴摩擦。求非线性运动方程、小振动周期与初值解、最低点速率及轴反力的求法。}{ex39}{例39题图}
% BEGIN DETAILED ANSWER ex39
\solution \Step{由重力矩建立非线性方程}
以图~\ref{fig:solution-ex39}~中杆中点重力与铰轴 $O$ 的力臂为依据，先写出图右侧第一行的几何或动力学关系，再代入其下方的计算式。

端点惯量 $I_O=mL^2/3$，质心距轴 $h=L/2$。以竖直向下为 $q=0$，重力的回复矩为 $-mgh\sin q$，因此
\[
\frac{mL^2}{3}\ddot q=-mg\frac L2\sin q,\qquad
\ddot q+\frac{3g}{2L}\sin q=0.
\]
\SolutionFigure{ex39}{例39的计算示意}
\Step{在线性近似中应用初值}
当 $|q_0|\ll1$，以弧度计有 $\sin q\simeq q$。令
$\omega_0=\sqrt{3g/(2L)}$，通解为 $q=A\cos\omega_0t+B\sin\omega_0t$。
由 $q(0)=q_0,\dot q(0)=0$ 得
\[
q(t)\simeq q_0\cos\omega_0t,\qquad
\tau_0=2\pi\sqrt{\frac{2L}{3g}}.
\]
对应的等效单摆长为 $2L/3$。小振动频率 $\omega_0$ 不是摆杆某瞬间的真实角速度 $\dot q$。
\Step{用精确能量求最低点速率}
以最低点势能为零，
\[
\frac12I_O\dot q^2+mgh(1-\cos q)=mgh(1-\cos q_0).
\]
所以 $\dot q^2=(3g/L)(\cos q-\cos q_0)$，到 $q=0$ 时
$|\dot q|=\sqrt{3g(1-\cos q_0)/L}$。初始 $q_0>0$ 的第一次回落取负根，$q_0<0$ 则取正根。
\Step{说明支反力的独立求法}
取径向单位矢量从轴指向质心、切向沿 $q$ 增大方向，
$\vect a_C=h\ddot q\vect e_t-h\dot q^2\vect e_r$。
由 $\vect R_O-mg\vect j=m\vect a_C$ 可求轴反力。
例如最低点径向向下、切向加速度为零，故 $R_y=mg+mh\dot q^2$，显然一般大于静重。

% END DETAILED ANSWER ex39
\example{悬点位置如何影响小振动周期}\label{ex:c3-40}

\QuestionStatement{刚体中心轴惯量为 $I_C=mk^2$，
平行悬轴距质心 $h>0$。求使小振动周期最短的位置。
悬轴方向保持不变，允许在刚体或刚性延长支架上选择悬点。求周期关于 $h$ 的函数、极小值及两端极限。}{ex40}{例40题图}
% BEGIN DETAILED ANSWER ex40
\solution \Step{写出依赖悬点的惯量与回复系数}
以图~\ref{fig:solution-ex40}~中悬轴 $O$ 与质心 $C$ 的距离 $h$为依据，先写出图右侧第一行的几何或动力学关系，再代入其下方的计算式。

悬轴方向不变时，中心轴惯量 $I_C=mk^2$ 为常量；距质心 $h$ 的平行悬轴有
$I_O=m(k^2+h^2)$。
对质心位于悬轴下方的小摆角 $q$，方程为
$m(k^2+h^2)\ddot q+mghq=0$，于是
\[
\tau_0(h)=2\pi\sqrt{\frac{k^2+h^2}{gh}},\qquad
\tau_0^2=\frac{4\pi^2}{g}\left(h+\frac{k^2}{h}\right).
\]
\SolutionFigure{ex40}{例40的计算示意}
\Step{驻点与最小值}
因周期为正，极小化 $\tau_0$ 等价于极小化
$F(h)=h+k^2/h$。有
\[
F'(h)=1-\frac{k^2}{h^2}=0\ \Longrightarrow\ h=k,\qquad
F''(h)=\frac{2k^2}{h^3}>0.
\]
又由 $h+k^2/h\ge2k$，可知它还是 $h>0$ 上的全局最小值：
\[
\tau_{\min}=2\pi\sqrt{\frac{2k}{g}},\qquad
\ell_{\text{等效,min}}=2k.
\]
\Step{边界与安装范围}
$h\to0^+$ 时回复矩趋零而 $I_C$ 非零，周期发散；$h\to\infty$ 时惯量随 $h^2$ 增长快于回复系数，周期也发散。
所以两种效应的竞争产生有限最优距离。
本题允许延长支架；若实际悬点只能在一个有限区间内选择，且该区间不包含 $k$，则应比较区间端点，而不能给出不可实现的悬点。

% END DETAILED ANSWER ex40
\subsubsection*{重力矩、回转半径与初值的联系}
令 $\vect i$ 水平向右、$\vect j$ 竖直向上，$\vect k$ 指向纸外。
从向下的竖线逆时针量有向角 $q$，定义
\[
\vect e_r=\sin q\,\vect i-\cos q\,\vect j,
\qquad
\vect e_t=\cos q\,\vect i+\sin q\,\vect j.
\]
于是 $\vect r_{OC}=h\vect e_r$，$\dot{\vect e}_r=\dot q\vect e_t$。
重力关于 $O$ 的矩为
\[
(\vect r_{OC}\times(-mg\vect j))_z=-mgh\sin q.
\]
这给出回复矩的号数：$q>0$ 时角加速度为负，$q<0$ 时角加速度为正。
用 $I_C=mk^2$ 和平行轴定理，有
\[
I_O=m(k^2+h^2),\qquad
\ddot q=-\frac{gh}{k^2+h^2}\sin q.
\]
同一质心距离 $h$ 下，回转半径越大，质量越难绕悬轴改变角速度。
小振幅近似只把 $\sin q$ 换成 $q$，惯量仍是完整的 $I_O$。

对于线性初值，写 $q=A\cos(\omega_0t-\delta)$，比较 $t=0$ 的位移和速度，得
\[
A=\sqrt{q_i^2+\frac{u_i^2}{\omega_0^2}},\qquad
A\cos\delta=q_i,\quad A\sin\delta=\frac{u_i}{\omega_0}.
\]
相位应由这两个有向条件共同决定，单独求反正切会丢失象限。
使用显式初值解可直接绕开这个问题。
如果 $h=0$，回复矩为零，方程变成 $I_O\ddot q=0$；
这时是自由匀速转动或静止，不能称为具有无穷周期的实际往返振动。

\subsubsection*{非线性积分、振动支路与时间参数}
将摆方程乘以 $\dot q$，积分得到
\[
\frac12I_O\dot q^2+mgh(1-\cos q)
=\frac12I_Ou_i^2+mgh(1-\cos q_i)=E.
\]
能量只给 $\dot q^2$，不能独自给速度方向。
一般初态的运动支路由 $E$ 与 $u_i$ 的方向共同选取：
$0<E<2mgh$ 时有两个转向点，形成往返振动；
$E>2mgh$ 时速度不为零，可连续越过直立位置；
$E=2mgh$ 是两者的分界，接近直立位置所需时间无限增加。
在静止下垂状态 $E=0$，平衡位置不随时间变化。

考虑 $q_0>0$ 的静止释放振动，先研究从 $q_0$ 回落到 $0$ 的四分之一周期。
此段应取
\[
\dot q=-\sqrt{2}\,\omega_0\sqrt{\cos q-\cos q_0},\qquad
\omega_0^2=\frac{mgh}{I_O}.
\]
因此到任意位置 $q$ 的时间为
\[
t=\frac1{\sqrt2\,\omega_0}
\int_q^{q_0}\frac{\dd s}{\sqrt{\cos s-\cos q_0}}.
\]
端点的被积函数趋于无穷，但当 $q_0<\pi$ 时是可积的平方根奇点。
令 $k=\sin(q_0/2)$、$\sin(s/2)=k\sin u$，则
\[
\cos s-\cos q_0=2k^2\cos^2u,\qquad
\dd s=\frac{2k\cos u}{\sqrt{1-k^2\sin^2u}}\dd u.
\]
代入后，两个 $k\cos u$ 因子相消，得
\[
t=\frac1{\omega_0}\left[K(k)-F(u,k)\right],
\quad F(u,k)=\int_0^u\frac{\dd\eta}{\sqrt{1-k^2\sin^2\eta}},
\quad \sin\frac q2=k\sin u.
\]
令 $q=0$ 即得到四分之一周期 $K(k)/\omega_0$。
借助运动的左右与时间反演对称性，四段相接成为完整周期。
若 $q_0<0$，用 $|q_0|$ 求周期，再把位移和速度整体镜像。
一般初值先由能量求振幅，再选与初速度同向的积分支路确定相位。

展开被积函数
\[
(1-k^2\sin^2u)^{-1/2}
=1+\frac12k^2\sin^2u+\frac38k^4\sin^4u+\cdots,
\]
得到
\[
K(k)=\frac\pi2\left(1+\frac{k^2}{4}+\frac{9k^4}{64}+\cdots\right),
\quad
\tau=\tau_0\left(1+\frac{q_0^2}{16}
+\frac{11q_0^4}{3072}+\cdots\right).
\]
有限振幅越大，周期越长。$k\to1$ 时积分在 $u=\pi/2$ 附近发散。
在分界能量上，一条通过下垂位置的有向轨道可写成
\[
q(t)=4\arctan\exp[\omega_0(t-t_c)]-\pi.
\]
它在 $t=t_c$ 通过 $q=0$，两端只在无限时间趋向直立姿态。
这条轨道与 $q=\pi,\dot q=0$ 的精确静止解是不同的初值解。

\subsubsection*{物理摆的轴反力必须由质心加速度求}
质心沿半径 $h$ 的圆周运动，故
\[
\vect a_C=-h\dot q^2\vect e_r+h\ddot q\vect e_t.
\]
设悬轴对刚体的平面合反力为 $\vect R$，由
$\vect R-mg\vect j=m\vect a_C$ 得
\begin{align*}
R_x&=mh(\ddot q\cos q-\dot q^2\sin q),\\
R_y&=mg+mh(\ddot q\sin q+\dot q^2\cos q).
\end{align*}
也可沿瞬时径向和切向投影：
\[
R_r=-m(h\dot q^2+g\cos q),\qquad
R_t=m(h\ddot q+g\sin q).
\]
其中负 $R_r$ 表示指向悬轴。
把摆方程代入切向式，可进一步写成
\[
R_t=mg\sin q\left(1-\frac{mh^2}{I_O}\right)
=mg\sin q\frac{I_C}{I_O}.
\]
点质量单摆的 $I_C=0$，反力沿悬线；
一般物理摆的 $I_C\ne0$，铰反力还可能有切向分量。
有回复矩不意味着悬轴反力恒等于重力，也不意味着反力始终沿 $OC$。

对均匀端点杆，$mh^2/I_O=3/4$。
静止释放的瞬间，
\[
R_x(0)=-\frac34mg\sin q_0\cos q_0,\qquad
R_y(0)=mg\left(1-\frac34\sin^2q_0\right).
\]
最低点的 $\ddot q=0$ 而向心加速度最大，故
\[
R_x=0,\qquad
R_y=mg+\frac{mL}{2}\dot q^2
=mg\left(\frac52-\frac32\cos q_0\right).
\]
这是到最低点的一瞬间、碰到任何障碍之前的反力。
如果悬轴由两个轴承实现，上式只给它们在摆动平面内的合反力；
单独分配到两支点还需轴向位置和横向力矩条件。

\subsubsection*{等效长度的优化与悬点互换}
由 $L_{\text{等效}}=h+k^2/h$，求导得到
\[
\frac{\dd L_{\text{等效}}}{\dd h}=1-\frac{k^2}{h^2},\qquad
\frac{\dd^2 L_{\text{等效}}}{\dd h^2}=\frac{2k^2}{h^3}>0.
\]
因此 $h=k$ 是唯一最短周期位置，但刚体上实际允许的悬点必须能达到这个距离。
若结构限制 $h$ 只能取某个区间，应在该区间比较端点与 $h=k$，不能忽略几何可实现性。
此外，取 $h'=k^2/h$，有
\[
h'+\frac{k^2}{h'}=h+\frac{k^2}{h}.
\]
所以位于质心两侧、距离分别为 $h,h'$ 的平行悬轴有相同的小振动周期，
两轴间距为 $h+h'=L_{\text{等效}}$。
这里比较的是小振动周期；大振幅还须指定相同的最大摆角。


\subsection{轴承附加压力的两个来源}\label{sub:c3-6.3}
固定轴有两个轴承 $A,B$，坐标分别为 $z=0,L$。
质心坐标为 $(x_C,y_C,z_C)$。
把某瞬间的本体横轴选作 $x,y$，角速度 $\omega\vect e_z$、
角加速度 $\alpha\vect e_z$，则
\[
(a_C)_x=-\omega^2x_C-\alpha y_C,\qquad
(a_C)_y=-\omega^2y_C+\alpha x_C.
\]
即使转速恒定，偏心质心也有向心加速度，需额外横向合力。
轴承与主动外力的合力必须满足 $m\vect a_C$。

另一方面，
$\vect L=(I_{xz}\omega,I_{yz}\omega,I_{zz}\omega)$，
其惯性系导数为
\[
(\dot L)_x=I_{xz}\alpha-I_{yz}\omega^2,\quad
(\dot L)_y=I_{yz}\alpha+I_{xz}\omega^2,\quad
(\dot L)_z=I_{zz}\alpha.
\]
非零惯性积造成额外横向力矩。
这是与质心偏心不同的第二个来源。

设两轴承对转子的力为 $\vect R_A,\vect R_B$。
把维持相同姿态的瞬时静态反力扣除，以 $\Delta\vect R$ 表示运动附加项，
且主动载荷不随转速变化。由上轴承力矩
$L\vect e_z\times\Delta\vect R_B$ 得
\begin{align}
\Delta R_{Bx}&=\frac{I_{yz}\alpha+I_{xz}\omega^2}{L},&
\Delta R_{By}&=\frac{-I_{xz}\alpha+I_{yz}\omega^2}{L},\\
\Delta R_{Ax}&=m(-\omega^2x_C-\alpha y_C)-\Delta R_{Bx},&
\Delta R_{Ay}&=m(-\omega^2y_C+\alpha x_C)-\Delta R_{By}.
\end{align}
这里 $I_{xz},I_{yz}$ 是带符号的实际矩阵元，
与取正的惯性积之间差一个负号。
若求轴承受到转子的压力，须把上述力整体反向。
仅由整体合力无法分配两个轴承的载荷，必须同时使用横向力矩方程。
定轴运动还有 $(a_C)_z=0$。
若主动外力的轴向分量为 $F_z$，则
\[
F_z+R_{Az}+R_{Bz}=0.
\]
这个方程只决定两轴承轴向反力之和。
它们各自承担多少载荷，要由哪一处设置止推轴承、轴向间隙及其他支承条件确定。

\TeachingFigure{c3-diagram-31}{偏心合力与惯性积力偶共同决定两个轴承的附加反力}
% BEGIN ORIGINAL FIGURES
\OriginalFigure[6.2cm]{c3-4-p013-bearing-forces}{两轴承约束力与主动外力；图中 $\vect N_A,\vect N_B$ 对应正文的 $\vect R_A,\vect R_B$，正文将原点移到 $A$}
% END ORIGINAL FIGURES



\subsubsection*{扣除静态载荷时究竟扣除了什么}
为求运动附加反力，先保持某一瞬间的姿态不变，
用相同的主动合力和横向主矩求其静态横向支承反力。
再从运动方程扣去这些静态关系。
这是一种反力分解方法，并没有把运动中的刚体当作静止。
例如水平轴上的转子受到固定竖直重力，静态反力是在空间中向上的固定矢量；
转子的本体 $x$ 轴随时间转动，所以这个静态反力的本体系分量一般随相位变化。
在某一相位取 $x$ 向上，不允许在整个周期仍把 $x$ 当作空间竖直方向。

将 $A$ 取作原点，$B$ 位于 $L\vect e_z$。
静态与动态方程相减后，横向合力为
\[
\Delta\vect R_A+\Delta\vect R_B
=m\alpha\vect e_z\times\vect e_C-m\omega^2\vect e_C,
\qquad
\vect e_C=x_C\vect e_x+y_C\vect e_y.
\]
附加横向矩为
\[
L\vect e_z\times\Delta\vect R_B
=(I_{xz}\alpha-I_{yz}\omega^2)\vect e_x
+(I_{yz}\alpha+I_{xz}\omega^2)\vect e_y.
\]
逐分量解出 $B$ 的附加力，再由合力求 $A$，得到前述四式。
第一式依赖偏心矢量，第二式依赖惯性积；
所以消除其中一个来源，不能自动消除另一个来源。

\subsubsection*{由质量分布直接看见合力与力偶}
% BEGIN ORIGINAL FIGURES
\OriginalFigure[6.2cm]{c3-4-p024-centrifugal-couple}{随体旋转参考系中的离心力偶；红箭头表示 $A$ 端附加反力，$B$ 端须有等大反向反力}
% END ORIGINAL FIGURES

定义两个本体系横向矢量
\[
\vect U=\int\vect e\,\dd m=m\vect e_C,
\qquad
\vect V=\int z\vect e\,\dd m
=-I_{xz}\vect e_x-I_{yz}\vect e_y,
\]
其中 $\vect e=x\vect e_x+y\vect e_y$，$z$ 从 $A$ 量起。
$\vect U$ 的量纲为质量乘长度，$\vect V$ 的量纲为质量乘长度平方。
质量元的横向加速度是
$\alpha\vect e_z\times\vect e-\omega^2\vect e$，
按质量积分产生附加合力。
在不同轴向位置再乘以 $z$，其一阶矩决定两个支点分配。
于是反力可写成更直观的矢量形式
\begin{align*}
\Delta\vect R_B&=\frac1L
\left(\alpha\vect e_z\times\vect V-\omega^2\vect V\right),\\
\Delta\vect R_A&=
\alpha\vect e_z\times\left(\vect U-\frac{\vect V}{L}\right)
-\omega^2\left(\vect U-\frac{\vect V}{L}\right).
\end{align*}
若只有一个偏心小质量 $m_0$，位于 $z=z_0$、横向位置 $\vect e$，
恒速时上式变为
\[
\Delta\vect R_A=-m_0\omega^2\left(1-\frac{z_0}{L}\right)\vect e,
\qquad
\Delta\vect R_B=-m_0\omega^2\frac{z_0}{L}\vect e.
\]
当 $0<z_0<L$，两力都指向轴，按两个支点到质量所在截面的距离分配。
当质量位于支承区间外，一个系数可为负，封闭支承需要提供反向力。
不能仅凭两支点反力的图形方向猜测所有结构都应同向。

若 $\vect U=0$ 而 $\vect V\ne0$，则
\[
\Delta\vect R_A=-\Delta\vect R_B.
\]
这是一对纯力偶，质心没有横向加速度，轴承仍可承受明显的周期载荷。
把这种力偶误画成两个同向力，会违反整体合力条件；
把偏心质心所需的向心合力误画成反向等长力，则会违反质心运动条件。

\subsubsection*{固定空间基下的整周期反力}
设恒速转动时 $q=q_0+\omega t$，固定横向基为 $\vect e_X,\vect e_Y$，
二维转动矩阵记作
\[
\mathsf Q(q)=
\begin{pmatrix}\cos q&-\sin q\\\sin q&\cos q\end{pmatrix}.
\]
取 $q_0=0$，将 $\vect U_0,\vect V_0$ 表示为初始固定基的分量，得
\[
\Delta\vect R_A(t)=-\omega^2\mathsf Q(\omega t)
\left(\vect U_0-\frac{\vect V_0}{L}\right),\qquad
\Delta\vect R_B(t)=-\frac{\omega^2}{L}\mathsf Q(\omega t)\vect V_0.
\]
两个附加力各自大小恒定，方向在横向平面内旋转。
如果一个支点的静态力为固定矢量 $\vect G$，附加力大小为 $D\ge0$，
令其与 $\vect G$ 的夹角为 $\chi(t)$，总反力满足
\[
R^2=G^2+D^2+2GD\cos\chi(t).
\]
经过一个完整周期，$\cos\chi$ 遍历 $[-1,1]$，所以
\[
R_{\max}=G+D,\qquad R_{\min}=|G-D|.
\]
最大值是两个矢量同向时的长度，最小值是反向时的长度。
因此不能把固定方向分量的最大值与总压力的最大值混淆。
两个支点附加力相反时，各自峰值通常出现在相差半周期的时刻。
对某个支点，反力矢量的端点在横向力平面中走一个圆，圆心为静态力的端点。

\TeachingFigure{c3-diagram-32}{恒速附加反力的方向旋转，总压力由它与固定静态力的合成决定}
% BEGIN ORIGINAL FIGURES
\OriginalFigure[6.2cm]{c3-4-p019-tilted-disk-3d}{倾斜圆盘及轴承反力的空间示意；$N$ 的分量对应正文 $R$ 的分量，右图为动画的一个时刻}
% END ORIGINAL FIGURES


非匀速时仍可使用前面的瞬时矢量式，但 $\alpha$ 也参与反力。
因为 $\vect e_z\times\vect V$ 与 $\vect V$ 正交，
\[
|\Delta\vect R_B|=\frac{|\vect V|}{L}\sqrt{\omega^4+\alpha^2}.
\]
同理，$A$ 的瞬时附加力大小是
$|\vect U-\vect V/L|\sqrt{\omega^4+\alpha^2}$。
这个结果展示了加速阶段的切向载荷和恒速阶段的向心载荷如何合成。
若静态载荷或转速也随时间改变，整段最大压力须在实际时间区间内求极值，
不能再直接套恒定 $G,D$ 的整周期结论。


\subsection{静平衡、动平衡与校正质量}\label{sub:c3-6.4}
若质心在轴上，$x_C=y_C=0$，称静平衡；
它消除了恒速转动时因质心偏心产生的附加合力。
若固定轴同时是关于质心的惯量主轴，
$I_{xz}=I_{yz}=0$，则称动平衡，
附加合力和附加力偶都消失。
静平衡并不自动保证动平衡。

\example{质心在轴上仍有轴承压力}\label{ex:c3-41}

\QuestionStatement{均匀细杆长 $l$、质量 $m$，中心位于固定竖直轴上，
杆与轴成角 $\beta$，并随轴恒速 $\omega$ 转动。
杆在 $xz$ 平面，轴承间距 $L$，忽略其他横向主动载荷。
取两轴承沿 $+z$ 由 $A$ 至 $B$，间距 $L$；$\vect e_x$ 随杆转动。求两轴承的附加横向反力，解释静平衡是否足够。}{ex41}{例41题图}
% BEGIN DETAILED ANSWER ex41
\solution \Step{由杆方向写出中心惯量张量}
以图~\ref{fig:solution-ex41}~中杆倾角 $\beta$ 与沿轴分开的两个轴承为依据，先写出图右侧第一行的几何或动力学关系，再代入其下方的计算式。

随杆转动的正交基中，$\vect n=(\sin\beta,0,\cos\beta)$，令 $J=ml^2/12$。细杆的中心张量为
\[
\mathsf I_C=J(\mathsf1-\vect n\vect n^{\mathsf T}),\quad
I_{xz}=-J\sin\beta\cos\beta,\quad I_{yz}=0.
\]
负号来自张量交叉元的定义 $I_{xz}=-\int xz\dd m$。
\SolutionFigure{ex41}{例41的计算示意}
\Step{用角动量的空间导数求所需力偶}
总角速度为 $\omega\vect e_z$，故
$\vect L_C=\omega(I_{xz}\vect e_x+I_{zz}\vect e_z)$。
分量恒定，但 $\dot{\vect e}_x=\omega\vect e_y$，所以
\[
\vect M_C=\frac{\dd\vect L_C}{\dd t}=I_{xz}\omega^2\vect e_y.
\]
质心不作横向运动，附加反力之和必须为零。
\Step{用轴承间距把力偶还原为两力}
设 $\Delta\vect R_B=R\vect e_x$，则 $\Delta\vect R_A=-R\vect e_x$。关于任一点，这对力的矩都是
$L\vect e_z\times(R\vect e_x)=LR\vect e_y$。
与所需力偶比较，有
\[
\Delta R_{Bx}=\frac{I_{xz}\omega^2}{L}
=-\frac{ml^2\omega^2}{12L}\sin\beta\cos\beta,\qquad
\Delta R_{Ax}=-\Delta R_{Bx}.
\]
其空间方向随轴一起转动。这里求的是附加横向反力，竖直静载分配需轴向支承资料。
$\beta=0$ 或 $\pi/2$ 时交叉元为零；$\omega\to0$ 时附加载荷按平方趋零。质心在轴上只消除偏心合力，并不自动消除这个力偶。

% END DETAILED ANSWER ex41
\example{点质量转子的单平面和双平面校正}\label{ex:c3-42}

\QuestionStatement{无质量轴上固定若干小质量 $m_i$，
横向位置矢量为 $\vect e_i$，轴向坐标为 $z_i$。
求恒速动平衡条件，说明为什么通常需要两个校正平面。
进一步给出两个不同校正平面 $z_a,z_b$ 上所需的质量半径矢量及其安装方位。}{ex42}{例42题图}
% BEGIN DETAILED ANSWER ex42
\solution \Step{区分零合力与零力偶条件}
以图~\ref{fig:solution-ex42}~中不同轴向坐标的两个校正平面为依据，先写出图右侧第一行的几何或动力学关系，再代入其下方的计算式。

在恒速转动中，每个小质量需要横向向心力 $-m_i\omega^2\vect e_i$。零附加合力要求
\[
\vect U=\sum_i m_i\vect e_i=0.
\]
对横截面原点取矩，轴向位置使这些力额外形成力偶，其消失要求
\[
\vect V=\sum_i m_i z_i\vect e_i=0.
\]
这两个横向矢量条件共四个标量。第一条件等价于质心在轴上；两者一起才使该轴成为中心主轴。
\SolutionFigure{ex42}{例42的计算示意}
\Step{建立两个校正平面的线性方程}
把校正质量与半径的乘积矢量记为 $\vect U_a,\vect U_b$，则
\[
\vect U_a+\vect U_b=-\vect U,\qquad
z_a\vect U_a+z_b\vect U_b=-\vect V.
\]
第二式减第一式的 $z_b$ 倍，得到
$(z_a-z_b)\vect U_a=z_b\vect U-\vect V$，因此
\[
\vect U_a=\frac{z_b\vect U-\vect V}{z_a-z_b},\qquad
\vect U_b=\frac{\vect V-z_a\vect U}{z_a-z_b}.
\]
\Step{回代、安装与单平面特例}
两式相加恰为 $-\vect U$，分别乘 $z_a,z_b$ 再相加恰为 $-\vect V$，故消元结果满足原条件。
选定安装半径 $r_a$ 时，校正质量 $m_a=|\vect U_a|/r_a$，方位由 $\vect U_a$ 的方向确定，另一平面同理。
单平面 $z_c$ 只能给出 $\vect V_c=z_c\vect U_c$，仅当原不平衡满足 $\vect V=z_c\vect U$ 才可一次消除。
$z_a\to z_b$ 时一般需要很大的校正量，反映两个平面趋于重合后失去独立调节力偶的能力。

% END DETAILED ANSWER ex42
\TeachingFigure{c3-diagram-33}{两个不同轴向位置的校正量可以分别调整合力与力偶}


\subsubsection*{静平衡与动平衡是两组独立质量矩条件}
静平衡要求 $\vect U=0$，等价于质心在轴上；
动平衡还要求 $\vect V=0$，等价于转轴的两项横向惯性积为零。
这两组条件一共含四个横向标量条件。
它们消除的是运动附加反力，重力、皮带拉力或其他主动载荷的静态支承力仍存在。
在某一瞬间恰好 $\omega=\alpha=0$ 时，任何质量分布都没有运动附加项，
这只能说明该瞬间的状态，不能说明转子结构已经平衡。

若改变轴上原点，从旧原点移到 $z=z_*$，则
\[
\vect V'=\sum_i m_i(z_i-z_*)\vect e_i
=\vect V-z_*\vect U.
\]
因此对于已经静平衡的转子，$\vect V'=\vect V$，
偶不平衡条件不依赖选择哪个轴上点作原点。
这也解释了为什么“质心在轴上，且该轴为质心主轴”是与原点无关的几何表述。

前面的斜杆给出一个可逐项核查的例子。
沿杆方向取参数 $s\in[-l/2,l/2]$，有
\[
x=s\sin\beta,\quad z=s\cos\beta,\quad
\dd m=\frac ml\dd s.
\]
因奇函数积分，$\int x\,\dd m=0$；但
\[
\int xz\,\dd m=
\frac ml\sin\beta\cos\beta\int_{-l/2}^{l/2}s^2\dd s
=\frac{ml^2}{12}\sin\beta\cos\beta.
\]
故实际 $I_{xz}$ 是它的负值。
杆中点在轴上仍有周期横向矩，轴承以反向等长力承担它。
杆平行轴或垂直轴时，上述积分为零，这两种极限都应消除该附加力偶。
由此可以同时核查惯性积号数和两支点力的方向。

\subsubsection*{单平面校正何时足够}
在 $z=z_c$ 的一个校正平面添加质量半径矢量 $\vect U_c$，其贡献满足
\[
\Delta\vect U=\vect U_c,\qquad
\Delta\vect V=z_c\vect U_c.
\]
要同时抵消原转子的 $\vect U,\vect V$，必须满足
\[
\vect U_c=-\vect U,\qquad \vect V=z_c\vect U.
\]
一般的两个二维矢量并不满足这种共线比例关系，所以单平面校正通常不够。
如果所有不平衡小质量原本就在同一个轴向平面，确有
$\vect V=z_c\vect U$，则在该平面加一个校正量即可。
若 $\vect U=0$ 而 $\vect V\ne0$，添加任何单平面校正量都会重新产生非零合力，
因而无法单独消除这个纯偶不平衡。

\subsubsection*{双平面校正的消元、安装与回代}
两校正平面坐标为 $z_a\ne z_b$。从
\[
\vect U_a+\vect U_b=-\vect U,\qquad
z_a\vect U_a+z_b\vect U_b=-\vect V
\]
出发，将第一式乘 $z_b$ 后与第二式相减，得到
\[
(z_a-z_b)\vect U_a=z_b\vect U-\vect V.
\]
再求 $\vect U_b=-\vect U-\vect U_a$，即得两项校正公式。
公式中的每一个分量都可独立消元，但最终安装必须使用矢量的方向与大小。
若允许安装的半径为 $r_a,r_b$，则
\[
m_a=\frac{|\vect U_a|}{r_a},\qquad
m_b=\frac{|\vect U_b|}{r_b}.
\]
安装方位分别与 $\vect U_a,\vect U_b$ 同向。
如果采用去除质量的工艺，等价的去除位置应在所需添加矢量的相反方位，
因为去除质量贡献的是负质量半径矢量。

以一组计算数据说明分量到实物的换算。
取固定的初始横向基 $\vect i,\vect j$，原不平衡为
\[
\vect U=(0.003\vect i+0.004\vect j)\,\mathrm{kg\,m},\qquad
\vect V=(0.0006\vect i-0.0003\vect j)\,\mathrm{kg\,m^2}.
\]
选择 $z_a=0,z_b=0.30\,\mathrm m$，则
\[
\vect U_b=-\frac{\vect V}{0.30}
=(-0.002\vect i+0.001\vect j)\,\mathrm{kg\,m},
\]
\[
\vect U_a=-\vect U-\vect U_b
=(-0.001\vect i-0.005\vect j)\,\mathrm{kg\,m}.
\]
在两处均取 $0.10\,\mathrm m$ 的安装半径时，质量分别为
\[
m_a\simeq0.05099\,\mathrm{kg},\qquad
m_b\simeq0.02236\,\mathrm{kg}.
\]
$\vect U_a$ 位于第三象限，$\vect U_b$ 位于第二象限；
不能只由分量比的反正切把它们装到相反方向。
回代可见 $\vect U+\vect U_a+\vect U_b=0$、
$\vect V+0\vect U_a+0.30\vect U_b=0$，两种不平衡同时消除。

当 $|z_a-z_b|$ 很小时，为形成同样的校正力偶，需要较大的质量半径矢量。
这正是分母中校正平面间距的物理含义。
若题目限制添加质量不能为负、半径有上限、安装方位有禁区，
首先求矢量目标，再检查几何与质量限制；必要时把一个矢量分解为多个允许方位的质量。
这些限制影响校正的实现，不改变整体合力和力偶应为零的基本条件。


\subsection{定轴算例与支反力的独立求解}\label{sub:c3-6.5}
\example{立方体绕地面水平棱倾倒}\label{ex:c3-43}

\QuestionStatement{均匀实心立方体边长 $a$、质量 $m$，
一条水平棱由理想铰固定在地面，其余接触可脱离。
底面抬起角 $\theta_0$ 后从静止释放，求回落时的角速度与角加速度。
先取 $0<\theta_0<\pi/4$，研究触地之前。
取底面相对地面的抬角 $\theta$ 为正；给出任意回落位置与即将触地时的有向角速度、角加速度。}{ex43}{例43题图}
% BEGIN DETAILED ANSWER ex43
\solution \Step{建立质心位置和固定棱惯量}
以图~\ref{fig:solution-ex43}~中抬起底面后质心 $C$ 的高度为依据，先写出图右侧第一行的几何或动力学关系，再代入其下方的计算式。

在垂直铰棱的截面内取 $O$ 为原点，底面抬角为 $\theta$。质心坐标为
\[
x_C=\frac a2(\cos\theta-\sin\theta),\quad
y_C=\frac a2(\sin\theta+\cos\theta).
\]
中心平行轴惯量是 $ma^2/6$，质心到棱的距离为 $a/\sqrt2$，所以
$I_O=ma^2/6+m(a/\sqrt2)^2=2ma^2/3$。
\SolutionFigure{ex43}{例43的计算示意}
\Step{用能量和初值求回落角速度}
铰反力对固定棱不做功，由 $\dot\theta(0)=0$，
\[
\frac12I_O\dot\theta^2+mg y_C(\theta)=mg y_C(\theta_0).
\]
因 $0<\theta<\theta_0<\pi/4$ 为回落支路，取负根：
\[
\dot\theta=-\sqrt{\frac{3g}{2a}
[\sin\theta_0+\cos\theta_0-\sin\theta-\cos\theta]}.
\]
令 $\theta=0$，得到即将触地的速度
$-\sqrt{3g(\sin\theta_0+\cos\theta_0-1)/(2a)}$。
\Step{用有向重力矩求角加速度}
重力矩为 $-mgx_C$，因此
\[
\ddot\theta=-\frac{3g}{4a}(\cos\theta-\sin\theta).
\]
在给定支路它为负，与回落方向一致；触地前为 $-3g/(4a)$。
当 $\theta_0\to0$ 时触地前速率趋零；$\theta_0=\pi/4$ 的精确静止初态位于不稳定平衡，理想方程不会自行选边。
碰撞后的角速度必须另给冲量或恢复系数模型，不能延用碰撞前的能量积分。

% END DETAILED ANSWER ex43
\TeachingFigure{c3-diagram-34}{立方体倾倒的质心高度随姿态变化}
% BEGIN ORIGINAL FIGURES
\OriginalFigure[6.2cm]{c3-4-p009-tipping-cube}{立方体绕地面棱转动的截面；图中 $\alpha$ 对应正文倾角 $\theta$，释放时为 $\theta_0$}
% END ORIGINAL FIGURES


\example{有惯量定滑轮的两重物系统}\label{ex:c3-44}

\QuestionStatement{质量为 $m$ 的均匀圆盘滑轮半径 $r$，
轴固定且光滑，轻绳不可伸长并与轮缘无滑动，
两端悬挂质量 $m$ 与 $2m$ 的物块。由静止释放，求加速度、张力和速度。
另求重物下降高度 $h$ 后的速率，并用全系统能量验证。}{ex44}{例44题图}
% BEGIN DETAILED ANSWER ex44
\solution \Step{分别隔离两物块与滑轮}
以图~\ref{fig:solution-ex44}~中左右两段绳上的不同张力为依据，先写出图右侧第一行的几何或动力学关系，再代入其下方的计算式。

轻物 $m$ 向上、重物 $2m$ 向下取正，共同加速度大小为 $a$，两张力记为 $T_1,T_2$。对两个物块：
\[
T_1-mg=ma,\qquad 2mg-T_2=2ma.
\]
绳与轮无滑动，轮的顺时针角加速度为 $a/r$。滑轮惯量 $I=mr^2/2$，因此
\[
(T_2-T_1)r=I\frac ar\quad\Longrightarrow\quad T_2-T_1=\frac12ma.
\]
\SolutionFigure{ex44}{例44的计算示意}
\Step{消元并回代求张力}
由前两式 $T_1=m(g+a)$、$T_2=2m(g-a)$，相减代入第三式：
$mg-3ma=ma/2$，故
\[
a=\frac{2g}{7},\qquad
T_1=\frac97mg,\quad T_2=\frac{10}{7}mg,\quad
\alpha_{\text{顺时针}}=\frac{2g}{7r}.
\]
两张力皆正，且 $mg<T_1<T_2<2mg$，与两物块的加速方向相容。由静止释放后 $v(t)=2gt/7$。
\Step{高度与速度的关系}
重物下降 $h$，轻物上升同样的 $h$，总重力势能减少 $mgh$；滑轮也有转动能：
\[
mgh=\frac12mv^2+\frac12(2m)v^2+\frac12\left(\frac12mr^2\right)\frac{v^2}{r^2}
=\frac74mv^2.
\]
所以 $v=\sqrt{4gh/7}$，与 $v^2=2ah$ 一致。
若忽略滑轮质量，相当于把有效惯性分母从 $7m/2$ 改为 $3m$，会误得 $a=g/3$；无滑动只约束运动，并不保证两侧张力相等。

% END DETAILED ANSWER ex44
\TeachingFigure{c3-diagram-35}{滑轮惯量使两侧张力产生差值}
% BEGIN ORIGINAL FIGURES
\OriginalFigure[6.2cm]{c3-4-p011-massive-pulley}{有质量滑轮两侧的张力、悬挂物与正转方向}
% END ORIGINAL FIGURES


\example{倾斜薄圆盘的最大轴承压力}\label{ex:c3-45}

\QuestionStatement{均匀薄圆盘质量 $m$、半径 $r$，
质心 $O$ 在水平轴 $AB$ 上，$AO=a,OB=b$。
盘法线与轴夹角 $\beta$，以恒定 $\omega$ 随轴转动。
取 $z$ 沿 $A$ 指向 $B$，瞬时盘法线为
$\vect n=\sin\beta\,\vect e_x+\cos\beta\,\vect e_z$，
$t=0$ 时 $\vect e_x$ 竖直向上。求两个轴承的反力和最大压力。
明确反力是轴承作用于盘轴系统的力；分别求最大模，并说明最大值是否同时出现。}{ex45}{例45题图}
% BEGIN DETAILED ANSWER ex45
\solution \Step{从主惯量写出实际轴向的张量}
以图~\ref{fig:solution-ex45}~中盘法线与实际转轴的夹角 $\beta$为依据，先写出图右侧第一行的几何或动力学关系，再代入其下方的计算式。

薄盘沿法线的主惯量为 $mr^2/2$，盘内为 $mr^2/4$。令 $J=mr^2/4$，有
\[
\mathsf I_O=J(\mathsf1+\vect n\vect n^{\mathsf T}),\quad
I_{xz}=J\sin\beta\cos\beta=\frac{mr^2}{8}\sin2\beta,
\quad I_{yz}=0.
\]
绕实际轴的 $I_{zz}=J(1+\cos^2\beta)$。即使质心在轴上，只要 $I_{xz}\ne0$，轴也不是主轴。
\SolutionFigure{ex45}{例45的计算示意}
\Step{先分配重力静载，再求动力力偶}
令 $\vect e_X$ 为固定向上方向。由静力平衡与关于 $O$ 的矩平衡，
\[
G_A+G_B=mg,\quad -aG_A+bG_B=0,
\quad G_A=\frac{mgb}{a+b},\ G_B=\frac{mga}{a+b}.
\]
恒速转动所需力偶为 $I_{xz}\omega^2\vect e_y(t)$。令动力附加反力为
$\Delta\vect R_B=D\vect e_x(t)$、$\Delta\vect R_A=-D\vect e_x(t)$，其力偶为 $(a+b)D\vect e_y(t)$，所以
\[
D=\frac{mr^2\omega^2\sin2\beta}{8(a+b)}.
\]
\Step{在固定空间方向合成整周期反力}
$t=0$ 时 $\vect e_x=\vect e_X$，并有 $\vect e_x(t)\cdot\vect e_X=\cos\omega t$。于是
\[
\vect R_A=G_A\vect e_X-D\vect e_x(t),\qquad
\vect R_B=G_B\vect e_X+D\vect e_x(t),
\]
\[
R_A^2=G_A^2+D^2-2G_AD\cos\omega t,\quad
R_B^2=G_B^2+D^2+2G_BD\cos\omega t.
\]
\Step{最大模与退化姿态}
余弦在 $[-1,1]$ 内变化，故两轴承各自最大压力为
$R_{A,\max}=G_A+|D|$、$R_{B,\max}=G_B+|D|$，若 $D\ne0$，最大值相位相差 $\pi$。
所有轴承中的最大值为 $\max(G_A,G_B)+|D|$。
当 $\beta=0$ 或 $\pi/2$，$D=0$，恢复纯静载；当 $\omega$ 加倍时，附加力增至四倍。压力取力的模，因此盘轴作用于轴承的反力虽方向相反，最大压力数值相同。

% END DETAILED ANSWER ex45
\TeachingFigure{c3-diagram-36}{倾斜圆盘质心在轴上，但轴通常不是主轴}
% BEGIN ORIGINAL FIGURES
\OriginalFigure[6.2cm]{c3-4-p018-tilted-disk}{倾斜薄圆盘的转轴、主轴与轴承反力；图中 $\alpha$ 对应正文 $\beta$，$N$ 的分量对应 $R$ 的分量}
% END ORIGINAL FIGURES


\example{T形双杆的瞬时角加速度与轴反力}\label{ex:c3-46}

\QuestionStatement{两均匀细杆各长 $L$、各质量 $m$，
水平杆 $OA$ 一端固定于轴 $O$，
另一端与竖直杆的中点 $A$ 刚接，组成 T 形刚体。
在 $OA$ 水平、竖杆竖直时，刚体顺时针角速度为 $\omega$。
求重力作用下的角加速度和轴反力。
取 $x$ 向右、$y$ 向上、逆时针为正；再代入 $m=8\,\mathrm{kg}$、$L=0.5\,\mathrm m$、$\omega=4\,\mathrm{s^{-1}}$、$g=9.8\,\mathrm{m\,s^{-2}}$。}{ex46}{例46题图}
% BEGIN DETAILED ANSWER ex46
\solution \Step{分别计算两根杆的惯量与总质心}
以图~\ref{fig:solution-ex46}~中两个杆中点的重力作用位置为依据，先写出图右侧第一行的几何或动力学关系，再代入其下方的计算式。

水平杆关于端点 $O$ 的惯量是 $mL^2/3$；竖杆中心在距 $O$ 为 $L$ 的 $A$，故其惯量为 $mL^2/12+mL^2$。相加得
\[
I_O=\frac{17}{12}mL^2,\quad
x_C=\frac{m(L/2)+mL}{2m}=\frac{3L}{4},\quad y_C=0.
\]
\SolutionFigure{ex46}{例46的计算示意}
\Step{以固定轴取矩求有向角加速度}
两重力关于 $O$ 的矩为 $-mg(L/2)-mgL=-3mgL/2$，因此
\[
\alpha=-\frac{3mgL/2}{17mL^2/12}=-\frac{18g}{17L}.
\]
记顺时针角加速度大小 $\beta=18g/(17L)$。给定瞬时角速度也为顺时针，但离心型加速度含 $\omega^2$，不受其符号影响。
\Step{由质心加速度求轴反力}
\[
\vect a_C=\alpha\vect k\times\left(\frac{3L}{4}\vect i\right)
-\omega^2\frac{3L}{4}\vect i
=-\frac{3L\omega^2}{4}\vect i-\frac{3L\beta}{4}\vect j.
\]
总质量为 $2m$，动量方程 $\vect R_O-2mg\vect j=2m\vect a_C$ 给出
\[
R_{Ox}=-\frac32mL\omega^2,\qquad
R_{Oy}=2mg-\frac32mL\beta=\frac7{17}mg.
\]
\Step{数值代入与量纲关系}
给定 $m=8,L=0.5,\omega=4,g=9.8$（均为 SI 单位），得
$\beta=20.75294\,\mathrm{s^{-2}}$，
$R_{Ox}=-96\,\mathrm N$、$R_{Oy}=32.28235\,\mathrm N$，
$|\vect R_O|=101.28253\,\mathrm N$。
向左反力提供质心的法向加速度，向下的质心切向加速度使向上轴反力小于总重；仅有轴力矩方程不能确定这两个反力分量。

% END DETAILED ANSWER ex46
\TeachingFigure{c3-diagram-37}{T形杆的质心偏在横杆中点之外}
% BEGIN ORIGINAL FIGURES
\OriginalFigure[6.2cm]{c3-4-p027-t-rod-loads}{T形双杆的尺寸、重力与顺时针角速度}
% END ORIGINAL FIGURES



\subsubsection*{立方体：先建立质心几何，再决定回落分支}
在立方体绕铰棱转动的垂直截面中，取铰点为原点。
未抬起时截面质心为 $(a/2,a/2)$，抬起 $\theta$ 后
\[
x_C=\frac a2(\cos\theta-\sin\theta),\qquad
 y_C=\frac a2(\sin\theta+\cos\theta).
\]
于是 $x_C^2+y_C^2=a^2/2$，质心沿半径 $a/\sqrt2$ 的圆运动。
关于与铰棱平行的中心轴，长方体公式给
$I_C=m(a^2+a^2)/12=ma^2/6$。
再用平行轴定理加上 $ma^2/2$，才得到 $I_O=2ma^2/3$。
不能只用截面面积惯性矩，也不能把立方体当成集中在质心的一个点质量。

选重力势能 $V=mgy_C$，对转角求导：
\[
\frac{\dd V}{\dd\theta}=\frac{mga}{2}(\cos\theta-\sin\theta)
=mgx_C.
\]
因而重力矩为 $-mgx_C$。
在 $0<\theta<\pi/4$，质心在铰点右侧，重力使转角减小；
在 $\pi/4<\theta<\pi/2$，质心在铰点左侧，重力使转角增大。
$\theta=\pi/4$ 时质心恰在铰点上方，势能达到最大值。
这些几何关系在求平方根之前就已确定回落方向。

从静止释放后，能量式可以整理为任意中间位置的结果
\[
\dot\theta^2=\frac{3g}{2a}
\bigl(\sin\theta_0+\cos\theta_0-\sin\theta-\cos\theta\bigr).
\]
左支路应取负平方根，右支路应取正平方根。
左支路的回落时间满足
\[
t(\theta)=\sqrt{\frac{2a}{3g}}
\int_{\theta}^{\theta_0}
\frac{\dd s}{\sqrt{\sin\theta_0+\cos\theta_0-\sin s-\cos s}}.
\]
这给出实际初值的隐式解，而角加速度由力矩方程直接给出。
若从能量式求导并除以 $\dot\theta$，只在非零速度区间得到同一角加速度；
静止释放点仍应由原力矩方程或连续极限确定。

\subsubsection*{立方体：触地前支反力与不稳定姿态}
由 $\dot x_C=-y_C\dot\theta$、$\dot y_C=x_C\dot\theta$，得
\[
(a_C)_x=-y_C\ddot\theta-x_C\dot\theta^2,\qquad
(a_C)_y=x_C\ddot\theta-y_C\dot\theta^2.
\]
理想铰的合反力因此为
\[
R_x=-m(y_C\ddot\theta+x_C\dot\theta^2),\qquad
R_y=mg+m(x_C\ddot\theta-y_C\dot\theta^2).
\]
将它与 $I_O\ddot\theta=-mgx_C$ 联立，才同时获得运动与支反力。
左支路释放瞬间 $\dot\theta=0$，反力为
\[
R_x(0)=\frac38mg\cos2\theta_0,
\qquad
R_y(0)=mg\left[1-\frac38(\cos\theta_0-\sin\theta_0)^2\right].
\]
在 $\theta=0$ 的触地前一瞬，令 $s_0=\sin\theta_0+\cos\theta_0$，得到
\[
R_x^-=mg\left(\frac98-\frac34s_0\right),\qquad
R_y^-=mg\left(\frac{11}{8}-\frac34s_0\right).
\]
这两个量不包含随后地面接触产生的冲量。
理想铰可承担有向反力；如果把它换成无铰的粗糙角点接触，
还须检查 $R_y\ge0$ 和 $|R_x|\le\mu_sR_y$，
否则原来假定的固定棱转动不能保持。
保持运动学条件的铰链和单侧地面接触是不同的力学模型。

在不稳定姿态附近令 $\theta=\pi/4+\varepsilon$，线性化得
\[
\ddot\varepsilon=\frac{3g}{2\sqrt2a}\varepsilon.
\]
扰动按指数增长，而不是作回复振动。
精确的 $\varepsilon=\dot\varepsilon=0$ 仍是静止解。
右支路首次在 $\theta=\pi/2$ 触地，角速度为同一根式的正值，
角加速度为 $+3g/(4a)$；相应水平反力与左支路镜像。
两支路的能量与边界姿态对称，但碰撞过程必须由另外的冲量模型处理。

\subsubsection*{定滑轮：三个对象的张力与转向}
设重物下降位移为 $x$，则轻物上升同样的 $x$，轻绳速度与轮缘速度大小均为 $\dot x$。
把轮惯量先保留为一般的 $I$，逐对象写
\[
T_1=mg+m\ddot x,\qquad T_2=2mg-2m\ddot x.
\]
上图中右端重物下降，滑轮实际顺时针转动。
若用 $\beta=\ddot x/r$ 表示顺时针角加速度大小，则
\[
(T_2-T_1)r=I\beta.
\]
若统一采用逆时针有向角加速度 $\alpha_z$，应写
\[
\alpha_z=-\frac{\ddot x}{r},\qquad
(T_1-T_2)r=I\alpha_z.
\]
两种写法是同一个有向方程，不能把正的大小 $\beta$ 与有向 $\alpha_z$ 混用。
代入两块的张力式得到
\[
mg-3m\ddot x=\frac I{r^2}\ddot x,
\qquad
\ddot x=\frac{mg}{3m+I/r^2}.
\]
对均匀圆盘 $I/r^2=m/2$，分母为 $7m/2$，故 $a=2g/7$。
再回代才得到两边张力，张力差为 $mg/7$，恰好提供轮的角加速度。
自由绳段的张力各自恒定，而绕在有惯量轮上的接触段有分布切向力，
所以轻绳条件不要求跨过有惯量滑轮后的两个自由段张力相等。

\subsubsection*{定滑轮的能量、功率与轴反力}
重块势能减少 $2mgx$，轻块势能增加 $mgx$，净减少 $mgx$。
全部动能为
\[
T=\frac12\left(3m+\frac I{r^2}\right)\dot x^2.
\]
由静止开始，$mgx=T$，故
\[
\dot x^2=\frac{2mgx}{3m+I/r^2}.
\]
在非零速度区间用 $\dot x\,\dd\dot x/\dd x=\ddot x$，
即可得到与逐对象动力学一致的常加速度。
初速度不为零时，只需把右边动能改为动能增量，动力学加速度不变。

张力对轻块、重块和轮的功率分别含
$+T_1\dot x$、$-T_2\dot x$、$(T_2-T_1)\dot x$，三者严格相消。
因此张力把能量在各对象之间传递，但不在这个理想系统中产生净耗散。
轮的转动能占总动能比例为
\[
\frac{I/r^2}{3m+I/r^2};
\]
圆盘时为 $1/7$。忽略轮惯量会同时改掉加速度、张力差和能量分配。

滑轮中心固定，其合力为零；两绳段对轮均向下，加上轮自重，轴反力为
\[
R_x=0,\qquad R_y=mg+T_1+T_2=\frac{26}{7}mg.
\]
对包括两个物块与轮的整个系统，总质量为 $4m$，
竖直总动量导数是 $ma-2ma=-ma$，所以
$R_y-4mg=-ma$，也给出同一个 $R_y$。
这个核查同时区分了轮的自重与两块的重量。

\subsubsection*{倾斜盘：从主惯量到实际张量}
薄圆盘沿法线的主惯量为 $I_\parallel=mr^2/2$，
盘面内任一直径的主惯量为 $I_\perp=mr^2/4$。
任何方向上的张量可写成
\[
\mathsf I=I_\perp\mathsf1
+(I_\parallel-I_\perp)\vect n\vect n^{\mathsf T}.
\]
代入 $\vect n=(\sin\beta,0,\cos\beta)$，完整矩阵为
\[
\mathsf I=\frac{mr^2}{4}
\begin{pmatrix}
1+\sin^2\beta&0&\sin\beta\cos\beta\\
0&1&0\\
\sin\beta\cos\beta&0&1+\cos^2\beta
\end{pmatrix}.
\]
因此实际 $I_{xz}$ 为正的 $mr^2\sin2\beta/8$。
盘法线倾向 $+x$，而薄盘沿法线惯量大于横向惯量，
二者共同决定这个非对角元的号数。
若换成沿自身长轴惯量较小的细杆，同方向倾斜时 $I_{xz}$ 为负，
这与前面斜杆的结果一致。

恒速转动时
\[
\vect L_O=I_{zz}\omega\vect e_z+I_{xz}\omega\vect e_x(t),
\qquad
\dot{\vect L}_O=I_{xz}\omega^2\vect e_y(t).
\]
质心 $O$ 在轴上且固定，故支承反力之和只能平衡重力。
关于 $O$ 取矩，$A$、$B$ 的臂分别为 $-a\vect e_z,b\vect e_z$，得到
\[
\vect R_A+\vect R_B=mg\vect e_X,
\]
\[
\vect e_z\times(-a\vect R_A+b\vect R_B)
=I_{xz}\omega^2\vect e_y(t).
\]
将静态反力与附加项分开，前式给出 $G_A+G_B=mg$，
后式的静态部分给出 $-aG_A+bG_B=0$。
附加部分因此是反向等长力，大小
$D=I_{xz}\omega^2/(a+b)$，方向分别为 $-\vect e_x(t)$ 与 $+\vect e_x(t)$。
这样同时得到两反力的合力条件和力偶条件。

\subsubsection*{倾斜盘：各相位的压力与模型边界}
取 $\vect e_Y=\vect e_z\times\vect e_X$，则
$\vect e_x(t)=\cos\omega t\,\vect e_X+\sin\omega t\,\vect e_Y$。
总反力分量为
\begin{align*}
R_{AX}&=G_A-D\cos\omega t,&R_{AY}&=-D\sin\omega t,\\
R_{BX}&=G_B+D\cos\omega t,&R_{BY}&=D\sin\omega t.
\end{align*}
在 $0<\beta<\pi/2$ 时 $D>0$。
$t=0$ 的反力分别是 $(G_A-D)\vect e_X$、$(G_B+D)\vect e_X$；
$t=\pi/|\omega|$ 时分别是 $(G_A+D)\vect e_X$、$(G_B-D)\vect e_X$。
四分之一周期处，两压力大小是 $\sqrt{G_A^2+D^2}$、$\sqrt{G_B^2+D^2}$。
这些不同相位的计算说明为什么仅算初始竖直分量无法得到每个轴承的峰值。

各自最低压力为 $|G_A-|D||$、$|G_B-|D||$。
当 $|D|>G_A$ 或 $|D|>G_B$，对应支承某些相位需要向下分量；
封闭轴承可以提供这种有向力，而仅从下方托住轴的单侧支座不能保证接触。
这时应重新分析支承类型，不能把负分量擅自改成零而仍维持原定轴运动。
附加项随 $\omega^2$ 增大，在 $\beta=\pi/4$ 取最大几何系数，
在 $\beta=0,\pi/2$ 消失。仅消除附加力偶不消除重力支承力。

本理想模型中支点速度为零，重力作用点固定，横向约束力偶垂直 $\vect\omega$。
所以即使轴承承受随相位变化的力，其瞬时功率仍为零，转动动能保持不变。
实际轴承的摩擦损耗需要另用轴向阻力矩计入能量方程。

\subsubsection*{T形杆：惯量、质心和有向加速度逐项求和}
% BEGIN ORIGINAL FIGURES
\OriginalFigure[6.2cm]{c3-4-p028-t-rod-accelerations}{T形双杆两质心的实际加速度方向；图中 $X_O,Y_O$ 对应正文的 $R_{Ox},R_{Oy}$}
% END ORIGINAL FIGURES

水平杆的质量元为 $\dd m=(m/L)\dd x$，$0\le x\le L$，故
\[
I_{O,\text{水平}}=\frac mL\int_0^Lx^2\dd x=\frac13mL^2.
\]
竖直杆沿 $x=L$、$-L/2\le y\le L/2$，其对 $O$ 的距离平方是 $L^2+y^2$，所以
\[
I_{O,\text{竖直}}=\frac mL\int_{-L/2}^{L/2}(L^2+y^2)\dd y
=mL^2+\frac1{12}mL^2.
\]
相加为 $17mL^2/12$。
这里 $m$ 是每根杆的质量，整体质量为 $2m$；
若把 $m$ 当总质量，质心计算与转矩都会发生两倍的错误。
质量加权给
\[
x_C=\frac{m(L/2)+mL}{2m}=\frac{3L}{4},\qquad y_C=0.
\]
总重力矩为 $-(mg)(L/2)-(mg)L=-3mgL/2$。
取逆时针有向角加速度为 $\alpha_z$，则
\[
\alpha_z=-\frac{18g}{17L}=-\beta.
\]
顺时针角速度虽然为负，其平方始终非负。
质心的加速度由两项明确构成
\[
\vect a_C=\underbrace{-\frac{3L}{4}\omega^2\vect i}_{\text{向心}}
+\underbrace{-\frac{3L}{4}\beta\vect j}_{\text{切向}}.
\]
代入整体合力式，得到轴反力，并可化简竖直分量为
\[
R_{Oy}=2mg-\frac{27}{17}mg=\frac7{17}mg.
\]
当前水平姿态的竖直反力与瞬时转速无关；水平反力则随 $\omega^2$ 改变。
这一点来自向心项水平、切向项竖直的特定几何关系，不能外推到任意姿态。

\subsubsection*{T形杆：数值回代与刚接条件}
给定数据下，
\[
I_O=\frac{17}{12}(8)(0.5)^2=\frac{17}{6}\,\mathrm{kg\,m^2},
\qquad |M_O|=\frac32(8)(9.8)(0.5)=58.8\,\mathrm{N\,m}.
\]
因此 $\beta=58.8/(17/6)=20.75294\,\mathrm{s^{-2}}$。
向心加速度为 $-6\,\mathrm{m\,s^{-2}}$，
切向加速度约为 $-7.78235\,\mathrm{m\,s^{-2}}$。
再乘整体质量 $16\,\mathrm{kg}$ 并加重力，分别给出
$R_{Ox}=-96\,\mathrm N$、$R_{Oy}=32.28235\,\mathrm N$。
压力大小是 $\sqrt{R_{Ox}^2+R_{Oy}^2}$，不是两分量大小相加。

“刚接”还意味着中点 $A$ 能传递力偶。
只把竖直杆作为研究对象，它的质心在 $A$，加速度为
$-L\omega^2\vect i-L\beta\vect j$。
水平杆对它的接点力是
\[
F_{Ax}=-mL\omega^2,\qquad F_{Ay}=mg-mL\beta.
\]
其中心轴惯量为 $mL^2/12$，重力对中心的矩为零，
故接点还须施加有向力偶
\[
K_A=-\frac{mL^2}{12}\beta.
\]
这个内部力偶在两根杆的整体方程中相消，但维持两杆同一角加速度不可缺少。
如果连接处改成光滑铰，竖直杆就不是整体刚体的一部分，
不能继续沿用上面的总惯量和单一转角方程。


\clearpage
\section{平面运动、滚动与耦合系统}\label{sec:c3-7}
\subsection{平面速度场、投影定理与瞬心}\label{sub:c3-7.1}
刚体各点始终在互相平行的平面内运动，称平面平行运动。
用质心平面坐标 $x_C,y_C$ 与姿态角 $q$ 描述，有三个自由度。
角速度垂直运动平面，$\vect\omega=\omega\vect k$。
% BEGIN ORIGINAL FIGURES
\OriginalFigure[6.2cm]{c3-5-p002-parallel-planes}{平面平行运动中的平行截面}
\OriginalFigure[6.2cm]{c3-5-p002-rotation-translation}{刚体的转动轴与平面内平移}
% END ORIGINAL FIGURES

对两固连点，
\[
\vect v_B=\vect v_A+\omega\vect k\times\vect r_{AB},\quad
\vect a_B=\vect a_A+\alpha\vect k\times\vect r_{AB}-\omega^2\vect r_{AB}.
\]
把速度差向 $AB$ 投影，叉积项为零，因此
\[
\vect v_B\cdot\vect e_{AB}=\vect v_A\cdot\vect e_{AB}.
\]
这称速度投影定理：两点速度在连线上的投影相同。
它给约束关系，但只有一个投影式一般不能确定全部速度与角速度。

当 $\omega\ne0$，存在唯一平面位置 $P$ 使速度场在该位置为零：
\[
\vect r_{CP}=\frac{\vect k\times\vect v_C}{\omega},\qquad
\vect v_A=\omega\vect k\times\vect r_{PA}.
\]
它称速度瞬心，可以位于刚体外部。
每一点的速度垂直于它与瞬心的连线，大小为 $|\omega|\,PA$。
求瞬心常用以下条件：已知一点速度及角速度，
在垂直该速度的方向上按有向距离定位；
已知两点速度方向且不平行，取两条速度法线交点；
已知无滑动接触固定轨道，接触处为速度瞬心。
若两点速度平行，要先判断它们的大小和转向；
若 $\omega=0$ 而平移速度非零，瞬心在无穷远，不能强找有限交点。

\TeachingFigure{c3-diagram-38}{平面速度场可以在一瞬间等效为绕瞬心转动}
% BEGIN ORIGINAL FIGURES
\OriginalFigure[6.2cm]{c3-5-p006-instant-center}{平面运动的速度瞬心与垂直距离；图中 $C$ 为瞬心，正文记为 $P$}
\OriginalFigure[6.2cm]{c3-5-p007-one-velocity}{由一点速度和角速度确定瞬心}
\OriginalFigure[6.2cm]{c3-5-p007-rolling-contact}{固定曲面上纯滚动时，接触点为速度瞬心}
\OriginalFigure[6.2cm]{c3-5-p008-nonparallel-velocities}{两条不平行速度的法线相交于瞬心}
% END ORIGINAL FIGURES



\subsubsection*{三个运动变量如何产生整个速度场}
取惯性平面基 $\vect i,\vect j$，$\vect k=\vect i\times\vect j$。
刚体上任一材料点的本体坐标为 $\vect b$，空间位置可写为
\[
\vect r=\vect r_C+\mathsf Q(q)\vect b.
\]
这里 $\vect r_C$ 给两个平移变量，$q$ 给一个姿态变量，
所以三者已经确定整个刚体的平面位置。
对时间求导，令 $\vect\rho=\mathsf Q(q)\vect b$，得到
\[
\vect v=\vect v_C+\omega\vect k\times\vect\rho,
\qquad \omega=\dot q.
\]
第二次求导时，$\dot{\vect\rho}=\omega\vect k\times\vect\rho$，因而
\[
\vect a=\vect a_C+\alpha\vect k\times\vect\rho-\omega^2\vect\rho.
\]
任选两个材料点相减即可消去质心量，得前面的基点合成式。
基点可以是任何固连点，选质心方便动力学，选铰接点方便连接约束，
选速度已知的点方便运动学。
所谓平移加转动是一种运动描述的分解，不表示刚体先平移再转动两个独立过程。
% BEGIN ORIGINAL FIGURES
\OriginalFigure[6.2cm]{c3-5-p003-planar-position}{基点位置与相对位置的矢量合成}
% END ORIGINAL FIGURES


设 $\vect r_{AB}=\Delta x\vect i+\Delta y\vect j$，速度分量关系是
\[
v_{Bx}-v_{Ax}=-\omega\Delta y,\qquad
v_{By}-v_{Ay}=\omega\Delta x.
\]
于是若两点完整速度已知，可直接求
\[
\omega=\frac{\bigl(\vect r_{AB}\times(\vect v_B-\vect v_A)\bigr)_z}
{|\vect r_{AB}|^2}.
\]
同时还必须检查两速度差在 $AB$ 上的投影是否为零，
否则给定数据本身不满足刚性长度约束。

\subsubsection*{速度投影来自长度约束，加速度投影含向心项}
固连两点的距离不变，即 $\vect r_{AB}\cdot\vect r_{AB}=l^2$。
第一次求导给
\[
\vect r_{AB}\cdot(\vect v_B-\vect v_A)=0.
\]
因此两点沿连线的速度分量相等。
这里比较的是有向分量，不能把投影后的代数值擅自换成速度大小。
例如一端沿水平轨道、另一端沿竖直轨道时，
连线方向的两投影可能同时为负，实际两端却沿不同方向运动。

再求一次导数，得到
\[
|\vect v_B-\vect v_A|^2
+\vect r_{AB}\cdot(\vect a_B-\vect a_A)=0.
\]
利用相对速度大小为 $|\omega|l$，有
\[
(\vect a_B-\vect a_A)\cdot\vect e_{AB}=-\omega^2l.
\]
所以加速度在连线上的投影并不相等；差值恰是相对向心加速度。
向连杆方向投影能消去 $\alpha$ 引起的切向项，
却不能消去 $\omega^2$ 项，这将在曲柄滑块的加速度计算中使用。

\subsubsection*{速度瞬心的有向坐标与唯一性}
在质心速度场中令 $\vect v_P=0$，写分量
\[
0=v_{Cx}-\omega(y_P-y_C),\qquad
0=v_{Cy}+\omega(x_P-x_C).
\]
当 $\omega\ne0$，唯一解为
\[
x_P=x_C-\frac{v_{Cy}}\omega,\qquad
 y_P=y_C+\frac{v_{Cx}}\omega.
\]
该点可能在实体外，但刚体的速度场可以几何延拓到这个平面位置。
消去 $\vect v_C$ 后，对任意点有
$\vect v_A=\omega\vect k\times\vect r_{PA}$。
因此速度法线必过 $P$，速度大小和到 $P$ 的距离成比例。
方向由 $\omega$ 的正负决定：逆时针正角速度使每条半径按正向旋转 $90^\circ$。

已知两点速度方向且其法线不平行时，两法线交点给出 $P$。
已知一点速度和角速度时，要用有向式定位，而不仅是距离 $|v|/|\omega|$。
例如一点速度向右，$\omega<0$，瞬心在该点下方；
若将它放在上方，就得到相反转向。

平行速度需要分情况讨论。
设 $A=(0,0),B=(d,0)$，两速度均沿竖直方向，
$v_{Ay}=u_A,v_{By}=u_B$，则
\[
\omega=\frac{u_B-u_A}{d},\qquad
 x_P=\frac{du_A}{u_A-u_B},\quad y_P=0.
\]
同向且大小不等时瞬心在两点之外；反向时位于两点之间；
% BEGIN ORIGINAL FIGURES
\OriginalFigure[6.2cm]{c3-5-p008-opposite-directions}{两点反向运动时，瞬心位于两点之间}
\OriginalFigure[6.2cm]{c3-5-p008-same-direction}{两点同向且不等速时，瞬心位于两点之外}
% END ORIGINAL FIGURES

一端速度为零时瞬心就是该端。
若 $u_A=u_B\ne0$，则 $\omega=0$，整个速度场是瞬时平移，有限瞬心不存在。
% BEGIN ORIGINAL FIGURES
\OriginalFigure[6.2cm]{c3-5-p009-instant-translation}{不同位置的两点速度相同，对应瞬时平移}
% END ORIGINAL FIGURES

若两个不同点速度都为零，则 $\omega=0,\vect v_C=0$，整个刚体瞬时静止，
任意位置的速度场都为零，因而不存在唯一的瞬心。

这些结论只描述一个时刻的速度。
瞬心没有因此成为固定铰，也不能仅凭它的存在就写固定轴的动力学方程。
速度瞬心法最直接的用途是求速度与角速度，
求加速度和力矩时还需处理它随时间改变的位置。


\subsection{极迹与瞬心加速度}\label{sub:c3-7.2}
瞬心在空间中的轨迹称空间极迹，在本体中的轨迹称本体极迹。
这两条曲线在运动过程中相切，并作相对无滑动的几何滚动。
直线轨道上的车轮，空间极迹为地面直线，本体极迹为轮缘圆。
极迹上的位置是依次改变的几何位置，通常不是同一材料点的轨迹。
% BEGIN ORIGINAL FIGURES
\OriginalFigure[6.2cm]{c3-5-p011-body-circle}{本体极迹为轮缘圆；接触材料点随滚动而改变}
\OriginalFigure[6.2cm]{c3-5-p011-contact-loci}{滚动过程中依次成为瞬心的空间位置与材料位置}
\OriginalFigure[6.2cm]{c3-5-p011-fixed-circle}{空间极迹为直线；蓝点标识当前接触位置}
% END ORIGINAL FIGURES


速度瞬心处恰好速度为零，不保证该处材料点加速度为零。
直线匀速纯滚动的轮子在接触处，向上的加速度为 $v_C^2/R$。
不能把 $\vect v_P=0$ 对时间直接求导并断言
$\vect a_P=0$，因为每一时刻选取的接触材料点都不同。

\example{等速拖动梯子的两条极迹}\label{ex:c3-47}

\QuestionStatement{梯子长 $l$，端点 $A$ 沿水平地面向右等速 $u$ 移动，
端点 $B$ 沿竖直墙向下滑动。取墙脚为原点，
$A=(x,0),B=(0,y)$，$x^2+y^2=l^2$。
求角速度、空间极迹和本体极迹。
取 $x,y>0$，逆时针为正，写出本体系的取向和两条极迹适用的弧段。}{ex47}{例47题图}
% BEGIN DETAILED ANSWER ex47
\solution \Step{由两端速度确定瞬时角速度中心}
以图~\ref{fig:solution-ex47}~中两端速度法线交点 $P$为依据，先写出图右侧第一行的几何或动力学关系，再代入其下方的计算式。

$A$ 的速度水平，故瞬心在过 $A$ 的竖线上；$B$ 的速度竖直，故瞬心在过 $B$ 的横线上。交点为 $P=(x,y)$。
由刚体速度式，$\vect v_A=\omega\vect k\times(0,-y)=\omega y\vect i$，所以
\[
\omega=\frac uy>0.
\]
对长度约束求导，$x\dot x+y\dot y=0$，故
$\vect v_B=-xu\vect j/y$，与绕 $P$ 的速度一致。
\SolutionFigure{ex47}{例47的计算示意}
\Step{消去时间得到空间极迹}
瞬心坐标为 $(x_P,y_P)=(x,y)$，而 $x^2+y^2=l^2$，因此
\[
x_P^2+y_P^2=l^2,\qquad x_P>0,\ y_P>0.
\]
空间极迹是墙脚为圆心、半径 $l$ 的第一象限圆弧；它描述的是各瞬时不同的速度零点在空间的位置。
\Step{投影到固定在梯子上的正交轴}
以 $A$ 为本体原点，令 $\vect e_1=(-x,y)/l$ 沿 $AB$，
$\vect e_2=(-y,-x)/l=\vect k\times\vect e_1$。因 $\vect r_{AP}=(0,y)$，
\[
q_1=\vect r_{AP}\cdot\vect e_1=\frac{y^2}{l},\qquad
q_2=\vect r_{AP}\cdot\vect e_2=-\frac{xy}{l}.
\]
代入约束得 $q_1^2+q_2^2=lq_1$，整理为
\[
(q_1-l/2)^2+q_2^2=l^2/4,\qquad q_2<0.
\]
本体极迹是以梯子中点为圆心、直径 $l$ 的下半圆弧。
$y\to0$ 时 $\omega=u/y$ 发散，说明等速拖动的理想模型在梯子即将放平时要求无界转速。

% END DETAILED ANSWER ex47
\TeachingFigure{c3-diagram-39}{梯子端点的速度法线与空间极迹}
% BEGIN ORIGINAL FIGURES
\OriginalFigure[6.2cm]{c3-5-p012-ladder-loci}{梯子的空间与本体极迹；课件以 $C$ 表示瞬心，正文以 $P$ 表示，课件 $\xi,\eta$ 与正文 $q_1,q_2$ 的轴向约定不同}
% END ORIGINAL FIGURES



\subsubsection*{空间极迹与本体极迹为何相切且等弧长}
设空间瞬心的位置为 $\vect p(t)$，在以材料点 $A$ 为本体原点的坐标中，
它的位置为 $\vect b_P(t)$。几何变换给
\[
\vect p=\vect r_A+\mathsf Q(q)\vect b_P.
\]
这里 $\vect b_P$ 随时间改变，因为不同的本体位置依次成为瞬心。
求导得到
\[
\dot{\vect p}
=\vect v_A+\omega\vect k\times(\vect p-\vect r_A)
+\mathsf Q\dot{\vect b}_P.
\]
前两项之和是当前几何瞬心位置的材料速度，按定义为零。
因此
\[
\dot{\vect p}=\mathsf Q\dot{\vect b}_P.
\]
在两条极迹均为正则曲线的位置，右边将本体极迹的切向转到空间方向，
恰好等于空间极迹切向；两个矢量长度也相等。
所以两曲线相切，而且对应的微小弧长相等。
这就是极迹之间无滑动几何滚动的推导。
若 $\dot{\vect p}=0$，两曲线的参数速度同时为零，
尖点或退化处的切线需要由邻近曲线的极限判断。

极迹滚动描述整个刚体的运动几何，
并不要求刚体轮廓真的与某个物面接触。
连杆的本体极迹甚至可以部分位于实体之外；
它仍是本体坐标系中一条有明确意义的曲线。

\subsubsection*{几何瞬心的速度与当前材料点的加速度}
用 $\vect v(\vect r,t)$ 表示刚体速度场。
瞬心条件是 $\vect v(\vect p(t),t)=0$。
沿几何瞬心的路径求导，链式法则给
\[
0=\frac{\partial\vect v}{\partial t}
+\omega\vect k\times\dot{\vect p}.
\]
材料点加速度为
$\partial\vect v/\partial t+(\vect v\cdot\nabla)\vect v$；
在瞬心位置 $\vect v=0$，第二项为零。因此当前材料点的加速度为
\[
\vect a_{P,\text{材料}}=-\omega\vect k\times\dot{\vect p}.
\]
几何瞬心本身的移动速度 $\dot{\vect p}$ 并不等于该材料点的速度，
二者不能互相替代。
对右行匀速滚轮，几何接触位置随轮心以 $v\vect i$ 移动，
而 $\omega=-v/R$，所以
\[
\vect a_{P,\text{材料}}
=-\left(-\frac vR\right)\vect k\times(v\vect i)
=\frac{v^2}{R}\vect j.
\]
这直接恢复向上的非零加速度，
也解释了错误地对 $\vect v_P=0$ 求导为什么会漏掉链式法则的移动项。

\subsubsection*{梯子：长度约束、角速度与角加速度}
取 $0<x<l$、$y>0$，由 $x^2+y^2=l^2$ 和 $\dot x=u>0$ 得
\[
x=x_i+ut,\qquad y=\sqrt{l^2-(x_i+ut)^2},\qquad
\dot y=-\frac{xu}{y}.
\]
有效时间区间终止于 $x=l$，不能把平方根公式继续用于地面交叉后的构形。
再求导可得
\[
\ddot y=-\frac{u^2}{y}-\frac{x^2u^2}{y^3}
=-\frac{l^2u^2}{y^3}.
\]
虽然 $A$ 等速，$B$ 的速度和加速度都非恒定。
令 $x=l\cos\theta,y=l\sin\theta$，$0<\theta<\pi/2$，
则 $\dot\theta=-u/(l\sin\theta)$。
沿 $A$ 到 $B$ 的杆方向角是 $\pi-\theta$，所以刚体有向角速度为
\[
\omega=-\dot\theta=\frac uy>0,
\qquad \alpha=\frac{xu^2}{y^3}>0.
\]
梯子随拖动在空间逆时针转动，角速度随 $y$ 下降而增大。
这里 $\theta$ 是几何锐角，$\dot\theta$ 与刚体有向角速度相反，
不应把二者混写。

若梯子均匀，质心位于 $(x/2,y/2)$，于是
\[
\vect v_C=\frac u2\vect i-\frac{xu}{2y}\vect j,
\qquad
\vect a_C=-\frac{l^2u^2}{2y^3}\vect j.
\]
空间瞬心 $P=(x,y)$ 的几何速度是
$u\vect i-(xu/y)\vect j$。
代入材料瞬心加速度式，得到
\[
\vect a_{P,\text{材料}}
=-\frac{xu^2}{y^2}\vect i-\frac{u^2}{y}\vect j.
\]
它显然非零。因此可以用速度法线求 $P$，
但不能把 $P$ 当成固定基点求各点加速度。

\subsubsection*{梯子：两条极迹及其覆盖区间}
空间极迹参数为 $P=(l\cos\theta,l\sin\theta)$，弧长微元为
$\dd s_{\text{空间}}=l|\dd\theta|$。
在 $A$ 原点的本体基 $\vect e_1,\vect e_2$ 中，瞬心坐标化为
\[
q_1=l\sin^2\theta=\frac l2(1-\cos2\theta),\qquad
q_2=-l\sin\theta\cos\theta=-\frac l2\sin2\theta.
\]
因此本体极迹的圆心是 $(l/2,0)$、半径是 $l/2$，
从 $\theta=\pi/2$ 到 $0$ 描出 $q_2\le0$ 的半圆。
相应弧长微元为
\[
\dd s_{\text{本体}}=\frac l2|\dd(2\theta)|=l|\dd\theta|.
\]
两条极迹弧长相等的结论在本例中得到显式验证。
空间轨迹的四分之一圆与本体轨迹的半圆具有相同总弧长 $\pi l/2$。
此外，均匀梯子质心与瞬心之距为
\[
PC=\sqrt{(x/2)^2+(y/2)^2}=l/2,
\]
在这个特定端点约束下恒定。
一般机构的瞬心到质心距离可能改变，但本例恰是恒定距离的一种特殊情况。

\TeachingFigure{c3-diagram-40}{均匀梯子的本体极迹以质心为圆心，覆盖下半圆}


\subsection{动力学、动能与移动瞬心取矩}\label{sub:c3-7.3}
平面动力学的基本式为
\begin{equation}\label{eq:planar}
m\ddot x_C=F_x,\qquad m\ddot y_C=F_y,\qquad
I_C\dot\omega=M_C.
\end{equation}
前三式决定平面内的三个自由度。
若刚体本身为一般三维质量分布，维持平面运动还需三个约束方程。
取质心为原点、本体 $z$ 轴垂直运动平面，
实际惯量矩阵的非对角元仍记 $I_{xz},I_{yz}$，则
\[
F_z=0,\qquad
M_{Cx}=I_{xz}\alpha-I_{yz}\omega^2,\qquad
M_{Cy}=I_{yz}\alpha+I_{xz}\omega^2.
\]
它们由质心的法向加速度为零，以及两个横向角动量导数得到。
如果 $z$ 是中心主轴，右边两横向矩为零；
否则约束须提供这些矩，不能只用平面三式就判断支承载荷。
选择刚体上三个点，其在运动平面的投影互不共线，并保持固定法向坐标，
可以约束法向平移及两个倾斜自由度；
约束点共线则仍可能绕该线倾转。
动能是
$T=\tfrac12mv_C^2+\tfrac12I_C\omega^2$。
若瞬心与质心距离为 $\rho$，
该瞬间 $v_C=|\omega|\rho$，所以
$T=\tfrac12(I_C+m\rho^2)\omega^2=\tfrac12I_P\omega^2$。
这只是瞬时动能表达，动力学取矩还需核查移动点项。
% BEGIN ORIGINAL FIGURES
\OriginalFigure[6.2cm]{c3-5-p041-moving-center}{移动瞬心 $A$ 与质心 $C$ 的相对位置}
% END ORIGINAL FIGURES


对于沿某几何轨迹移动的参考点 $P(t)$，
用惯性速度定义角动量，满足
\[
\frac{\dd\vect L_P}{\dd t}
=\vect M_P-\vect v_{P,\text{几何}}\times m\vect v_C.
\]
若 $P$ 始终是速度瞬心，$\vect L_P=I_P\omega\vect k$。
令 $\vect\rho=\vect r_C-\vect r_P$，
$\vect v_C=\omega\vect k\times\vect\rho$，
$\vect v_{P,\text{几何}}=\vect v_C-\dot{\vect\rho}$，
则
$\vect v_{P,\text{几何}}\times m\vect v_C
=-m\omega\rho\dot\rho\,\vect k$。
又 $\dot I_P=2m\rho\dot\rho$，故
\begin{equation}\label{eq:instanttorque}
I_P\dot\omega=M_P-m\rho\dot\rho\,\omega.
\end{equation}
当 $\rho$ 恒定时，才可简化为 $I_P\dot\omega=M_P$。
固定轨道上的圆轮纯滚动，$\rho$ 恒为轮半径，满足这个条件；
一般连杆或梯子的瞬心到质心距离可能变化，不能直接套简式。
% BEGIN ORIGINAL FIGURES
\OriginalFigure[6.2cm]{c3-5-p015-rolling-centers}{斜面上各瞬间的接触位置依次成为速度瞬心}
% END ORIGINAL FIGURES

以当前材料点取矩也可写
$M_P=I_P\alpha+m(\vect r_{PC}\times\vect a_{P,\text{材料}})_z$，
同样揭示瞬心材料加速度的重要性。


\subsubsection*{平面内三式和维持平面运动的三式}
由质心运动定理，合外力的平面分量满足
$m\ddot x_C=F_x,m\ddot y_C=F_y$。
由关于质心的角动量定理，
$\dot{\vect L}_C=\vect M_C$。
即使质心在空间移动，也没有额外的
$-\vect v_C\times m\vect v_C$ 项，因为这个叉积恒为零。
将 $\vect\omega=\omega\vect e_z$ 代入实际中心惯量张量，再按惯性空间求导，得到
\begin{align*}
M_{Cx}&=I_{xz}\alpha-I_{yz}\omega^2,\\
M_{Cy}&=I_{yz}\alpha+I_{xz}\omega^2,\\
M_{Cz}&=I_{zz}\alpha.
\end{align*}
因此平面内的转动方程使用 $I_C=I_{zz}$，
横向两式用于求约束矩，另有 $F_z=m(a_C)_z=0$。
同样的六个整体关系在不同约束下承担不同任务：
平面内三个决定运动，法向一个和横向矩两个保证运动没有偏离规定平面。

若刚体是理想薄片，质量分布位于 $z=0$，则 $I_{xz}=I_{yz}=0$。
对于一般三维体，即使质心只在 $xy$ 平面运动，也不必有这两个惯性积为零。
例如厚度不对称的构件可能需要横向约束力偶，
但其 $x_C,y_C,q$ 的运动仍可由平面内三式描述。
不能由“运动为平面运动”倒推“质量分布也必是平面薄片”。

三个法向约束力如何维持平面运动，可以用线性代数看清。
设它们沿 $z$ 作用于三个固连点，投影坐标为 $(x_i,y_i)$，有向大小为 $N_i$。
它们产生的合力与横向矩为
\[
\begin{pmatrix}F_z^{(\text{约束})}\\M_x^{(\text{约束})}\\M_y^{(\text{约束})}\end{pmatrix}
=
\begin{pmatrix}
1&1&1\\y_1&y_2&y_3\\-x_1&-x_2&-x_3
\end{pmatrix}
\begin{pmatrix}N_1\\N_2\\N_3\end{pmatrix}.
\]
该矩阵行列式的绝对值等于投影三角形面积的两倍。
投影非共线时，三个有向法向力可独立提供一个法向合力和两个横向矩。
投影共线时矩阵降秩，仍可能绕该线倾转。
如果约束只是单侧接触，还必须检查所解 $N_i$ 的允许符号；
有向铰约束和只能受压的接触不能按同一可行性条件处理。

\subsubsection*{动能分解与瞬心表达的严格含义}
令每个质元相对质心的位置为 $\vect\rho_i$，则
\[
\vect v_i=\vect v_C+\omega\vect k\times\vect\rho_i.
\]
将平方展开并按质量求和，交叉项含
$\sum m_i\vect\rho_i=0$，因此
\[
T=\frac12mv_C^2+\frac12I_C\omega^2.
\]
若 $P$ 是速度瞬心，则 $\vect v_C=\omega\vect k\times\vect r_{PC}$。
令 $\rho=PC$，得到 $T=\tfrac12(I_C+m\rho^2)\omega^2$。
$I_P=I_C+m\rho^2$ 是这一瞬间关于 $P$ 的惯量；
当 $P$ 在本体中的位置改变时，$I_P$ 也随时间变化。
所以这个动能式本身不能保证 $I_P$ 是常数，
也不能保证取矩方程具有固定轴形式。

关于瞬心的外力矩仍给出瞬时总功率。
因为所有点的速度场为 $\omega\vect k\times\vect r_{Pi}$，有
\[
\mathcal P=\sum_i\vect F_i\cdot\vect v_i
=\omega\sum_i(\vect r_{Pi}\times\vect F_i)_z=M_P\omega,
\]
纯力偶的功率也按同一个式子加入。
从动能求导，得
\[
\dot T=I_P\omega\alpha+m\rho\dot\rho\,\omega^2.
\]
与功率式比较，在非零角速度区间即得到移动瞬心修正项。
但要在 $\omega=0$ 或不能使用唯一有限瞬心的时刻正确理解，
仍需下面不依赖除以 $\omega$ 的角动量推导。

\subsubsection*{移动取矩点的运输项逐质元推导}
选择任意空间路径 $\vect p(t)$ 作为取矩点，使用质元的惯性速度定义
\[
\vect L_P=\sum_i(\vect r_i-\vect p)\times m_i\vect v_i.
\]
对时间求导，
\begin{align*}
\dot{\vect L}_P
&=\sum_i(\vect v_i-\dot{\vect p})\times m_i\vect v_i
 +\sum_i(\vect r_i-\vect p)\times m_i\vect a_i\\
&=-\dot{\vect p}\times m\vect v_C+\vect M_P.
\end{align*}
第一项中 $\vect v_i\times\vect v_i=0$；
第二项使用牛顿方程，刚体内部成对的中心作用力矩相消。
因此取矩点的几何移动产生运输项。
这里的 $\dot{\vect p}$ 是参考位置的空间速度，不能用当前材料点的零速度替代。

现在令 $P$ 始终为速度瞬心，并记
$\vect\rho=\vect r_C-\vect p$。
质心轨道角动量加中心角动量给
\[
\vect L_P=I_C\omega\vect k+\vect\rho\times m\vect v_C
=(I_C+m\rho^2)\omega\vect k=I_P\omega\vect k.
\]
又有 $\dot{\vect p}=\vect v_C-\dot{\vect\rho}$，
$\vect v_C=\omega\vect k\times\vect\rho$，故
\begin{align*}
\dot{\vect p}\times m\vect v_C
&=-\dot{\vect\rho}\times m\omega(\vect k\times\vect\rho)\\
&=-m\omega(\dot{\vect\rho}\cdot\vect\rho)\vect k
=-m\omega\rho\dot\rho\vect k.
\end{align*}
而 $\dot I_P=2m\rho\dot\rho$，所以两侧分别为
\[
\dot L_P=I_P\alpha+2m\rho\dot\rho\omega,
\qquad M_P-\dot{\vect p}\times m\vect v_C
=M_P+m\rho\dot\rho\omega.
\]
移项后仅剩一个修正项：
\[
I_P\alpha=M_P-m\rho\dot\rho\omega.
\]
出现一个而不是两个，是因为惯量导数给出的两项中有一项由运输项抵消。
只对 $I_P\omega$ 求导而把右边直接写 $M_P$，会错误地留下两倍修正。

\subsubsection*{固定材料取矩点的加速度式与特殊简化}
在当前瞬心位置取一个固连于刚体的材料点或几何延拓点，
每个质元的加速度满足
\[
\vect a_i=\vect a_{P,\text{材料}}
+\alpha\vect k\times\vect r_{Pi}-\omega^2\vect r_{Pi}.
\]
乘以 $\vect r_{Pi}$ 作叉积并按质量求和，
径向项叉积为零，得到
\[
M_P=I_P\alpha+m(\vect r_{PC}\times\vect a_{P,\text{材料}})_z.
\]
这与移动几何瞬心的运输式描述相同的物理事实。
固定轨道圆轮的 $PC$ 是恒定半径，故 $\rho\dot\rho=0$，
即使当前接触材料点加速度不为零，也可简化为 $I_P\alpha=M_P$。
此时材料瞬心加速度与 $PC$ 平行，叉积为零。
只要求相关叉积或修正项为零，不要求整个加速度矢量为零。

一个非恒定 $\rho$ 的具体核查来自等长曲柄滑块。
均匀连杆的质心和瞬心分别为
\[
C=(\tfrac32l\cos\varphi,\tfrac12l\sin\varphi),\qquad
P=(2l\cos\varphi,2l\sin\varphi).
\]
于是
\[
\rho^2=\frac{l^2}{4}(\cos^2\varphi+9\sin^2\varphi),\qquad
\rho\dot\rho=2l^2\Omega\sin\varphi\cos\varphi.
\]
恒速曲柄下，连杆 $\omega=-\Omega,\alpha=0$，但仍需要
\[
M_P=m\rho\dot\rho\omega
=-2ml^2\Omega^2\sin\varphi\cos\varphi.
\]
这不是说连杆不受外力矩，而是外力矩恰用于改变移动瞬心角动量中的轨道部分。
若误套 $I_P\alpha=M_P$，会把这个非零外矩错误地删去。


\subsection{斜面滚动、退绕落盘与滚动阻力}\label{sub:c3-7.4}
\example{圆柱沿粗糙斜面滚下：三种解法}\label{ex:c3-48}

\QuestionStatement{圆柱质量 $m$、半径 $a$、中心轴惯量 $I_C=mk^2$，
沿倾角 $\beta$ 的固定粗糙斜面无滑动滚下。
从静止开始，忽略滚动阻力。求质心加速度和接触反力。
设坡面向右下倾斜；同时求保证纯滚动所需的最小静摩擦系数，并检验实心圆柱特例。}{ex48}{例48题图}
% BEGIN DETAILED ANSWER ex48
\solution \Step{明确模型与计算依据}
以图~\ref{fig:solution-ex48}~中同一个接触点出发的法向与切向反力为依据，先写出图右侧第一行的几何或动力学关系，再代入其下方的计算式。
 第一种用能量。下移斜距 $x$ 后，
\[
mgx\sin\beta=\frac12m\dot x^2+\frac12mk^2\frac{\dot x^2}{a^2},
\quad
\ddot x=\frac{g\sin\beta}{1+k^2/a^2}.
\]
能量法确定加速度，却不直接确定法向力与静摩擦。

第二种用质心动力学。沿斜面向下为正，角加速度 $\alpha$ 按逆时针取正，摩擦 $f$ 向上，
\[
m\ddot x=mg\sin\beta-f,\quad N=mg\cos\beta,\quad
mk^2\alpha=-fa,\quad \ddot x=-a\alpha.
\]

\SolutionFigure{ex48}{例48的计算示意}
\Step{逐项计算与求解}
因此
\[
f=\frac{mgk^2}{a^2+k^2}\sin\beta,\quad
N=mg\cos\beta,\quad
\mu_s\ge\frac{k^2}{a^2+k^2}\tan\beta.
\]
第三种以接触瞬心 $P$ 取矩，$\rho=a$ 恒定，
\[
I_P\alpha=-mga\sin\beta,\quad I_P=m(k^2+a^2),
\]

\Step{结果与适用条件}
与前两法一致。对均匀实心圆柱 $k^2=a^2/2$，
\[
\ddot x=\frac23g\sin\beta,\quad f=\frac13mg\sin\beta,\quad
\mu_s\ge\frac13\tan\beta.
\]
静摩擦在固定斜面上的接触瞬时速度为零，故总功为零。
它仍把重力输入的一部分分配给转动；
不能从总功为零推出摩擦力为零。
\Step{说明能量求导与初始静止的相容性} 能量式在 $\dot x\ne0$ 的运动段求导后得到加速度；再由运动方程的连续性延伸至释放瞬间，避免在 $t=0$ 直接除以零。实心圆柱由静止开始还有 $x(t)=g\sin\beta\,t^2/3$，与常加速度 $2g\sin\beta/3$ 一致。

% END DETAILED ANSWER ex48
\TeachingFigure{c3-diagram-41}{无滑滚动的摩擦方向由平移与转动共同决定}
% BEGIN ORIGINAL FIGURES
\OriginalFigure[6.2cm]{c3-5-p017-incline-forces}{圆柱沿斜面纯滚动的受力；原图倾角 $\alpha$ 对应正文 $\beta$，图示斜面上升方向与正文坐标选择相反}
% END ORIGINAL FIGURES


\example{固定轻绳上退绕下落的圆盘}\label{ex:c3-49}

\QuestionStatement{均匀圆盘质量 $m$、半径 $r$，轻绳上端固定，
下段沿竖直方向与圆盘边缘相切，绳与盘无滑动。
圆盘从静止释放，圆心仅竖直下降。求加速度与绳张力。
绳在圆盘左侧与轮缘相切，取逆时针为正；给出动力学、能量和瞬心三种求解，并核查约束方向。}{ex49}{例49题图}
% BEGIN DETAILED ANSWER ex49
\solution \Step{从固定绳段确定无滑动约束}
以图~\ref{fig:solution-ex49}~中圆盘左侧切点处的向上张力为依据，先写出图右侧第一行的几何或动力学关系，再代入其下方的计算式。

绳沿圆盘左侧竖直相切。圆心向下速率记为 $v$，逆时针角速度为 $\omega$。
接触材料点相对圆心为 $-r\vect i$，绝对速度为 $-(v+r\omega)\vect j$。
绳固定，故 $v=-r\omega$，求导得 $a=-r\alpha$。
\SolutionFigure{ex49}{例49的计算示意}
\Step{列平动与绕质心转动方程}
下降取正，绳张力向上，且关于质心产生顺时针矩：
\[
ma=mg-T,\qquad I_C\alpha=-Tr,\qquad I_C=\frac12mr^2.
\]
将 $\alpha=-a/r$ 代入第二式，得 $T=I_Ca/r^2=ma/2$；再代入第一式：
\[
\frac32ma=mg\quad\Longrightarrow\quad a=\frac23g,\quad T=\frac13mg,
\quad\alpha=-\frac{2g}{3r}.
\]
\Step{功与能量的计算}
接触点速度为零，张力对盘的总功率为零。下降高度 $x$ 后
\[
mgx=\frac12mv^2+\frac12I_C\frac{v^2}{r^2}=\frac34mv^2.
\]
由此 $v^2=4gx/3$，再用 $v^2=2ax$ 得到相同的 $a$。
\Step{瞬心取矩的适用条件}
这一几何中接触点到质心距离恒为 $r$，质心仅竖直平移，关于接触瞬心的方程可写为
$-mgr=(I_C+mr^2)\alpha$，故仍有 $\alpha=-2g/(3r)$。
接触瞬心取矩的简式依赖本题的约束，不能无条件用于加速移动的任意取矩点。
本题只有一条绳段承重；若换成两绳段的动滑轮，受力及约束都须重新建立。

% END DETAILED ANSWER ex49
\TeachingFigure{c3-diagram-42}{固定绳段使绳盘接触点成为速度瞬心}
% BEGIN ORIGINAL FIGURES
\OriginalFigure[6.2cm]{c3-5-p019-unwinding-disc}{固定轻绳上退绕下落的圆盘；原图重量 $P$ 对应正文 $mg$}
% END ORIGINAL FIGURES


滚动阻力与上述静摩擦不同。
接触变形使法向合力的作用线相对圆心向前偏移距离 $b$，
产生阻碍转动的矩 $M_{\text{滚}}=Nb$。
% BEGIN ORIGINAL FIGURES
\OriginalFigure[6.2cm]{c3-5-p020-rolling-resistance}{接触变形使法向力作用线向前偏移，产生滚动阻力矩}
% END ORIGINAL FIGURES

$b$ 有长度量纲，等效切向阻力为 $Nb/R$；
若另定义无量纲滚阻系数，则是 $b/R$。
滚阻一般耗散机械能，不能把它与不做总功的理想静摩擦混为一谈。
在存在驱动、滑动与接触变形时，
要同时保留实际切向力和滚阻矩，避免重复计数同一经验阻力。


\subsubsection*{斜面三法：同一约束下逐步消元}
令 $\vect e_s$ 沿斜面向右下，$\vect e_n$ 为向上法向，
仍取 $\vect k$ 指向纸外。圆心速度为 $v\vect e_s$，
接触位置相对圆心为 $-a\vect e_n$，因此
\[
\vect v_P=v\vect e_s+\omega\vect k\times(-a\vect e_n)
=(v+a\omega)\vect e_s.
\]
固定斜面纯滚动要求 $v+a\omega=0$，即
$\omega=-v/a,\alpha=-\dot v/a$。
圆柱向右下运动时有向角速度为负，所有三种方法都使用这个同一约束。
记 $\kappa=k^2/a^2=I_C/(ma^2)$。

第一法从整个刚体的能量入手。
质心下移斜距 $x$，势能减少 $mgx\sin\beta$；
平动能与转动能之和为
\[
T=\frac12mv^2+\frac12mk^2\frac{v^2}{a^2}
=\frac12m(1+\kappa)v^2.
\]
所以由静止释放后
\[
v^2=\frac{2g\sin\beta}{1+\kappa}x,
\qquad \ddot x=\frac{g\sin\beta}{1+\kappa}.
\]
转动能占总动能的比例为 $\kappa/(1+\kappa)$，
平动能占 $1/(1+\kappa)$。
相同质量与半径下，惯量较大的滚动物体把更多能量分给转动，
因而相同下降高度时质心速度较小。

第二法分别列质心平移和中心转动。
假定摩擦有向大小 $f$ 向上坡，方程为
\[
m\ddot x=mg\sin\beta-f,
\qquad I_C\alpha=-fa,
\qquad \ddot x=-a\alpha.
\]
后两式给 $f=I_C\ddot x/a^2=m\kappa\ddot x$。
代入平移式得到 $(1+\kappa)\ddot x=g\sin\beta$，
同时得到 $f=mg\kappa\sin\beta/(1+\kappa)$。
计算结果 $f>0$ 验证了预设的上坡方向。
沿法向圆心没有加速度，故 $N=mg\cos\beta$。
法向没有加速度这一结论针对直斜面；曲槽内滚动时不能继续使用它。

第三法关于瞬心取矩。
接触力 $N,f$ 都作用于 $P$，对 $P$ 的矩为零。
重力关于 $P$ 的矩是 $-mga\sin\beta$，
$I_P=I_C+ma^2=ma^2(1+\kappa)$。
因 $PC=a$ 恒定，移动瞬心修正项为零，得到
\[
ma^2(1+\kappa)\alpha=-mga\sin\beta.
\]
再用 $\ddot x=-a\alpha$，恢复相同线加速度。
瞬心法最少使用接触力未知量，但接触力仍须回到质心合力式求出，
以检查所假定纯滚动能否由实际摩擦维持。

\subsubsection*{静摩擦的零总功与非零能量传递}
在惯性系中，静摩擦施力的接触材料点瞬时速度为零，所以总功率为零。
若分拆平动和转动，摩擦对质心平动的功率是 $-fv$，
对中心转动的功率则为
\[
(-fa)\omega=(-fa)(-v/a)=fv.
\]
两项相消，但各自不为零。
因此静摩擦降低重力全部转为平动能的比例，同时增加转动能，
并没有消耗理想滚动系统的总机械能。
“静摩擦不做总功”指作用点的速度与功率，不能理解为摩擦力不存在。

静摩擦实际取值由动力学决定，只需落在允许区间 $|f|\le\mu_sN$。
故条件是
\[
\mu_s\ge\frac{\kappa}{1+\kappa}\tan\beta.
\]
等号仅表示临界状态，并非每次滚动都取满最大静摩擦。
对环、均匀实心圆柱和均匀实心球，$\kappa$ 依次为 $1,1/2,2/5$，
质心加速度分别为
\[
\frac12g\sin\beta,\qquad\frac23g\sin\beta,\qquad\frac57g\sin\beta,
\]
静摩擦所需系数分别至少为
$\tfrac12\tan\beta,\tfrac13\tan\beta,\tfrac27\tan\beta$。
这些比较以相同的无滑动条件、忽略滚阻和同一斜面为前提。

若可用静摩擦不足，应重新采用有滑动的接触方程。
例如从静止开始，下坡接触滑动时 $f=\mu_kmg\cos\beta$，有
\[
\dot v=g(\sin\beta-\mu_k\cos\beta),\qquad
\alpha=-\frac{\mu_kg\cos\beta}{\kappa a}.
\]
接触相对滑动速度是 $u=v+a\omega$，其增长率为
\[
\dot u=g\left[\sin\beta-\mu_k\cos\beta
\left(1+\frac1\kappa\right)\right].
\]
若这个量为正，所选下坡滑动方向自洽，摩擦始终向上坡。
机械能消耗率是 $f|u|$，而不再是零。
不能把纯滚动的 $a\omega=-v$ 约束与正在滑动的动摩擦模型同时强加。

\subsubsection*{退绕盘三法：固定绳的速度约束与张力功率}
在上图中，绳切点位于盘心左侧。圆心下降速度为 $v>0$，
$\vect v_C=-v\vect j$、$\vect r_{CP}=-r\vect i$。
固定绳段的速度为零，接触无滑动要求
\[
0=-v\vect j+\omega\vect k\times(-r\vect i)
=-(v+r\omega)\vect j.
\]
故 $\omega=-v/r$，下降加速度 $a$ 与有向角加速度满足 $\alpha=-a/r$。
保留一般 $I_C=\kappa mr^2$，质心和转动方程为
\[
ma=mg-T,\qquad I_C\alpha=-Tr.
\]
后一式与约束给 $T=\kappa ma$，前式于是给
\[
a=\frac{g}{1+\kappa},\qquad
T=\frac{\kappa}{1+\kappa}mg.
\]
圆盘 $\kappa=1/2$，即恢复 $2g/3$ 和 $mg/3$。
它比自由落体慢，是因为重力输入同时增加转动动能，
并不是盘上承受两段绳的双倍张力。

能量法逐项为
\[
mgx=\frac12mv^2+\frac12I_C\frac{v^2}{r^2}
=\frac12m(1+\kappa)v^2.
\]
固定绳对盘的总功率为零：
张力对质心平移的功率是 $-Tv$，对转动的功率是
$(-Tr)\omega=Tv$，二者相消。
因此理想绳不耗散机械能，但通过张力改变平动与转动的分配。

瞬心法中 $\vect r_{PC}=r\vect i$，
重力矩为 $-mgr$，张力矩为零。
$I_P=I_C+mr^2$、$PC=r$ 恒定，所以
\[
I_P\alpha=-mgr,
\]
再用 $a=-r\alpha$ 即得同一结果。
接触材料点的加速度却为
\[
\vect a_{P,\text{材料}}=
-a\vect j+\alpha\vect k\times(-r\vect i)-\omega^2(-r\vect i)
=\frac{v^2}{r}\vect i.
\]
两个竖直项相消，剩余向右的法向加速度。
它与 $\vect r_{PC}$ 平行，所以材料点取矩式的加速度叉积为零。
这再次给出可以使用简化瞬心力矩式、但材料瞬心加速度不为零的例子。

\subsubsection*{滚动阻力矩如何进入三种方程}
设圆柱正下坡滚动，法向合力向前偏移 $b>0$，
等效阻转力偶为逆时针的 $Nb$，与顺时针 $\omega$ 相反。
保留切向接触力 $f$，有
\[
m\ddot x=mg\sin\beta-f,\qquad
I_C\alpha=-fa+Nb,\qquad \ddot x=-a\alpha.
\]
因此
\[
(1+\kappa)\ddot x=g\left(\sin\beta-\frac ba\cos\beta\right),
\]
\[
f=\frac{mg}{1+\kappa}
\left(\kappa\sin\beta+\frac ba\cos\beta\right).
\]
关于瞬心取矩也给
$I_P\alpha=-mga\sin\beta+Nb$，与上式一致。
阻转力偶的功率为 $Nb\omega=-Nbv/a$，
所以圆心滚过斜距 $x$ 时，相应耗散为 $Nbx/a$。
$b/a$ 是无量纲比值，$b$ 本身具有长度量纲。

在这一方向的持续滚动中，要有正的下坡加速度，需要 $\tan\beta>b/a$；
纯滚动还须检查
\[
\mu_s\ge\frac{\kappa\tan\beta+b/a}{1+\kappa}.
\]
临界坡度以下能否静止，还需滚动阻力在静止时允许的力矩区间，
不能把运动中的恒定反向力偶公式无条件延伸到速度为零或转向改变以后。


\subsection{轮缘速度、加速度与行星轮}\label{sub:c3-7.5}
\example{变速纯滚动轮的四点运动}\label{ex:c3-50}

\QuestionStatement{半径 $R$ 的轮在固定水平面纯滚动，
轮心瞬时速度 $v\vect i$、加速度 $a\vect i$。
轮缘 $A$ 在底部、$B$ 在左侧、$C$ 在顶部、$D$ 在右侧。
取逆时针角速度为正，求四点的速度与加速度。
四点均指当前位于指定位置的材料点；使用统一的基点速度和加速度公式逐点代入。}{ex50}{例50题图}
% BEGIN DETAILED ANSWER ex50
\solution \Step{用纯滚动条件确定角运动}
以图~\ref{fig:solution-ex50}~中当前处于四个方位的材料点为依据，先写出图右侧第一行的几何或动力学关系，再代入其下方的计算式。

取 $\vect i$ 向右、$\vect j$ 向上、$\vect k$ 向外。底部点瞬时无滑动，
$0=v\vect i+\omega\vect k\times(-R\vect j)$，故
$\omega=-v/R$、$\alpha=-a/R$。
对任意轮缘材料点 $P$，统一使用
\[
\vect v_P=v\vect i+\omega\vect k\times\vect r_{OP},\quad
\vect a_P=a\vect i+\alpha\vect k\times\vect r_{OP}-\omega^2\vect r_{OP}.
\]
\SolutionFigure{ex50}{例50的计算示意}
\Step{逐点代入相对位置求速度}
四点的相对位置依次为 $-R\vect j,-R\vect i,R\vect j,R\vect i$，所以
\[
\vect v_A=0,\quad \vect v_B=v(\vect i+\vect j),\quad
\vect v_C=2v\vect i,\quad \vect v_D=v(\vect i-\vect j).
\]
例如左点 $B$ 的自转速度向上，和轮心的向右速度合成，故两分量均为正。
\Step{分别计算平移、切向和法向加速度}
以 $B$ 为例：轮心项为 $a\vect i$，切向项为 $a\vect j$，法向项为 $(v^2/R)\vect i$。四点结果为
\[
\vect a_A=\frac{v^2}{R}\vect j,\qquad
\vect a_B=\left(a+\frac{v^2}{R}\right)\vect i+a\vect j,
\]
\[
\vect a_C=2a\vect i-\frac{v^2}{R}\vect j,\qquad
\vect a_D=\left(a-\frac{v^2}{R}\right)\vect i-a\vect j.
\]
\Step{瞬心与匀速极限}
$a=0$ 时轮心无加速度，但轮缘仍有大小 $v^2/R$ 的指向轮心的法向加速度。
底部点的水平平移项和切向项相消，竖直项保留，因此它速度为零而加速度不为零。
不能对“每时刻位于底部的不同材料点速度为零”直接求导，来代替跟踪同一个材料点的加速度。

% END DETAILED ANSWER ex50
\TeachingFigure{c3-diagram-43}{速度瞬心在底部，其材料加速度仍向上}
% BEGIN ORIGINAL FIGURES
\OriginalFigure[6.2cm]{c3-5-p005-wheel-relative}{匀速纯滚动车轮各点的相对速度与向心加速度；此图左右轮缘分别记为 $D,B$，与本节正文的 $B,D$ 对调}
\OriginalFigure[6.2cm]{c3-5-p032-wheel-bottom-acceleration}{轮底接触材料点的水平加速度分量相消，法向分量向上}
\OriginalFigure[6.2cm]{c3-5-p033-wheel-left-acceleration}{轮缘左侧点的相对切向与法向加速度；图中浅色 $\omega$ 为角速度，$\varepsilon$ 对应正文角加速度 $\alpha$}
\OriginalFigure[6.2cm]{c3-5-p034-wheel-velocity-field}{纯滚动车轮的速度场与底部瞬心}
% END ORIGINAL FIGURES


固定曲线轨道上的圆轮平面纯滚动，同样满足
$|v_C|=R|\omega|$：接触点速度为零，而圆心到瞬心仍为 $R$。
轨道曲率改变圆心轨迹与加速度，未改变这个瞬时大小关系。
若轨道本身运动，则须使用接触点速度匹配，
不能再把接触处绝对速度设为零。

\example{固定轮外侧的行星轮}\label{ex:c3-51}

\QuestionStatement{半径 $R$ 的轮固定，半径 $r$ 的行星轮与它外接且纯滚动，
中心 $A$ 绕固定中心 $O$ 以角速度 $\Omega\vect k$ 运动。
求行星轮绝对自转角速度及轮缘上
$AM\perp OA$ 的点 $M$ 的速度大小。
在两个满足垂直条件的轮缘点中，指定 $\vect r_{AM}=r\vect e_t$，其中 $\vect e_t=\vect k\times\vect e_r$。同时给出相对转臂的自转率。}{ex51}{例51题图}
% BEGIN DETAILED ANSWER ex51
\solution \Step{由中心轨迹求平移速度}
以图~\ref{fig:solution-ex51}~中两轮接触点与垂直于 $OA$ 的半径 $AM$为依据，先写出图右侧第一行的几何或动力学关系，再代入其下方的计算式。

外接使 $OA=R+r$。令 $\vect e_r$ 从 $O$ 指向 $A$，$\vect e_t=\vect k\times\vect e_r$，则
$\vect v_A=(R+r)\Omega\vect e_t$。
\SolutionFigure{ex51}{例51的计算示意}
\Step{在真实接触点施加无滑动条件}
行星轮上的接触点 $P$ 满足 $\vect r_{AP}=-r\vect e_r$。设其绝对自转为 $\omega\vect k$，固定大轮接触点速度为零，因此
\[
0=\vect v_P=\vect v_A+\omega\vect k\times(-r\vect e_r)
=[(R+r)\Omega-r\omega]\vect e_t.
\]
故 $\omega=(1+R/r)\Omega$，与转臂同向。
转臂架自身以 $\Omega$ 转动，故相对自转率是
$\omega_{\text{相对}}=\omega-\Omega=(R/r)\Omega$，两者不可混用。
\Step{计算题定轮缘点的速度}
题设选 $\vect r_{AM}=r\vect e_t$，由 $\vect k\times\vect e_t=-\vect e_r$，
\[
\vect v_M=(R+r)\Omega\vect e_t-r\omega\vect e_r
=(R+r)\Omega(\vect e_t-\vect e_r).
\]
两单位方向正交，故 $|\vect v_M|=\sqrt2(R+r)|\Omega|$。
\Step{接触点、对侧点与等半径极限}
把 $\vect r_{AP}=-r\vect e_r$ 代回速度式确为零；对侧点 $+r\vect e_r$ 的速度为 $2\vect v_A$。
若 $R=r$，绝对自转率为 $2\Omega$，相对转臂仅为 $\Omega$。
这些特殊位置说明不能只根据轮心轨迹或周长比，忽略参考架转动来推算绝对自转。

% END DETAILED ANSWER ex51
\TeachingFigure{c3-diagram-44}{绝对自转包含转臂的转动，不能只用相对弧长比}
% BEGIN ORIGINAL FIGURES
\OriginalFigure[6.2cm]{c3-5-p023-planet-arc}{外滚行星轮的接触弧长与相对转角；绝对自转还包含转臂转角}
\OriginalFigure[6.2cm]{c3-5-p023-planet-mechanism}{固定轮外侧行星轮的机构；原图转臂角速度 $\omega$ 对应正文 $\Omega$}
\OriginalFigure[6.2cm]{c3-5-p024-planet-velocities}{行星轮指定点 $M$ 的速度合成；原图 $\omega^{\prime}$ 表示绝对自转角速度}
% END ORIGINAL FIGURES



\subsubsection*{四个轮缘点：把三个加速度来源分别列出}
为避免轮缘顶部 $C$ 与质心符号混淆，本段把轮心记为 $O$。
基点加速度为 $a\vect i$，角速度为 $-v/R$，角加速度为 $-a/R$。
每一点都要加上基点、相对切向和相对法向三项。
相对位置和各相对项逐点为
\[
\begin{array}{c|c|c|c|c}
\text{点}&\vect r_{O\text{点}}&\omega\vect k\times\vect r
&\alpha\vect k\times\vect r&-\omega^2\vect r\\ \hline
A&-R\vect j&-v\vect i&-a\vect i&(v^2/R)\vect j\\
B&-R\vect i&v\vect j&a\vect j&(v^2/R)\vect i\\
C&R\vect j&v\vect i&a\vect i&-(v^2/R)\vect j\\
D&R\vect i&-v\vect j&-a\vect j&-(v^2/R)\vect i
\end{array}
\]
给第三列加 $v\vect i$ 得四点速度；
给最后两列的和加 $a\vect i$ 得四点加速度。
例如底部点的两水平项相消，故即使 $a\ne0$，其材料加速度仍纯竖直向上。
顶部点则两水平项相加，水平加速度为 $2a$，另有向下的向心项。
左点和右点的向心项方向相反，切向项也方向相反。
不能根据两点速度大小相同，就判断它们的加速度也同向或同大小。

在匀速条件 $a=0$ 下，四点加速度都只含相对轮心的向心项：
底部向上、左侧向右、顶部向下、右侧向左，大小均为 $v^2/R$。
此时顶部材料点速度大小为 $2v$，却不能用 $(2v)^2/R$ 求其加速度，
因为它在空间的轨迹不是半径 $R$、中心固定的圆。
$R$ 是相对轮心的位置半径，空间轨迹需要另行计算。

\subsubsection*{跟踪同一个材料点得到摆线}
四点表列的是当前位于四个特殊位置的材料点。
为了追踪同一个点，令 $\gamma$ 为从初始底部位置开始的顺时针累计转角。
无滑动条件给轮心水平位移 $R\gamma$。
初始底部的那个材料点在空间中的坐标为
\[
x=R(\gamma-\sin\gamma),\qquad
 y=R(1-\cos\gamma).
\]
对时间求导，得到
\begin{align*}
\vect v&=R\dot\gamma
\bigl[(1-\cos\gamma)\vect i+\sin\gamma\vect j\bigr],\\
\vect a&=R\ddot\gamma
\bigl[(1-\cos\gamma)\vect i+\sin\gamma\vect j\bigr]
+R\dot\gamma^2
\bigl(\sin\gamma\vect i+\cos\gamma\vect j\bigr).
\end{align*}
在 $\gamma=0$，该点速度为零，加速度为
$R\dot\gamma^2\vect j=v^2\vect j/R$。
在它到达顶部的 $\gamma=\pi$ 时，速度为 $2R\dot\gamma\vect i$，
加速度为 $2R\ddot\gamma\vect i-R\dot\gamma^2\vect j$。
这些结果与瞬时四点表完全一致。
摆线的尖点正是材料点与地面接触的时刻，
尖点处的零速度不表示该点在附近时间保持静止。

任意当前轮缘位置也可统一求。
若 $\vect r=R(\cos\chi\vect i+\sin\chi\vect j)$，则
\[
\vect v=v(1+\sin\chi)\vect i-v\cos\chi\vect j,
\]
\[
\vect a=\left[a(1+\sin\chi)-\frac{v^2}{R}\cos\chi\right]\vect i
+\left[-a\cos\chi-\frac{v^2}{R}\sin\chi\right]\vect j.
\]
代入 $\chi=-\pi/2,\pi,\pi/2,0$ 即分别得到 $A,B,C,D$。
这既是四点结果的核查，也能处理其他轮缘位置。

\subsubsection*{曲线轨道上应区分轮心轨道角速度与绝对自转}
固定轨道纯滚动时，接触处速度为零，
$\vect v_O=-\omega\vect k\times\vect r_{OP}$，故大小总有 $|v_O|=R|\omega|$。
轨道弯曲改变的是接触半径方向和轮心加速度；
$\omega$ 仍指轮子在惯性空间的绝对自转角速度。
若圆心绕某轨道中心转动，则其公转角速度还要乘轮心轨迹半径才能得到 $|v_O|$。
所以公转角速度与自转角速度一般不相等。
移动轨道的接触材料点速度不为零，无滑动关系要写成
\[
\vect v_{P,\text{轮}}=\vect v_{P,\text{轨道}},
\]
而不是把轮接触点的绝对速度设为零。

\subsubsection*{行星轮：绝对转动与相对转动的分别计数}
设 $d=R+r$，$\vect e_r$ 沿 $OA$，$\vect e_t=\vect k\times\vect e_r$。
转臂的角速度为 $\Omega$，故 $\vect v_A=d\Omega\vect e_t$。
接触点 $P$ 在行星轮内侧，$\vect r_{AP}=-r\vect e_r$。
速度合成给
\[
\vect v_P=d\Omega\vect e_t-r\omega\vect e_t.
\]
固定轮在接触位置的速度为零，故 $\omega=d\Omega/r$。
转臂坐标架也以 $\Omega$ 旋转，因此相对于转臂的自转角速度为
\[
\omega_{\text{相对}}=\omega-\Omega=\frac Rr\Omega.
\]
两个结果的差正是一项转臂角速度，不能把相对角速度直接放进惯性基点公式。
转臂转一周后，相对于转臂行星轮转 $R/r$ 周，
相对于固定空间则转 $1+R/r$ 周。
当 $R=r$ 时，相对转臂一周，绝对转动两周。
这里转数指有向累计角位移，不涉及轮上某个标记是否恰好回到某个接触位置。

\subsubsection*{行星轮：任意轮缘点和指定点的速度}
在当前瞬间，把轮缘点写成
\[
\vect r_{AM}=r(\cos\chi\vect e_r+\sin\chi\vect e_t).
\]
用 $\omega=d\Omega/r$ 代入基点式，得到
\[
\vect v_M=d\Omega\bigl[(1+\cos\chi)\vect e_t-\sin\chi\vect e_r\bigr].
\]
所以
\[
|v_M|=2d|\Omega|\left|\cos\frac\chi2\right|.
\]
内侧接触点 $\chi=\pi$ 给零速度，外侧点 $\chi=0$ 给 $2d|\Omega|$，
$AM\perp OA$ 的 $\chi=\pm\pi/2$ 给 $\sqrt2d|\Omega|$。
题图采用 $\chi=\pi/2$，故速度向 $\vect e_t-\vect e_r$。
另一个垂直位置的速度向 $\vect e_t+\vect e_r$，大小相同，方向不同。
指定垂直条件是当前瞬间的位置条件；同一个材料点以后不一定继续满足它。

若 $\Omega$ 随时间变化，角速度约束给
$\alpha=(d/r)\dot\Omega$，轮心加速度为
\[
\vect a_A=d\dot\Omega\vect e_t-d\Omega^2\vect e_r.
\]
任意点加速度必须用
\[
\vect a_M=\vect a_A+\alpha\vect k\times\vect r_{AM}
-\omega^2\vect r_{AM}.
\]
例如恒定 $\Omega$ 且当前 $\chi=\pi/2$ 时，
\[
\vect a_M=-d\Omega^2\vect e_r-\frac{d^2}{r}\Omega^2\vect e_t.
\]
虽然公转与绝对自转的角速度均恒定，材料点仍有两项向心加速度。
使用相对转臂角速度 $R\Omega/r$ 替代绝对 $\omega$ 会漏掉转臂运动，
相对坐标法只有把参考系运动的全部加速度项补齐才与上式一致。


\subsection{曲面滚动与有摩擦的动力学}\label{sub:c3-7.6}
\example{圆槽内圆柱的方程、反力与振动}\label{ex:c3-52}

\QuestionStatement{均匀实心圆柱半径 $r$、质量 $m$，
在半径 $R>r$ 的固定圆槽内无滑动滚动，轴线平行槽轴。
令 $q=R-r$，圆心与槽心连线从竖直向下向右偏转角 $\varphi$。
初始 $\varphi=\varphi_0,\dot\varphi=0$，
取 $0<|\varphi_0|<\pi/2$，忽略滚阻。求运动、接触力及小振动周期。
写出初值能量积分、纯滚动的摩擦条件及接触保持条件；小振动解应满足给定初值。}{ex52}{例52题图}
% BEGIN DETAILED ANSWER ex52
\solution \Step{明确模型与计算依据}
以图~\ref{fig:solution-ex52}~中圆槽心 $O$ 至圆柱心 $C$ 的距离 $q$为依据，先写出图右侧第一行的几何或动力学关系，再代入其下方的计算式。
 取
$\vect e_\rho=(\sin\varphi,-\cos\varphi)$、
$\vect e_t=(\cos\varphi,\sin\varphi)$。
逆时针自转角 $\theta$ 为正，
无滑约束是 $r\dot\theta=-q\dot\varphi$。
切向力 $F$ 按 $\vect e_t$ 取有向值，法向力 $N$ 指向槽心。
\[
mq\ddot\varphi=F-mg\sin\varphi,\quad
mq\dot\varphi^2=N-mg\cos\varphi,\quad
\frac12mr^2\ddot\theta=rF.
\]
消去自转与摩擦，得
\[
\frac32q\ddot\varphi+g\sin\varphi=0,\qquad
F=\frac13mg\sin\varphi.
\]

\SolutionFigure{ex52}{例52的计算示意}
\Step{逐项计算与求解}
总动能为 $\tfrac34mq^2\dot\varphi^2$，能量积分给出
\[
\dot\varphi^2=\frac{4g}{3q}(\cos\varphi-\cos\varphi_0),\qquad
N=\frac{mg}{3}(7\cos\varphi-4\cos\varphi_0).
\]
在允许区间 $|\varphi|\le|\varphi_0|$，
$|F|/N=|\sin\varphi|/(7\cos\varphi-4\cos\varphi_0)$
的最大值在端点，因此全程纯滚要求
\[
\mu_s\ge\frac13\tan|\varphi_0|.
\]

\Step{结果与适用条件}
$\varphi_0=\pi/2$ 时起点法向力为零而切向力非零，
有限摩擦不能维持这个从静止纯滚动的模型。
小振动方程为
$\ddot\varphi+[2g/(3q)]\varphi=0$，故
\[
\tau_0=2\pi\sqrt{\frac{3(R-r)}{2g}}.
\]
大振幅需保留正弦项，周期为
$4\sqrt{3q/(2g)}K(\sin(|\varphi_0|/2))$。
圆心的轨迹半径是 $R-r$，而不是槽半径 $R$。
\Step{小振动初值与接触条件} 线性近似给出 $\varphi(t)\simeq\varphi_0\cos(\sqrt{2g/(3q)}\,t)$。释放端 $N=mg\cos\varphi_0>0$，运动中 $\cos\varphi\ge\cos\varphi_0$，故接触持续受压。最大摩擦比在端点取得，保证全程纯滚动的条件须在释放前核查。

% END DETAILED ANSWER ex52
\TeachingFigure{c3-diagram-45}{槽内滚动的法向力、切向力与圆心轨迹}
% BEGIN ORIGINAL FIGURES
\OriginalFigure[6.2cm]{c3-5-p025-groove-geometry}{圆槽与圆柱的接触几何，圆心轨迹半径为 $R-r$}
\OriginalFigure[6.2cm]{c3-5-p026-groove-forces}{圆槽内圆柱的切向与法向加速度及接触力；原图自转正向按顺时针，正文按逆时针}
% END ORIGINAL FIGURES


\example{双轮提升装置}\label{ex:c3-53}

\QuestionStatement{两个均匀圆盘 $A,B$ 的轴固定，重量为 $P_1,P_2$，
半径为 $r_1,r_2$。轻绳绕过 $B$ 后由 $A$ 收卷，
绳端物体 $C$ 重 $P_3$。绳不可伸长且无滑动，
有效卷绕半径恒定，轮 $A$ 受恒定驱动矩 $M_1$。
求物体的向上加速度及张力。
分别记两轮之间张力为 $T$、悬挂段为 $S$；要求写出提升的驱动力矩条件。}{ex53}{例53题图}
% BEGIN DETAILED ANSWER ex53
\solution \Step{分别隔离两轮与被提升物体}
以图~\ref{fig:solution-ex53}~中两轮之间和悬挂段的两种张力为依据，先写出图右侧第一行的几何或动力学关系，再代入其下方的计算式。

两轮轴固定但有惯量。取收绳、提升对应的角加速度为正，轮间张力为 $T$、悬挂段为 $S$，则
\[
I_A\alpha_1=M_1-Tr_1,\quad
I_B\alpha_2=(T-S)r_2,\quad
\frac{P_3}{g}a=S-P_3.
\]
无伸长、无滑动使 $a=r_1\alpha_1=r_2\alpha_2$。
\SolutionFigure{ex53}{例53的计算示意}
\Step{按有效平动惯性消元}
第一式除以 $r_1$、第二式除以 $r_2$，代入约束得到
\[
T=\frac{M_1}{r_1}-\frac{I_A}{r_1^2}a,\qquad
T-S=\frac{I_B}{r_2^2}a,\quad S=P_3+\frac{P_3}{g}a.
\]
相加消去两张力，有
\[
a=\frac{M_1/r_1-P_3}{P_3/g+I_A/r_1^2+I_B/r_2^2}.
\]
\Step{代入均匀圆盘惯量并求张力}
$I_A=P_1r_1^2/(2g)$、$I_B=P_2r_2^2/(2g)$，所以
\[
a=\frac{2g(M_1/r_1-P_3)}{P_1+P_2+2P_3},\quad
T=\frac{M_1}{r_1}-\frac{P_1a}{2g},\quad S=P_3(1+a/g).
\]
从静止向上加速要求 $M_1>P_3r_1$；等号时静止或匀速均可，具体由初值决定。
\Step{能量关系与参数极限}
若提升 $x$，速度 $v$，由静止开始有
\[
(M_1/r_1-P_3)x=\frac12(P_3/g+I_A/r_1^2+I_B/r_2^2)v^2.
\]
与 $v^2=2ax$ 一致。$r_2$ 消去是均匀圆盘的 $I_B\propto r_2^2$ 所致，第二轮仍通过 $P_2$ 增大有效惯性。
两轮惯量趋零时，退化为 $a=g(M_1/(P_3r_1)-1)$，并有 $T=S$。

% END DETAILED ANSWER ex53
\TeachingFigure{c3-diagram-46}{两个固定轮的转动惯量共同进入提升加速度}
% BEGIN ORIGINAL FIGURES
\OriginalFigure[6.2cm]{c3-5-p028-two-wheel-mechanism}{双轮提升装置的原始绕绳形式}
% END ORIGINAL FIGURES


\example{墙角圆柱的停止转动时间}\label{ex:c3-54}

\QuestionStatement{均匀实心圆柱半径 $r$、重量 $Q$，同时接触左侧竖墙和水平地面。
初始质心静止，角速度为逆时针 $\omega_0>0$。
两处动摩擦系数均为 $\mu>0$，忽略滚阻，求首次停止时间。
取向右、向上为正，画明四个接触力，检查双接触假设并说明停止后方程是否仍适用。}{ex54}{例54题图}
% BEGIN DETAILED ANSWER ex54
\solution \Step{按初始转向确定滑动和摩擦方向}
以图~\ref{fig:solution-ex54}~中墙面和地面的四个接触力为依据，先写出图右侧第一行的几何或动力学关系，再代入其下方的计算式。

逆时针转动时，墙接触处材料点向下滑，故墙摩擦 $f_A$ 向上；底部点向右滑，故地面摩擦 $f_B$ 向左。
墙法向 $N_A$ 向右、地面法向 $N_B$ 向上。两接触保持时圆心被几何约束固定。
\SolutionFigure{ex54}{例54的计算示意}
\Step{联立平动力平衡与动摩擦律}
\[
N_A-f_B=0,\quad N_B+f_A-Q=0,\quad
f_A=\mu N_A,\quad f_B=\mu N_B.
\]
先由水平式 $N_A=\mu N_B$，再代入竖直式：
$(1+\mu^2)N_B=Q$。因此
\[
N_B=\frac Q{1+\mu^2},\quad N_A=\frac{\mu Q}{1+\mu^2},\quad
f_A=\frac{\mu^2Q}{1+\mu^2},\quad f_B=\frac{\mu Q}{1+\mu^2}.
\]
$\mu>0$ 时两法向力都为正，双接触假设成立。
\Step{用摩擦力矩积分角速度}
两摩擦都产生顺时针矩，均匀圆柱惯量 $I_C=Qr^2/(2g)$，所以
\[
I_C\dot\omega=-r(f_A+f_B),\quad
\dot\omega=-\frac{2g\mu(1+\mu)}{r(1+\mu^2)}=-c.
\]
在仍滑动的阶段 $\omega(t)=\omega_0-ct$。首次到零的时刻为
\[
t_* =\frac{r\omega_0(1+\mu^2)}{2g\mu(1+\mu)}.
\]
\Step{停止后的接触状态}
$t=t_*$ 后不能继续用动摩擦律外推成反向加速。静止状态可取 $N_B=Q$，其余三力为零。
$\mu\to0$ 时 $t_*\to\infty$；若初始改为顺时针，地面摩擦将向右，水平平衡需要负墙反力，因而原双接触解失效，须另分析脱离墙面的运动。

% END DETAILED ANSWER ex54
\TeachingFigure{c3-diagram-47}{初始转向决定两个接触是否能够持续保持}
% BEGIN ORIGINAL FIGURES
\OriginalFigure[6.2cm]{c3-5-p029-corner-contacts}{墙角圆柱的两个接触点与初始逆时针转向}
\OriginalFigure[6.2cm]{c3-5-p029-corner-forces}{墙角圆柱的法向力与摩擦力}
% END ORIGINAL FIGURES



\subsubsection*{圆槽：两个不同半径和自转方向的推导}
槽心到圆柱心的距离为 $q=R-r$，而圆柱心到接触点的距离为 $r$。
前者决定公转平移，后者决定自转与摩擦力矩，二者不能互换。
用 $\vect r_{OC}=q\vect e_\rho$、$\dot{\vect e}_\rho=\dot\varphi\vect e_t$，得
\[
\vect v_C=q\dot\varphi\vect e_t,\qquad
\vect a_C=q\ddot\varphi\vect e_t-q\dot\varphi^2\vect e_\rho.
\]
接触点在圆心的外侧，即 $\vect r_{CP}=r\vect e_\rho$，
由于 $\vect k\times\vect e_\rho=\vect e_t$，接触速度为
\[
\vect v_P=(q\dot\varphi+r\dot\theta)\vect e_t.
\]
所以固定槽纯滚动给 $r\dot\theta=-q\dot\varphi$。
当圆柱从右侧向槽底运动，$\dot\varphi<0$，实际自转逆时针 $\dot\theta>0$。
这是内侧滚动的方向关系，与固定轮外侧的行星轮不同。

重力的切向分量为 $-mg\sin\varphi$，径向外分量为 $mg\cos\varphi$。
法向 $N$ 指向槽心，即 $-\vect e_\rho$；有向切向力 $F$ 沿 $+\vect e_t$。
因此质心两式与中心转动式逐项为
\[
mq\ddot\varphi=F-mg\sin\varphi,
\qquad -mq\dot\varphi^2=mg\cos\varphi-N,
\]
\[
I_C\ddot\theta=rF.
\]
后一式的正号来自 $\vect e_\rho\times\vect e_t=\vect k$。
令一般中心惯量为 $I_C=\kappa mr^2$，用
$\ddot\theta=-(q/r)\ddot\varphi$ 消元，得到
\[
F=-\kappa mq\ddot\varphi,\qquad
(1+\kappa)q\ddot\varphi+g\sin\varphi=0.
\]
对实心圆柱 $\kappa=1/2$，即得到前面的 $3q/2$ 系数。
也可由 $F$ 与角加速度关系核查：右侧回落时 $\ddot\varphi<0$，
$F>0$ 指向上坡，提供逆时针自转加速度。
摩擦不是沿圆心运动方向的阻力系数，而是满足无滑动所需的有向接触力。

\subsubsection*{圆槽：非线性积分与大振幅初值解}
总动能为
\[
T=\frac12mq^2\dot\varphi^2
+\frac12\kappa mr^2\frac{q^2}{r^2}\dot\varphi^2
=\frac12m(1+\kappa)q^2\dot\varphi^2.
\]
以槽底为势能零点，$V=mgq(1-\cos\varphi)$。
固定槽接触处速度为零，法向和静摩擦均不做总功，所以
\[
\frac12m(1+\kappa)q^2\dot\varphi^2
+mgq(1-\cos\varphi)=mgq(1-\cos\varphi_0).
\]
从 $\varphi_0>0$ 静止回落的第一段应取
\[
\dot\varphi=-\sqrt{\frac{2g}{(1+\kappa)q}
(\cos\varphi-\cos\varphi_0)}.
\]
圆柱的时间积分为
\[
t=\sqrt{\frac{3q}{4g}}
\int_{\varphi}^{\varphi_0}
\frac{\dd s}{\sqrt{\cos s-\cos\varphi_0}}.
\]
令 $k=\sin(\varphi_0/2)$、$\sin(\varphi/2)=k\sin u$，转换后为
\[
t=\sqrt{\frac{3q}{2g}}[K(k)-F(u,k)].
\]
这给出从初始位置到当前角度的完整非线性时间关系。
四个对称区间相接，得到有限振幅周期
\[
\tau=4\sqrt{\frac{3q}{2g}}K(k),\qquad
\tau=\tau_0\left(1+\frac{\varphi_0^2}{16}+\cdots\right).
\]
不同于自由物理摆可讨论 $|q_0|$ 接近 $\pi$，
槽内这一从静止纯滚模型还受到接触和摩擦条件限制，不能任意取大摆角。
初态 $\varphi_0<0$ 对应镜像解；$\varphi_0=0$ 是静止槽底平衡。
自转角可随公转角积分为
\[
\theta(t)-\theta(0)=-\frac qr[\varphi(t)-\varphi_0].
\]
因此能量与时间积分决定圆心运动后，也同时决定材料点的自转姿态。

\subsubsection*{圆槽：全程接触与最大静摩擦比}
一般惯量下的接触力为
\[
F=\frac{\kappa}{1+\kappa}mg\sin\varphi,
\quad
N=\frac{mg}{1+\kappa}
\bigl[(\kappa+3)\cos\varphi-2\cos\varphi_0\bigr].
\]
圆柱时恢复 $F=mg\sin\varphi/3$、
$N=mg(7\cos\varphi-4\cos\varphi_0)/3$。
在 $|\varphi|\le|\varphi_0|<\pi/2$ 内，法向最小值在转向点，
\[
N_{\min}=mg\cos\varphi_0>0,
\qquad
N_{\max}=\frac{mg}{3}(7-4\cos\varphi_0).
\]
所以整个振动保持单侧接触。
最低点的摩擦恰为零，但法向最大，
这两个极值发生在不同位置，不应独立取摩擦最大值与法向最小值后无条件相除。

令 $u=|\varphi|$、$c_0=\cos|\varphi_0|>0$。
圆柱所需摩擦比为
\[
h(u)=\frac{\sin u}{7\cos u-4c_0},\qquad
h'(u)=\frac{7-4c_0\cos u}{(7\cos u-4c_0)^2}>0.
\]
因此最大比确实在 $u=|\varphi_0|$，给出
\[
\max\frac{|F|}{N}=\frac13\tan|\varphi_0|.
\]
若给定 $\mu_s$，允许的最大静止释放角为
$|\varphi_0|\le\arctan(3\mu_s)$，并仍要求 $|\varphi_0|<\pi/2$。
在转向点摩擦非零、圆心速度为零，静摩擦仍能提供转动加速度。
角度趋于 $\pi/2$ 时，所需比值发散；
恰在 $\pi/2$ 静止释放，$N=0$ 却需要 $|F|=mg/3$，有限摩擦无法满足。
更大的静止释放角甚至要求负法向力，单侧槽接触立即失效。
若摩擦条件不满足，应分别求滑动平移和自转，再检查脱离条件 $N=0$；
不能继续保留纯滚动约束来求接触力。

\subsubsection*{双轮：统一逆时针正方向后的三张受力图}
对图示上方水平绳段，提升物块时该段向左运动。
轮 $A$ 顶部向左意味着逆时针转动；轮 $B$ 右侧向上也意味着逆时针转动。
所以在这幅几何中，两轮均采用逆时针正的角加速度：
\[
\alpha_1=a/r_1,\qquad\alpha_2=a/r_2,
\]
其中 $a$ 是物块向上的有向加速度。
轮 $A$ 顶部的张力指向右，其矩为 $-Tr_1$，驱动矩为 $+M_1$。
轮 $B$ 顶部张力指向左，其矩为 $+Tr_2$；
右侧悬挂段张力向下，其矩为 $-Sr_2$。
各轮重量作用在固定中心，对中心矩为零。
这些叉积号数明确后，才可以联立加速度约束。

\TeachingFigure{c3-diagram-48}{双轮逐对象取矩：顶部两绳力对两轮的矩符号相反}
% BEGIN ORIGINAL FIGURES
\OriginalFigure[6.2cm]{c3-5-p028-two-wheel-forces}{双轮和悬挂物的受力；$T,T^{\prime}$ 大小相等，悬挂段张力为 $S$。原图 $\varepsilon_1,\varepsilon_2$ 分别以各轮提升重物时的转向为正，正文两轮统一取逆时针为正}
% END ORIGINAL FIGURES


\subsubsection*{双轮：张力消元与有效平移惯量}
先把各轮转动方程除以相应半径，再用 $a=r_i\alpha_i$，得到
\[
\frac{I_A}{r_1^2}a=\frac{M_1}{r_1}-T,\quad
\frac{I_B}{r_2^2}a=T-S,\quad
\frac{P_3}{g}a=S-P_3.
\]
三式相加，两个内部张力相消，得
\[
\left(\frac{I_A}{r_1^2}+\frac{I_B}{r_2^2}+\frac{P_3}{g}\right)a
=\frac{M_1}{r_1}-P_3.
\]
括号内三个量都具有质量量纲，称沿绳运动坐标的有效惯量。
任意轮的形状只通过 $I_i/r_i^2$ 进入，
均匀盘的这个比为其质量的一半，才使半径 $r_2$ 在最后式中消去。
若轮 $B$ 换成薄环，同样重量下 $I_B/r_2^2=P_2/g$，
相应加速度会变小，不能继续使用圆盘结果。

求出 $a$ 后，从物块式求 $S=P_3(1+a/g)$，
从轮 $B$ 式求 $T=S+I_Ba/r_2^2$，
再回代轮 $A$ 核查驱动矩。
等价地也可从轮 $A$ 式先求 $T$，但只有把三式全部检查才确认符号相容。
当 $M_1>P_3r_1$，从静止开始向上加速；
当 $M_1=P_3r_1$，$a=0$，可以保持静止或已有的匀速提升；
当 $0\le M_1<P_3r_1$，$a<0$，若初始向上运动则减速，若从静止释放则下降。
有向加速度为负不自动意味着轻绳松弛，仍应检查 $T,S>0$。
在此非负驱动矩与正重量模型中，$a>-g$，悬挂段的 $S$ 保持正值。

\subsubsection*{双轮的能量与合并系统角动量}
令物块向上位移为 $x$，两轮转角增量分别为 $x/r_1,x/r_2$。
驱动做功 $M_1x/r_1$，物块势能增加 $P_3x$。
由静止开始，总能量关系为
\[
\left(\frac{M_1}{r_1}-P_3\right)x
=\frac12\left(\frac{I_A}{r_1^2}+\frac{I_B}{r_2^2}+\frac{P_3}{g}\right)v^2.
\]
这个式子求导即给同一加速度。
张力对三对象的功率为 $-Tv,(T-S)v,Sv$，恰好相消；
两轮中心不动，轮自重没有功率。

也可只合并轮 $B$ 和物块 $C$，对固定轮心 $O_B$ 取矩。
物块沿距轮心横向距离 $r_2$ 的竖直线运动，其角动量为
$(P_3/g)r_2v$，轮角动量为 $I_Bv/r_2$，所以
\[
L_{O_B}=\left(\frac{I_B}{r_2}+\frac{P_3}{g}r_2\right)v.
\]
轮与物块之间悬挂绳力 $S$ 的总矩相消，外矩只剩 $(T-P_3)r_2$，得到
\[
\left(\frac{I_B}{r_2^2}+\frac{P_3}{g}\right)a=T-P_3.
\]
这正是轮 $B$ 与物块平移两式的和。
这里积分的是角动量，不是动能；物块也没有刚接在轮缘随轮作圆周运动。
采用固定点角动量时只需其直线轨道的横向力臂恒为 $r_2$。

\subsubsection*{墙角圆柱的滑动方向与单侧法向}
逆时针 $\omega>0$ 时，地面接触点速度为 $r\omega\vect i$，
墙接触点速度为 $-r\omega\vect j$。
因此地面摩擦向左、墙摩擦向上，二者对圆心的矩都顺时针。
法向分别向上与向右，方向由两个单侧接触面决定。
圆心若保持固定，平面合力必须为零；于是
\[
N_A=f_B,\qquad N_B=Q-f_A.
\]
再用 $f_A=\mu N_A,f_B=\mu N_B$，得到
$(1+\mu^2)N_B=Q$，其余三力依次求出。
所求法向均非负，故在这段转向不变的滑动过程中双接触自洽。

转动方程给常角减速度，记其正的大小为
\[
c=\frac{2g\mu(1+\mu)}{r(1+\mu^2)}.
\]
于是 $\omega=\omega_0-ct$，首次停止时间为 $\omega_0/c$，
停止前转过的角度为
\[
\Delta\theta=\frac{\omega_0^2}{2c}.
\]
初始转动能全部耗散：
\[
\frac12I_C\omega_0^2=r(f_A+f_B)\Delta\theta.
\]
每处法向与对应接触速度正交，不做功；
两处动摩擦功率分别为 $-f_Ar\omega$、$-f_Br\omega$。
虽然质心没有平动，材料接触点仍滑动，所以摩擦仍耗散转动能。
若两面系数不同，换成 $\mu_A,\mu_B$，同样消元给
\[
N_B=\frac Q{1+\mu_A\mu_B},\quad
N_A=\frac{\mu_BQ}{1+\mu_A\mu_B},\quad
\dot\omega=-\frac{2g\mu_B(1+\mu_A)}{r(1+\mu_A\mu_B)}.
\]
它清楚表明地面摩擦先产生墙面压紧力，再使墙摩擦参与制动。

\subsubsection*{墙角圆柱：停止以后与相反初始转向}
在首次 $\omega=0$ 后，两个接触处相对速度均为零，
原来的动摩擦定向关系不再适用。
此时可取 $N_B=Q,N_A=f_A=f_B=0$，圆柱保持静止。
墙仍与圆柱几何相切，但法向载荷可为零。
把停止前的恒负角加速度继续积分，会产生并不存在的自发反向转动。
静摩擦应由停后的平衡需要决定，而不是自动等于 $\mu N$。

若初始顺时针角速度为 $-\omega_0$，地面点向左、墙点向上滑，
摩擦分别向右、向下。
保持固定圆心就需 $N_A+f_B=0$，这要求负墙法向，单侧接触不能实现。
因此墙面立即失去载荷，圆柱向右离墙，问题切换为仅与地面接触的滑动转动。
在这个新阶段，$N_B=Q$、地面摩擦向右，均匀实心圆柱满足
\[
\dot v_C=\mu g,\qquad \dot\omega=\frac{2\mu g}{r},
\qquad v_C(0)=0,\quad\omega(0)=-\omega_0.
\]
地面接触相对速度为
\[
u=v_C+r\omega=-r\omega_0+3\mu gt.
\]
首次无滑动发生于
\[
t_1=\frac{r\omega_0}{3\mu g},\qquad
v_C(t_1)=\frac{r\omega_0}{3},\qquad
\omega(t_1)=-\frac{\omega_0}{3}.
\]
它此后向右匀速纯滚动，并没有停下来。
离墙的圆心位移为 $r^2\omega_0^2/(18\mu g)$，为正，证实分离后的几何自洽。
初始转动能为 $mr^2\omega_0^2/4$，最终总动能为 $mr^2\omega_0^2/12$，
所以这一相反转向模型只耗散 $mr^2\omega_0^2/6$。
初始方向改变的是接触状态与最终运动，不能只在停止时间公式中换一个角速度负号。


\subsection{连杆机构与接触状态的过渡}\label{sub:c3-7.7}
\example{等长曲柄滑块：基点法、投影法与瞬心法}\label{ex:c3-55}

\QuestionStatement{曲柄 $OA$ 与连杆 $AB$ 均长 $l$，
$O$ 固定，滑块 $B$ 在过 $O$ 的水平导轨上运动，
曲柄以恒定逆时针角速度 $\Omega$ 转动。
取右侧装配支路，当前曲柄与导轨夹角 $\varphi=\pi/4$，
求滑块速度与连杆角速度。
分别用基点速度、沿连杆投影和瞬心几何核验，说明装配支路的作用。}{ex55}{例55题图}
% BEGIN DETAILED ANSWER ex55
\solution \Step{先用几何关系直接求导}
以图~\ref{fig:solution-ex55}~中曲柄端 $A$ 与滑块 $B$ 的同杆约束为依据，先写出图右侧第一行的几何或动力学关系，再代入其下方的计算式。

右侧装配支路有 $A=(l\cos\varphi,l\sin\varphi)$，
$B=(2l\cos\varphi,0)$，且 $\dot\varphi=\Omega$。
故 $\vect v_B=-2l\Omega\sin\varphi\vect i$；连杆从 $A$ 向 $B$ 的方位角为 $-\varphi$，所以 $\omega_{AB}=-\Omega$。
在 $\varphi=\pi/4$ 时，$\vect v_B=-\sqrt2l\Omega\vect i$，连杆顺时针转。
\SolutionFigure{ex55}{例55的计算示意}
\Step{以曲柄端为基点复核两个分量}
\[
\vect v_A=l\Omega(-\sin\varphi\vect i+\cos\varphi\vect j),\quad
\vect r_{AB}=l(\cos\varphi\vect i-\sin\varphi\vect j).
\]
在 $\vect v_B=\vect v_A+\omega_{AB}\vect k\times\vect r_{AB}$ 中，竖直分量为
$0=l(\Omega+\omega_{AB})\cos\varphi$，故在本位置 $\omega_{AB}=-\Omega$；水平分量继而给出同一滑块速度。
\Step{沿连杆投影消去相对转动}
令连杆方向 $\vect n=(\cos\varphi,-\sin\varphi)$。
刚性约束要求两端速度沿连杆的投影相等，故
\[
v_{Bx}\cos\varphi=\vect v_A\cdot\vect n
=-2l\Omega\sin\varphi\cos\varphi.
\]
得到 $v_{Bx}=-2l\Omega\sin\varphi$。投影式消去了连杆转速，仍需速度差求它。
\Step{瞬心的距离比例}
过 $A$ 作 $\vect v_A$ 的垂线，即 $OA$ 延长线；过 $B$ 作竖线，两线交于 $P$。
$45^\circ$ 时 $AP=l,BP=\sqrt2l$，而 $|v_A|=l|\Omega|$，所以 $|\omega_{AB}|=|\Omega|$、$|v_B|=\sqrt2l|\Omega|$。
方向由 $A$ 点的速度判定为顺时针，与向量结果一致。另一装配支路或死点位置须重新检查几何与可除因子。

% END DETAILED ANSWER ex55
\TeachingFigure{c3-diagram-49}{等长机构在 $45^\circ$ 位置的速度瞬心}
% BEGIN ORIGINAL FIGURES
\OriginalFigure[6.2cm]{c3-5-p035-equal-link-mechanism}{等长曲柄滑块的机构与曲柄端点速度}
\OriginalFigure[6.2cm]{c3-5-p035-equal-link-velocities}{基点法中的连杆相对速度与滑块绝对速度}
\OriginalFigure[6.2cm]{c3-5-p036-equal-link-center}{等长连杆的瞬心 $P$ 及两端速度}
% END ORIGINAL FIGURES

当曲柄或连杆通过退化直线位置时，
用除以 $\cos\varphi$ 的局部公式须重新判断装配支路；
几何微分与原向量方程比盲目外推更可靠。

\example{不等长曲柄滑块的加速度}\label{ex:c3-56}

\QuestionStatement{曲柄 $OA=R$ 以恒定角速度 $\Omega$ 转动，
当前与水平导轨成 $45^\circ$。
连杆 $AB$ 向右下且与导轨成 $30^\circ$，
$B$ 为水平滑块，求 $a_B$ 和连杆角加速度。
取 $x$ 向右、$y$ 向上，逆时针为正；先由几何和速度约束求 $AB$ 长度及角速度，再求加速度。}{ex56}{例56题图}
% BEGIN DETAILED ANSWER ex56
\solution \Step{用同一高度关系求连杆长度}
以图~\ref{fig:solution-ex56}~中曲柄的 $45^\circ$ 位置和连杆的 $30^\circ$ 位置为依据，先写出图右侧第一行的几何或动力学关系，再代入其下方的计算式。

点 $A$ 高度为 $R\sin45^\circ$，连杆竖直投影为 $L\sin30^\circ$，故
$L=\sqrt2R$。取曲柄逆时针速率为 $\Omega$，
\[
\vect v_A=\frac{R\Omega}{\sqrt2}(-\vect i+\vect j),\quad
\vect r_{AB}=\frac{L}{2}(\sqrt3\vect i-\vect j).
\]
\SolutionFigure{ex56}{例56的计算示意}
\Step{先由速度约束求连杆转速}
由 $v_{By}=0$，
$0=R\Omega/\sqrt2+\omega_{AB}L\sqrt3/2$，所以
$\omega_{AB}=-\Omega/\sqrt3$。
后续加速度中的法向项必须使用这个瞬时转速。
\Step{分别列加速度的两个分量}
曲柄匀速，故 $\vect a_A=-R\Omega^2(\vect i+\vect j)/\sqrt2$。完整基点式为
\[
\vect a_B=\vect a_A+\alpha_{AB}\vect k\times\vect r_{AB}
-\omega_{AB}^2\vect r_{AB}.
\]
竖直约束给出
$0=-R\Omega^2/\sqrt2+\alpha_{AB}L\sqrt3/2+\omega_{AB}^2L/2$，从而
\[
\alpha_{AB}=\frac{2\Omega^2}{3\sqrt3}.
\]
再代入水平分量：
\[
a_{Bx}=-\frac{R\Omega^2}{\sqrt2}+\frac L2\alpha_{AB}
-\frac{\sqrt3L}{2}\frac{\Omega^2}{3}
=-\left(\frac{\sqrt2}{2}+\frac{\sqrt6}{18}\right)R\Omega^2.
\]
\Step{方向与投影关系}
连杆角加速度逆时针，滑块加速度向左。其大小为\[|a_B|\simeq0.8431895R\Omega^2.\]
沿连杆投影时切向项消失，得
\[|a_B|\cos30^\circ=R\Omega^2\cos75^\circ+\frac{\sqrt2}{3}R\Omega^2,\]与分量法一致。
即使曲柄角加速度为零，连杆的角加速度也不为零；漏掉相对法向项还会给出错误数值。

% END DETAILED ANSWER ex56
\TeachingFigure{c3-diagram-50}{位置角给出连杆长度，速度约束先确定连杆角速度}
% BEGIN ORIGINAL FIGURES
\OriginalFigure[6.2cm]{c3-5-p037-unequal-link-mechanism}{不等长曲柄滑块在 $45^\circ$、$30^\circ$ 位置的几何}
\OriginalFigure[6.2cm]{c3-5-p038-unequal-link-acceleration}{不等长连杆的加速度分解；相对切向箭头为假定正向，计算结果可为负值}
% END ORIGINAL FIGURES


\example{球与活动板由滑动过渡到纯滚动}\label{ex:c3-57}

\QuestionStatement{均匀实心球质量 $m$、半径 $r$，
放在质量 $M$ 的水平长板上，板下支承光滑。
初始球质心静止且无自转，板以速度 $v>0$ 向右运动。
球板间动摩擦系数为 $\mu>0$，忽略滚阻，
板足够长，使球在滑动至首次相对纯滚动的研究时段内始终保持在板上。
求开始相对于板纯滚动的时间、此时各速度以及机械能损失。
求解时同时检查总水平动量和摩擦耗散，明确球的接触点速度与质心速度的区别。}{ex57}{例57题图}
% BEGIN DETAILED ANSWER ex57
\solution \Step{分清球、板受力与接触速度}
以图~\ref{fig:solution-ex57}~中球底部接触点与活动板的相对滑动为依据，先写出图右侧第一行的几何或动力学关系，再代入其下方的计算式。

球底部相对板初始向左滑，球受摩擦向右、板向左。地面光滑，竖直无加速度，故接触法向力为 $mg$，滑动摩擦大小 $f=\mu mg$。
逆时针角速度为正，球底部点速度为 $v_C+r\omega$，并非 $v_C-r\omega$。
\SolutionFigure{ex57}{例57的计算示意}
\Step{分别积分球和平板的动力学方程}
实心球 $I_C=2mr^2/5$，有
\[
\dot v_C=\mu g,\quad \dot v_{\text{板}}=-\frac{\mu mg}{M},\quad
\dot\omega=\frac{fr}{I_C}=\frac{5\mu g}{2r}.
\]
用初值 $v_C(0)=\omega(0)=0,v_{\text{板}}(0)=v$，得
$v_C=\mu gt$、$v_{\text{板}}=v-\mu mgt/M$、$r\omega=5\mu gt/2$。
\Step{由首次无相对滑动求过渡时刻}
接触相对速度为
$u=v_C+r\omega-v_{\text{板}}=-v+\mu g(7/2+m/M)t$，其第一次为零时
\[
t_* =\frac{2Mv}{(7M+2m)\mu g}.
\]
记 $D=7M+2m$，此时
\[
v_C^*=\frac{2Mv}{D},\qquad
v_{\text{板}}^*=\frac{7Mv}{D},\qquad r\omega^*=\frac{5Mv}{D}.
\]
它们满足接触约束，且 $mv_C^*+Mv_{\text{板}}^*=Mv$，检查了总水平动量守恒。
\Step{从摩擦功率计算耗散}
球板系统的摩擦总功率为 $f(v_C+r\omega-v_{\text{板}})=fu\le0$。
$|u|$ 从 $v$ 线性减到零，滑动距离为 $vt_*/2$，故
\[
\Delta E=f\frac{vt_*}{2}=\frac{mMv^2}{7M+2m}.
\]
将三个末速度代入 $T_* =mv_C^{*2}/2+I_C\omega^{*2}/2+Mv_{\text{板}}^{*2}/2$，也得到 $Mv^2/2-T_* =\Delta E$。
过渡后无外驱动，可取静摩擦为零，三速度均保持不变；板足够长的条件保证这一时段不发生越界。

% END DETAILED ANSWER ex57
\TeachingFigure{c3-diagram-51}{球与板有相对运动，必须作为两个动力学研究对象}
% BEGIN ORIGINAL FIGURES
\OriginalFigure[6.2cm]{c3-5-p039-ball-board}{移动板与球的接触；球与板须分别列动力学方程}
% END ORIGINAL FIGURES



\subsubsection*{等长机构：装配支路与三个有向速度算法}
取局部装配区间 $-\pi/2<\varphi<\pi/2$，此时 $\cos\varphi>0$。
曲柄端点 $A=(l\cos\varphi,l\sin\varphi)$，
连杆长约束给
\[
(x_B-l\cos\varphi)^2+l^2\sin^2\varphi=l^2.
\]
因此 $x_B=l\cos\varphi\pm l\cos\varphi$。
右侧展开的装配支路取加号，才得到 $x_B=2l\cos\varphi$；
另一根是滑块与 $O$ 重合的折叠构形，不能与本题运动支路混用。
连杆向右下的方向角是 $-\varphi$，故其角速度为 $-\Omega$。

第一种基点算法在两个坐标分量中展开：
\[
\begin{pmatrix}v_B\\0\end{pmatrix}
=l\Omega\begin{pmatrix}-\sin\varphi\\\cos\varphi\end{pmatrix}
+l\omega_{AB}\begin{pmatrix}\sin\varphi\\\cos\varphi\end{pmatrix}.
\]
竖直分量立即给 $\omega_{AB}=-\Omega$，
水平分量再给 $v_B=-2l\Omega\sin\varphi$。
$\varphi=\pi/4$、$\Omega>0$ 时，滑块向左，连杆顺时针。

第二种投影算法取 $\vect e_{AB}=\cos\varphi\vect i-\sin\varphi\vect j$。
曲柄端速度的有向投影为
\[
\vect v_A\cdot\vect e_{AB}
=l\Omega(-\sin\varphi\cos\varphi-\cos\varphi\sin\varphi)
=-2l\Omega\sin\varphi\cos\varphi.
\]
滑块投影为 $v_B\cos\varphi$，令二者相等得到相同 $v_B$。
投影法消去了相对转动项，却没有单独确定 $\omega_{AB}$；
还需由两速度差或另一分量补求它。
在 $45^\circ$ 恰有速度大小关系 $|v_B|=|v_A|/\cos\varphi$，
该简单大小关系不是所有角度都成立。

第三种瞬心算法由两条速度法线构造。
$A$ 的速度垂直 $OA$，所以其法线为直线 $OA$；
$B$ 的速度水平，所以其法线为过 $B$ 的竖线。
交点为
\[
P=(2l\cos\varphi,2l\sin\varphi),\qquad
AP=l,\quad BP=2l|\sin\varphi|.
\]
$\omega_{AB}$ 的大小由 $|v_A|/AP=|\Omega|$ 得到，方向由 $A$ 的实际速度确定。
再用 $\vect v_B=\omega_{AB}\vect k\times\vect r_{PB}$ 求有向滑块速度。
在 $\varphi=0$，$P=B$，滑块瞬时速度为零，但连杆仍以 $-\Omega$ 转动。
若正好经过 $\cos\varphi=0$ 的折叠位置，装配支路和几何约束需重新选择，
不能把除以 $\cos\varphi$ 的式子当作那里仍独立有效的求解关系。

\subsubsection*{等长机构：恒速曲柄不等于零滑块加速度}
在上述装配支路，连杆姿态角为 $-\varphi$。
当 $\dot\varphi=\Omega$ 恒定时，$\alpha_{AB}=0$，
但几何坐标二次求导给
\[
\vect a_B=-2l\Omega^2\cos\varphi\vect i.
\]
也可用完整基点式核查。
曲柄端加速度为 $-l\Omega^2(\cos\varphi\vect i+\sin\varphi\vect j)$；
连杆相对法向加速度为
$-l\Omega^2(\cos\varphi\vect i-\sin\varphi\vect j)$。
两竖直项相消，两水平项相加，正好得到上式。
在 $45^\circ$ 位置，$\vect a_B=-\sqrt2l\Omega^2\vect i$。
在 $\varphi=0$，滑块虽然速度为零，却有 $-2l\Omega^2\vect i$ 的加速度。
这再次展示速度瞬心与加速度零点之间没有必然对应关系。

\subsubsection*{不等长机构：先求当前几何和速度}
以逆时针 $\Omega$ 为曲柄有向角速度，连杆有向姿态角记 $\lambda$。
当前 $\varphi=\pi/4,\lambda=-\pi/6$，由两端高度差
\[
R\sin\varphi+L\sin\lambda=0
\]
得到 $L=\sqrt2R$。水平坐标为
\[
x_B=R\cos\varphi+L\cos\lambda
=\frac{1+\sqrt3}{\sqrt2}R.
\]
当前相对位置为
\[
\vect r_{AB}=\frac{\sqrt6}{2}R\vect i-\frac{\sqrt2}{2}R\vect j,
\qquad
\vect v_A=\frac{R\Omega}{\sqrt2}(-\vect i+\vect j).
\]
完整速度分量式是
\[
v_B=-\frac{R\Omega}{\sqrt2}+\frac{R}{\sqrt2}\omega_{AB},\qquad
0=\frac{R\Omega}{\sqrt2}+\frac{\sqrt6R}{2}\omega_{AB}.
\]
所以
\[
\omega_{AB}=-\frac\Omega{\sqrt3},\qquad
v_B=-\frac{R\Omega}{\sqrt2}\left(1+\frac1{\sqrt3}\right).
\]
必须先求 $\omega_{AB}$，才能在加速度式中代入它的平方。
连杆的向心项使用 $\omega_{AB}^2=\Omega^2/3$，并不使用曲柄的 $\Omega^2$。

\subsubsection*{不等长机构：基点法的两个加速度分量}
恒速曲柄的 $\alpha_{OA}=0$，因此
$\vect a_A=-R\Omega^2(\vect i+\vect j)/\sqrt2$。
将叉积项和向心项逐个写入，得到
\begin{align*}
a_B&=-\frac{R\Omega^2}{\sqrt2}
+\frac R{\sqrt2}\alpha_{AB}
-\frac{\sqrt6R}{2}\omega_{AB}^2,\\
0&=-\frac{R\Omega^2}{\sqrt2}
+\frac{\sqrt6R}{2}\alpha_{AB}
+\frac R{\sqrt2}\omega_{AB}^2.
\end{align*}
竖直式里相对向心项为正，是因为连杆从 $A$ 指向 $B$ 的竖直分量为负。
先代入 $\omega_{AB}^2=\Omega^2/3$，得
\[
\frac{\sqrt6R}{2}\alpha_{AB}
=\frac{2R\Omega^2}{3\sqrt2},
\qquad \alpha_{AB}=\frac{2\Omega^2}{3\sqrt3}>0.
\]
再代入水平式，得到
\[
a_B=-\frac{R\Omega^2}{\sqrt2}
+\frac{2R\Omega^2}{3\sqrt6}-\frac{\sqrt6R\Omega^2}{6}
=-\left(\frac{\sqrt2}{2}+\frac{\sqrt6}{18}\right)R\Omega^2.
\]
当前连杆顺时针转动而角加速度逆时针，表示其顺时针角速度的绝对值正在减小。
滑块的速度和加速度都向左；这两个方向需要分别由计算确定。

\subsubsection*{不等长机构的投影法与几何微分法}
将加速度差投影到连杆方向
$\vect e_{AB}=(\sqrt3\vect i-\vect j)/2$，切向项消去，得
\[
a_B\cos30^\circ
=\vect a_A\cdot\vect e_{AB}-\omega_{AB}^2L
=-R\Omega^2\cos75^\circ-\frac{\sqrt2}{3}R\Omega^2.
\]
因此 $a_B<0$，并给同样的 $0.8431895R\Omega^2$ 大小。
由这个投影先求滑块加速度后，再回代任一分量求角加速度，
便是第二种完整算法。

第三种算法始终使用几何约束而不预设加速度方向。
从一般位置关系
\[
R\sin\varphi+L\sin\lambda=0,
\qquad x_B=R\cos\varphi+L\cos\lambda
\]
分别求一、二阶导数：
\begin{align*}
0&=R\Omega\cos\varphi+L\omega_{AB}\cos\lambda,\\
0&=-R\Omega^2\sin\varphi
+L\alpha_{AB}\cos\lambda-L\omega_{AB}^2\sin\lambda,\\
a_B&=-R\Omega^2\cos\varphi
-L\alpha_{AB}\sin\lambda-L\omega_{AB}^2\cos\lambda.
\end{align*}
代入当前两角和 $L=\sqrt2R$，逐式得到相同的角速度、角加速度和滑块加速度。
这一方法也能直接处理曲柄非恒速的情况，只需在二阶关系中保留曲柄角加速度项。

\subsubsection*{用移动瞬心速度式求加速度时必须连同几何一起求导}
当前不等长机构的瞬心仍是 $OA$ 直线与过 $B$ 竖线的交点，
一般位置的高度为 $h_P=x_B\tan\varphi$。
因此有向滑块速度可写为 $v_B=\omega_{AB}h_P$。
正确求导是
\[
a_B=\alpha_{AB}h_P+\omega_{AB}\dot h_P,
\qquad
\dot h_P=v_B\tan\varphi+x_B\sec^2\varphi\,\Omega.
\]
不能只保留 $\alpha_{AB}h_P$，因为瞬心的高度也在变化。
当前 $45^\circ$ 位置，
\[
h_P=\frac{1+\sqrt3}{\sqrt2}R,
\qquad
\dot h_P=\frac{R\Omega}{\sqrt2}\left(1+\frac5{\sqrt3}\right).
\]
代入已求的 $\omega_{AB},\alpha_{AB}$，再一次得到
\[
a_B=-\frac{R\Omega^2}{\sqrt2}\left(1+\frac1{3\sqrt3}\right).
\]
它与前面根式完全相同。
瞬心方法在速度层面非常简洁，但若用到加速度，
必须求几何瞬心的移动或材料瞬心的加速度，不能把它当成固定中心。

\subsubsection*{球板：接触相对速度决定摩擦，不能只看球心速度}
板向右，球初始不动且不转，故球底材料点相对板向左滑。
球的动摩擦向右，其关于球心的矩为 $+fr$，所以球逆时针加速。
竖直方向球心加速度为零，$N=mg$，球板间作用与反作用的水平力为 $+f,-f$。
板下的光滑支承不给水平外力，因此球加板的水平总动量守恒。

分别积分动力学，得到三个线性变化的速度。
用 $u$ 表示球底材料点相对板的有向滑动速度，则
\[
u=v_C+r\omega-v_{\text{板}}
=-v+\mu g\left(\frac72+\frac mM\right)t.
\]
在 $0\le t<t_*$ 内，$u<0$，所选摩擦向右一直自洽。
首次 $u=0$ 给出 $t_*$，只表示接触处无相对滑动，
并不表示球心与板同速，也不表示球已经绕固定地面纯滚动。
把终态关系写成 $v_C=r\omega$ 会使用错误的参考系和接触约束。

终态也可用总摩擦冲量 $J=\int_0^{t_*}f\dd t$ 单独核查。
平移冲量和中心角冲量分别给
\[
v_C^*=\frac Jm,\qquad
v_{\text{板}}^*=v-\frac JM,\qquad
r\omega^*=\frac{r^2J}{I_C}=\frac{5J}{2m}.
\]
代入首次无滑动条件，
\[
J\left(\frac1m+\frac{5}{2m}+\frac1M\right)=v,
\qquad J=\frac{2mMv}{7M+2m}.
\]
因此三个终态速度并不依赖动摩擦系数；
摩擦系数决定实现同一所需冲量的时间和位移。
只有 $\mu>0$、接触持续且板够长时，才实际能达到这个终态。
$\mu=0$ 时 $J$ 始终为零，不能把有限摩擦终态公式当成零摩擦的有限时间结果。

\subsubsection*{球板：分对象功率与总机械能耗散}
球的平动功率为 $fv_C$，转动功率为 $fr\omega$，
板的功率为 $-fv_{\text{板}}$。
相加得到
\[
\dot E_{\text{机械}}=f(v_C+r\omega-v_{\text{板}})=fu=-f|u|.
\]
所以两个接触力虽然是系统内部力，仍可以因相对滑动消耗机械能。
动量守恒要求内部力相消，能量变化还要比较两施力材料点的速度。
用球心位移单独求 $fx_C$，只能求球平动能的增加，不能表示总热耗散。

因为 $|u|$ 从 $v$ 匀减到零，累计接触滑动距离为
\[
s_{\text{滑}}=\int_0^{t_*}|u|\dd t=\frac12vt_*.
\]
于是 $\Delta E=fs_{\text{滑}}$，给出前面的耗散结果。
再直接核查终态动能。令 $D=7M+2m$，则
\begin{align*}
T_{\text{球平动}}^*&=\frac{2mM^2v^2}{D^2},\\
T_{\text{球转动}}^*&=\frac{5mM^2v^2}{D^2},\\
T_{\text{板}}^*&=\frac{49M^3v^2}{2D^2}.
\end{align*}
相加得 $T^*=7M^2v^2/(2D)$，初始能量是 $Mv^2/2$，故
\[
T(0)-T^*=\frac12Mv^2-\frac{7M^2v^2}{2D}
=\frac{mMv^2}{D}.
\]
动量方法、摩擦功率方法和动能差三种核查得到相同终态与耗散。
因为系统总水平动量不变，换到另一个匀速惯性参考系时，
两时刻能量同时改变相同的平动修正，总能量损失仍相同。

\subsubsection*{球板：球心相对位移与接触滑动距离并不相同}
令初始球心和板的水平坐标各取其初值为零，则滑动阶段
\[
x_C=\frac12\mu gt^2,\qquad
x_{\text{板}}=vt-\frac12\frac{\mu mg}{M}t^2.
\]
球心相对于板的位移为
\[
\Delta x_{\text{心/板}}
=-vt+\frac12\mu g\left(1+\frac mM\right)t^2.
\]
其速度为 $v_C-v_{\text{板}}$，与接触滑动速度 $u$ 相差 $r\omega$。
在整个滑动阶段球心相对于板始终向左，
因为 $v_C-v_{\text{板}}$ 达到零的时间大于 $t_*$。
代入 $t_*$ 得首次无滑动前相对左移的大小
\[
d_*=-\Delta x_{\text{心/板}}(t_*)
=\frac{2M(6M+m)v^2}{(7M+2m)^2\mu g}.
\]
球的累计自转角为 $\theta_*=\tfrac12[5\mu g/(2r)]t_*^2$，故
\[
r\theta_*=
\frac{5M^2v^2}{(7M+2m)^2\mu g},
\qquad
\Delta x_{\text{心/板}}(t_*)+r\theta_*=-s_{\text{滑}}.
\]
这把三种不同的位移清楚联系起来：球心相对板位移、自转弧长和接触滑动距离。
检查是否越过板端，使用球心在板上的接触投影位置，
应比较 $d_*$ 与初始到左端的余量，而不是只比较 $s_{\text{滑}}$。

若初始接触投影到板左端的距离为 $d_0$，
要使球在到达首次相对纯滚动前仍留在板上，需 $d_0>d_*$；
等号表示过渡与到达板端同时发生，边缘接触另有几何问题。
若 $d_0<d_*$，球在滑动阶段就离板，前述终态不能实际达到。
首次无滑动后，静摩擦可取零，球心相对板速度保持
\[
v_C^*-v_{\text{板}}^*=-\frac{5Mv}{7M+2m}<0.
\]
因此有限长板不可能在无限时间内始终托住这个球。
当 $d_0>d_*$ 时，从过渡到球到达左端的额外时间为
\[
\Delta t_{\text{留板}}=
\frac{(d_0-d_*)(7M+2m)}{5Mv}.
\]
题设板足够长应理解为覆盖所研究的滑动与首次过渡时段；
若继续研究离板后的运动，需另给板端几何与后续支承环境。

\subsubsection*{球板：质量极限与无滑动后的接触切换}
在 $M\gg m$ 的极限，板速几乎不变，结果趋于
\[
t_*\to\frac{2v}{7\mu g},\quad
v_C^*\to\frac{2v}{7},\quad
r\omega^*\to\frac{5v}{7},\quad
\Delta E\to\frac{mv^2}{7}.
\]
此时仍有 $v_C^*+r\omega^*=v$，因为参考接触面本身向右运动。
它与静止地面上的球由初始自转过渡到纯滚动是不同的初值模型。
从首次 $u=0$ 起，没有水平外驱动，
设静摩擦为零即使三个速度保持恒定，约束自然继续满足。
若以后给板施加外力或使支承运动，需重新求所需静摩擦并检查上限；
不能在无滑动阶段继续保持原来的定值动摩擦。


\addtocontents{toc}{\protect\newpage}
\clearpage
\section{定点运动、欧拉方程与自由转动}\label{sec:c3-8}
\subsection{瞬时轴、极面与点的加速度}\label{sub:c3-8.1}
固定点 $O$ 的位置不变，姿态有三个自由度。
每一瞬间 $\vect v=\vect\omega\times\vect r$，
所有位于瞬时角速度轴上的体内点速度为零。
瞬时轴在空间扫出的锥面称空间极面，
在本体中扫出的锥面称本体极面。
两极面是几何构造，不能据此认为存在真实接触支承或摩擦。

点的加速度为
\[
\vect a=\vect\alpha\times\vect r
+\vect\omega\times(\vect\omega\times\vect r)
=\vect\alpha\times\vect r-\omega^2\vect r_\perp.
\]
第一项称转动加速度，第二项指向瞬时轴。
若轴方向变化，$\vect\alpha$ 不一定平行 $\vect\omega$，
转动加速度也不一定沿点的实际轨道切线。
% BEGIN ORIGINAL FIGURES
\OriginalFigure[6.2cm]{c3-6-p004-perpendicular-radius}{点到瞬时转轴的垂直矢量；原图 $\vect R$ 对应正文 $\vect r_\perp$}
% END ORIGINAL FIGURES

在 $\omega=0$ 时仍可使用双叉积公式，
不必定义不存在唯一方向的瞬时轴。


\subsubsection*{从姿态矩阵构造瞬时轴和两个极面}
取固定点为原点，体内点的空间位置是
$\vect r(t)=\mathsf R(t)\vect\rho$，其中 $\vect\rho$ 不随时间变。
正交性给出
\[
\dot{\mathsf R}\mathsf R^{\mathsf T}
+\mathsf R\dot{\mathsf R}^{\mathsf T}=0.
\]
因此 $\dot{\mathsf R}\mathsf R^{\mathsf T}$ 是反对称矩阵，可以且只能写成
\[
[\vect\omega]_\times=
\begin{pmatrix}
0&-\omega_z&\omega_y\\
\omega_z&0&-\omega_x\\
-\omega_y&\omega_x&0
\end{pmatrix},\qquad
\dot{\mathsf R}=[\vect\omega]_\times\mathsf R.
\]
这里 $[\vect\omega]_\times\vect r=\vect\omega\times\vect r$。
当 $\omega\ne0$ 时，速度为零的条件是
$\vect\omega\times\vect r=0$，即
$\vect r=\lambda\vect\omega$。这是一根过 $O$ 的直线，
不是一块由静止物质点组成的有限平面。
轴外一点到该直线的垂距为
\[
r_\perp=\left|\vect r-(\vect r\cdot\widehat{\vect\omega})
\widehat{\vect\omega}\right|,
\qquad v=\omega r_\perp.
\]
在相邻时刻，瞬时轴一般改变方向。令
$\vect n_S(t)=\vect\omega(t)/\omega(t)$，
$\vect n_B(t)=\mathsf R^{\mathsf T}(t)\vect n_S(t)$。
空间极面和本体极面分别是参数曲面
\[
\vect x_S(\lambda,t)=\lambda\vect n_S(t),\qquad
\vect x_B(\lambda,t)=\lambda\vect n_B(t).
\]
两曲面的母线都通过顶点 $O$；它们未必是圆锥，也未必为闭曲面。
因为
\[
\dot{\vect n}_B
=\mathsf R^{\mathsf T}
\bigl(\dot{\vect n}_S-\vect\omega\times\vect n_S\bigr)
=\mathsf R^{\mathsf T}\dot{\vect n}_S,
\]
在当前共同母线处，本体极面的相邻母线映入空间后，与空间极面的相邻母线有相同的一阶变化。
当 $\dot{\vect n}_S\ne0$ 时，两曲面的公共切平面由
$\vect n_S$ 和 $\dot{\vect n}_S$ 张成。
这一几何滚动描述由姿态运动产生，并没有额外引入接触力。
若方向恒定，极面退化为直线；若 $\omega=0$，瞬时轴方向不唯一，应回到姿态矩阵。

\subsubsection*{为什么瞬时轴上的零速度不保证零加速度}
对 $\vect v=\vect\omega\times\vect r$ 取惯性系导数，
必须同时求导 $\vect\omega$ 和 $\vect r$：
\[
\dot{\vect v}=\dot{\vect\omega}\times\vect r
+\vect\omega\times\dot{\vect r}
=\vect\alpha\times\vect r
+\vect\omega\times(\vect\omega\times\vect r).
\]
用三重积恒等式，
\[
\vect\omega\times(\vect\omega\times\vect r)
=\vect\omega(\vect\omega\cdot\vect r)-\omega^2\vect r
=-\omega^2\vect r_\perp.
\]
若某点此刻在瞬时轴上，第二项为零，但
$\vect\alpha\times\vect r$ 可以不为零。
它下一瞬间就可能离开新的瞬时轴。
只有恒定方向的固定轴情形，轴上的每一个物质点才在整个运动中保持静止。
另一个退化例是刚体从静止开始：$\omega=0$，
而受外力矩后 $\vect\alpha\ne0$；点的初加速度为
$\vect\alpha\times\vect r$，完全可以非零。

若还存在基点平动，且 $A$ 是刚体固连点，
相同求导给出
\[
\vect v_B=\vect v_A+\vect\omega\times\vect r_{AB},\qquad
\vect a_B=\vect a_A+\vect\alpha\times\vect r_{AB}
+\vect\omega\times(\vect\omega\times\vect r_{AB}).
\]
这正是下面螺旋桨题所用的形式。
不能把飞机轨迹中心 $O$ 和桨心 $A$ 混成同一基点。

\example{飞机转弯时螺旋桨点的速度与加速度}\label{ex:c3-58}

\QuestionStatement{飞机以恒速 $u$ 在半径 $R$ 的水平圆周飞行，
螺旋桨中心为 $A$，桨轴沿飞行切向，以恒定相对自转速率 $\omega_s$ 转动。
点 $B$ 距桨中心为 $l$。
取随飞机转弯的右手轴组：
$\vect i$ 径向向外、$\vect j$ 沿飞行切向、$\vect k$ 竖直向上。
令 $\vartheta$ 从竖直向上朝径向向外增加，
$\vect r_{AB}=l(\sin\vartheta\,\vect i+\cos\vartheta\,\vect k)$，
$\dot\vartheta=\omega_s$。忽略其他俯仰与滚转，求 $B$ 的运动。
假设 $A$ 即沿给定半径 $R$ 的圆周运动，$\vartheta(0)=\vartheta_0$；写出任意时刻的速度、加速度，并检验直线飞行极限。}{ex58}{例58题图}
% BEGIN DETAILED ANSWER ex58
\solution \Step{写明运动轴组的导数}
以图~\ref{fig:solution-ex58}~中桨点 $B$ 在 $\vect i,\vect k$ 平面内的半径为依据，先写出图右侧第一行的几何或动力学关系，再代入其下方的计算式。

转弯角速率 $\Omega=u/R$，随飞机转弯的基矢满足
$\dot{\vect i}=\Omega\vect j$、$\dot{\vect j}=-\Omega\vect i$、$\dot{\vect k}=0$。
桨的总角速度为 $\vect\omega=\omega_s\vect j+\Omega\vect k$，故即使两速率恒定，仍有
$\vect\alpha=-\Omega\omega_s\vect i$。
\SolutionFigure{ex58}{例58的计算示意}
\Step{从桨点相对位置计算绝对速度}
$\vartheta=\vartheta_0+\omega_st$，$\vect r_{AB}=l(\sin\vartheta\vect i+\cos\vartheta\vect k)$。
由 $\vect v_B=u\vect j+\vect\omega\times\vect r_{AB}$，逐项计算得到
\[
\vect v_B=\omega_sl\cos\vartheta\vect i
+(u+\Omega l\sin\vartheta)\vect j-\omega_sl\sin\vartheta\vect k.
\]
第一、三项是绕桨轴的相对速度，切向项包含飞机中心速度及桨点随转弯的牵连速度。
\Step{对空间速度求导，保留基矢变化}
$i$ 分量的求导给出 $-\omega_s^2l\sin\vartheta\vect i+\Omega\omega_sl\cos\vartheta\vect j$；
$j$ 分量的求导给出 $\Omega\omega_sl\cos\vartheta\vect j-\Omega(u+\Omega l\sin\vartheta)\vect i$；
$k$ 分量的求导给出 $-\omega_s^2l\cos\vartheta\vect k$。合并：
\[
\begin{aligned}
\vect a_B={}&-\left[\frac{u^2}{R}+l(\omega_s^2+\Omega^2)\sin\vartheta\right]\vect i\\
&+2\Omega\omega_sl\cos\vartheta\vect j-\omega_s^2l\cos\vartheta\vect k.
\end{aligned}
\]
\Step{刚体加速度公式与参数极限}
同样的结果可由 $\vect a_A+\vect\alpha\times\vect r_{AB}+\vect\omega\times(\vect\omega\times\vect r_{AB})$ 得到。
两个不同来源的切向交叉项相加，产生系数 $2$。
$R\to\infty$ 且 $u$ 固定时，只剩桨自转的向心加速度；$\omega_s=0$ 时，恢复固定在飞机上的点随转弯运动的加速度。

% END DETAILED ANSWER ex58
\TeachingFigure{c3-diagram-52}{飞机随动轴组与桨上点的相对位置}
% BEGIN ORIGINAL FIGURES
\OriginalFigure[6.2cm]{c3-6-p005-propeller-axes}{螺旋桨计算用的随动基矢、叶尖 $B$ 与轨道中心 $O$}
\OriginalFigure[6.2cm]{c3-6-p005-propeller-motion}{螺旋桨自转与飞机沿圆轨道转弯}
% END ORIGINAL FIGURES


\subsubsection*{螺旋桨例题的独立分项计算}
以下把上例中的每一项分别算出。随飞机转弯的三基矢满足
\[
\dot{\vect i}=\Omega\vect j,\quad
\dot{\vect j}=-\Omega\vect i,\quad
\dot{\vect k}=0,\qquad \Omega=u/R.
\]
这是仅随飞机转弯的参考系，不随桨叶绕 $\vect j$ 自转。
桨叶相对该系的角速度是 $\omega_s\vect j$，
该系相对地面的角速度是 $\Omega\vect k$。
两者都是瞬时角速度，故可以相加。
虽然 $\omega_s$ 和 $\Omega$ 的大小恒定，
$\omega_s\vect j$ 的空间方向却变化，所以
\[
\vect\alpha=
\frac{\dd}{\dd t}(\omega_s\vect j+\Omega\vect k)
=-\Omega\omega_s\vect i.
\]

速度先分成基点、自转与转弯三部分：
\[
\begin{aligned}
\vect v_A&=u\vect j,\\
\omega_s\vect j\times\vect r_{AB}
&=\omega_s l(\cos\vartheta\,\vect i-\sin\vartheta\,\vect k),\\
\Omega\vect k\times\vect r_{AB}
&=\Omega l\sin\vartheta\,\vect j.
\end{aligned}
\]
所用叉积为 $\vect j\times\vect i=-\vect k$、
$\vect j\times\vect k=\vect i$、$\vect k\times\vect i=\vect j$。
把三式相加得到上例的速度，不需要先猜一个速度方向。

加速度中的基点项为
\[
\vect a_A=u\dot{\vect j}=-u\Omega\vect i=-u^2\vect i/R.
\]
角加速度项单独给出
\[
\vect\alpha\times\vect r_{AB}
=(-\Omega\omega_s\vect i)\times
l(\sin\vartheta\,\vect i+\cos\vartheta\,\vect k)
=\Omega\omega_s l\cos\vartheta\,\vect j.
\]
再令 $\vect\omega_s=\omega_s\vect j$、
$\vect\Omega=\Omega\vect k$，将双叉积展开为四项：
\[
\begin{aligned}
\vect\omega\times(\vect\omega\times\vect r_{AB})
={}&\vect\omega_s\times(\vect\omega_s\times\vect r_{AB})
+\vect\omega_s\times(\vect\Omega\times\vect r_{AB})\\
&+\vect\Omega\times(\vect\omega_s\times\vect r_{AB})
+\vect\Omega\times(\vect\Omega\times\vect r_{AB}).
\end{aligned}
\]
各项结果是
\[
\begin{aligned}
\vect\omega_s\times(\vect\omega_s\times\vect r_{AB})
&=-\omega_s^2l(\sin\vartheta\,\vect i+\cos\vartheta\,\vect k),\\
\vect\omega_s\times(\vect\Omega\times\vect r_{AB})
&=\omega_s\vect j\times(\Omega l\sin\vartheta\,\vect j)=0,\\
\vect\Omega\times(\vect\omega_s\times\vect r_{AB})
&=\Omega\omega_s l\cos\vartheta\,\vect j,\\
\vect\Omega\times(\vect\Omega\times\vect r_{AB})
&=-\Omega^2l\sin\vartheta\,\vect i.
\end{aligned}
\]
这里两种交叉项并不相等：一个为零，另一个非零。
总的切向交叉贡献还包括前面的角加速度项，因此最终是
$2\Omega\omega_s l\cos\vartheta\,\vect j$。
若先把 $\vect\alpha$ 错当成零，就会恰好丢掉其中一半。

也可在转弯参考系中计算。相对速度和相对加速度为
\[
\vect v_{\mathrm{rel}}=\omega_s l
(\cos\vartheta\,\vect i-\sin\vartheta\,\vect k),\qquad
\vect a_{\mathrm{rel}}=-\omega_s^2\vect r_{AB}.
\]
因 $\Omega$ 恒定，运输加速度公式给出
\[
\vect a_B=\vect a_A+\vect a_{\mathrm{rel}}
+2\vect\Omega\times\vect v_{\mathrm{rel}}
+\vect\Omega\times(\vect\Omega\times\vect r_{AB}).
\]
其中科里奥利项正好是上面两个非零交叉贡献之和。
两种算法只是选择的中间参考系不同，最终加速度必须相同。

可用三个极限检查：$\Omega=0$ 时得到固定桨轴上的圆周运动；
$\omega_s=0$ 时点 $B$ 固连于飞机，只有转弯的运输加速度；
$l=0$ 时回到桨心的 $u\vect j$ 与 $-u^2\vect i/R$。
在 $\vartheta=0$ 时，向轴项仅含竖直的
$-\omega_s^2l\vect k$，但仍有水平基点向心项及切向交叉项。


\subsection{本体主轴上的欧拉动力学}\label{sub:c3-8.2}
关于固定点 $O$ 的主惯量为 $I_1,I_2,I_3$。
本体主轴中的角动量分量为 $I_i\omega_i$，
由运输公式
\[
\vect M_O=\left(\frac{\dd\vect L_O}{\dd t}\right)_B
+\vect\omega\times\vect L_O
\]
得到
\begin{equation}\label{eq:eulerdyn}
\begin{aligned}
I_1\dot\omega_1+(I_3-I_2)\omega_2\omega_3&=M_1,\\
I_2\dot\omega_2+(I_1-I_3)\omega_3\omega_1&=M_2,\\
I_3\dot\omega_3+(I_2-I_1)\omega_1\omega_2&=M_3.
\end{aligned}
\end{equation}
本体密度可以不均匀，只要随时间不变；
关于固定点的本体张量就恒定。
若没有预先选择主轴，更一般的矩阵式为
$\mathsf I\dot{\vect\omega}_B+\vect\omega\times(\mathsf I\vect\omega)=\vect M_B$。

与欧拉运动学联立，方程确定姿态与角速度。
% BEGIN ORIGINAL FIGURES
\OriginalFigure[6.2cm]{c3-6-p009-euler-portrait}{莱昂哈德·欧拉}
% END ORIGINAL FIGURES

需给初始姿态和角速度，并把外力矩换成本体分量。
空间中恒定的力矩，其本体分量未必恒定。
理想固定点反力对 $O$ 不产生矩；
反力本身则由
\[
\vect N=m\vect a_C-\vect F_{\text{主动}}
\]
求出。固定点摩擦偶矩若存在仍必须保留。

三式分别乘 $\omega_i$ 后相加，惯量交叉项互相消去：
\[
\dot T=\vect M\cdot\vect\omega,\qquad
T=\frac12\sum_i I_i\omega_i^2.
\]
主动外力保守、势能无显含时间且约束不做功时，$T+V$ 守恒。
固定点反力因作用点静止而不做功；
驱动矩和阻矩则可以改变机械能。


\subsubsection*{运输公式的逐分量推导}
在本体正交轴组中写
$\vect L=L_1\vect e_1+L_2\vect e_2+L_3\vect e_3$。
每根基矢的空间导数是
\[
\dot{\vect e}_1=\omega_3\vect e_2-\omega_2\vect e_3,
\quad
\dot{\vect e}_2=\omega_1\vect e_3-\omega_3\vect e_1,
\quad
\dot{\vect e}_3=\omega_2\vect e_1-\omega_1\vect e_2.
\]
因此角动量的惯性系导数不仅包含 $\dot L_i$，还包括基矢的导数：
\[
\begin{aligned}
\left(\frac{\dd\vect L}{\dd t}\right)_S
={}&(\dot L_1+\omega_2L_3-\omega_3L_2)\vect e_1\\
&+(\dot L_2+\omega_3L_1-\omega_1L_3)\vect e_2\\
&+(\dot L_3+\omega_1L_2-\omega_2L_1)\vect e_3.
\end{aligned}
\]
$\dot L_i$ 是本体分量函数的普通时间导数。
取关于同一固定点的本体主轴，$L_i=I_i\omega_i$，
且 $I_i$ 不随时间变，第一分量就成为
\[
M_1=I_1\dot\omega_1+
\omega_2I_3\omega_3-\omega_3I_2\omega_2
=I_1\dot\omega_1+(I_3-I_2)\omega_2\omega_3.
\]
第二、第三分量依次是
\[
\begin{aligned}
M_2&=I_2\dot\omega_2+
\omega_3I_1\omega_1-\omega_1I_3\omega_3,\\
M_3&=I_3\dot\omega_3+
\omega_1I_2\omega_2-\omega_2I_1\omega_1.
\end{aligned}
\]
这给出式\eqref{eq:eulerdyn}，也明确了循环下标与符号。
若采用非主轴，惯量矩阵仍在本体中恒定，但
$L_i=\sum_j I_{ij}\omega_j$，不可逐分量替成 $I_i\omega_i$。

角加速度的本体分量恰为 $\dot\omega_i$：
\[
\vect\alpha=
\left(\frac{\dd\vect\omega}{\dd t}\right)_B
+\vect\omega\times\vect\omega
=\sum_i\dot\omega_i\vect e_i.
\]
这并不表示任意矢量的空间导数都等于本体分量导数；
这里只是运输叉积对 $\vect\omega$ 自身为零。

\subsubsection*{功率、能量与支承反力的分别求法}
分别将三条欧拉方程乘以 $\omega_1,\omega_2,\omega_3$。
惯量交叉项的总系数为
\[
(I_3-I_2)+(I_1-I_3)+(I_2-I_1)=0,
\]
而惯量固定时
\[
\sum_i I_i\omega_i\dot\omega_i
=\frac{\dd}{\dd t}\left(\frac12\sum_i I_i\omega_i^2\right).
\]
故总功率就是 $\sum_i M_i\omega_i=\vect M\cdot\vect\omega$。
也可直接对质元计算外力功率：
\[
\int\vect F\cdot\vect v\,\dd m
=\int\vect F\cdot(\vect\omega\times\vect r)\,\dd m
=\vect\omega\cdot\int\vect r\times\vect F\,\dd m.
\]
在这一写法中 $\vect F$ 表示单位质量所受力；若直接用质元力 $\dd\vect F$，
末式相应写成 $\vect\omega\cdot\int\vect r\times\dd\vect F$。
固定点理想反力的作用点速度为零，因而既不贡献关于该点的矩，也不贡献功率。
它却仍可能承担相当大的载荷。

姿态解出后，应以质心加速度求反力。例如只有重力为主动合力时，
\[
\vect r_{OC}=\mathsf R\vect\rho_{OC},\qquad
\vect a_C=\vect\alpha\times\vect r_{OC}
+\vect\omega\times(\vect\omega\times\vect r_{OC}),
\]
\[
\vect N=m\vect a_C-m\vect g.
\]
这里不能用 $\vect N=-m\vect g$，除非质心加速度恰为零。
轴承若还施加偶矩，其偶矩必须进入 $\vect M$；
“固定点”只约束位置，并不自动保证理想无摩擦。

若保守外力的势能为 $V(\text{姿态})$，无显含时间，
其沿真实运动的导数满足 $\dot V=-\vect M_{\text{保守}}\cdot\vect\omega$。
把总矩分成保守矩与其他矩，得到
\[
\frac{\dd}{\dd t}(T+V)
=\vect M_{\text{其他}}\cdot\vect\omega.
\]
这同时说明机械能守恒所需的条件，和给定驱动矩时怎样累计机械能变化。

\subsection{无力矩运动的守恒量与稳定性}\label{sub:c3-8.3}
若关于 $O$ 的外力矩恒为零，则空间角动量 $\vect L$ 恒定，
动能也恒定。在本体分量中，
\[
\sum I_i\omega_i^2=2T,\qquad
\sum I_i^2\omega_i^2=L^2.
\]
两二次曲面的交线约束角速度在本体中的轨迹。
本体角速度分量变化，并不违反空间角动量守恒。
只有关于所取点的总外力矩为零时，才可归为无力矩转动。
均匀重力关于该点的力矩为 $\vect r_{OC}\times m\vect g$。
取矩点位于质心，或 $OC$ 与重力方向共线时，重力矩为零；
其他外矩也可能恰好抵消重力矩。

\example{三根主轴转动的线性稳定性}\label{ex:c3-59}

\QuestionStatement{设 $I_1<I_2<I_3$。绕主轴以速率 $W$ 纯转动，
考察小横向角速度扰动。
在无外力矩且主惯量互不相同的条件下，分别对绕三根主轴的解线性化，求横向扰动频率或增长率并判断稳定性。}{ex59}{例59题图}
% BEGIN DETAILED ANSWER ex59
\solution \Step{从无外力矩的欧拉方程开始}
以图~\ref{fig:solution-ex59}~中三根不同主惯量轴及横向扰动为依据，先写出图右侧第一行的几何或动力学关系，再代入其下方的计算式。

在本体主轴系中
\[
I_1\dot\omega_1=(I_2-I_3)\omega_2\omega_3,\quad
I_2\dot\omega_2=(I_3-I_1)\omega_3\omega_1,\quad
I_3\dot\omega_3=(I_1-I_2)\omega_1\omega_2.
\]
绕任一主轴的纯转动为精确解。线性化时保留 $W$ 与横向小扰动的乘积，舍去两个小扰动之积。
\SolutionFigure{ex59}{例59的计算示意}
\Step{绕最大惯量轴的扰动}
令 $\omega_3=W$，$\omega_1=u,\omega_2=v$，一阶方程为
$I_1\dot u=-(I_3-I_2)Wv$、$I_2\dot v=(I_3-I_1)Wu$。
对第一式求导并代入第二式，得
\[
\ddot u+\nu_3^2u=0,\qquad
\nu_3^2=\frac{(I_3-I_2)(I_3-I_1)}{I_1I_2}W^2>0.
\]
例如 $u(0)=u_0,v(0)=v_0$ 时，
$u=u_0\cos\nu_3t-[(I_3-I_2)W/(I_1\nu_3)]v_0\sin\nu_3t$，扰动有界。
\Step{最小惯量轴与中间惯量轴}
绕第一轴同理得到
\[
\nu_1^2=\frac{(I_2-I_1)(I_3-I_1)}{I_2I_3}W^2>0,
\]
仍为振荡。绕第二轴，令 $\omega_2=W,\omega_1=u,\omega_3=v$，则
$I_1\dot u=-(I_3-I_2)Wv$、$I_3\dot v=-(I_2-I_1)Wu$，于是
\[
\ddot u-\lambda^2u=0,\qquad
\lambda^2=\frac{(I_2-I_1)(I_3-I_2)}{I_1I_3}W^2>0.
\]
通解含 $e^{\lambda t}$，一般小扰动会增长。
\Step{说明稳定性的含义与范围}
最大、最小惯量轴的自由转动在横向角速度的线性意义下稳定，中间惯量轴不稳定。
沿原轴的微小扰动只改变常转速，不产生上述横向恢复或增长。
当扰动增长到不再很小时须解非线性方程；重根、有外力矩或受轴承约束的转子均须另行分析。

% END DETAILED ANSWER ex59
\subsubsection*{三个主轴附近的线性初值问题}
上面的指数/振荡判别可以展开为带初值的完整解。
在一根主轴上取无扰动速率 $W\ne0$，保留横向小分量的一阶项。
沿主轴的速率扰动在一阶下为常量，只有横向的二乘二系统决定稳定性。
设 $I_1<I_2<I_3$。

绕第一轴时，令 $u=\delta\omega_2$、$v=\delta\omega_3$，有
\[
\dot u=a_1v,\qquad \dot v=-b_1u,
\quad
 a_1=\frac{(I_3-I_1)W}{I_2},\quad
 b_1=\frac{(I_2-I_1)W}{I_3}.
\]
特征值为 $\lambda=\pm\mathrm i\sigma_1$，其中
\[
\sigma_1=|W|\sqrt{\frac{(I_3-I_1)(I_2-I_1)}{I_2I_3}}.
\]
给定 $u(0)=u_0,v(0)=v_0$，解为
\[
\begin{aligned}
u(t)&=u_0\cos\sigma_1t+\frac{a_1}{\sigma_1}v_0\sin\sigma_1t,\\
v(t)&=v_0\cos\sigma_1t-\frac{b_1}{\sigma_1}u_0\sin\sigma_1t.
\end{aligned}
\]
$a_1,b_1$ 保留 $W$ 的符号，$\sigma_1$ 则是非负角频率。
因此改变主转动方向会改变相位旋转方向，但不会把稳定轴变成不稳定轴。

绕第三轴时，令 $u=\delta\omega_1$、$v=\delta\omega_2$，
\[
\dot u=-a_3v,\quad\dot v=b_3u,\qquad
 a_3=\frac{(I_3-I_2)W}{I_1},\quad
 b_3=\frac{(I_3-I_1)W}{I_2}.
\]
\[
\sigma_3=|W|\sqrt{\frac{(I_3-I_2)(I_3-I_1)}{I_1I_2}},
\qquad\lambda=\pm\mathrm i\sigma_3,
\]
\[
\begin{aligned}
u(t)&=u_0\cos\sigma_3t-\frac{a_3}{\sigma_3}v_0\sin\sigma_3t,\\
v(t)&=v_0\cos\sigma_3t+\frac{b_3}{\sigma_3}u_0\sin\sigma_3t.
\end{aligned}
\]
两种稳定轴的横向轨迹一般是椭圆，而不是两个分量等幅的圆。

绕第二轴时，令 $u=\delta\omega_1$、$v=\delta\omega_3$，则
\[
\dot u=-a_2v,\quad\dot v=-b_2u,\qquad
 a_2=\frac{(I_3-I_2)W}{I_1},\quad
 b_2=\frac{(I_2-I_1)W}{I_3}.
\]
特征值为实数 $\pm\lambda_2$，
\[
\lambda_2=|W|\sqrt{\frac{(I_3-I_2)(I_2-I_1)}{I_1I_3}}.
\]
初值解为
\[
\begin{aligned}
u(t)&=u_0\cosh\lambda_2t-\frac{a_2}{\lambda_2}v_0\sinh\lambda_2t,\\
v(t)&=v_0\cosh\lambda_2t-\frac{b_2}{\lambda_2}u_0\sinh\lambda_2t.
\end{aligned}
\]
增长项系数为 $[u_0-a_2v_0/\lambda_2]/2$。
某个精确配比的初值可以只激发衰减模态，但任意微小偏离都会激发增长模态，
这正是“不稳定”的含义。
当扰动长到与 $W$ 同量级时，线性解失效，不能用指数式预测无限大的角速度。
完整欧拉方程仍受恒定 $T,L$ 约束。

例如取 $(I_1,I_2,I_3)=(1,2,3)\,\mathrm{kg\,m^2}$，
$W=10\,\mathrm{s^{-1}}$，则
\[
\sigma_1=10/\sqrt3,\quad\sigma_3=10,
\quad\lambda_2=10/\sqrt3\quad (\mathrm{s^{-1}}).
\]
中间轴的增长时间尺度为 $1/\lambda_2$；它随转速增加而缩短。

\subsubsection*{用守恒量解释为什么中间轴特殊}
在角动量分量空间中，恒定 $L$ 给出球面
$L_1^2+L_2^2+L_3^2=L^2$，而
\[
T=\frac12\left(\frac{L_1^2}{I_1}
+\frac{L_2^2}{I_2}+\frac{L_3^2}{I_3}\right).
\]
在第 $j$ 主轴附近消去 $L_j^2$，得到
\[
T=\frac{L^2}{2I_j}
+\frac12\sum_{k\ne j}
\left(\frac1{I_k}-\frac1{I_j}\right)L_k^2.
\]
$I_1$ 轴是这个球面上的能量极大，$I_3$ 轴是极小；
守恒能量限制运动只能留在相应极值附近。
$I_2$ 轴的两个二次系数异号，形成鞍点，故允许沿一方向离开。
这里能量极大也能稳定，因为讨论的是无耗散、固定角动量的运动，
不能把“能量越低越稳定”的直觉无条件搬来。
当两个主惯量相等时，某一特征频率退化为零，
须重新处理对称性；上述互不相等的三轴判别不能直接套用。

\subsection{对称自由陀螺的完整解析解}\label{sub:c3-8.4}
设 $I_1=I_2=A,I_3=C$，外矩为零。
轴向分量恒定，记 $\omega_3=s$。
令有向频率 $n=(C-A)s/A$，
前两条欧拉方程化为
$\dot\omega_1=-n\omega_2,\dot\omega_2=n\omega_1$，
故
\[
\omega_1=w\cos(nt+\delta),\quad
\omega_2=w\sin(nt+\delta),\quad \omega_3=s.
\]
$w,\delta,s$ 由初值决定。
$C>A$ 与 $C<A$ 对应不同的本体分量旋转方向，
$C=A$ 时所有角速度分量恒定。

取空间 $\zeta$ 轴沿恒定角动量 $\vect L$。
由本体投影，
\[
L\sin\theta\sin\psi=A\omega_1,\quad
L\sin\theta\cos\psi=A\omega_2,\quad
L\cos\theta=Cs.
\]
因此 $\theta=\theta_0$，
$L\sin\theta_0=Aw$、$L\cos\theta_0=Cs$。
联立运动学，得到
\begin{equation}\label{eq:freegyro}
\theta=\theta_0,\qquad
\varphi=\varphi_0+\frac LA t,\qquad
\psi=\psi_0-\frac{C-A}{A}s\,t.
\end{equation}
以 $L/A$ 写空间进动速率，可避免
$s=0,\cos\theta_0=0$ 时用 $Cs/(A\cos\theta_0)$ 出现 $0/0$。
实际轴向角速度 $s$ 与欧拉角自转率 $\dot\psi$ 仍不同。

令 $\alpha$ 为 $\vect\omega$ 与正对称轴的夹角，
$\beta$ 为 $\vect L$ 与正对称轴的夹角，
对于轴向分量为正的普通分支，
\[
\tan\alpha=\frac ws,\qquad
\tan\beta=\frac{Aw}{Cs}=\frac AC\tan\alpha.
\]
两矢量夹角为 $|\alpha-\beta|$。
$C>A$ 时 $\beta<\alpha$，$\vect L$ 位于对称轴与角速度之间；
$C<A$ 时相反。
一般符号情形应用双参数反正切确定象限，
不能对所有惯量都画同一顺序的三根矢量。

\TeachingFigure{c3-diagram-53}{自由对称陀螺中角动量与角速度的夹角顺序取决于惯量比}
% BEGIN ORIGINAL FIGURES
\OriginalFigure[6.2cm]{c3-6-p015-free-top-disc}{自由对称刚体的角速度与角动量方向；原图角动量记为 $\vect J$，正文记为 $\vect L$}
\OriginalFigure[6.2cm]{c3-6-p015-free-top-section}{自由转动的剖切视图；三矢量夹角顺序须由主惯量比判断}
% END ORIGINAL FIGURES


\example{给定惯量和角速度的自由陀螺}\label{ex:c3-60}

\QuestionStatement{关于固定点的主惯量为
$A=2\,\mathrm{kg\,m^2},C=3\,\mathrm{kg\,m^2}$，
初始轴向分量 $s=10\,\mathrm{s^{-1}}$，横向大小 $w=5\,\mathrm{s^{-1}}$。
求本体频率、空间进动速率与三个方向的夹角。
陀螺轴对称，$I_1=I_2=A,I_3=C$，无外力矩；取初始横向分量沿本体 $\vect e_1$。求本体有向旋转频率、空间进动率及 $\vect\omega,\vect L,\vect e_3$ 两两夹角，并求欧拉角自转率。}{ex60}{例60题图}
% BEGIN DETAILED ANSWER ex60
\solution \Step{积分轴对称自由陀螺的本体方程}
以图~\ref{fig:solution-ex60}~中位于同一子午面内的 $\vect\omega,\vect L,\vect e_3$为依据，先写出图右侧第一行的几何或动力学关系，再代入其下方的计算式。

令 $I_1=I_2=A,I_3=C$。无矩时 $\omega_3=s$ 为常数，横向方程为
\[
\dot\omega_1=-n\omega_2,\quad\dot\omega_2=n\omega_1,\quad
n=\frac{C-A}{A}s.
\]
初始横向分量取 $(w,0)$，故
$\omega_1=w\cos nt,\omega_2=w\sin nt$。
代入数据 $n=(3-2)10/2=5\,\mathrm{s^{-1}}$，这是本体分量的有向旋转率。
\SolutionFigure{ex60}{例60的计算示意}
\Step{由守恒角动量求空间进动率}
\[
L=\sqrt{A^2w^2+C^2s^2}=\sqrt{10^2+30^2}=10\sqrt{10}\,\mathrm{kg\,m^2\,s^{-1}}.
\]
取空间 $\zeta$ 轴沿 $\vect L$。由
$\vect L=A\vect\omega+(C-A)s\vect e_3$，叉乘 $\vect e_3$，得到
$\dot{\vect e}_3=\vect\omega\times\vect e_3=(\vect L/A)\times\vect e_3$。
所以 $\dot\varphi=L/A=5\sqrt{10}\simeq15.81139\,\mathrm{s^{-1}}$。
\Step{求三组夹角，避免混用角速度与角动量}
令 $\alpha=\angle(\vect\omega,\vect e_3)$、$\theta_0=\angle(\vect L,\vect e_3)$，则
\[
\tan\alpha=\frac ws=\frac12,\quad
\tan\theta_0=\frac{Aw}{Cs}=\frac13.
\]
故 $\alpha=26.56505^\circ$、$\theta_0=18.43495^\circ$。
本题 $A,C,s,w$ 均为正，$\vect L$ 与 $\vect\omega$ 在 $\vect e_3$ 的同一侧，二者夹角为
$\alpha-\theta_0=8.13010^\circ$。
\Step{欧拉角自转率}
$\omega_3=\dot\psi+\dot\varphi\cos\theta_0=s$，而
$\dot\varphi\cos\theta_0=Cs/A=15\,\mathrm{s^{-1}}$，所以
$\dot\psi=10-15=-5\,\mathrm{s^{-1}}=-n$。
欧拉角自转率为负而实际轴向分量为正，两者关系由进动的轴向投影保证，没有矛盾。
这里 $\theta$ 恒定、角动量空间固定；角速度仍随体轴绕它进动。

% END DETAILED ANSWER ex60
\subsubsection*{本体解的初相位与欧拉角的衔接}
令 $\omega_1(0)=a$、$\omega_2(0)=b$、$\omega_3(0)=s$。
对称自由陀螺的前两条欧拉方程为
\[
\frac{\dd}{\dd t}(\omega_1+\mathrm i\omega_2)
=\mathrm i n(\omega_1+\mathrm i\omega_2),\qquad n=\frac{C-A}{A}s.
\]
所以不需要先引入三角初相位也可以写出
\[
\begin{aligned}
\omega_1(t)&=a\cos nt-b\sin nt,\\
\omega_2(t)&=a\sin nt+b\cos nt,\qquad \omega_3(t)=s.
\end{aligned}
\]
若 $w=\sqrt{a^2+b^2}>0$，令
$\delta=\operatorname{atan2}(b,a)$，就恢复上面的幅相式。
$w=0$ 时 $\delta$ 没有必要定义，运动是沿对称主轴纯转动。
角速度大小 $\sqrt{w^2+s^2}$ 恒定，但其方向一般并不恒定。

取空间轴沿 $\vect L$，在本体中的投影给出
\[
\sin\psi=\omega_1/w,\qquad\cos\psi=\omega_2/w.
\]
因而对于非退化分支
\[
\psi(t)=\frac\pi2-\delta-nt\pmod{2\pi}.
\]
这说明 $\psi_0$ 与本体横向角速度的初相位不是两个可以独立任意指定的量。
$L\cos\theta_0=Cs$、$L\sin\theta_0=Aw$ 给出
\[
L=\sqrt{A^2w^2+C^2s^2},\qquad
\theta_0=\operatorname{atan2}(Aw,Cs).
\]
对 $0<\theta_0<\pi$，由欧拉运动学中
\[
\omega_1=\dot\varphi\sin\theta_0\sin\psi,
\quad
\omega_2=\dot\varphi\sin\theta_0\cos\psi
\]
立刻有 $\dot\varphi=L/A$。
再由 $s=\dot\psi+\dot\varphi\cos\theta_0$ 得
$\dot\psi=-(C-A)s/A$，而不是把 $s$ 当成 $\dot\psi$。
$\varphi_0$ 决定初始轴在空间中的方位；选空间轴沿角动量只是选择坐标，
不把初始姿态的信息抹去。

\subsubsection*{同一运动在空间和本体中的两个旋转频率}
对称惯量张量可写成
\[
\vect L=A\vect\omega+(C-A)s\vect e_z,
\qquad
\vect\omega=\frac{\vect L}{A}-n\vect e_z.
\]
刚体轴的惯性系导数为
\[
\dot{\vect e}_z=\vect\omega\times\vect e_z
=\frac{\vect L}{A}\times\vect e_z.
\]
所以本体轴在空间中以角速率 $L/A$ 绕恒定角动量进动。
对上一式再求导，$s,n$ 均恒定，得
\[
\dot{\vect\omega}=-n\dot{\vect e}_z
=\frac{\vect L}{A}\times\vect\omega.
\]
空间角速度也绕 $\vect L$ 转动，角频率仍为 $L/A$；
这不是本体中横向分量绕第三轴转动的 $n$。

给定初始单位轴 $\vect e_{z0}$，记
$\vect\ell=\vect L/L$、$\Omega_p=L/A$。
不使用欧拉角奇点也能写出完整空间解：
\[
\begin{aligned}
\vect e_z(t)={}&\vect e_{z0}\cos\Omega_pt
+(\vect\ell\times\vect e_{z0})\sin\Omega_pt\\
&+\vect\ell(\vect\ell\cdot\vect e_{z0})(1-\cos\Omega_pt),\\
\vect\omega(t)={}&\Omega_p\vect\ell-n\vect e_z(t).
\end{aligned}
\]
式中保持 $\vect e_z\cdot\vect\ell=\cos\theta_0$。
当 $s=0$ 时 $n=0$、$\vect\omega=\vect L/A$ 恒定，
但对称轴仍可绕它转动；当 $C=A$ 时惯量各向同性，
所有角速度分量都恒定。$w=0$ 时三根矢量共线，
通常的圆锥进动图退化，欧拉角的 $\varphi,\psi$ 也不必分别可观测。

\subsubsection*{夹角关系的矢量推导与数值计算}
$\vect L$ 与 $\vect\omega$ 的横向投影平行，故两矢量与对称轴共面。
不管 $s$ 的符号如何，直接使用
\[
\alpha=\operatorname{atan2}(w,s),\qquad
\beta=\operatorname{atan2}(Aw,Cs)
\]
确定两角的象限；$A,C>0$ 保证 $\beta=\theta_0$。
也可以独立用点积算夹角 $\eta$：
\[
\cos\eta=
\frac{Aw^2+Cs^2}{\sqrt{w^2+s^2}\sqrt{A^2w^2+C^2s^2}}.
\]
分子等于 $\vect\omega\cdot\vect L=2T>0$，
所以真实 $\eta$ 总在 $[0,\pi/2)$，不应由反正切相减得到负的无向角。
叉积检查为
\[
\sin\eta=\frac{|C-A|\,|s|w}
{\sqrt{w^2+s^2}\sqrt{A^2w^2+C^2s^2}}.
\]
上例 $A=2,C=3,w=5,s=10$ 中，
\[
2T=Aw^2+Cs^2=350\,\mathrm J,
\quad |\vect\omega|=5\sqrt5\,\mathrm{s^{-1}},
\quad \sin\eta=\frac{50}{(5\sqrt5)(10\sqrt{10})}.
\]
得到 $\eta\simeq8.130^\circ$，与 $|\alpha-\beta|$ 相同。
空间进动一周的时间为 $2\pi A/L$，
本体横向分量一周的时间为 $2\pi/|n|$；本例分别约
$0.3974\,\mathrm s$ 和 $1.2566\,\mathrm s$。
这两段时间测量的是不同对象，不能互相代换。

\subsection{潘索几何构造与两条极迹}\label{sub:c3-8.5}
对于一般无力矩刚体，$T$ 与空间 $\vect L$ 固定。
取参考惯量 $I_*$，画附着于本体的惯量椭球
$\vect q^{\mathsf T}\mathsf I\vect q=I_*$。
令瞬时接触位置
\[
\vect q_P=\sqrt{\frac{I_*}{2T}}\vect\omega.
\]
它落在椭球上。椭球在该点的法线为
$\mathsf I\vect q_P=\sqrt{I_*/(2T)}\,\vect L$，
所以切平面垂直于固定的空间角动量。
切平面方程为
\[
\vect L\cdot\vect q=\sqrt{2TI_*}.
\]
平面离中心的距离 $d=\sqrt{2TI_*}/L$ 恒定，
称不可变平面。
这是由动力学守恒量构造出的几何平面，不是实际地面。

接触处的椭球材料点速度与
$\vect\omega\times\vect q_P$ 成正比，为零。
所以附着椭球在不可变平面上作无滑动的几何滚动。
椭球上的接触轨迹称本体极迹，
平面上的轨迹称空间极迹。
一般非对称体两轨迹不是圆；
对称自由陀螺的特殊情形可形成圆形极迹。

\TeachingFigure{c3-diagram-54}{潘索构造：角速度沿接触半径，角动量垂直切平面}
% BEGIN ORIGINAL FIGURES
\OriginalFigure[6.2cm]{c3-6-p016-poinsot-tangent}{惯量椭球、接触半径与不可变平面的距离}
\OriginalFigure[6.2cm]{c3-6-p017-poinsot-ellipsoid}{惯量椭球在不可变平面上作无滑动几何滚动}
% END ORIGINAL FIGURES


\subsubsection*{法向、切平面距离和无滑动的严格对应}
对惯量椭球定义
$F(\vect q)=\vect q^{\mathsf T}\mathsf I\vect q-I_*$。
其梯度是 $\nabla F=2\mathsf I\vect q$。
令 $c_* =\sqrt{I_*/(2T)}$，则
\[
F(c_*\vect\omega)=c_*^2(2T)-I_*=0,
\qquad\mathsf I\vect q_P=c_*\vect L.
\]
切平面由一阶变分 $\dd F=0$ 给出：
\[
(\mathsf I\vect q_P)\cdot(\vect q-\vect q_P)=0.
\]
代入 $\mathsf I\vect q_P=c_*\vect L$，并用
$\vect L\cdot\vect q_P=c_*(2T)$，得到
$\vect L\cdot\vect q=\sqrt{2TI_*}$。
因为 $\vect L$ 的空间方向和 $T$ 都恒定，这个平面在空间中固定。
沿法向 $\vect\ell=\vect L/L$，从 $O$ 到平面的距离为
\[
d=\vect q_P\cdot\vect\ell
=c_*\frac{2T}{L}=\frac{\sqrt{2TI_*}}{L}.
\]
$I_*$ 只是选定的归一化尺度。取 $I_*=1$ 的数值书写时，
也必须保留惯量空间的尺度含义，不能把这个 $d$ 直接当作真实支杆长度。

接触处固连于椭球的材料点，瞬时速度是
$\vect\omega\times\vect q_P=c_*\vect\omega\times\vect\omega=0$。
而每一时刻重新确定的几何接触位置可以运动：
\[
\dot{\vect q}_P=c_*\vect\alpha.
\]
两者不矛盾：前者追踪当前材料点，后者追踪不断换材料点的接触位置。
又因 $\vect L\cdot\vect\alpha=\dot T=0$，
\[
\vect L\cdot\dot{\vect q}_P=0,
\]
几何接触轨迹始终留在不可变平面上。
在本体中追踪 $\vect q_{P,B}$ 得本体极迹，
在空间中追踪 $\vect q_{P,S}$ 得空间极迹；
它们不是同一固定质点在两个坐标系中的普通轨迹。

\subsubsection*{由守恒曲面得到极迹及其退化}
在本体中，$\vect q=c_*\vect\omega$ 把两个守恒曲面化为
\[
\vect q^{\mathsf T}\mathsf I\vect q=I_*,\qquad
\vect q^{\mathsf T}\mathsf I^2\vect q=
\frac{I_*L^2}{2T}.
\]
它们的交线就是本体极迹。两曲面一般不是同心球，
所以不能先验地把交线画成平面圆。
对称情形中 $\omega_3=s$、$\omega_1^2+\omega_2^2=w^2$，
故本体极迹在 $q_3=c_*s$ 的平面上，半径为 $c_*w$。
当 $n\ne0,w>0$ 时它按本体相位完整走过这一圆。

空间中 $|\vect\omega|$ 恒定且它绕 $\vect L$ 转动，
因此空间极迹也是圆，其圆心是 $O$ 在不可变平面上的正交投影。
圆半径为
\[
\rho=|\vect q_P\times\vect\ell|
=c_*\frac{|C-A|\,|s|w}{L}
=d\tan\eta,
\qquad\eta=\angle(\vect\omega,\vect L).
\]
当 $C=A$、$s=0$ 或 $w=0$ 使 $\rho=0$ 时，空间圆退化为点。
其中 $s=0$ 时本体横向相位也停止，不能仍说极迹以有限周期绕圆。
对一般非对称刚体，角速度大小本身会变；
不能用这一个对称特例替代全部潘索运动。
最后，若 $T=L=0$，刚体静止，以上含 $1/\sqrt T$ 或 $1/L$ 的归一化构造不适用。


\clearpage
\section{重陀螺、章动与定常进动}\label{sec:c3-9}
\subsection{轴对称重刚体的模型与重力矩}\label{sub:c3-9.1}
轴对称刚体固定在对称轴上的一点 $O$，
关于 $O$ 的主惯量为 $A,A,C$。
质心在 $+\vect e_z$ 方向，$OC=l$，
空间 $\zeta$ 轴竖直向上，重力为 $-mg\vect e_\zeta$。
固定点理想反力对 $O$ 不产生力矩。
这称拉格朗日重陀螺模型；固定点与质心重合时则退化为无重力矩的自由转动。

由 $\vect r_C=l\vect e_z$，
\[
\vect M_O=-mgl\,\vect e_z\times\vect e_\zeta
=mgl\sin\theta(\cos\psi\,\vect e_x-\sin\psi\,\vect e_y).
\]
势能为 $V=mgl\cos\theta$，零点任意。
与欧拉方程联立，
\[
\begin{aligned}
A\dot\omega_1+(C-A)\omega_2\omega_3&=mgl\sin\theta\cos\psi,\\
A\dot\omega_2+(A-C)\omega_3\omega_1&=-mgl\sin\theta\sin\psi,\\
C\dot\omega_3&=0.
\end{aligned}
\]
本体对称轴分量的角动量守恒，整个角动量矢量却不守恒，
因为重力矩一般不为零。

\TeachingFigure{c3-diagram-55}{重力作用线不通过固定点时，重力对固定点产生力矩}
% BEGIN ORIGINAL FIGURES
\OriginalFigure[6.2cm]{c3-6-p018-heavy-top}{重陀螺的倾斜对称轴、自转与重力}
% END ORIGINAL FIGURES



\subsubsection*{重力矩的方向与欧拉角方程的逐项转换}
以下记 $G=mgl>0$，它的量纲是力矩或能量；它不是万有引力常数。
在已经约定的右手本体轴中，竖直向上单位矢量的分量为
\[
\vect e_\zeta=\sin\theta\sin\psi\,\vect e_x
+\sin\theta\cos\psi\,\vect e_y+\cos\theta\,\vect e_z.
\]
逐项叉乘得到
\[
-G\vect e_z\times\vect e_\zeta
=G\sin\theta\cos\psi\,\vect e_x
-G\sin\theta\sin\psi\,\vect e_y.
\]
重力矩垂直于竖直方向和对称轴，故其竖直投影与轴向投影均为零。
倾角增大意味着质心降低，$V=G\cos\theta$ 减小，因而没有自转时重力应使
$0<\theta<\pi$ 的倾角加速增大；这也是核查正负号的一条直接方法。

令 $s=\omega_3$，由第三条欧拉方程得 $\dot s=0$。
为了把另两式化为倾角方程，先使用运动学的两个组合：
\[
\dot\theta=\omega_1\cos\psi-\omega_2\sin\psi,
\qquad \dot\varphi\sin\theta=\omega_1\sin\psi+\omega_2\cos\psi.
\]
对第一式求导，不得漏掉 $\psi$ 的导数：
\[
\ddot\theta=\dot\omega_1\cos\psi-\dot\omega_2\sin\psi
-\dot\psi(\omega_1\sin\psi+\omega_2\cos\psi).
\]
设 $n=(C-A)s/A$，欧拉方程给出
\[
\dot\omega_1=-n\omega_2+\frac GA\sin\theta\cos\psi,
\qquad
\dot\omega_2=n\omega_1-\frac GA\sin\theta\sin\psi.
\]
所以
\[
\ddot\theta=\frac GA\sin\theta
-(n+\dot\psi)\dot\varphi\sin\theta.
\]
再由 $\dot\psi=s-\dot\varphi\cos\theta$，有
$n+\dot\psi=P/A-\dot\varphi\cos\theta$，从而
\[
A\ddot\theta=A\dot\varphi^2\sin\theta\cos\theta
-P\dot\varphi\sin\theta+G\sin\theta.
\]
第一项来自运动学和惯性，第二项包含轴向角动量，第三项才是重力作用。
$\dot\varphi$ 随时间变化时仍须使用这条方程，不能用定常进动条件替代。

重力的瞬时功率也可直接核对：
\[
\vect M\cdot\vect\omega
=G\sin\theta(\omega_1\cos\psi-\omega_2\sin\psi)
=G\sin\theta\,\dot\theta=-\dot V.
\]
这一步明确了能量积分中势能的符号。
固定点的理想反力作用于速度为零的点，不进入上述功率；
若有轴承摩擦偶矩或驱动矩，其分量与功率必须另行加入。


\subsection{三个守恒量与初始条件}\label{sub:c3-9.2}
令轴向角速度 $s=\omega_3$，
轴向角动量 $P=Cs$，竖直角动量 $K=L_\zeta$。
由于 $M_3=0$，$P$ 恒定；
由于 $\vect M\cdot\vect e_\zeta=0$，$K$ 恒定。
能量为
\begin{equation}\label{eq:heavyintegrals}
\begin{aligned}
P&=C(\dot\psi+\dot\varphi\cos\theta),\\
K&=A\dot\varphi\sin^2\theta+P\cos\theta,\\
E&=\frac A2(\dot\theta^2+\dot\varphi^2\sin^2\theta)
+\frac{P^2}{2C}+mgl\cos\theta.
\end{aligned}
\end{equation}
三个动能平方项都须有 $1/2$。
给定初始 $\theta_0,\dot\theta_0,\dot\varphi_0,\dot\psi_0$，
先由这三式计算常数 $P,K,E$，
再求后续运动。初始姿态角 $\varphi_0,\psi_0$ 还决定空间方位和盘上标记点。

在非极点处，
\[
\dot\varphi=\frac{K-P\cos\theta}{A\sin^2\theta},\qquad
\dot\psi=\frac PC-\dot\varphi\cos\theta.
\]
轴向分量恒定不等于 $\dot\psi$ 恒定；
$\theta$ 变化时进动在轴上的投影也在变化。


\subsubsection*{从动力学逐项建立三个积分}
第一积分直接来自 $C\dot\omega_3=0$：
\[
P=C\omega_3=C(\dot\psi+\dot\varphi\cos\theta)=\text{常量}.
\]
第二积分要注意竖直方向在本体系中也在运动。
将竖直单位矢量的本体分量记作 $\vect e_{\zeta,B}$，
角动量的本体分量为
\[
\vect L_B=(A\omega_1,A\omega_2,C\omega_3).
\]
于是
\[
\dot K=\dot{\vect L}_B\cdot\vect e_{\zeta,B}
+\vect L_B\cdot\dot{\vect e}_{\zeta,B}.
\]
运输公式分别给出
\[
\dot{\vect L}_B=\vect M_B-\vect\omega\times\vect L_B,
\qquad
\dot{\vect e}_{\zeta,B}=-\vect\omega\times\vect e_\zeta.
\]
两项含角速度的标量三重积相消，因此
\[
\dot K=\vect M\cdot\vect e_\zeta=0.
\]
用上面的竖直投影和运动学组合，有
\[
\begin{aligned}
K&=A\sin\theta(\omega_1\sin\psi+\omega_2\cos\psi)
+P\cos\theta\\
&=A\dot\varphi\sin^2\theta+P\cos\theta.
\end{aligned}
\]
这里的守恒是空间固定方向上的投影守恒；
它不表示本体横向分量 $A\omega_1,A\omega_2$ 恒定。
直接求导还得到另一条欧拉角动力学方程：
\[
A\sin^2\theta\,\ddot\varphi
+(2A\dot\varphi\cos\theta-P)\sin\theta\,\dot\theta=0.
\]

第三积分由功率方程和 $\dot T=-\dot V$ 得到。
先逐项展开横向角速度平方：
\[
\begin{aligned}
\omega_1^2+\omega_2^2
={}&(\dot\varphi\sin\theta\sin\psi+\dot\theta\cos\psi)^2\\
&+(\dot\varphi\sin\theta\cos\psi-\dot\theta\sin\psi)^2\\
={}&\dot\varphi^2\sin^2\theta+\dot\theta^2.
\end{aligned}
\]
两交叉项大小相等、符号相反。
于是 $T=A(\dot\theta^2+\dot\varphi^2\sin^2\theta)/2+P^2/(2C)$，
加上 $G\cos\theta$ 得常量 $E$。
这些积分分别利用轴对称性、竖直方向的转动对称性和时间平移对称性；
改变模型条件时应重新检查哪一个对称性仍然存在。

\subsubsection*{由任意初值确定同一条运动}
给定 $0<\theta_0<\pi$ 和三个初始欧拉角速率，先计算
\[
\begin{aligned}
P&=C(\dot\psi_0+\dot\varphi_0\cos\theta_0),\\
K&=A\dot\varphi_0\sin^2\theta_0+P\cos\theta_0,\\
E'&=\frac A2(\dot\theta_0^2+\dot\varphi_0^2\sin^2\theta_0)
+G\cos\theta_0.
\end{aligned}
\]
然后计算 $u_0=\cos\theta_0$、
$\dot u_0=-\sin\theta_0\dot\theta_0$。
$P,K,E'$ 决定允许的章动区间，$\dot u_0$ 决定最初向哪一个转向点运动；
$\varphi_0,\psi_0$ 则给出两相位常数。
若 $\dot u_0=0$，必须由二阶方程或转向点邻域展开判断出发方向，
不能仅在平方根式中随意选择正号。
极点处节点线不存在，应改用轴方向矢量和本体角速度给初值，再取光滑极限。


\subsection{有效势、章动区间与积分解}\label{sub:c3-9.3}
扣除恒定轴向动能，令 $E'=E-P^2/(2C)$，得到
\begin{equation}\label{eq:heavypotential}
\frac A2\dot\theta^2+U(\theta)=E',\qquad
U(\theta)=\frac{(K-P\cos\theta)^2}{2A\sin^2\theta}+mgl\cos\theta.
\end{equation}
第一项由竖直与本体轴向角动量的约束产生，
第二项为重力势能。允许区间须满足 $E'\ge U(\theta)$；
边界 $E'=U(\theta)$ 是章动转向点。
若 $U$ 在极点发散，轴方向不能到达该极点；
若分子同时趋零，须取极限而不能只看分母。

令 $u=\cos\theta$，则
\[
\dot u^2=\frac2A(E'-mgl\,u)(1-u^2)-\frac{(K-Pu)^2}{A^2}
=:f(u).
\]
右边是三次多项式。实际解满足 $-1\le u\le1$ 与 $f(u)\ge0$。
一般运动由积分
\[
t-t_0=\pm\int_{u_0}^u\frac{\dd v}{\sqrt{f(v)}}
\]
确定，再对 $\dot\varphi,\dot\psi$ 积分。
到转向点应换 $\dot u$ 的符号以继续运动。
一般解涉及椭圆积分，不能任意把 $\theta$ 当常量来代替章动。

若章动在 $\theta_{\min},\theta_{\max}$ 间进行，
周期和每周期的进动量为
\[
\tau_\theta=2\sqrt{\frac A2}
\int_{\theta_{\min}}^{\theta_{\max}}\frac{\dd\theta}{\sqrt{E'-U(\theta)}},
\]
\[
\Delta\varphi=2\sqrt{\frac A2}
\int_{\theta_{\min}}^{\theta_{\max}}
\frac{K-P\cos\theta}{A\sin^2\theta}
\frac{\dd\theta}{\sqrt{E'-U(\theta)}}.
\]
这些式子区分章动周期与平均进动速率
$\langle\dot\varphi\rangle=\Delta\varphi/\tau_\theta$。
它们一般不是三个欧拉角都以常速变化的运动。


\subsubsection*{三次多项式的系数、根与允许区间}
由 $\dot u=-\sin\theta\dot\theta$，把能量式乘以 $2(1-u^2)/A$，
再代入竖直角动量约束，逐项展开为
\[
\begin{aligned}
f(u)={}&\frac{2G}{A}u^3
-\left(\frac{2E'}A+\frac{P^2}{A^2}\right)u^2\\
&+\left(\frac{2KP}{A^2}-\frac{2G}A\right)u
+\frac{2E'}A-\frac{K^2}{A^2}.
\end{aligned}
\]
这里各系数的量纲均为 $\mathrm{s^{-2}}$，$u$ 无量纲。
两端值为
\[
f(1)=-\frac{(K-P)^2}{A^2},\qquad
f(-1)=-\frac{(K+P)^2}{A^2}.
\]
对 $G>0$、运动区间不含极点且有两个不同转向点的普通情形，
可把三个实根编号为
\[
-1<u_1<u_2<1<u_3,
\qquad
f(u)=\frac{2G}{A}(u-u_1)(u_2-u)(u_3-u).
\]
若根恰落在极点，严格不等号应改为相应等号并另取极限。
根的编号不是任意的：最高次系数为正，两个物理端点之间 $f>0$，
而 $f(1)<0$，因此第三根位于其右侧。
物理章动范围为
\[
u_1\le u\le u_2,\qquad
\theta_{\min}=\arccos u_2,\quad
\theta_{\max}=\arccos u_1.
\]
转向点满足 $\dot\theta=0$，一般仍有 $\dot\varphi\ne0$；
章动停止一瞬间，不等于刚体停止转动。
求根时应先核对 $f(u_0)=\dot u_0^2$，再排除 $|u|>1$ 的根作为实际倾角端点。

\subsubsection*{带初值的三个求积}
在一段没有越过转向点的运动中，令 $\sigma=\operatorname{sgn}\dot u$，则
\[
\begin{aligned}
t-t_0&=\sigma\int_{u_0}^{u}\frac{\dd v}{\sqrt{f(v)}},\\
\varphi-\varphi_0&=\sigma\int_{u_0}^{u}
\frac{K-Pv}{A(1-v^2)\sqrt{f(v)}}\,\dd v,\\
\psi-\psi_0&=\frac PC(t-t_0)
-\sigma\int_{u_0}^{u}
\frac{v(K-Pv)}{A(1-v^2)\sqrt{f(v)}}\,\dd v.
\end{aligned}
\]
这三个式子一起才给出姿态。
在 $u_1,u_2$ 处换向后，应把已累积的时间和角度作为新积分常数，
继续下一段，而不能把方位角重新归零。
等价地，可使用连续展开的椭圆积分振幅，自动记录经过了多少个半周期。

\subsubsection*{第一类与第三类椭圆积分的明确形式}
定义
\[
D=u_2-u_1,\quad k^2=\frac{D}{u_3-u_1},\quad
\nu=\sqrt{\frac{G(u_3-u_1)}{2A}},\qquad
u=u_1+D\sin^2\chi.
\]
这里左侧最后一个 $u$ 是倾角变量，$\nu$ 是有量纲的频率参数。
在从 $u_1$ 出发向 $u_2$ 运动的半周期内，$0\le\chi\le\pi/2$，
因而
\[
\frac{\dd u}{\sqrt{f(u)}}
=\frac{\dd\chi}{\nu\sqrt{1-k^2\sin^2\chi}}.
\]
定义第一类不完全椭圆积分及第三类不完全椭圆积分
\[
F(\chi,k)=\int_0^\chi\frac{\dd a}{\sqrt{1-k^2\sin^2a}},
\]
\[
\Pi(n;\chi,k)=\int_0^\chi
\frac{\dd a}{(1-n\sin^2a)\sqrt{1-k^2\sin^2a}}.
\]
取 $t_*$ 为 $u=u_1$ 的时刻，得到
\[
F(\chi,k)=\nu(t-t_*),\qquad
u(t)=u_1+D\,\operatorname{sn}^2\bigl(\nu(t-t_*)\mid k^2\bigr).
\]
上式的 $u(t)$ 表示倾角余弦；$\operatorname{sn}$ 是由
$\operatorname{sn}(F(\chi,k)\mid k^2)=\sin\chi$ 定义的 Jacobi 函数。
为了避免频率参数 $\nu$ 与变量 $u$ 混淆，计算中可始终把后者记作 $\cos\theta(t)$。
章动周期是
\[
\tau_\theta=\frac{2\mathcal K(k)}{\nu},\qquad
\mathcal K(k)=F(\pi/2,k).
\]
初始振幅由 $\sin^2\chi_0=(u_0-u_1)/D$ 决定；
选择 $\chi_0$ 所在半周期时要同时匹配 $\dot u_0$ 的符号。
这比只对 $\arcsin\sqrt{(u_0-u_1)/D}$ 取主值多了一项不可省略的初值信息。

为求方位角，先做部分分式：
\[
\frac{K-Pu}{1-u^2}
=\frac{K-P}{2(1-u)}+\frac{K+P}{2(1+u)}.
\]
记
\[
n_-=\frac{D}{1-u_1},\quad n_+=-\frac{D}{1+u_1},\qquad
b_-=\frac{K-P}{2A\nu(1-u_1)},\quad
b_+=\frac{K+P}{2A\nu(1+u_1)}.
\]
从同一时刻 $t_*$ 出发，有
\[
\varphi(t)-\varphi_*
=b_-\Pi(n_-;\chi,k)+b_+\Pi(n_+;\chi,k).
\]
对轴向欧拉角再使用
\[
\frac{u(K-Pu)}{1-u^2}
=P+\frac{K-P}{2(1-u)}-\frac{K+P}{2(1+u)},
\]
于是
\[
\begin{aligned}
\psi(t)-\psi_*={}&\left(\frac PC-\frac PA\right)(t-t_*)\\
&-b_-\Pi(n_-;\chi,k)+b_+\Pi(n_+;\chi,k).
\end{aligned}
\]
任意初值的解只需在这些表达式中减去 $\chi_0$ 对应的积分值。
多周期运动的 $\chi$ 要连续展开，不能每次折回 $[0,\pi/2]$。
每个完整章动周期的方位增量为
\[
\Delta\varphi=2b_-\Pi(n_-;\pi/2,k)+2b_+\Pi(n_+;\pi/2,k).
\]
轴方向只有在 $\Delta\varphi/(2\pi)$ 为有理数时才会经过有限多个章动周期闭合；
一般情形章动周期相同，而球面轨迹不闭合。

\subsubsection*{极点与重根为什么必须另行处理}
若 $K=P$，$\theta\to0$ 时的约束项为
\[
\frac{P^2(1-\cos\theta)^2}{2A\sin^2\theta}
=\frac{P^2}{2A}\tan^2\frac\theta2,
\]
没有发散。若 $K=-P$，在 $\theta\to\pi$ 处也有同样的可约去因子。
此时应先约分，再计算极限；欧拉角本身仍可能有坐标奇点。
如果两个转向根合并，解可能是定常进动，或者是渐近接近不稳定特殊姿态的分离轨道；
不能继续把双根当成两个有限周期的普通端点。
当 $k\to1$ 且相应积分端点可达时，$\mathcal K(k)$ 发散，
对应无限长的渐近运动时间。
$G=0$ 时三次项消失，回到自由陀螺，宜直接用自由转动的解析解，
不应机械沿用含 $1/G$ 的参数化。


\subsection{定常进动与直立自转稳定性}\label{sub:c3-9.4}
对于 $0<\theta_0<\pi$ 的定常进动，令
$\theta=\theta_0,\dot\varphi=\Omega$。
从 $\theta$ 方程
$A\ddot\theta=A\dot\varphi^2\sin\theta\cos\theta
-P\dot\varphi\sin\theta+mgl\sin\theta$
可得
\begin{equation}\label{eq:regularprecession}
A\Omega^2\cos\theta_0-P\Omega+mgl=0.
\end{equation}
当 $\cos\theta_0\ne0$，
\[
\Omega_\pm=\frac{P\pm\sqrt{P^2-4Amgl\cos\theta_0}}
{2A\cos\theta_0}.
\]
须有非负判别式，并使初始角速度匹配某一支。
水平轴 $\cos\theta_0=0$ 时应使用原一次方程，
$\Omega=mgl/P$，要求 $P\ne0$。
若 $P>0$、$\cos\theta_0>0$ 且
$P^2\gg4Amgl\cos\theta_0$，两支近似为
\[
\Omega_{\text{慢}}\simeq\frac{mgl}{P},\qquad
\Omega_{\text{快}}\simeq\frac{P}{A\cos\theta_0}-\frac{mgl}{P}.
\]
高速陀螺常用的是慢进动近似；一般陀螺不必处于这一分支。

定常进动也是有效势极值。
在固定 $P,K$ 的小章动扰动下，
\[
\nu_\theta^2=\frac{U''(\theta_0)}A
=\left(\frac PA-2\Omega\cos\theta_0\right)^2
+\Omega^2\sin^2\theta_0.
\]
非极点定常分支在此约化的意义下为稳定极小值；
整体方位的变化还包含旋转对称性所对应的中性方向。

直立自转 $\theta=0$ 属于欧拉角奇点，需要直接取极限。
此时 $K=P$，有效势在小角度展开为
\[
U(\theta)=mgl+
\left(\frac{P^2}{8A}-\frac{mgl}{2}\right)\theta^2+O(\theta^4).
\]
因此稳定条件为
\[
P^2>4Amgl,\qquad
|s|>\frac{2\sqrt{Amgl}}{C}.
\]
低于阈值时直立姿态的小扰动增长；等号是临界情形，
不能仅由二次项得出严格稳定结论。
与此相比，质心在支点下方的倒向轴分支具有重力回复效应。


\subsubsection*{定常分支的选取、展开与小章动初值解}
给定 $P$ 与倾角并不直接规定同一个初态；
两支定常解还分别要求 $\dot\theta_0=0$ 和对应的 $\dot\varphi_0=\Omega$，
并据此取 $K=A\Omega\sin^2\theta_0+P\cos\theta_0$。
$\dot\psi=s-\Omega\cos\theta_0$ 随分支改变。
对于 $P\ne0$、判别式 $D_0=P^2-4AG\cos\theta_0\ge0$，
与高速极限连续相接的慢分支宜写成有理化形式
\[
\Omega_{\text{慢}}=
\frac{2G}{P+\operatorname{sgn}(P)\sqrt{D_0}},
\]
它避免在 $|P|$ 很大时相减两个近似相等的数。
当 $\cos\theta_0\ne0$，另一支由根和得到
\[
\Omega_{\text{快}}=\frac{P}{A\cos\theta_0}-\Omega_{\text{慢}}.
\]
对于 $\cos\theta_0<0$，判别式自动为正，两支符号也可能不同；
“快”和“慢”在此表示高速展开中的大小，并不表示固定的正负转向。
若 $P=0$，定常非极点解要求 $\cos\theta_0<0$，
且 $\Omega=\pm\sqrt{-G/(A\cos\theta_0)}$。
$\theta_0=\pi/2$ 时原方程只保留 $-P\Omega+G=0$，
不存在可以由除以 $\cos\theta_0$ 得到的第二个有限分支。

令 $\varepsilon=4AG\cos\theta_0/P^2$。
在 $|\varepsilon|\ll1$ 的条件下，逐项展开平方根，有
\[
\Omega_{\text{慢}}=\frac GP
\left(1+\frac\varepsilon4+\frac{\varepsilon^2}{8}
+O(\varepsilon^3)\right).
\]
首项的相对修正量约为 $\varepsilon/4$，因此慢进动近似的误差可由实际参数估计。
慢分支只是一种经过初值匹配的定常运动，
不能因 $|P|$ 很大便把所有运动都改写为 $\dot\varphi=G/P$。

稳定性推导须在固定 $P,K$ 时求导。
记 $\Omega(\theta)=(K-P\cos\theta)/(A\sin^2\theta)$，直接求导得
\[
\Omega'(\theta)=\frac{P/A-2\Omega\cos\theta}{\sin\theta}.
\]
有效势的一阶导数可整理为
\[
U'(\theta)=\sin\theta\,[P\Omega-A\Omega^2\cos\theta-G].
\]
在定常点括号为零，二阶导数为
\[
\begin{aligned}
\frac{U''(\theta_0)}A
&=\sin\theta_0\left[
\left(\frac PA-2\Omega\cos\theta_0\right)\Omega'
+\Omega^2\sin\theta_0\right]\\
&=\left(\frac PA-2\Omega\cos\theta_0\right)^2
+\Omega^2\sin^2\theta_0=\nu_\theta^2>0.
\end{aligned}
\]
因此所有存在的非极点定常分支都是固定这些积分时的严格局部极小值。
设小扰动 $\delta=\theta-\theta_0$，给定 $\delta(0)=\delta_0$、
$\dot\delta(0)=d_0$，线性解为
\[
\delta(t)=\delta_0\cos\nu_\theta t+
\frac{d_0}{\nu_\theta}\sin\nu_\theta t.
\]
记 $a_0=\Omega'(\theta_0)$，相应的方位扰动为
\[
\delta\varphi(t)=\delta\varphi(0)+a_0
\left[\frac{\delta_0}{\nu_\theta}\sin\nu_\theta t
+\frac{d_0}{\nu_\theta^2}(1-\cos\nu_\theta t)\right].
\]
轴向相位扰动由
\[
\delta\dot\psi=(-a_0\cos\theta_0+\Omega\sin\theta_0)\delta
\]
积分得到。若扰动还改变 $P$ 或 $K$，则应先移动到邻近定常解，
再描述围绕新倾角的章动；相位随时间积累的差异也不能误认为倾角指数发散。

\subsubsection*{直立临界点的四次势与完整倾斜方程}
直立时 $K=P$，有效势的可约分形式为
\[
U(\theta)=\frac{P^2}{2A}\tan^2\frac\theta2+G\cos\theta.
\]
利用 $\tan^2(\theta/2)=\theta^2/4+\theta^4/24+O(\theta^6)$，
得到
\[
U(\theta)=G+\left(\frac{P^2}{8A}-\frac G2\right)\theta^2
+\left(\frac{P^2}{48A}+\frac G{24}\right)\theta^4+O(\theta^6).
\]
在 $P^2=4AG$ 时，固定 $K=P$ 的约化势四次项为 $G\theta^4/8>0$。
它描述受这组精确积分限制的非线性回复，二次近似已经不能给出非零小振动频率。

为看清一般倾斜扰动，避开欧拉角极点，写轴方向
$\vect n=(x,y,1)+O(x^2+y^2)$，取空间竖直为 $\vect k$。
对轴对称体，精确关系为
\[
\vect L=A\vect n\times\dot{\vect n}+P\vect n,
\qquad \vect M=G\vect k\times\vect n.
\]
取一阶项，得到
\[
A\ddot x+P\dot y-Gx=0,\qquad
A\ddot y-P\dot x-Gy=0.
\]
记 $z=x+\mathrm i y$，合写为
\[
A\ddot z-\mathrm iP\dot z-Gz=0.
\]
取 $z\propto e^{\mathrm i\sigma t}$，特征频率满足
\[
A\sigma^2-P\sigma+G=0,\qquad
\sigma_\pm=\frac{P\pm\sqrt{P^2-4AG}}{2A}.
\]
阈值以上两频率为实数，小倾斜由两个旋转模态叠加而成。
对初值 $z_0,\dot z_0$，若 $\sigma_+\ne\sigma_-$，
\[
\begin{aligned}
z(t)={}&c_+e^{\mathrm i\sigma_+t}+c_-e^{\mathrm i\sigma_-t},\\
c_+&=\frac{\dot z_0-\mathrm i\sigma_-z_0}
{\mathrm i(\sigma_+-\sigma_-)},\qquad
c_-=\frac{\mathrm i\sigma_+z_0-\dot z_0}
{\mathrm i(\sigma_+-\sigma_-)}.
\end{aligned}
\]
阈值以下，频率出现虚部，因而存在增长模态。
恰在阈值处，重根给出
\[
z(t)=(c_1+c_2t)e^{\mathrm iPt/(2A)}.
\]
因此完整线性化的一般初值可含随 $t$ 增长的项。
这个线性项不能外推为精确非线性运动无限增大；
同样，固定 $P,K$ 约化势的正四次项也不能替代一般扰动的完整分析。
严格大于阈值的条件才给出这里使用的通常线性稳定结论。


\example{定常进动的两支角速率}\label{ex:c3-61}

\QuestionStatement{取 $A=1\,\mathrm{kg\,m^2}$、$C=.5\,\mathrm{kg\,m^2}$、
$mgl=1\,\mathrm{N\,m}$，
轴向分量 $s=8\,\mathrm{s^{-1}}$，
定常倾角 $\theta_0=\pi/3$。求两支进动率与欧拉角自转率。
模型为轴对称重陀螺，$I_1=I_2=A,I_3=C$。列出定常进动存在条件，比较慢进动近似及精确值。}{ex61}{例61题图}
% BEGIN DETAILED ANSWER ex61
\solution \Step{列出定常进动的代数条件}
以图~\ref{fig:solution-ex61}~中重力矩与固定倾角 $\theta_0$为依据，先写出图右侧第一行的几何或动力学关系，再代入其下方的计算式。

固定倾角 $\theta_0$、恒定轴向分量 $s$ 和进动率 $\Omega$ 必须满足倾角方程
\[
A\Omega^2\cos\theta_0-Cs\Omega+mgl=0.
\]
若 $\cos\theta_0\ne0$，实根存在的条件是
$(Cs)^2-4Amgl\cos\theta_0\ge0$。题给 $P=Cs=4\,\mathrm{kg\,m^2\,s^{-1}}$，代入得
$\Omega^2/2-4\Omega+1=0$。
\SolutionFigure{ex61}{例61的计算示意}
\Step{解二次方程并区分两支}
\[
\Omega_\pm=\frac{4\pm\sqrt{16-2}}{1}
=4\pm\sqrt{14},
\]
即快支 $7.74166\,\mathrm{s^{-1}}$、慢支 $0.25834\,\mathrm{s^{-1}}$。
两支均需相应的初始进动速度；同一个完整初态不会任意选择其一。
\Step{计算欧拉角自转率}
由于 $s=\dot\psi+\Omega\cos\theta_0$，
\[
\dot\psi_+=8-\frac{7.74166}{2}=4.12917\,\mathrm{s^{-1}},\qquad
\dot\psi_-=8-\frac{0.25834}{2}=7.87083\,\mathrm{s^{-1}}.
\]
把它们与对应 $\Omega/2$ 相加都回到给定的轴向分量 $8$，可检查符号。
\Step{评估慢进动近似}
若忽略慢支方程中的二次项，$\Omega_{\text{慢}}\simeq mgl/P=0.25\,\mathrm{s^{-1}}$。
它相对精确慢支低估约 $(0.2583426-0.25)/0.2583426=3.23\%$。
小参数 $Amgl\cos\theta_0/P^2=1/32$ 控制这一近似；快支不可用此式计算。

% END DETAILED ANSWER ex61
\subsection{单向轨迹、回环与尖点}\label{sub:c3-9.5}
刚体对称轴的单位端点在球面上的速度平方为
\[
|\dot{\vect e}_z|^2=\dot\theta^2+\dot\varphi^2\sin^2\theta.
\]
在非极点章动区间内，
$\dot\varphi$ 的符号只取决于 $K-P\cos\theta$。
若该量不变号，方位角单调，得到伴随章动的单向波状轨迹；
若它在区间内部过零，进动方向反转，产生回环；
若它在章动端点为零，$\dot\theta$ 同时为零，
轴端点瞬时停止，可形成尖点。
这些是轴方向在单位球面上的轨迹，
不是盘上某一材料点或质心的空间轨迹。

\example{只有轴向自转的倾斜释放}\label{ex:c3-62}

\QuestionStatement{陀螺初始轴倾角满足 $0<\theta_0<\pi$，
初始无章动、无进动，但轴向角速度为 $s\ne0$。
理想固定支点不做功，除均匀恒定重力外没有其他外力矩。
求守恒量，说明为何开始时出现尖点而不是直接匀速进动。
其中 $A,C$ 是固定支点处的横向、轴向主惯量，质心位于 $+\vect e_3$ 方向距离 $l$。求另一章动转向点的条件和高速自转近似的最大下倾量。}{ex62}{例62题图}
% BEGIN DETAILED ANSWER ex62
\solution \Step{由初值求三个积分常数}
以图~\ref{fig:solution-ex62}~中倾斜释放时的轴线与重力方向为依据，先写出图右侧第一行的几何或动力学关系，再代入其下方的计算式。

采用 $\theta$ 从竖直向上量起，令 $P=L_3$、$K=L_\zeta$，并从总能量中扣去恒定自转能 $P^2/(2C)$，剩余能量记为 $E'$。初始 $\dot\theta=\dot\varphi=0,\omega_3=s$，故
\[
P=Cs,\qquad K=P\cos\theta_0,\qquad E'=mgl\cos\theta_0.
\]
\SolutionFigure{ex62}{例62的计算示意}
\Step{求初始倾角加速度与进动率}
进动率由 $K=A\dot\varphi\sin^2\theta+P\cos\theta$ 给出：
\[
\dot\varphi=\frac{P(\cos\theta_0-\cos\theta)}{A\sin^2\theta}.
\]
初始确为零。倾角动力学为
$A\ddot\theta=A\dot\varphi^2\sin\theta\cos\theta-P\dot\varphi\sin\theta+mgl\sin\theta$，所以
$\ddot\theta(0)=mgl\sin\theta_0/A>0$，陀螺开始下倾。
\Step{建立章动转向点方程}
令 $u=\cos\theta,u_0=\cos\theta_0$。约化能量写为
\[
\frac A2\dot\theta^2+\frac{P^2(u_0-u)^2}{2A(1-u^2)}+mglu=mglu_0.
\]
乘以 $2(1-u^2)/A$，得到
\[
\dot u^2=(u_0-u)\left[\frac{2mgl}{A}(1-u^2)-\frac{P^2}{A^2}(u_0-u)\right].
\]
除初始根 $u_0$ 外，另一允许根 $u_1<u_0$ 满足
$2Amgl(1-u_1^2)=P^2(u_0-u_1)$。最大倾角是 $\arccos u_1$，并须检验 $-1\le u_1<u_0$ 及运动区间非负性。
\Step{解释尖点与高速小章动}
令 $\theta=\theta_0+\delta$，在 $|\delta|\ll1$ 时有效势能增量为
$P^2\delta^2/(2A)-mgl\sin\theta_0\delta$。
转向点满足其为零，故
$\delta_{\max}\simeq2Amgl\sin\theta_0/P^2$，要求 $P^2\gg Amgl$、初始轴远离两极，并使此量远小于1。
初始轴轨迹速度为零，随后 $\delta\sim t^2$ 而 $\dot\varphi\sim t^2$，形成尖点。高速条件下平均慢进动约 $mgl/P$，但初始瞬间的进动率仍是零，不能直接替换为定常进动。

% END DETAILED ANSWER ex62
\subsubsection*{纯自转释放的准确转向角与早期尖点}
以下明确取 $P\ne0$、$0<\theta_0<\pi$，且初始
$\dot\theta_0=\dot\varphi_0=0$。
令 $u_0=\cos\theta_0$。除了根 $u_0$，另两根来自
\[
2AGu^2-P^2u+(P^2u_0-2AG)=0.
\]
因此
\[
\begin{aligned}
u_1&=\frac{P^2-\sqrt{P^4-8AGP^2u_0+16A^2G^2}}{4AG},\\
u_3&=\frac{P^2+\sqrt{P^4-8AGP^2u_0+16A^2G^2}}{4AG}.
\end{aligned}
\]
普通倾斜释放中，$u_1<u_0<u_3$，实际章动在 $u_1$ 与 $u_0$ 之间，
所以准确的最大下倾角为 $\theta_{\max}=\arccos u_1$。
高速时第一个根的表达式可能发生相减消去，可改用
\[
u_1=\frac{2(P^2u_0-2AG)}
{P^2+\sqrt{P^4-8AGP^2u_0+16A^2G^2}}.
\]
分母非零；极限 $G\to0$ 应由原方程单独处理。

令时间零点为释放时刻。
由二阶倾角方程，最初的展开为
\[
\theta-\theta_0=\frac{G\sin\theta_0}{2A}t^2+O(t^4).
\]
再从角动量积分展开，而不是假定进动率恒定：
\[
\begin{aligned}
\dot\varphi&=\frac{P(\cos\theta_0-\cos\theta)}{A\sin^2\theta}
=\frac{PG}{2A^2}t^2+O(t^4),\\
\varphi-\varphi_0&=\frac{PG}{6A^2}t^3+O(t^5).
\end{aligned}
\]
于是局部经纬展开图中纵向位移为 $t^2$ 量级，横向位移为 $t^3$ 量级。
消去时间得到半三次尖点关系
\[
(\varphi-\varphi_0)^2=
\frac{2P^2}{9AG\sin^3\theta_0}(\theta-\theta_0)^3
+\text{更高阶项}.
\]
这也给出尖点的时间方向：释放后倾角先增大，进动方向随后由 $P$ 的符号确定。
若 $P=0$，系统退化为在竖直平面内摆动的轴，
横向三次项消失，不能再把轨迹称为普通的半三次进动尖点。
若 $\sin\theta_0=0$，释放条件位于特殊直立或悬垂姿态，
也须使用相应的稳定性分析。

\subsubsection*{高速释放的有时间依赖的章动与进动}
令 $\delta=\theta-\theta_0$ 很小。
固定 $P,K=P\cos\theta_0$ 时，初态附近
\[
U'(\theta_0)=-G\sin\theta_0,\qquad
U''(\theta_0)=\frac{P^2}{A}-G\cos\theta_0.
\]
保留势的二次项，得到一个受常值驱动的振子：
\[
\ddot\delta+\omega_\theta^2\delta=\frac GA\sin\theta_0,
\qquad
\omega_\theta^2=\frac{P^2}{A^2}-\frac GA\cos\theta_0.
\]
若 $\omega_\theta^2>0$，且初始 $\delta=\dot\delta=0$，解为
\[
\delta(t)=\frac{AG\sin\theta_0}{P^2-AG\cos\theta_0}
(1-\cos\omega_\theta t).
\]
相应的一阶进动率和方位角为
\[
\dot\varphi\simeq\frac{P\delta}{A\sin\theta_0}
=\frac{PG}{P^2-AG\cos\theta_0}(1-\cos\omega_\theta t),
\]
\[
\varphi-\varphi_0\simeq\frac{PG}{P^2-AG\cos\theta_0}
\left(t-\frac{\sin\omega_\theta t}{\omega_\theta}\right).
\]
在更强的 $P^2\gg AG$ 条件下，$\omega_\theta\simeq|P|/A$，
最大章动幅度化为 $2AG\sin\theta_0/P^2$，平均进动率化为 $G/P$。
瞬时进动却在零与约 $2G/P$ 之间变化。
上述近似还要求计算所得章动幅度远小于一，
且不逼近使 $\sin\theta_0$ 消失的极点。
它补足了“平均慢进动”所省略的启动过程，同时保留了初始进动率为零的条件。


下面用同一个模型 $A=1,C=.5,mgl=1,P=4$，
初始 $\theta_0=.6,\dot\theta_0=0$，
分别取 $\dot\varphi_0=.25,-.15,0$。
曲线由上述运动方程积分得到，每幅截取两个完整章动周期。
为清楚显示回环和尖点，将单位球面的方位、极角方向展开：
横轴为 $\varphi$，纵向向下表示 $\theta$ 增大。
三种轨迹对应不同的初始竖直角动量，不是同一初态的三种任意画法。
\TeachingFigure{c3-diagram-56}{重陀螺方程积分得到的两个章动周期：单向、回环与尖点；箭头表示时间方向}
% BEGIN ORIGINAL FIGURES
\OriginalFigure[6.2cm]{c3-6-p021-nutation-tracks}{单位球面上的重陀螺轴迹：单向、回环与尖点；原图上、下边界分别为 $\theta_2=\theta_{\min}$、$\theta_1=\theta_{\max}$}
% END ORIGINAL FIGURES

轨迹的每一点同时满足同一组 $P,K,E$。
示意图的纵向尺度分别归一化以显示形状，
它不意味着三种章动幅度完全相同。

\clearpage
\section{拉莫尔进动与核磁共振}\label{sec:c3-10}
\subsection{带电圆环与陀螺磁比}\label{sub:c3-10.1}
质量 $m$、半径 $r$、总电荷 $q$ 的均匀薄圆环
绕法线 $\vect n$ 以有向角速率 $\omega$ 转动。
电流按正电荷运动定义，$q$ 必须保留符号：
\[
I=\frac{q\omega}{2\pi},\qquad
\vect\mu=I\pi r^2\vect n=\frac12qr^2\omega\vect n,\qquad
\vect L=mr^2\omega\vect n.
\]
因此
\[
\vect\mu=\gamma\vect L,\qquad \gamma=\frac q{2m}.
\]
$\gamma$ 称陀螺磁比。
若处处有同一荷质比 $\dd q/\dd m=q/m$，关于同一原点，
\[
\vect\mu=\frac12\int\vect r\times\vect v\,\dd q
=\frac q{2m}\int\vect r\times\vect v\,\dd m
=\frac q{2m}\vect L.
\]
若电荷与质量分布不成比例，不能仅用总 $q/m$ 给出这个关系。
电子自旋与核自旋也不是经典圆环的几何模型，
它们的磁矩须用各自的陀螺磁比或 $g$ 因子。

\TeachingFigure{c3-diagram-57}{均匀带电圆环的磁矩与角动量；负电荷时两者反向}
% BEGIN ORIGINAL FIGURES
\OriginalFigure[6.2cm]{c3-6-p022-magnetic-ring}{磁场中的电流圆环及其磁矩方向}
% END ORIGINAL FIGURES



\subsubsection*{磁矩与角动量必须使用同一系统和同一原点}
薄圆环上任意一段满足 $\dd q=(q/m)\dd m$。
这段粒子的机械角动量为
$\dd\vect L=\vect r\times\vect v\,\dd m$，
其在局域电流磁矩中的贡献为
$\dd\vect\mu=\vect r\times\vect v\,\dd q/2$。
先在每一段使用荷质比，再积分，才可把 $q/(2m)$ 提到积分号外。
若两个同轴环的荷质比分别为 $q_1/m_1,q_2/m_2$，则
\[
\vect\mu=\frac{q_1}{2m_1}\vect L_1+
\frac{q_2}{2m_2}\vect L_2.
\]
一般不能把它写成
$(q_1+q_2)(\vect L_1+\vect L_2)/[2(m_1+m_2)]$。
例如两个半径不同、质量相同、带等量异号电荷的环同速转动，
总电荷为零，但两磁矩未必抵消。
这说明决定磁矩的是电流的空间分布，而不是单独的总电荷。

这里 $\vect L$ 是关于选定原点的机械角动量。
在外磁场中引入矢势之后，正则角动量一般与机械角动量不同；
两者的关系将在轨道拉莫尔变换中明确写出。
对于电子自旋，通常写
$\vect\mu_S=-g_S\mu_B\vect S/\hbar$，
$\mu_B=e\hbar/(2m_e)$、$e>0$。
它的 $g_S\simeq2$ 不是由一个均匀圆环的质量和电荷积分推出来的。
核自旋的磁矩也须用具体核的 $\gamma_n$，
因此同一种进动方程可以用于不同磁偶极模型，但参数的物理来源并不相同。


\subsection{闭合电流的磁力矩：积分证明}\label{sub:c3-10.2}
稳恒闭合细线圈在均匀磁场中，
$\dd\vect F=I\,\dd\vect s\times\vect B$，
合力为 $I\oint\dd\vect s\times\vect B=0$。
磁矩的正确线积分形式为
\begin{equation}
\vect\mu=\frac I2\oint\vect r\times\dd\vect s,
\end{equation}
体电流则写
$\vect\mu=\tfrac12\int\vect r\times\vect j\,\dd V$。
$1/2$ 必须保留，因为
$\tfrac12\oint\vect r\times\dd\vect s$
等于平面线圈的有向面积。

力矩为
\[
\vect M=I\oint\vect r\times(\dd\vect s\times\vect B)
=I\oint[(\vect r\cdot\vect B)\dd\vect s
-(\vect r\cdot\dd\vect s)\vect B].
\]
第二项的闭环积分为零，因为
$\oint\vect r\cdot\dd\vect s=\tfrac12\oint\dd(r^2)=0$。
于是 $\vect M=I\oint(\vect B\cdot\vect r)\dd\vect s$。
另一方面，
\[
\vect\mu\times\vect B
=\frac I2\oint[(\vect r\cdot\vect B)\dd\vect s
-(\dd\vect s\cdot\vect B)\vect r].
\]
由闭合线积分分部，
\[
0=\oint\dd[(\vect B\cdot\vect r)\vect r]
=\oint(\vect B\cdot\dd\vect s)\vect r
+\oint(\vect B\cdot\vect r)\dd\vect s,
\]
可见后两积分互为相反数，并非各自都为零。
最终
\begin{equation}\label{eq:magtorque}
\vect M=\vect\mu\times\vect B.
\end{equation}
恒定外场、稳恒闭合电流是这个推导的条件。
时变电荷分布、真实带电刚体的全部倾转电流，
以及非均匀场的磁力，需要重新核查模型。
高速旋转的弱场近似或自旋磁偶极模型可使用下节的进动方程，
不应把两种不同层次的模型自动视为完全等价。


\subsubsection*{闭环力矩的三分量展开}
取 $\dd\vect s=(\dd x,\dd y,\dd z)$，均匀场为 $(B_x,B_y,B_z)$。
三个微元力为
\[
\begin{aligned}
\dd F_x&=I(B_z\dd y-B_y\dd z),\\
\dd F_y&=I(B_x\dd z-B_z\dd x),\\
\dd F_z&=I(B_y\dd x-B_x\dd y).
\end{aligned}
\]
因闭合路径上 $\oint\dd x=\oint\dd y=\oint\dd z=0$，合力的三个分量分别为零。
不能把这个结论推广到带坐标权重的积分；力矩正包含这样的权重。
例如第一分量为
\[
\begin{aligned}
M_x&=\oint(y\dd F_z-z\dd F_y)\\
&=I\oint[B_y y\dd x+B_z z\dd x-B_x(y\dd y+z\dd z)]\\
&=I\left(B_y\oint y\dd x+B_z\oint z\dd x\right).
\end{aligned}
\]
最后一步只消掉了全微分
$\oint(y\dd y+z\dd z)=\oint\dd(y^2+z^2)/2=0$。
相同计算给出
\[
\begin{aligned}
M_y&=I\left(B_z\oint z\dd y+B_x\oint x\dd y\right),\\
M_z&=I\left(B_x\oint x\dd z+B_y\oint y\dd z\right).
\end{aligned}
\]
定义有向面积矢量 $\vect{\mathcal A}$ 的三个分量：
\[
\begin{aligned}
\mathcal A_x&=\frac12\oint(y\dd z-z\dd y)=\oint y\dd z,\\
\mathcal A_y&=\frac12\oint(z\dd x-x\dd z)=\oint z\dd x,\\
\mathcal A_z&=\frac12\oint(x\dd y-y\dd x)=\oint x\dd y.
\end{aligned}
\]
等号使用 $\oint\dd(yz)=\oint\dd(zx)=\oint\dd(xy)=0$。
对于平面简单闭环，$\vect{\mathcal A}$ 等于几何面积乘右手法向；
对于非平面闭环，它的三个分量分别是三个坐标平面上的有向投影面积，
不必假定整条线圈有同一个平面法向。
因此
\[
\begin{aligned}
M_x&=I(\mathcal A_yB_z-\mathcal A_zB_y),\\
M_y&=I(\mathcal A_zB_x-\mathcal A_xB_z),\\
M_z&=I(\mathcal A_xB_y-\mathcal A_yB_x),
\end{aligned}
\]
逐分量恢复 $\vect M=(I\vect{\mathcal A})\times\vect B$，
即 $\vect\mu=I\vect{\mathcal A}$。
面积线积分中的 $1/2$ 与最终力矩中没有额外 $1/2$ 同时成立。

\subsubsection*{积分分部、原点平移与适用边界}
对任意恒定 $\vect B$，闭合路径满足
\[
\oint\dd[(\vect B\cdot\vect r)\vect r]=0.
\]
乘积求导产生两个积分，正是把力矩积分变成磁矩积分所需的关系。
只有它们的和为零，而不是每一项各为零。
这与一维分部积分完全相同，只是被积分函数为矢量。

如果原点移到 $O'=O+\vect a$，则 $\vect r'=\vect r-\vect a$，
\[
\vect\mu'=\frac I2\oint(\vect r-\vect a)\times\dd\vect s
=\vect\mu-\frac I2\vect a\times\oint\dd\vect s=\vect\mu.
\]
同理 $\vect M_{O'}=\vect M_O-\vect a\times\vect F=\vect M_O$。
均匀场中的合力为零保证了力矩不随取矩原点改变。
这是一对相互独立的检验：磁矩的闭环公式与一般力系的矩迁移公式给出同一结果。

若场随位置显著改变，$\vect B$ 不能提出积分号，合力未必为零；
对于尺寸远小于场变化尺度的永久小磁偶极，最低阶近似才可写
$\vect F=\nabla(\vect\mu\cdot\vect B)$。
若电流由外部电源维持，讨论总能量时还要把电源与场包含在系统中。
这些情况不能仅凭均匀静场的闭环积分直接决定全部运动。


\example{检验一个不能直接消掉的闭环积分}\label{ex:c3-63}

\QuestionStatement{半径 $r$ 的圆环位于 $xy$ 平面，
$\vect B=B\vect i$。计算
$\oint(\vect B\cdot\dd\vect s)\vect r$。
位置矢量从圆心量起，积分方向从 $+z$ 看为逆时针。写出参数化与各分量积分，并解释闭合路径为何不保证此加权积分为零。}{ex63}{例63题图}
% BEGIN DETAILED ANSWER ex63
\solution \Step{按给定方向参数化圆周}
以图~\ref{fig:solution-ex63}~中圆周的有向切线 $\dd\vect s$ 与径矢 $\vect r$为依据，先写出图右侧第一行的几何或动力学关系，再代入其下方的计算式。

取 $0\le\chi\le2\pi$，从 $+z$ 看逆时针，则
\[
\vect r=r(\cos\chi\vect i+\sin\chi\vect j),\qquad
\dd\vect s=r(-\sin\chi\vect i+\cos\chi\vect j)\dd\chi.
\]
题设 $\vect B=B\vect i$，所以 $\vect B\cdot\dd\vect s=-Br\sin\chi\dd\chi$。
\SolutionFigure{ex63}{例63的计算示意}
\Step{逐分量保留位置权重}
\[
\oint(\vect B\cdot\dd\vect s)\vect r
=-Br^2\left[\vect i\int_0^{2\pi}\sin\chi\cos\chi\dd\chi
+\vect j\int_0^{2\pi}\sin^2\chi\dd\chi\right].
\]
第一积分的原函数为 $\sin^2\chi/2$，两端均为零；第二个用 $\sin^2\chi=(1-\cos2\chi)/2$ 得到 $\pi$。故结果为
\[
-\pi Br^2\vect j.
\]
\Step{解释闭合路径不能消去哪一项}
无权积分 $\oint\vect B\cdot\dd\vect s=B\oint\dd x=0$；本题则是 $B\oint\vect r\dd x$，$\vect r$ 沿路径变化，不能从积分中提出。
$x$ 分量是 $B\oint x\dd x=0$，$y$ 分量是 $B\oint y\dd x=-B\pi r^2$，亦可由有向面积核验。
反向绕行会使答案变号，$r\to0$ 时按面积趋零；量纲为磁场乘面积，与每项积分相符。

% END DETAILED ANSWER ex63
\subsection{拉莫尔方程、转向与旋转参考系}\label{sub:c3-10.3}
在磁偶极模型中，若 $\vect\mu=\gamma\vect L$，
磁场均匀、恒定，则
\[
\dot{\vect L}=\gamma\vect L\times\vect B
=\vect\Omega_L\times\vect L,\qquad
\vect\Omega_L=-\gamma\vect B.
\]
它表明角动量绕磁场进动。由叉积正交性，
$L^2$ 与 $\vect L\cdot\vect B$ 都恒定，
进动圆锥夹角不变。
频率大小与周期为
\[
\omega_L=|\gamma|B,\quad f_L=\frac{|\gamma|B}{2\pi},\quad
\tau_L=\frac{2\pi}{|\gamma|B}.
\]
$\omega_L$ 是角频率，$f_L$ 是普通频率，不能漏 $2\pi$。
对于经典圆环 $\gamma=q/(2m)$；
电子轨道取 $q=-e,e>0$，故
$\vect\Omega_L=+e\vect B/(2m)$。
若把 $e$ 已定义成带符号电荷，则公式应同时改变解释。

取 $\vect B=B\vect k$、初始横向角动量沿 $\vect i$，
令有向 $\Omega=-\gamma B$，解为
\[
\vect L(t)=L_\perp[\cos(\Omega t)\vect i+\sin(\Omega t)\vect j]
+L_\parallel\vect k.
\]
在以 $\vect\Omega_L$ 转动的参考系中，
运输项恰抵消真实进动导数，$\vect L$ 的分量恒定。
该参考系不自动等同于实际刚体本体系，
也不能把 $\vect\Omega_L$ 当成物体全部的机械角速度。

\TeachingFigure{c3-diagram-58}{进动的角动量端点位于固定圆锥上，转向取决于陀螺磁比符号}
% BEGIN ORIGINAL FIGURES
\OriginalFigure[6.2cm]{c3-6-p025-magnetic-precession}{不同初始取向的磁偶极进动；竖直轴为磁场方向}
% END ORIGINAL FIGURES



\subsubsection*{任意初始角动量的分量解与旋转系运输}
先取 $\vect B=B\vect k$，$B\ge0$，并保留 $\gamma$ 的符号。
进动方程逐分量为
\[
\dot L_x=\gamma B L_y,\qquad
\dot L_y=-\gamma B L_x,\qquad \dot L_z=0.
\]
因此 $\ddot L_x+\gamma^2B^2L_x=0$。
对于任意初值 $L_{x0},L_{y0},L_{z0}$，
\[
\begin{aligned}
L_x(t)&=L_{x0}\cos(\gamma Bt)+L_{y0}\sin(\gamma Bt),\\
L_y(t)&=L_{y0}\cos(\gamma Bt)-L_{x0}\sin(\gamma Bt),\\
L_z(t)&=L_{z0}.
\end{aligned}
\]
或者记 $L_+=L_x+\mathrm iL_y$，便有
$L_+(t)=L_+(0)e^{-\mathrm i\gamma Bt}$。
在初始横向矢量沿正 $x$ 方向时，$\dot L_y(0)=-\gamma BL_{x0}$；
由这条局部导数即可核对图上的进动箭头。
频率大小相同而 $\gamma$ 变号时，圆周方向反转。

对任意方向的恒定磁场，令 $\vect b=\vect B/B$、
$\vect L_{\parallel0}=(\vect L_0\cdot\vect b)\vect b$，
$\vect L_{\perp0}=\vect L_0-\vect L_{\parallel0}$。
Rodrigues 旋转公式给出完整解
\[
\vect L(t)=\vect L_{\parallel0}
+\cos(\gamma Bt)\vect L_{\perp0}
-\sin(\gamma Bt)\vect b\times\vect L_{\perp0}.
\]
若 $\vect L_0$ 平行磁场，横向分量为零，实际矢量静止；
$|\gamma|B$ 仍是横向小扰动的进动频率，但没有可见的进动圆周。
若 $B=0$ 或 $\gamma=0$，直接取 $\dot{\vect L}=0$，不必定义进动轴。

在以 $\vect\Omega_L=-\gamma\vect B$ 转动的坐标系 $R$ 中，运输公式为
\[
\left(\frac{\dd\vect L}{\dd t}\right)_S
=\left(\frac{\dd\vect L}{\dd t}\right)_R
+\vect\Omega_L\times\vect L.
\]
将空间导数的动力学表达式代入，得到
$({\dd\vect L}/{\dd t})_R=0$。
“磁矩在旋转系中静止”是分量恒定的陈述，
不要求所有带电材料点都在这个旋转系中静止。
若另外存在机械自转、章动或阻尼，它们须由各自方程描述。
理想进动方程既保持 $L$ 的大小，也保持与磁场的夹角，
因此它自身不会使一组倾斜磁矩逐渐排列到磁场方向。
产生平衡净磁化的热交换与弛豫是后面引入的另一层动力学。


\example{弱磁场中的经典轨道拉莫尔变换}\label{ex:c3-64}

\QuestionStatement{带电粒子在中心势 $U(r)$ 中运动，
$m\ddot{\vect r}=-\nabla U+q\dot{\vect r}\times\vect B$，
外磁场恒定且弱。说明以
$\vect\Omega=-q\vect B/(2m)$ 转动的参考系为何能消去一阶磁场项。
所有旋转率保留电荷 $q$ 的符号；写出未截断的旋转系方程和被忽略项，说明弱场近似的尺度条件。}{ex64}{例64题图}
% BEGIN DETAILED ANSWER ex64
\solution \Step{在同一空间中写速度与加速度变换}
以图~\ref{fig:solution-ex64}~中中心力轨道与磁场方向为依据，先写出图右侧第一行的几何或动力学关系，再代入其下方的计算式。

恒定旋转角速度取 $\vect\Omega=-q\vect B/(2m)$，与磁场平行。旋转系中的导数记为撇号，则
\[
\vect v=\vect v'+\vect\Omega\times\vect r,\quad
\vect a=\vect a'+2\vect\Omega\times\vect v'
+\vect\Omega\times(\vect\Omega\times\vect r).
\]
中心势具有旋转不变性，所以 $-\nabla U$ 在旋转坐标中保持相同函数形式。
\SolutionFigure{ex64}{例64的计算示意}
\Step{逐项代回并消去一阶速度项}
用 $q\vect B=-2m\vect\Omega$，有
\[
q\vect v'\times\vect B=2m\vect\Omega\times\vect v',
\]
恰与左侧科里奥利项相消。由牵连速度产生的洛伦兹项则为
\[
q(\vect\Omega\times\vect r)\times\vect B
=2m\vect\Omega\times(\vect\Omega\times\vect r).
\]
再扣除左侧一个离心型项，得到准确的旋转系方程
\[
m\vect a'=-\nabla U+m\vect\Omega\times(\vect\Omega\times\vect r).
\]
\Step{标明弱场近似的阶次与尺度}
最后一项为 $O(B^2)$，因此保留至 $B$ 的一阶时有 $m\vect a'\simeq-\nabla U$。
若无场运动特征频率为 $\omega_{\text{轨道}}$、尺度为 $r_0$，二阶项相对主加速度的量级为
$\Omega^2/\omega_{\text{轨道}}^2$，要求 $|\Omega|\ll\omega_{\text{轨道}}$。
\Step{还原物理含义与方向}
旋转系中轨道在一阶近似下保持无磁场形式；空间系观察到它以 $-q\vect B/(2m)$ 缓慢旋转。
负电荷的进动方向与正电荷相反，不能先取 $|q|$ 再推导。
自由粒子没有可比较的中心力轨道频率，故不能把此弱场结论用于将准确回旋频率误减半。

% END DETAILED ANSWER ex64
\subsubsection*{拉莫尔变换中的二阶磁场项}
在静止的空间系中，以中心势 $U(r)$ 为势能，
\[
m\vect a=-\nabla U+q\vect v\times\vect B.
\]
旋转系角速度 $\vect\Omega=-q\vect B/(2m)$ 恒定，原点不移动。
把速度与加速度的全部项代入，得到
\[
\begin{aligned}
m\vect a'+2m\vect\Omega\times\vect v'
+m\vect\Omega\times(\vect\Omega\times\vect r)
={}&-\nabla U+q\vect v'\times\vect B\\
&+q(\vect\Omega\times\vect r)\times\vect B.
\end{aligned}
\]
利用 $q\vect B=-2m\vect\Omega$，逐项比较：
\[
q\vect v'\times\vect B=2m\vect\Omega\times\vect v',
\]
\[
q(\vect\Omega\times\vect r)\times\vect B
=2m\vect\Omega\times(\vect\Omega\times\vect r).
\]
一阶项完全相消，二阶项却只相消一半，故准确的旋转系方程为
\[
m\vect a'=-\nabla U+m\vect\Omega\times(\vect\Omega\times\vect r).
\]
把位置分解为平行和垂直 $\vect B$ 的两部分，
$\vect\Omega\times(\vect\Omega\times\vect r)=-\Omega^2\vect r_\perp$，
从而
\[
m\vect a'=-\nabla U_{\text{eff}},\qquad
U_{\text{eff}}=U(r)+\frac12m\Omega^2r_\perp^2
=U(r)+\frac{q^2B^2}{8m}r_\perp^2.
\]
这是向磁轴的附加回复项，符号不能与单独的离心项混淆。
$U_{\text{eff}}$ 在旋转系中不显含时间，故
$E_R=m|\vect v'|^2/2+U_{\text{eff}}$ 守恒。
中心势的旋转不变性使 $U(r)$ 在该系中仍是同一个函数。
若外势固定在空间而各向异性，其旋转系表达式一般显含时间，
一阶磁项的代数消去并不意味着恢复原来的静止无场问题。

只有当二阶项相对于原中心力和所研究的时间精度可以忽略时，
才能说旋转系中恢复无场运动。
例如对典型半径 $r_c$、原动力学频率 $\omega_c$，
通常要求 $|q|B/(2m)\ll\omega_c$。
非常长时间内二阶频率差积累的相位仍可能可见；
“弱场”不等于在任意长时间内误差都为零。

\subsubsection*{正则动量说明同一个二阶项}
取对称规范 $\vect A=\vect B\times\vect r/2$。
带电粒子的正则动量为 $\vect p=m\vect v+q\vect A$，
所以空间哈密顿量逐项展开为
\[
\begin{aligned}
H&=\frac{(\vect p-q\vect A)^2}{2m}+U(r)\\
&=\frac{p^2}{2m}+U(r)-\frac q{2m}\vect B\cdot\vect L_{\text{can}}
+\frac{q^2B^2}{8m}r_\perp^2,
\end{aligned}
\]
其中 $\vect L_{\text{can}}=\vect r\times\vect p$。
中间项使用 $\vect p\cdot(\vect B\times\vect r)=\vect B\cdot(\vect r\times\vect p)$。
旋转系的哈密顿量减去 $\vect\Omega\cdot\vect L_{\text{can}}$，
恰好抵消一阶项，剩下同一个二阶势。
机械角动量则为
\[
\vect L_{\text{mech}}=\vect r\times m\vect v
=\vect L_{\text{can}}-\frac q2\vect r\times(\vect B\times\vect r).
\]
二者差为磁场的一阶量。
在严格的二阶计算中不能把两个角动量符号直接互换，
否则会改变 $B^2$ 项的系数。
一阶轨道塞曼项与二阶抗磁项的区分也由这个展开直接给出。
\footnote{旋转系中保留二阶磁场项的力学形式参见
\href{https://www.damtp.cam.ac.uk/user/tong/relativity/six.pdf}{David Tong：Non-Inertial Frames，Larmor Precession}。}

\example{各向同性谐振子对二阶项的检验}\label{ex:c3-65}

\QuestionStatement{设 $U=m\omega_c^2r^2/2$，磁场沿 $z$，$\omega_c>0$。
旋转系中横向方程准确化为
\[
\ddot x'+(\omega_c^2+\Omega^2)x'=0,\qquad
\ddot y'+(\omega_c^2+\Omega^2)y'=0,
\]
纵向仍是 $\ddot z'+\omega_c^2z'=0$。
此处 $\omega_c$ 是无磁场谐振子固有频率，$\Omega=-qB/(2m)$ 为有向拉莫尔率。求惯性系两个有向横向模态频率，展开至二阶，并解释自由粒子极限。}{ex65}{例65题图}
% BEGIN DETAILED ANSWER ex65
\solution \Step{把两个横向方程合成一个复方程}
以图~\ref{fig:solution-ex65}~中谐振子横向圆偏振与磁场方向为依据，先写出图右侧第一行的几何或动力学关系，再代入其下方的计算式。

令 $w=x+\mathrm i y$、$b=qB/m$，由洛伦兹力的两个分量有
\[
\ddot x+\omega_c^2x=b\dot y,\qquad
\ddot y+\omega_c^2y=-b\dot x,
\]
因此 $\ddot w+\mathrm i b\dot w+\omega_c^2w=0$。
\SolutionFigure{ex65}{例65的计算示意}
\Step{准确代入拉莫尔旋转}
取 $w=e^{\mathrm i\Omega t}w'$，则
$\dot w=e^{\mathrm i\Omega t}(\dot w'+\mathrm i\Omega w')$，
$\ddot w=e^{\mathrm i\Omega t}(\ddot w'+2\mathrm i\Omega\dot w'-\Omega^2w')$。
代入后一阶导数系数为 $\mathrm i(2\Omega+b)$；令 $\Omega=-b/2$，得到
\[
\ddot w'+(\omega_c^2+\Omega^2)w'=0.
\]
纵向没有磁力，仍为 $\ddot z'+\omega_c^2z'=0$。
\Step{还原惯性系模态与二阶展开}
旋转系两模态为 $w'\propto e^{\pm\mathrm i\nu t}$，$\nu=\sqrt{\omega_c^2+\Omega^2}$，所以惯性系的有向横向频率为
\[
\sigma_\pm=\Omega\pm\sqrt{\omega_c^2+\Omega^2}.
\]
在 $|\Omega|\ll\omega_c$ 时，
$\sqrt{\omega_c^2+\Omega^2}=\omega_c+\Omega^2/(2\omega_c)+O(\Omega^4/\omega_c^3)$。
一阶共同变化来自拉莫尔旋转，回复频率的二阶改变仍保留，故此变换不是任意磁场下的完全消磁。
\Step{无场与自由粒子极限}
$B=0$ 时两有向频率为 $\pm\omega_c$。
准确式取 $\omega_c=0$ 时，两频率之一为零、另一为 $2\Omega=-qB/m$，恢复静止导引中心和完整回旋运动。
这个极限须用准确式，不能使用已经要求 $|\Omega|\ll\omega_c$ 的展开。

% END DETAILED ANSWER ex65
\subsection{磁偶极能、塞曼分裂与自旋轨道项}\label{sub:c3-10.4}
磁偶极相互作用能为
\[
U_Z=-\vect\mu\cdot\vect B.
\]
其取向导数与 $\vect\mu\times\vect B$ 的力矩一致。
这个完整能量没有额外的 $1/2$。
$\vect\Omega_L$ 只描述角动量方向的进动，
不能代入刚体的 $T=\vect\omega\cdot\vect L/2$
来定义所谓进动动能。
静磁场对一个运动电荷的洛伦兹力与其速度垂直，
$\vect v\cdot q(\vect v\times\vect B)=0$；
磁偶极的有效相互作用能与单个电荷的瞬时机械功
也必须在相应系统模型中解释，不能混用。

在量子角动量模型中，磁矩为 $\gamma\vect J$，
磁场沿 $z$ 时能级
$E_{m_J}=-\gamma\hbar m_JB$。
相邻磁量子数相差一，
能级间隔为 $\hbar|\gamma|B$，称塞曼分裂。
自旋 $1/2$ 的两个能级是
$E_\pm=\mp\hbar\gamma B/2$；
这里的 $1/2$ 来自允许的自旋投影 $\pm\hbar/2$，
并不表示相互作用能公式被除以二。
\footnote{磁偶极能与共振频率可参见
\href{https://ocw.mit.edu/courses/22-920-a-hands-on-introduction-to-nuclear-magnetic-resonance-january-iap-1997/38a44e04e8f48b1a8969b3a2e5fb6363_lecture_1.pdf}{MIT：Nuclear Spin and Magnetic Resonance}。}

自旋与轨道角动量的耦合属于另一个问题。
对中心势能 $U(r)$，标准非相对论单电子近似中，
包含 Thomas 相对论修正的自旋轨道项为
\[
H_{\text{SO}}=\frac{1}{2m_e^2c^2r}\frac{\dd U}{\dd r}\,
\vect L\cdot\vect S.
\]
该半因子来自连续洛伦兹变换产生的参考轴进动修正，
不能归因于经典磁进动动能。
这一项涉及自旋与相对论，超出纯经典刚体模型；
在此列出是为了明确经典进动、塞曼能和自旋轨道耦合各自的含义。
\footnote{Thomas 因子及中心势耦合形式参见
\href{https://ocw.u-tokyo.ac.jp/lecture_files/engin_06/6/notes/en/4-1E.PDF}{东京大学：A New Degree of Freedom in Quantum Mechanics: Spin}。}


\subsubsection*{由虚转角得到磁偶极能，再区分三种半因子}
对大小固定的磁矩作无穷小转动 $\delta\vect\phi$，
$\delta\vect\mu=\delta\vect\phi\times\vect\mu$。
若势能为 $U_Z=-\vect\mu\cdot\vect B$，则
\[
\delta U_Z=-\vect B\cdot(\delta\vect\phi\times\vect\mu)
=-\delta\vect\phi\cdot(\vect\mu\times\vect B).
\]
按虚功定义 $\delta W=\vect M\cdot\delta\vect\phi=-\delta U_Z$，
正好得到 $\vect M=\vect\mu\times\vect B$。
若错误地把这个势能除以二，就会把由同一势能计算的力矩也除以二，
与闭环受力积分冲突。

取 $\vect\mu=\gamma\vect J$，量子哈密顿量为 $H_Z=-\gamma BJ_z$。
使用角动量对易关系，可直接核对方向：
\[
\dot J_x=\frac{\mathrm i}{\hbar}[H_Z,J_x]=\gamma BJ_y,
\qquad
\dot J_y=-\gamma BJ_x,\qquad \dot J_z=0.
\]
期望值因而满足与经典磁偶极相同的进动方程。
能量本征值由 $J_z=\hbar m_J$ 给出，
而不是由进动角速度代入刚体动能公式给出。
若 $\gamma>0$，正投影态能量较低；若 $\gamma<0$，负投影态能量较低。
一个自旋 $1/2$ 的两能级间隔仍为 $\hbar|\gamma|B$。

这里出现的三个半因子来自不同操作：
磁矩积分中的 $1/2$ 来自有向面积的线积分；
自旋 $1/2$ 的能级投影为 $\pm\hbar/2$；
Thomas 半因子则修正自旋轨道相互作用。
它们都不能给完整偶极势能额外乘一个 $1/2$。
理想静场中的进动满足 $\dot U_Z=0$，因为夹角恒定。
使磁矩逐渐转到低能取向需要能量交换，理想无耗散进动不会自行完成这一过程。

\subsubsection*{中心电场中的自旋轨道项与 Thomas 修正步骤}
为明确模型，考虑一个电子在静止中心电势中运动，
令 $q=-e$、$U(r)=q\Phi(r)$ 为电子势能，
$\vect E=-\nabla\Phi=-U'(r)\vect e_r/q$。
以下保留 $v/c$ 的最低非零相对论修正，并取电子 $g_S=2$ 的近似。
电子瞬时共动系中的一阶磁场为
\[
\vect B'=-\frac{\vect v\times\vect E}{c^2}
=-\frac{U'(r)}{q m_e c^2r}\vect L,
\qquad \vect L=m_e\vect r\times\vect v.
\]
这个符号可由 $\vect v\times\vect e_r=-\vect L/(m_er)$ 核对。
电子自旋磁矩在此近似为 $\vect\mu_S=q\vect S/m_e$，
所以只做场变换会给出
\[
H_{\text{naive}}=-\vect\mu_S\cdot\vect B'
=\frac{U'(r)}{m_e^2c^2r}\vect L\cdot\vect S.
\]
但是连续改变速度的瞬时共动轴并不是一组在空间中无转动的轴。
在低速极限，Thomas 进动矢量为
\[
\vect\Omega_T=\frac{\vect a\times\vect v}{2c^2},\qquad
\vect a=-\frac{U'(r)}{m_e}\vect e_r,
\]
故
\[
\vect\Omega_T=-\frac{U'(r)}{2m_e^2c^2r}\vect L.
\]
当把自旋运动转回空间描述时，这个附加轴进动对应
$H_T=\vect\Omega_T\cdot\vect S$，恰抵消朴素磁相互作用的一半。
二者合起来才是
\[
H_{\text{SO}}=H_{\text{naive}}+H_T
=\frac{U'(r)}{2m_e^2c^2r}\vect L\cdot\vect S.
\]
\footnote{低速 Thomas 进动与自旋轨道项的半因子参见
\href{https://galileoandeinstein.phys.virginia.edu/Elec_Mag/18_7420/18_7420_20_Thomas_Precession.pdf}{弗吉尼亚大学：Thomas Precession}。}
因此该修正来源于洛伦兹变换的运动学，
与外加磁场下完整偶极势能的系数是两件独立的事。
此推导还假定中心场、单电子有效模型和 $v\ll c$；
它没有包含电子磁矩的反常修正、一般多电子相互作用或强相对论效应。
电子自旋也没有被解释为一个经典刚体表面的机械自转。

为了看出自旋轨道项如何造成能级分裂，令总角动量
$\vect J=\vect L+\vect S$。平方后有
\[
\vect L\cdot\vect S=\frac12(J^2-L^2-S^2).
\]
在共同的 $j,l,s$ 态中，其本征值为
\[
\frac{\hbar^2}{2}[j(j+1)-l(l+1)-s(s+1)].
\]
电子 $s=1/2$ 时，$j=l+1/2$ 对应 $\hbar^2l/2$，
$j=l-1/2$ 对应 $-\hbar^2(l+1)/2$；后一分支要求 $l\ge1$。
能量修正还须对径向系数 $U'(r)/(2m_e^2c^2r)$ 取该态的期望值。
外加磁场弱于已有自旋轨道分裂时，可在一个固定 $j$ 多重态内讨论有效陀螺磁比；
若外场足以混合不同耦合态，简单的固定 $\gamma\vect J$ 模型需重新检查。
这些条件把经典轨道拉莫尔定理、自旋轨道分裂和核磁共振所使用的磁矩参数区分开来。


\subsection{核磁共振的频率与射频激发}\label{sub:c3-10.5}
具有自旋的原子核有磁矩 $\vect\mu=\gamma_n\vect J$，
$\gamma_n$ 须用该核的陀螺磁比，不能直接用圆环的 $q/(2m)$。
在静磁场 $B_0\vect k$ 中，
\[
\omega_0=|\gamma_n|B_0,\qquad
f_0=\frac{|\gamma_n|B_0}{2\pi}.
\]
对自旋 $1/2$，能级间隔 $\Delta E=\hbar\omega_0$。
静场建立量子能级与净磁化，横向射频场在匹配频率附近
可以驱动磁化方向改变或能级间跃迁。

以 $\gamma_n>0$ 为例，取圆偏振射频场
\[
\vect B_{\text{RF}}
=B_1[\cos\omega_{\text{RF}}t\,\vect i
-\sin\omega_{\text{RF}}t\,\vect j],
\]
它沿磁化自然进动方向旋转。
在角速度 $-\omega_{\text{RF}}\vect k$ 的旋转系中，
\[
\vect B_{\text{有效}}
=B_1\vect i'+(B_0-\omega_{\text{RF}}/\gamma_n)\vect k.
\]
当 $\omega_{\text{RF}}=\gamma_nB_0$，
纵向有效场消失，磁化绕横向有效场转动，
翻转角为 $\beta=|\gamma_n|B_1\tau$。
持续时间 $\tau_{\pi/2}=\pi/(2|\gamma_n|B_1)$ 为理想 $90^\circ$ 脉冲，
$\tau_\pi=\pi/(|\gamma_n|B_1)$ 为理想 $180^\circ$ 脉冲。
这里 $B_1$ 是旋转分量的幅度；
线偏振余弦场分成两旋转分量时各有一半幅度，
不能未经说明沿用同一 $B_1$。

实际信号来自宏观净磁化，
受纵向恢复时间 $T_1$、横向退相干时间 $T_2$
及局部场和化学环境影响。
刚体进动提供频率与旋转系的经典直观，
核磁共振的能级与跃迁则需要量子理论。
\footnote{旋转系与射频脉冲参见
\href{https://ocw.mit.edu/courses/22-920-a-hands-on-introduction-to-nuclear-magnetic-resonance-january-iap-1997/2c3c010568e1561a18c5b472db1d84ec_lecture_2.pdf}{MIT：The Rotating Frame, RF Pulses and the Bloch Equations}。}

\TeachingFigure{c3-diagram-59}{能级间隔确定共振频率，横向有效场产生脉冲翻转}


\subsubsection*{旋转基矢、失谐符号与有效场的逐项推导}
以下仍明确取 $\gamma_n>0$，并把 $\omega_{\text{RF}}>0$ 定义为载频大小。
选择与上述圆偏振场一起转动的正交基矢
\[
\begin{aligned}
\vect i'&=\cos\omega_{\text{RF}}t\,\vect i
-\sin\omega_{\text{RF}}t\,\vect j,\\
\vect j'&=\sin\omega_{\text{RF}}t\,\vect i
+\cos\omega_{\text{RF}}t\,\vect j,\qquad \vect k'=\vect k.
\end{aligned}
\]
它们满足 $\vect i'\times\vect j'=\vect k$，
$\dot{\vect i}'=-\omega_{\text{RF}}\vect j'$、
$\dot{\vect j}'=\omega_{\text{RF}}\vect i'$，
因此轴组的实际角速度为 $-\omega_{\text{RF}}\vect k$。
宏观磁化的无弛豫方程为
$\dot{\vect M}_S=\gamma_n\vect M\times\vect B$。
用运输公式，旋转系中的导数为
\[
\begin{aligned}
\dot{\vect M}_R
&=\gamma_n\vect M\times\vect B
+\omega_{\text{RF}}\vect k\times\vect M\\
&=\gamma_n\vect M\times
\left(\vect B-\frac{\omega_{\text{RF}}}{\gamma_n}\vect k\right).
\end{aligned}
\]
这就是有效场的负号来源。
定义带符号失谐和带符号射频角频率
\[
\Delta=\gamma_nB_0-\omega_{\text{RF}},\qquad
\omega_1=\gamma_nB_1.
\]
于是
\[
\vect B_{\text{eff}}=\frac{\omega_1}{\gamma_n}\vect i'
+\frac{\Delta}{\gamma_n}\vect k.
\]
共振是 $\Delta=0$；不同教材如果把失谐定义成相反顺序，
所有含失谐的一阶符号也应相应反向。
对于 $\gamma_n<0$，自然进动反转，应选择相反旋向的射频分量和参考轴组；
不能仅把当前载频大小公式中的绝对值去掉而保留同一条圆偏振波。

把磁化写成 $M_x\vect i'+M_y\vect j'+M_z\vect k$，三条方程为
\[
\dot M_x=\Delta M_y,\qquad
\dot M_y=\omega_1M_z-\Delta M_x,\qquad
\dot M_z=-\omega_1M_y.
\]
这里下标 $x,y$ 表示旋转轴分量，$M_z$ 同时也是空间竖直分量。
$M_x^2+M_y^2+M_z^2$ 保持不变，且
$\omega_1M_x+\Delta M_z$ 恒定；
磁化沿着以有效场为轴的圆锥转动，而不是总沿原静场转动。

\subsubsection*{任意恒定失谐的矩形脉冲解}
定义
\[
\Omega_R=\sqrt{\omega_1^2+\Delta^2},\qquad
\vect n_R=\frac{\omega_1\vect i'+\Delta\vect k}{\Omega_R}.
\]
在脉冲内 $B_1$、失谐和 RF 相位均恒定，且忽略弛豫时，
磁化绕 $\vect n_R$ 的有向转角为 $-\Omega_Rt$。
由 Rodrigues 公式，任意初值的解为
\[
\begin{aligned}
\vect M(t)={}&\vect M_0\cos\Omega_Rt
-(\vect n_R\times\vect M_0)\sin\Omega_Rt\\
&+\vect n_R(\vect n_R\cdot\vect M_0)(1-\cos\Omega_Rt).
\end{aligned}
\]
若初态为热平衡磁化 $\vect M_0=M_0\vect k$，逐项投影得到
\[
\begin{aligned}
M_x(t)&=M_0\frac{\omega_1\Delta}{\Omega_R^2}(1-\cos\Omega_Rt),\\
M_y(t)&=M_0\frac{\omega_1}{\Omega_R}\sin\Omega_Rt,\\
M_z(t)&=M_0\left(\frac{\Delta^2}{\Omega_R^2}
+\frac{\omega_1^2}{\Omega_R^2}\cos\Omega_Rt\right).
\end{aligned}
\]
可直接将 $t=0$ 代回核对 $M_x=M_y=0,M_z=M_0$，
并核对初始导数 $\dot M_y(0)=\omega_1M_0$。
失谐存在时，磁化围绕倾斜的有效场转动，
即便把脉冲时间取为 $\pi/\Omega_R$，也未必能完全反转：
\[
\frac{M_z(\pi/\Omega_R)}{M_0}
=\frac{\Delta^2-\omega_1^2}{\Delta^2+\omega_1^2}.
\]
若 $|\Delta|\gg|\omega_1|$，横向激发很弱。
这解释了给定幅度的脉冲为何只能有效激发有限的频率范围。
当 $\omega_1=\Delta=0$ 时磁化恒定，不需要定义 $\vect n_R$。

\subsubsection*{共振的 $90^\circ$、$180^\circ$ 脉冲与 RF 相位}
在 $\Delta=0$、$B_1>0$ 的当前约定下，令 $\beta=\omega_1t$，
\[
M_x=0,\qquad M_y=M_0\sin\beta,\qquad M_z=M_0\cos\beta.
\]
于是 $90^\circ$ 脉冲给出 $(0,M_0,0)$，
$180^\circ$ 脉冲给出 $(0,0,-M_0)$。
磁化绕正 $\vect i'$ 方向的有向转角实际上为 $-\beta$，
与 $\dot{\vect M}=\gamma_n\vect M\times\vect B$ 的叉积顺序一致。
常说的“翻转角为 $90^\circ$”是转角大小，不能用它独自判定转向。

若 RF 相位使共振有效场沿
$\cos\alpha\,\vect i'+\sin\alpha\,\vect j'$，对应空间场为
\[
\vect B_{\text{RF}}=B_1[
\cos(\omega_{\text{RF}}t-\alpha)\vect i
-\sin(\omega_{\text{RF}}t-\alpha)\vect j].
\]
从同一初态出发，磁化分量成为
\[
M_x=-M_0\sin\alpha\sin\beta,\quad
M_y=M_0\cos\alpha\sin\beta,\quad
M_z=M_0\cos\beta.
\]
因而相位决定 $90^\circ$ 脉冲后横向磁化的方向。
相同幅度和持续时间的脉冲可以产生不同的横向相位，
这不是共振频率改变，而是旋转轴改变。
把旋转系分量换回空间系时，应使用
\[
M_{x,S}=M_x\cos\omega_{\text{RF}}t+M_y\sin\omega_{\text{RF}}t,
\]
\[
M_{y,S}=-M_x\sin\omega_{\text{RF}}t+M_y\cos\omega_{\text{RF}}t.
\]
脉冲终点的空间方向还包含载频累积的相位。

\subsubsection*{线偏振分解与短脉冲近似的条件}
若实验给出的横向场是
$\vect B_{\text{lin}}=B_a\cos\omega_{\text{RF}}t\,\vect i$，准确恒等式为
\[
\begin{aligned}
\vect B_{\text{lin}}={}&\frac{B_a}{2}
(\cos\omega_{\text{RF}}t\,\vect i-\sin\omega_{\text{RF}}t\,\vect j)\\
&+\frac{B_a}{2}
(\cos\omega_{\text{RF}}t\,\vect i+\sin\omega_{\text{RF}}t\,\vect j).
\end{aligned}
\]
第一项与正 $\gamma_n$ 磁化同向进动，在选定旋转系中静止；
第二项在这个系中仍以约两倍载频转动。
若 $\omega_0$ 远大于 RF 耦合率和失谐，快速旋转分量的最低阶平均影响可忽略，
这时脉冲公式应取 $B_1=B_a/2$。
圆偏振场则不存在这个幅度再除以二的步骤。
强激发或极高精度下，反向旋转分量会产生更高阶修正，
不能把它在所有参数范围内都宣称为严格为零。

上述无弛豫矩形脉冲解还要求 $\tau\ll T_1,T_2$，
脉冲幅度和相位在所用时间内近似不变，且局部失谐可以视为恒定。
若要让一组略有频率分散的自旋近似产生同一翻转角，
还要求 $|\omega_1|$ 大于该组的典型失谐范围。
脉冲时间只决定理想翻转的一部分，空间幅度不均匀和失谐都会改变结果。

\example{从给定陀螺磁比计算载频和脉冲时间}\label{ex:c3-66}

\QuestionStatement{某自旋 $1/2$ 核的参数在本题中给定为
$\gamma_n=2.68\times10^8\,\mathrm{rad\,s^{-1}\,T^{-1}}$，
取 $B_0=3.00\,\mathrm T$、圆偏振旋转幅度
$B_1=10.0\,\mathrm{\mu T}$。忽略弛豫并调到共振。
求共振角频率、普通频率、共振翻转角速率及理想 $90^\circ$、$180^\circ$ 脉冲时间；说明若 $B_1$ 改指线偏振峰值，答案应如何变化。}{ex66}{例66题图}
% BEGIN DETAILED ANSWER ex66
\solution \Step{统一磁场幅度与频率单位}
以图~\ref{fig:solution-ex66}~中纵向静场与横向共振旋转场为依据，先写出图右侧第一行的几何或动力学关系，再代入其下方的计算式。

本题给定 $\gamma_n$ 的单位为 $\mathrm{rad\,s^{-1}\,T^{-1}}$，因此先得到角频率；普通频率必须另除 $2\pi$。
$B_1=10.0\,\mathrm{\mu T}=1.00\times10^{-5}\,\mathrm T$ 已是同向共振旋转分量，不须先减半。
\SolutionFigure{ex66}{例66的计算示意}
\Step{计算静场的共振频率}
\[
\omega_0=|\gamma_n|B_0=(2.68\times10^8)(3.00)
=8.04\times10^8\,\mathrm{rad\,s^{-1}},
\]
\[
f_0=\frac{\omega_0}{2\pi}=1.27960574\times10^8\,\mathrm{Hz}
\simeq127.96\,\mathrm{MHz}.
\]
若误把 $8.04\times10^8$ 直接当作 Hz，会多出 $2\pi$ 倍。
\Step{在共振旋转系中积分翻转角}
忽略弛豫和失谐，横向有效场导致角速率
$\omega_1=|\gamma_n|B_1=2680\,\mathrm{rad\,s^{-1}}$。
恒幅矩形脉冲的翻转角为 $\beta=\omega_1\tau$，故
\[
\tau_{90}=\frac{\pi}{2\omega_1}=5.86118032\times10^{-4}\,\mathrm s
\simeq0.5861\,\mathrm{ms},
\]
\[
\tau_{180}=\frac{\pi}{\omega_1}=1.17223606\times10^{-3}\,\mathrm s
\simeq1.1722\,\mathrm{ms}.
\]
两时间之比为 $1:2$，且与 $B_1$ 成反比。
\Step{线偏振峰值的约定}
线偏振场 $B_{\text{pk}}\cos\omega t\,\vect i$ 可分为两个反向旋转、幅度各为 $B_{\text{pk}}/2$ 的分量。
若 $10.0\,\mathrm{\mu T}$ 改指该峰值，在旋转波近似下有效 $B_1$ 为 $5.00\,\mathrm{\mu T}$，脉冲时间分别变为 $1.1722$、$2.3445\,\mathrm{ms}$，载频不变。
静场决定共振载频，射频旋转分量决定共振翻转速率；失谐时须改用倾斜有效场，不能直接沿用这些时间。

% END DETAILED ANSWER ex66
\subsubsection*{净磁化来自占据数差，而不是全部自旋同向}
考虑数密度 $n_s$ 的互不耦合自旋 $1/2$ 核，$\gamma_n>0$。
两能级为 $E_+=-\hbar\gamma_nB_0/2$ 与
$E_- =+\hbar\gamma_nB_0/2$。
热平衡时玻尔兹曼权重给出
\[
\frac{N_+}{N_-}=\exp\left(\frac{\hbar\gamma_nB_0}{k_BT}\right),
\qquad
\frac{N_+-N_-}{N_++N_-}
=\tanh\left(\frac{\hbar\gamma_nB_0}{2k_BT}\right).
\]
每个核的磁矩投影为 $\pm\hbar\gamma_n/2$，因此磁化密度为
\[
M_0=n_s\frac{\hbar\gamma_n}{2}
\tanh\left(\frac{\hbar\gamma_nB_0}{2k_BT}\right).
\]
在 $\hbar\gamma_nB_0\ll k_BT$ 的高温极限，
\[
M_0\simeq\frac{n_s\hbar^2\gamma_n^2B_0}{4k_BT}.
\]
所以热平衡净磁化通常只是两个大量占据群体之间的小差值，
并非每一个核都已沿磁场排列。
RF 脉冲旋转的是这个统计平均磁化；
单个核的能级与跃迁仍由量子态描述。

\subsubsection*{Bloch 弛豫方程与脉冲后信号}
仅靠理想进动无法恢复热平衡。
在宏观描述中，加入纵向恢复时间 $T_1$、横向弛豫时间 $T_2$，
旋转系中的 Bloch 方程为
\[
\begin{aligned}
\dot M_x&=\Delta M_y-\frac{M_x}{T_2},\\
\dot M_y&=\omega_1M_z-\Delta M_x-\frac{M_y}{T_2},\\
\dot M_z&=-\omega_1M_y-\frac{M_z-M_0}{T_1}.
\end{aligned}
\]
这些弛豫项是描述环境作用的宏观模型，不由一个保守刚体的欧拉方程产生。
$T_1$ 描述纵向磁化向 $M_0$ 的恢复，涉及自旋系统与环境的能量交换；
$T_2$ 描述横向相干的衰减，不能把它简单等同于单个磁矩大小减小。

RF 在时刻 $t=\tau$ 关闭后，记 $s=t-\tau\ge0$。
在空间系中令 $M_+=M_{x,S}+\mathrm iM_{y,S}$，解为
\[
M_+(t)=M_+(\tau)\exp[-s/T_2-\mathrm i\gamma_nB_0s],
\]
\[
M_z(t)=M_0+[M_z(\tau)-M_0]e^{-s/T_1}.
\]
因此理想 $90^\circ$ 脉冲之后
\[
|M_+|=M_0e^{-s/T_2},\qquad
M_z=M_0(1-e^{-s/T_1});
\]
理想 $180^\circ$ 脉冲之后则
\[
M_z=M_0(1-2e^{-s/T_1}).
\]
后一式在 $s=T_1\ln2$ 处通过零点，是反转恢复测量的基本关系。
横向磁化的旋转在接收线圈中产生随时间衰减的信号，称自由感应衰减。
信号的频率来自拉莫尔进动，包络衰减则来自另一类过程。

实际样品不同位置还可能有不同的静场偏移 $\delta\omega$。
集合信号于是为
\[
\langle M_+(s)\rangle=M_+(0)e^{-s/T_2-\mathrm i\omega_0s}
\int p(\delta\omega)e^{-\mathrm i\delta\omega s}\,\dd(\delta\omega),
\]
其中 $p$ 是归一化的频率分布。
这额外的相位分散使观测衰减常快于纯 $T_2$ 衰减，通常用 $T_2^*$ 表征。
只有当分布和衰减都满足相应的指数模型时，才可把额外衰减率简单相加；
一般分布会给出不同包络，不能把所有 $T_2^*$ 情形都固定成一个普适的加法公式。

若持续施加恒幅 RF，而不满足短脉冲条件，
令三个导数为零可求出稳态。
先由第一式得 $M_x=\Delta T_2M_y$，
由第三式得 $M_z=M_0-\omega_1T_1M_y$。
代入第二式，定义
$D_B=1+\Delta^2T_2^2+\omega_1^2T_1T_2$，得到
\[
M_y=\frac{M_0\omega_1T_2}{D_B},\qquad
M_x=\frac{M_0\Delta\omega_1T_2^2}{D_B},\qquad
M_z=\frac{M_0(1+\Delta^2T_2^2)}{D_B}.
\]
共振时 $\omega_1^2T_1T_2$ 很大，纵向占据数差明显减小，表现为饱和。
这与一个短时间内近似保持 $|\vect M|$ 的理想 $90^\circ$ 脉冲属于不同实验条件。
把频率、脉冲旋转和弛豫三个环节连起来，才得到宏观核磁共振信号的基本描述。


\clearpage\section{练习与解答}\label{sec:c3-11}
\subsection{运动学与平衡}\label{sub:c3-11.1}
\exercise{距离不变是否意味着相对位置矢量不变？}\label{test:c3-1}

\QuestionStatement{设刚体上两个材料点 $A,B$ 的距离恒为 $l>0$。
写出该条件对相对速度的限制，并以平面定轴转动说明它为什么不要求相对速度为零。}{tr01}{训练题1题图}
% BEGIN DETAILED ANSWER tr01
\solution \Step{距离约束的一阶导数}
以图~\ref{fig:solution-tr01}~中连线 $AB$ 与垂直于它的相对速度为依据，先写出图右侧第一行的几何或动力学关系，再代入其下方的计算式。
 令 $\vect r=\vect r_B-\vect r_A$。刚性约束为
\[
\vect r\cdot\vect r=l^2.
\]
在同一惯性系中求导，得到
\[
2\vect r\cdot(\vect v_B-\vect v_A)=0.
\]
\SolutionFigure{tr01}{训练题1的计算示意}
\Step{区分长度与方向} 所以相对速度在连线方向的投影为零，而垂直连线的分量可以存在。
例如 $A$ 固定，$B$ 以恒定角速度 $\omega$ 绕 $A$ 转动，则
\[
\vect r=l(\cos\omega t\,\vect i+\sin\omega t\,\vect j),\qquad
\dot{\vect r}=l\omega(-\sin\omega t\,\vect i+\cos\omega t\,\vect j).
\]
两式的内积为零，但 $\dot{\vect r}$ 的大小是 $l|\omega|$。
\Step{加速度约束} 进一步求导可得
\[
\vect r\cdot(\vect a_B-\vect a_A)
=-|\vect v_B-\vect v_A|^2.
\]
这给出相对加速度沿连线方向的分量，说明长度保持不变仍可具有向心加速度。
刚性约束限制的是距离；只有纯平动时，各材料点之间的位置矢量在空间也保持不变。
\Step{加速度关系与反例}
对给出的匀速圆周运动，$\vect a_B-\vect a_A=-\omega^2\vect r$，左侧内积为 $-\omega^2l^2$；右侧 $-|\vect v_B-\vect v_A|^2$ 也是 $-\omega^2l^2$。
长度约束的一阶式排除沿杆伸缩，二阶式要求相对速度非零时具有适当的向心加速度。若只写 $\vect v_A=\vect v_B$，便把刚性条件错误缩小为纯平动条件。

% END DETAILED ANSWER tr01
\exercise{同样沿圆周运动的座舱和刚性轮辐有何区别？}\label{test:c3-2}

\QuestionStatement{一座舱的悬挂点沿半径 $R$ 的圆周运动，座舱由机构保持方向恒定；
另一物体是一根与转轮固连的轮辐。两者都以圆周角速度 $\Omega$ 运行，
分别写出任意材料点的速度，并判断运动类型。}{tr02}{训练题2题图}
% BEGIN DETAILED ANSWER tr02
\solution \Step{座舱的固定姿态}
以图~\ref{fig:solution-tr02}~中相同圆周轨迹下的固定姿态座舱为依据，先写出图右侧第一行的几何或动力学关系，再代入其下方的计算式。
 对座舱，取悬挂点 $A$ 为基点，任意点 $P$ 满足
\[
\vect r_P=\vect r_A+\vect b_P,
\]
其中 $\vect b_P$ 在空间固定。因此
\[
\vect v_P=\vect v_A,\qquad \vect a_P=\vect a_A,\qquad \vect\omega=\vect0.
\]
各点轨迹可以都是圆，但圆心相差固定矢量 $\vect b_P$，
这属于曲线平动。若圆周运行匀速，各点加速度的大小均为 $R\Omega^2$，
方向也相同。
\SolutionFigure{tr02}{训练题2的计算示意}
\Step{轮辐的速度场} 对轮辐，取轮心 $O$ 固定、角速度 $\vect\omega=\Omega\vect k$，
则
\[
\vect v_P=\Omega\vect k\times\vect r_{OP},\qquad
\vect a_P=-\Omega^2\vect r_{OP}
\]
适用于轮平面中的点。距轮心不同的点，速度和加速度大小均不同。
轮辐上任意两点的连线随时间转动，因此是定轴转动。

\Step{以轮辐上的两点作直接比较}
若两点距轮心分别为 $r_1,r_2$，同一瞬时的速率为 $|\Omega|r_1,|\Omega|r_2$；当 $r_1\ne r_2$ 时显然不同。
座舱上两点之差则为固定 $\vect b_P-\vect b_Q$，其导数为零，任意两点速度相同。
所以应检验连线方向是否随时间改变。即使座舱悬挂点和某轮辐端点沿同一个圆周运行，它们也只是在那一个点上共享轨迹与速度，不能据此判断整个物体的运动。

% END DETAILED ANSWER tr02
\exercise{能否直接把两个有限转角当矢量相加？}\label{test:c3-3}

\QuestionStatement{先绕固定 $x$ 轴转 $90^\circ$，再绕固定 $y$ 轴转 $90^\circ$；
与顺序颠倒的结果比较。再说明无穷小转动为何可以在一阶精度下相加。}{tr03}{训练题3题图}
% BEGIN DETAILED ANSWER tr03
\solution \Step{按次序写转动矩阵}
以图~\ref{fig:solution-tr03}~中标记 $\vect e_z$ 经两种转动顺序的结果为依据，先写出图右侧第一行的几何或动力学关系，再代入其下方的计算式。
 采用主动转动，作用于列矢量的矩阵为
\[
R_x=\begin{pmatrix}1&0&0\\0&0&-1\\0&1&0\end{pmatrix},\qquad
R_y=\begin{pmatrix}0&0&1\\0&1&0\\-1&0&0\end{pmatrix}.
\]
取初始矢量 $\vect e_z=(0,0,1)^{\mathsf T}$，先 $x$ 后 $y$ 得
\[
R_yR_x\vect e_z=(0,-1,0)^{\mathsf T},
\]
先 $y$ 后 $x$ 则得
\[
R_xR_y\vect e_z=(1,0,0)^{\mathsf T}.
\]
\SolutionFigure{tr03}{训练题3的计算示意}
\Step{比较同一标记方向} 两个姿态不相同，而通常的矢量加法与顺序无关，故有限转动不能直接作矢量相加。
\Step{无穷小极限} 若两次小转动分别为 $\delta\vect\alpha,\delta\vect\beta$，
记 $[\vect a]_\times\vect r=\vect a\times\vect r$，则
\[
R_\alpha=\mathsf1+[\delta\vect\alpha]_\times+O(\delta^2),\qquad
R_\beta R_\alpha
=\mathsf1+[\delta\vect\alpha+\delta\vect\beta]_\times+O(\delta^2).
\]
次序的影响包含在二阶项中。一阶小转动可相加；
除以同一时间间隔并取极限，得到瞬时角速度的矢量合成。
这一步的极限不能用于把一段时间内的有限转角直接相加。
\Step{用二阶项定位不交换的来源}
写成 $R_x(\epsilon)=\mathsf1+\epsilon T_x+O(\epsilon^2)$，对另一轴同理，则两种次序之差为
\[
R_yR_x-R_xR_y=\epsilon^2(T_yT_x-T_xT_y)+O(\epsilon^3).
\]
这说明不是转动角大小造成一阶矛盾，而是二阶耦合记录了次序。
取 $\epsilon\sim\dd t$ 后除以 $\dd t$，差异趋于零，才允许瞬时角速度按矢量合成；对 $90^\circ$ 不存在这种截断依据。

% END DETAILED ANSWER tr03
\exercise{已知两点速度方向平行，能否必然找到有限瞬心？}\label{test:c3-4}

\QuestionStatement{平面刚体上 $A=(0,0)$、$B=(d,0)$，其中 $d>0$。
已知 $\vect v_A=u\vect j$、$\vect v_B=w\vect j$。
求瞬时角速度和速度为零的点；讨论 $u=w\ne0$ 的情形。}{tr04}{训练题4题图}
% BEGIN DETAILED ANSWER tr04
\solution \Step{求角速度}
以图~\ref{fig:solution-tr04}~中$A,B$ 两点沿 $y$ 的速度差为依据，先写出图右侧第一行的几何或动力学关系，再代入其下方的计算式。
 由两点速度关系，
\[
(w-u)\vect j=\omega\vect k\times(d\vect i)=\omega d\vect j,
\qquad \omega=\frac{w-u}{d}.
\]
\SolutionFigure{tr04}{训练题4的计算示意}
\Step{求有限瞬心位置} 设瞬心 $P=(x_P,y_P)$，利用 $\vect v_A=\omega\vect k\times\vect r_{PA}$，
\[
u\vect j=\omega(y_P\vect i-x_P\vect j).
\]
当 $w\ne u$ 时，$\omega\ne0$，故
\[
y_P=0,\qquad x_P=-\frac{u}{\omega}=-\frac{ud}{w-u}.
\]
这样才能由大小和几何位置确定有限瞬心。
\Step{讨论退化速度} 若 $u=w\ne0$，则 $\omega=0$，所有材料点瞬时具有相同非零速度。
不存在有限的速度零点；随着 $w-u\to0$，以上 $x_P$ 发散。
若 $u=w=0$，则本瞬时全体点速度为零，也不能从速度数据选出唯一瞬心。
平行速度既可能来自转动，也可能来自瞬时平动。作垂线相交时，还要处理重合、平行与零速度等退化情形。
\Step{数值特例}
若 $w=2u\ne0$，则 $\omega=u/d$、$x_P=-d$。点 $A$ 距瞬心 $d$、点 $B$ 距瞬心 $2d$，故其速度大小比确为 $1:2$，且同向。
若 $w=-u$，则 $x_P=d/2$，两点在瞬心两侧、速度反向。
“速度方向平行”只确定法线方向，必须加上速度大小和刚体关系才能定位瞬心；瞬时平动正是有限瞬心不存在的极限。

% END DETAILED ANSWER tr04
\exercise{欧拉角 $\theta=0$ 时，哪些角仍可确定？}\label{test:c3-5}

\QuestionStatement{使用 $z$--$x$--$z$ 欧拉角约定
$R=R_z(\varphi)R_x(\theta)R_z(\psi)$。
在 $\theta=0$ 时求姿态矩阵，并说明角变量为何不唯一。}{tr05}{训练题5题图}
% BEGIN DETAILED ANSWER tr05
\solution \Step{合并重合轴的转动}
以图~\ref{fig:solution-tr05}~中重合的 $z,\zeta$ 轴和水平面内的合角为依据，先写出图右侧第一行的几何或动力学关系，再代入其下方的计算式。
 当 $\theta=0$，中间矩阵为单位矩阵，因此
\[
R=R_z(\varphi)R_z(\psi)=R_z(\varphi+\psi).
\]
\SolutionFigure{tr05}{训练题5的计算示意}
\Step{说明角度非唯一性} 两个 $z$ 轴重合，绕它们的转动只通过总角
$\chi=\varphi+\psi$ 影响姿态。对任意函数 $f(t)$，变换
\[
\varphi\longmapsto\varphi+f(t),\qquad
\psi\longmapsto\psi-f(t)
\]
不改变 $R$。所以此时不能分别测定 $\varphi,\psi$。
\Step{角速度与穿越奇异点} 若一段运动始终保持 $\theta=0$，则
\[
\vect\omega=\dot\chi\,\vect e_z
=(\dot\varphi+\dot\psi)\vect e_z.
\]
如果运动只是瞬时通过 $\theta=0$，角速度仍应由
$\dot R R^{\mathsf T}$ 求出，不能把穿越时的横向角速度也删去。
欧拉角反演公式中的 $1/\sin\theta$ 表示坐标退化，而不是刚体运动本身发散。
\Step{写出可测量的矩阵并回代}
令 $\chi=\varphi+\psi$，则
\[
R=\begin{pmatrix}\cos\chi&-\sin\chi&0\\\sin\chi&\cos\chi&0\\0&0&1\end{pmatrix}.
\]
因此可由 $\chi=\operatorname{atan2}(R_{21},R_{11})$ 求总方位，却无法再从该矩阵分离两角。
例如 $(\varphi,\psi)=(0,\chi)$ 和 $(\chi/2,\chi/2)$ 产生完全相同的姿态。若时间变化的 $f$ 同时加到一角、减到另一角，两角率也发生变化而总角速度不变。

% END DETAILED ANSWER tr05
\exercise{为什么轴向角速度不一定等于欧拉角自转率？}\label{test:c3-6}

\QuestionStatement{设刚体以倾角 $\theta$ 进动，方位率 $\dot\varphi$，
绕节点轴的倾角变化率 $\dot\theta$，自转角率 $\dot\psi$。
求总角速度在本体 $z$ 轴上的分量，并给出匀速进动时的区别。}{tr06}{训练题6题图}
% BEGIN DETAILED ANSWER tr06
\solution \Step{合成三次瞬时转动}
以图~\ref{fig:solution-tr06}~中进动方向在本体轴 $\vect n$ 上的投影为依据，先写出图右侧第一行的几何或动力学关系，再代入其下方的计算式。
 总角速度由三次瞬时转动合成：
\[
\vect\omega=\dot\varphi\,\vect e_\zeta+
\dot\theta\,\vect e_N+\dot\psi\,\vect e_z.
\]
节点轴同时垂直于空间 $\zeta$ 轴与本体 $z$ 轴，
故 $\vect e_N\cdot\vect e_z=0$，而
$\vect e_\zeta\cdot\vect e_z=\cos\theta$。
\SolutionFigure{tr06}{训练题6的计算示意}
\Step{投影至本体轴} 两侧点乘 $\vect e_z$，得
\[
\omega_3=\dot\psi+\dot\varphi\cos\theta.
\]
\Step{匀速进动与特例} 若 $\dot\theta=0$、$\dot\varphi=\Omega$，则保持轴向角速度 $s$ 所需的自转角率为
\[
\dot\psi=s-\Omega\cos\theta.
\]
例如 $\theta=60^\circ$、$\Omega=2\,\mathrm{s^{-1}}$、
$s=10\,\mathrm{s^{-1}}$，相应 $\dot\psi=9\,\mathrm{s^{-1}}$。
这个差值来自进动角速度在本体轴上的投影。
当 $\theta=90^\circ$ 时进动投影为零，两者相等；
更一般地，$\omega_3=\dot\psi$ 的条件是 $\dot\varphi\cos\theta=0$，
所以 $\dot\varphi=0$ 时也有这一等式。
在始终保持 $\theta=0$ 的运动中，姿态及轴向角速度分别由
$\varphi+\psi$ 及 $\dot\varphi+\dot\psi$ 确定，不能将两角率分别测定。
\Step{总角速率与轴向投影}
因为节点轴与另外两根角速度方向都垂直，另两根方向夹角为 $\theta$，有
\[
\omega^2=\dot\theta^2+\dot\varphi^2\sin^2\theta+
(\dot\psi+\dot\varphi\cos\theta)^2.
\]
最后一项正是轴向分量的平方；第一、二项是横向分量的平方。
将匀速进动数值例代入，轴向为 $10$，横向为 $2\sin60^\circ=\sqrt3$，因此总角速率为 $\sqrt{103}\,\mathrm{s^{-1}}$，同样不能拿 $\dot\psi=9$ 替代。

% END DETAILED ANSWER tr06
\exercise{平面三点取矩为零为什么要求三点不共线？}\label{test:c3-7}

\QuestionStatement{设平面力系对三个点 $A,B,D$ 的力矩都为零。
证明三个点不共线时力系平衡，并构造三个点共线时的反例。}{tr07}{训练题7题图}
% BEGIN DETAILED ANSWER tr07
\solution \Step{改变取矩点}
以图~\ref{fig:solution-tr07}~中由不共线三点形成的有向面积为依据，先写出图右侧第一行的几何或动力学关系，再代入其下方的计算式。
 令合力为 $\vect F=F_x\vect i+F_y\vect j$。
取矩点平移关系为
\[
M_B=M_A-(\vect r_{AB}\times\vect F)\cdot\vect k,\qquad
M_D=M_A-(\vect r_{AD}\times\vect F)\cdot\vect k.
\]
\SolutionFigure{tr07}{训练题7的计算示意}
\Step{利用非零行列式} 三个矩都为零，意味着
\[
\begin{pmatrix}-y_{AB}&x_{AB}\\-y_{AD}&x_{AD}\end{pmatrix}
\begin{pmatrix}F_x\\F_y\end{pmatrix}
=\begin{pmatrix}0\\0\end{pmatrix}.
\]
系数行列式为
$x_{AB}y_{AD}-y_{AB}x_{AD}$，其绝对值等于三角形 $ABD$ 面积的两倍。
不共线时行列式非零，故 $F_x=F_y=0$，再结合 $M_A=0$，力系平衡。
\Step{构造退化反例} 若三点在 $x$ 轴上，一个非零力 $F\vect i$ 的作用线也在 $x$ 轴上，
它对三点的力矩都为零，但合力不为零。
所以仅有三条取矩方程还不够，取矩点的几何独立性不可省略。
\Step{把几何独立性与方程独立性联系}
不共线保证 $\vect r_{AB}$、$\vect r_{AD}$ 张成整个平面；同一个合力若同时平行于这两个不同方向，只能是零矢量。
共线时两条“相对矩”方程实际给出同一个力分量限制，不能确定沿该直线的合力。
反例中的力沿 $x$ 轴，即使再增加若干同轴取矩点，它们的矩仍全部为零，故增加退化条件的数量不能补足缺少的独立条件。

% END DETAILED ANSWER tr07
\exercise{有静摩擦的无滑滚动为何可以守恒机械能？}\label{test:c3-8}

\QuestionStatement{半径 $a$ 的圆轮在固定轨道上纯滚动，接触处静摩擦为 $\vect f$。
把摩擦的平动功率与转动功率分别写出，说明两项如何相消。}{tr08}{训练题8题图}
% BEGIN DETAILED ANSWER tr08
\solution \Step{写无滑动约束}
以图~\ref{fig:solution-tr08}~中接触点的静摩擦和纯滚动约束为依据，先写出图右侧第一行的几何或动力学关系，再代入其下方的计算式。
 在固定轨道上无滑动时，接触材料点 $P$ 的瞬时速度为零：
\[
\vect v_P=\vect v_C+\vect\omega\times\vect r_{CP}=\vect0.
\]
\SolutionFigure{tr08}{训练题8的计算示意}
\Step{展开摩擦的两项功率} 摩擦力对刚体整体的功率为
\[
\mathcal P_f=\vect f\cdot\vect v_C+
(\vect r_{CP}\times\vect f)\cdot\vect\omega
=\vect f\cdot(\vect v_C+\vect\omega\times\vect r_{CP})=0.
\]
等式使用了标量三重积的循环关系。
以轮向下滚动为例，沿坡向上的静摩擦对质心平动做负功，
但对转动做同样大小的正功，二者之和为零。
因此静摩擦能够重新分配平动和转动能量而不耗散总机械能。
\Step{明确守恒条件} 这一结论须同时满足固定支承、无滑动、无滚动阻力以及无其他耗散。
先解出所需 $f$，再核查 $|f|\le\mu_sN$；如果不满足，必须改用滑动模型。
若轨道本身运动，接触速度等于轨道速度，摩擦功率一般不再为零。
\Step{沿斜面的有向功率}
令轮心下坡速率为 $v$、向上静摩擦大小为 $f$，则质心平动功率为 $-fv$。
摩擦力矩大小为 $fa$，与实际滚动角速度同向，转动功率为 $(fa)|\omega|=fv$，使用了 $v=a|\omega|$。
两项相加为零。若发生滑动，接触速度不再为零，总摩擦功率等于 $\vect f\cdot\vect v_P<0$，机械能随之损失。
所以先验证无滑动条件，再选择机械能守恒方程，逻辑次序不能颠倒。

% END DETAILED ANSWER tr08
\exercise{取矩中心改动时主矢和主矩怎样变化？}\label{test:c3-9}

\QuestionStatement{力系含若干力 $\vect F_i$，作用点为 $P_i$，并可含自由力偶 $\vect M_c$。
由定义推导从 $O$ 改为 $A$ 的力系化简关系。}{tr09}{训练题9题图}
% BEGIN DETAILED ANSWER tr09
\solution \Step{定义主矢与主矩}
以图~\ref{fig:solution-tr09}~中从 $O,A$ 指向同一作用点 $P_i$ 的位置关系为依据，先写出图右侧第一行的几何或动力学关系，再代入其下方的计算式。
 主矢定义为
\[
\vect F=\sum_i\vect F_i,
\]
它不涉及取矩点，因此保持不变。主矩为
\[
\vect M_O=\sum_i\vect r_{OP_i}\times\vect F_i+\vect M_c.
\]
\SolutionFigure{tr09}{训练题9的计算示意}
\Step{展开位置矢量差} 由于 $\vect r_{AP_i}=\vect r_{OP_i}-\vect r_{OA}$，
\[
\begin{aligned}
\vect M_A
&=\sum_i(\vect r_{OP_i}-\vect r_{OA})\times\vect F_i+\vect M_c\\
&=\vect M_O-\vect r_{OA}\times\vect F.
\end{aligned}
\]
\Step{力偶与有向力矩} 所以当 $\vect F=\vect0$ 时，主矩对任意取矩点都相同，
这正是力偶的性质。若 $\vect F\ne\vect0$，沿合力方向移动取矩点也不改变主矩，
但一般方向的移动会改变主矩。
作为符号检查，设向上的力 $F\vect j$ 作用于 $x=b$，
原点矩为 $bF\vect k$。把取矩点移到 $A=(b,0)$，
公式给出 $M_A=bF-bF=0$，与作用点取矩一致。
\Step{与原点无关的轴向不变量}
将主矩变换式点乘主矢，有
\[
\vect F\cdot\vect M_A=\vect F\cdot\vect M_O-
\vect F\cdot(\vect r_{OA}\times\vect F)=\vect F\cdot\vect M_O.
\]
三重积中两个矢量相同，故最后一项为零。这是空间力系不能任意消去轴向力偶的原因。
主矩会随取矩点变化，但它沿主矢的投影保持不变；本题的自由力偶本身则完全不依赖取矩点。

% END DETAILED ANSWER tr09
\exercise{三力定理可以用于含一个独立力偶的物体吗？}\label{test:c3-10}

\QuestionStatement{一个刚体承受三个非零、非平行力，另外还承受独立力偶。
说明平衡时三个力的作用线是否必然共点，并给出平面反例。}{tr10}{训练题10题图}
% BEGIN DETAILED ANSWER tr10
\solution \Step{回顾纯三力条件}
以图~\ref{fig:solution-tr10}~中不共点的三个力及抵消其合矩的独立力偶为依据，先写出图右侧第一行的几何或动力学关系，再代入其下方的计算式。
 只有三个力时，任取两条相交作用线的交点取矩，
这两力的矩为零；第三力若要满足矩平衡，其作用线必须经过交点。
\SolutionFigure{tr10}{训练题10的计算示意}
\Step{加入独立力偶} 加上独立力偶后，矩方程变成
\[
\vect r\times\vect F_3+\vect M_c=\vect0,
\]
第三力的矩可以由力偶抵消，因此不再要求它通过交点。
\Step{逐项构造反例} 例如两力 $\vect F_1=F\vect i$、$\vect F_2=F\vect j$
都通过原点；第三力 $\vect F_3=-F(\vect i+\vect j)$
作用于 $(a,0)$。三力的合力为零，而合矩为
\[
(a\vect i)\times[-F(\vect i+\vect j)]=-aF\vect k.
\]
加入力偶 $\vect M_c=aF\vect k$ 后物体平衡，但第三力的作用线不通过原点。
故三力共点定理的适用条件必须先核对，不能把力偶漏记后作几何判断。
\Step{三条作用线的非共点性}
前两力的作用线是两坐标轴，它们唯一交于原点。第三力方向为 $(-1,-1)$，且过 $(a,0)$，其作用线方程为 $y=x-a$；当 $a\ne0$，它不经过原点。
因此这不是仅靠画图猜测的反例：合力逐分量为零，三力矩明确为 $-aF$，独立力偶为 $+aF$，三条线又确实不共点，平衡条件和不共点结论同时得到验证。

% END DETAILED ANSWER tr10
\subsection{惯量、定轴与平面运动}\label{sub:c3-11.2}
\exercise{细杆斜轴的方向惯量}\label{test:c3-11}

\QuestionStatement{均匀细杆质量 $m$、长度 $L$，通过中点的转轴与杆夹角为 $\beta$。
直接由积分求惯量，再用惯量张量验证。}{tr11}{训练题11题图}
% BEGIN DETAILED ANSWER tr11
\solution \Step{用垂距建立积分}
以图~\ref{fig:solution-tr11}~中过质心的转轴与杆的夹角 $\beta$为依据，先写出图右侧第一行的几何或动力学关系，再代入其下方的计算式。
 用沿杆坐标 $s\in[-L/2,L/2]$，
线密度 $\lambda=m/L$。坐标为 $s$ 的质元到转轴的垂距为 $|s\sin\beta|$，
因此
\[
I_\beta=\int_{-L/2}^{L/2}s^2\sin^2\beta\,\lambda\,\dd s
=\frac{m}{L}\sin^2\beta\left[\frac{s^3}{3}\right]_{-L/2}^{L/2}
=\frac{mL^2}{12}\sin^2\beta.
\]
\SolutionFigure{tr11}{训练题11的计算示意}
\Step{建立质心张量} 若令杆沿 $z$ 轴，则其中心惯量张量为
\[
\mathsf I_C=\operatorname{diag}\left(\frac{mL^2}{12},
\frac{mL^2}{12},0\right).
\]
\Step{二次型与极限} 选轴向单位矢量
$\vect n=(\sin\beta,0,\cos\beta)^{\mathsf T}$，
有 $I_\beta=\vect n^{\mathsf T}\mathsf I_C\vect n$，
与积分相同。$\beta=0$ 时惯量为零是理想细杆横截面被忽略的结果；
实际有截面的杆沿纵轴仍有小而非零的惯量。
$\beta=90^\circ$ 时得到通常的中点横轴惯量，且
$I_\beta$ 的量纲始终为质量乘长度平方。
\Step{平行偏置轴与端点斜轴}
本题轴通过中点。若另取平行的端点轴，两轴最短距离为 $(L/2)|\sin\beta|$，因此
\[
I_{\text{端}}=I_\beta+m\left(\frac L2\sin\beta\right)^2
=\frac{mL^2}{3}\sin^2\beta.
\]
这与从 $s=0$ 至 $L$ 的直接积分一致，也解释了常见的 $1/12$ 与 $1/3$ 系数分别对应什么原点。
不能把端点到质心的距离 $L/2$ 全部作为轴距，除非转轴与杆垂直。

% END DETAILED ANSWER tr11
\exercise{球的平行轴惯量}\label{test:c3-12}

\QuestionStatement{均匀实心球半径 $a$、质量 $m$。一条转轴与过球心的某直径平行，
两轴的垂距为 $d$。从平行轴定理的积分证明求其惯量。}{tr12}{训练题12题图}
% BEGIN DETAILED ANSWER tr12
\solution \Step{展开偏移轴的垂距}
以图~\ref{fig:solution-tr12}~中平行球心轴和偏置轴之间的垂距 $d$为依据，先写出图右侧第一行的几何或动力学关系，再代入其下方的计算式。
 取过球心的轴为 $z$ 轴，外轴位于 $x=d,y=0$。
质元到外轴的距离平方为 $(x-d)^2+y^2$，因此
\[
\begin{aligned}
I&=\int[(x-d)^2+y^2]\,\dd m\\
&=\int(x^2+y^2)\,\dd m-2d\int x\,\dd m+d^2\int\dd m.
\end{aligned}
\]
\SolutionFigure{tr12}{训练题12的计算示意}
\Step{计算中心轴积分} 球心就是质心，$\int x\,\dd m=0$，故
$I=I_C+md^2$。均匀实心球的中心轴惯量为
\[
I_C=\rho\int_0^a r^4\,\dd r
\int_0^\pi\sin^3\theta\,\dd\theta
\int_0^{2\pi}\dd\varphi
=\frac{8\pi\rho a^5}{15}=\frac25ma^2,
\]
其中 $\rho=3m/(4\pi a^3)$。所以
\[
I=\frac25ma^2+md^2.
\]
\Step{平行轴结果与特例} 当 $d=a$ 时为切线轴，$I=7ma^2/5$。
这里的 $d$ 是两条平行轴的最短距离；
若给出的点间距离包含轴向分量，必须先取其垂直于轴的部分。
\Step{质心条件与动能交叉项}
一般物体若选原点不是质心，$\int x\dd m$ 并不一定为零，展开后不能直接省去 $-2d\int x\dd m$。
本题球心与质心重合，才可得到 $I=I_C+md^2$。
$\partial I/\partial d=2md$ 表明在平行轴族中中心轴惯量最小；$d=0$ 恢复 $2ma^2/5$，$d=a$ 恢复切线轴的 $7ma^2/5$，两种极限同时核验了平移方向和符号。

% END DETAILED ANSWER tr12
\exercise{均匀实心半球从球心移到质心，哪个惯量不变？}\label{test:c3-13}

\QuestionStatement{半球占据 $x^2+y^2+z^2\le a^2,\ z\ge0$，质量 $m$。
求其质心和三条质心主轴的惯量。}{tr13}{训练题13题图}
% BEGIN DETAILED ANSWER tr13
\solution \Step{计算质心一阶矩}
以图~\ref{fig:solution-tr13}~中半球质心沿对称轴的移动为依据，先写出图右侧第一行的几何或动力学关系，再代入其下方的计算式。
 密度为 $\rho=3m/(2\pi a^3)$。
高度 $z$ 的圆截面面积为 $\pi(a^2-z^2)$，故
\[
z_C=\frac{\rho\pi}{m}\int_0^a z(a^2-z^2)\,\dd z
=\frac{\rho\pi a^4}{4m}=\frac{3a}{8}.
\]
\SolutionFigure{tr13}{训练题13的计算示意}
\Step{计算球心惯量} 由对称性 $x_C=y_C=0$。过球心 $O$ 的三个正交轴惯量都为
$2ma^2/5$：相同密度整球的积分在 $z\ge0$ 与 $z\le0$ 各占一半，
而半球质量也为整球的一半。对 $z$ 轴也可用截面圆盘验证：
\[
I_{Oz}=\frac{\rho\pi}{2}\int_0^a(a^2-z^2)^2\,\dd z
=\frac{4\pi\rho a^5}{15}=\frac25ma^2.
\]
\Step{移点至质心} 将坐标原点沿 $z$ 移到 $C$，平行轴定理给出
\[
I_{Cx}=I_{Ox}-mz_C^2
=\left(\frac25-\frac9{64}\right)ma^2=\frac{83}{320}ma^2,
\]
同理 $I_{Cy}=83ma^2/320$。中心对称轴与原 $z$ 轴是同一条直线，
不存在两轴间的垂距，所以
\[
I_{Cz}=\frac25ma^2.
\]
因此移点只改变横向轴惯量，沿对称轴的惯量保持不变。
\Step{张量平移中的交叉元}
令 $\vect d=\overrightarrow{OC}=(0,0,3a/8)$，则
\[
\mathsf I_C=\mathsf I_O-m(d^2\mathsf1-\vect d\vect d^{\mathsf T}).
\]
被减去的矩阵为 $m(3a/8)^2\operatorname{diag}(1,1,0)$，所以仅两个横向对角元减少，交叉元仍为零。
半球没有关于 $z=0$ 的质量对称性，故不能以“球的一半”为理由断言球心就是质心；应先计算一阶矩，再讨论惯量的平移。

% END DETAILED ANSWER tr13
\exercise{重根主惯量的主轴是否唯一？}\label{test:c3-14}

\QuestionStatement{设质心惯量张量在某组正交坐标中为
$\operatorname{diag}(A,A,C)$，且 $A\ne C$。
求出其主轴，并说明在重根平面中使用斜交方向会发生什么。}{tr14}{训练题14题图}
% BEGIN DETAILED ANSWER tr14
\solution \Step{求特征子空间}
以图~\ref{fig:solution-tr14}~中重根平面内可以旋转的正交方向为依据，先写出图右侧第一行的几何或动力学关系，再代入其下方的计算式。
 特征方程为
\[
(A-\lambda)^2(C-\lambda)=0.
\]
对 $\lambda=C$，特征矢量沿 $z$ 轴；
对 $\lambda=A$，任意矢量 $(u,v,0)$ 都是特征矢量。
\SolutionFigure{tr14}{训练题14的计算示意}
\Step{选正交主轴组} 可以任选角 $\chi$，建立单位正交基
\[
\vect e'_1=(\cos\chi,\sin\chi,0),\qquad
\vect e'_2=(-\sin\chi,\cos\chi,0),\qquad \vect e'_3=\vect e_z.
\]
这组三轴均为主轴，且转动动能始终为
\[
T=\frac12A(\omega_1'^2+\omega_2'^2)+\frac12C\omega_3'^2.
\]
\Step{斜交基底中的展开} 如果在特征平面内选两条夹角为 $\eta\ne90^\circ$ 的单位方向
$\vect u,\vect v$，虽然两者各自都是特征矢量，
但令 $\vect\omega=p\vect u+q\vect v$ 时，
\[
T=\frac A2(p^2+q^2+2pq\cos\eta)
\]
出现交叉项。原因是 $p,q$ 为斜交基展开系数，不能当作正交投影分量。
重根给出的是主轴选择自由，而不是取消正交坐标的要求。
若进一步 $A=C$，任意正交三轴组都可以作为主轴组。
\Step{斜交基中的交叉项与惯量对称性}
若 $\vect\omega=p\vect u+q\vect v$，两方向夹角为 $\eta$，则
$|\vect\omega|^2=p^2+q^2+2pq\cos\eta$。
在重根横向平面内，$T=A|\vect\omega|^2/2$，所以交叉项来自基底度量，和选择之前的张量对称性并无矛盾。
选取互相垂直的两根特征方向后，$\cos\eta=0$，才恢复无交叉项的主轴动能形式。

% END DETAILED ANSWER tr14
\exercise{均匀圆环与圆盘滚下同一斜面}\label{test:c3-15}

\QuestionStatement{同半径 $a$ 的均匀圆环和圆盘，从静止开始沿倾角 $\beta$ 的固定斜面纯滚动。
求加速度、所需静摩擦和保证纯滚动的最小静摩擦系数。}{tr15}{训练题15题图}
% BEGIN DETAILED ANSWER tr15
\solution \Step{列三个动力学条件}
以图~\ref{fig:solution-tr15}~中两个圆轮与同一斜面的接触为依据，先写出图右侧第一行的几何或动力学关系，再代入其下方的计算式。
 沿坡向下为正，设静摩擦大小 $f$ 沿坡向上。
令 $I_C=kma^2$，由平动、转动和无滑约束，
\[
m\ddot x=mg\sin\beta-f,\qquad
I_C|\alpha|=fa,\qquad \ddot x=a|\alpha|.
\]
\SolutionFigure{tr15}{训练题15的计算示意}
\Step{消元求一般结果} 后两式给出 $f=km\ddot x$，代入第一式，
\[
\ddot x=\frac{g\sin\beta}{1+k},\qquad
f=\frac{k}{1+k}mg\sin\beta,\qquad N=mg\cos\beta.
\]
\Step{分别代入两种惯量} 圆环 $k=1$，因此
\[
\ddot x_{\text{环}}=\frac12g\sin\beta,\qquad
f_{\text{环}}=\frac12mg\sin\beta,\qquad
\mu_s\ge\frac12\tan\beta.
\]
圆盘 $k=1/2$，因此
\[
\ddot x_{\text{盘}}=\frac23g\sin\beta,\qquad
f_{\text{盘}}=\frac13mg\sin\beta,\qquad
\mu_s\ge\frac13\tan\beta.
\]
两者都满足各自摩擦条件时，圆盘较早到达同一坡下位置。
也可由 $mgx\sin\beta=\frac12m(1+k)\dot x^2$ 求导得到相同加速度。
若斜面摩擦只够圆盘而不够圆环，上述圆环纯滚动解就不能用来比较实际运动。
\Step{用同一坡长比较到达时间}
从静止沿坡下移距离 $s$，时间为 $t=\sqrt{2s/\ddot x}$。
两物体均满足各自的纯滚动条件时，
\[
\frac{t_{\text{环}}}{t_{\text{盘}}}
=\sqrt{\frac{\ddot x_{\text{盘}}}{\ddot x_{\text{环}}}}=\sqrt{\frac43}>1.
\]
结果不依赖质量，来自不同的无量纲惯量 $I_C/(ma^2)$。
若 $\mu_s$ 位于 $\tan\beta/3$ 与 $\tan\beta/2$ 之间，圆盘可以纯滚，圆环则必须另用滑动摩擦方程，不能继续使用以上时间比。

% END DETAILED ANSWER tr15
\exercise{物理摆悬点过质心时会怎样？}\label{test:c3-16}

\QuestionStatement{刚体绕水平固定轴运动，质心到轴的距离为 $h$，轴惯量为 $I_O$。
从动力学方程解释 $h=0$ 与 $h\to0^+$ 的区别。}{tr16}{训练题16题图}
% BEGIN DETAILED ANSWER tr16
\solution \Step{建立精确动力学}
以图~\ref{fig:solution-tr16}~中悬点至质心的距离 $h$为依据，先写出图右侧第一行的几何或动力学关系，再代入其下方的计算式。
 以质心位于悬点下方为 $\theta=0$，
重力的有向力矩为 $-mgh\sin\theta$，故
\[
I_O\ddot\theta+mgh\sin\theta=0.
\]
\SolutionFigure{tr16}{训练题16的计算示意}
\Step{取小振幅极限} 对 $h>0$ 的小振幅运动，$\sin\theta\simeq\theta$，得到
\[
\omega_{\text{摆}}^2=\frac{mgh}{I_O},\qquad
T_{\text{摆}}=2\pi\sqrt{\frac{I_O}{mgh}}.
\]
当悬点向质心趋近，$I_O=I_C+mh^2$。
若 $I_C>0$，则 $h\to0^+$ 时回复系数趋零，周期发散。
\Step{严格令轴过质心} 在严格的 $h=0$ 条件下，方程直接变为
\[
I_C\ddot\theta=0,\qquad
\theta(t)=\theta(0)+\dot\theta(0)t.
\]
初始静止时可以停在任意姿态；初始有角速度时保持匀速转动。
这里不存在围绕唯一平衡姿态的小摆动，所以不应将无限周期解释成某种有限振荡。
如果附加扭簧或其他恢复力矩，则须把该项计入，原结论也随之改变。
\Step{从能量形状解释两个极限}
约去势能常数后 $U=mgh(1-\cos\theta)$。$h>0$ 时最低点附近有
$U\simeq mgh\theta^2/2$，曲率决定小振动回复强度。
严格 $h=0$ 时整个势能与姿态无关，所有角度同为平衡位置，不能再选出唯一的振动中心。
当 $h$ 很小而未为零，仍存在很弱的回复力矩和很长的周期；把“周期趋无穷”与“没有振荡”区分开，须回到未除以 $h$ 的原始动力学方程。

% END DETAILED ANSWER tr16
\exercise{静平衡是否保证动平衡？}\label{test:c3-17}

\QuestionStatement{转子绕固定 $z$ 轴以常角速度 $\Omega$ 转动。
设质心在轴上，但该轴不是中心主轴。
由角动量的空间导数证明仍需附加横向力偶。}{tr17}{训练题17题图}
% BEGIN DETAILED ANSWER tr17
\solution \Step{区分偏心与惯性积}
以图~\ref{fig:solution-tr17}~中轴上质心与倾斜杆造成的轴承力偶为依据，先写出图右侧第一行的几何或动力学关系，再代入其下方的计算式。
 质心在轴上意味着横向偏心距为零，
因此没有由质心绕轴公转引起的离心合力。
不过在随体坐标中，
\[
\vect L_C=\Omega(I_{xz}\vect e_x+I_{yz}\vect e_y+I_{zz}\vect e_z).
\]
\SolutionFigure{tr17}{训练题17的计算示意}
\Step{对角动量求空间导数} 各分量虽为常数，横向基矢却以 $\Omega$ 转动。
由运输公式，
\[
\vect M_C=\vect\omega\times\vect L_C
=-I_{yz}\Omega^2\vect e_x+I_{xz}\Omega^2\vect e_y.
\]
只要 $I_{xz}$ 或 $I_{yz}$ 非零，就需轴承提供横向力偶。
\Step{构造静平衡反例} 例如取均匀细杆，其质心在轴上，杆方向
$\vect n=(\sin\beta,0,\cos\beta)$。
它的中心张量是
\[
\mathsf I_C=\frac{mL^2}{12}(\mathsf1-\vect n\vect n^{\mathsf T}),
\qquad I_{xz}=-\frac{mL^2}{12}\sin\beta\cos\beta.
\]
当 $0<\beta<90^\circ$，$I_{xz}\ne0$，这是静平衡但不满足动平衡的明确例子。
要消除随转动变化的附加轴承载荷，
还须令固定轴成为中心主轴，即同时满足 $I_{xz}=I_{yz}=0$。
\Step{把所需力偶换成可测量的轴承力}
若两轴承间距为 $d$，它们的附加横向力合为零。设 $\vect R_B=X\vect e_x+Y\vect e_y$、$\vect R_A=-\vect R_B$，则
\[
d\vect e_z\times\vect R_B=d(X\vect e_y-Y\vect e_x)=\vect M_C.
\]
逐分量比较得到 $X=I_{xz}\Omega^2/d$、$Y=I_{yz}\Omega^2/d$。
因此轴承附加力幅度为 $\Omega^2\sqrt{I_{xz}^2+I_{yz}^2}/d$，即使质心不公转，也可产生随转速平方增大的载荷。

% END DETAILED ANSWER tr17
\exercise{轮子匀速纯滚动时底部点加速度}\label{test:c3-18}

\QuestionStatement{半径 $R$ 的轮向右以质心速率 $v$ 纯滚动。
取当前接触地面的材料点 $P$，求其速度和加速度，
并说明为何不能对瞬心速度零直接求导。}{tr18}{训练题18题图}
% BEGIN DETAILED ANSWER tr18
\solution \Step{求接触材料点速度}
以图~\ref{fig:solution-tr18}~中当前底部材料点 $P$ 与轮心 $C$为依据，先写出图右侧第一行的几何或动力学关系，再代入其下方的计算式。
 取 $\vect i$ 向右、$\vect j$ 向上，
有 $\vect v_C=v\vect i$，
$\vect\omega=-(v/R)\vect k$、$\vect r_{CP}=-R\vect j$。
因而
\[
\vect v_P=v\vect i+
\left(-\frac vR\vect k\right)\times(-R\vect j)=\vect0.
\]
\SolutionFigure{tr18}{训练题18的计算示意}
\Step{求同一材料点加速度} 匀速运动中 $\vect a_C=\vect0$、$\vect\alpha=\vect0$，所以
\[
\begin{aligned}
\vect a_P
&=\vect\omega\times(\vect\omega\times\vect r_{CP})\\
&=-\frac{v^2}{R^2}\vect r_{CP}
=\frac{v^2}{R}\vect j.
\end{aligned}
\]
\Step{说明不能直接求导的原因} 这里 $P$ 是本瞬时碰巧位于底部的同一个材料点，它随后将离开接触位置。
下一时刻的底部材料点已是另一个点。
“接触点速度恒为零”是在不断更换材料点的意义下成立，
不能由此推出某一材料点的加速度也为零。
结果还可由摆线的二阶导数核查：最低点的切向速度为零，但曲线具有非零向上加速度。
\Step{材料点的摆线轨迹}
选 $t=0$ 时恰在底部的材料点，轮心向右匀速，令 $\nu=v/R$，该点坐标可写为
\[
x_P=vt-R\sin\nu t,\qquad y_P=R-R\cos\nu t.
\]
在 $t=0$，两一阶导数都为零；但
$\ddot x_P(0)=0$、$\ddot y_P(0)=R\nu^2=v^2/R$。
这是对同一材料点的轨迹连续求导，严格恢复了刚体加速度公式，也说明零速度本身并不能推出零加速度。

% END DETAILED ANSWER tr18
\exercise{外滚行星轮半径相等时转速比}\label{test:c3-19}

\QuestionStatement{固定轮半径 $R$，外滚行星轮半径 $r$，
转臂绕固定轮心以角速度 $\Omega>0$ 逆时针转动。
求行星轮的绝对自转率，以及相对于转臂的自转率。}{tr19}{训练题19题图}
% BEGIN DETAILED ANSWER tr19
\solution \Step{求轮心速度}
以图~\ref{fig:solution-tr19}~中外滚轮接触点相对轮心的向内半径为依据，先写出图右侧第一行的几何或动力学关系，再代入其下方的计算式。
 行星轮心 $C$ 到固定轮心的距离为 $R+r$，
故其切向速度为
\[
\vect v_C=(R+r)\Omega\vect e_\theta.
\]
接触点相对轮心的位置为 $-r\vect e_r$。
\SolutionFigure{tr19}{训练题19的计算示意}
\Step{用接触点无滑动求绝对率} 设行星轮的绝对角速度为 $\omega\vect k$，
固定轮上的接触点速度为零；无滑条件给出
\[
\vect0=(R+r)\Omega\vect e_\theta+
\omega\vect k\times(-r\vect e_r)
=[(R+r)\Omega-r\omega]\vect e_\theta.
\]
因此
\[
\omega=\frac{R+r}{r}\Omega.
\]
\Step{减去转臂参考架} 转臂坐标也以 $\Omega$ 转动，所以相对自转率为
\[
\omega_{\text{相对}}=\omega-\Omega=\frac Rr\Omega.
\]
特别是 $R=r$ 时，绝对率为 $2\Omega$，相对率为 $\Omega$。
转臂完成一周，行星轮在空间转两周，而相对于转臂只转一周。
应先注明观察参考系，再使用“自转圈数”，否则容易漏掉转臂本身的一周。
\Step{一周运动的转角积分}
设转臂角 $\varphi$ 从0增至 $2\pi$，绝对自转角增量为
$\Delta\theta=(1+R/r)2\pi$，相对转臂转角为 $\Delta\theta-2\pi=2\pi R/r$。
等半径时分别为 $4\pi$ 与 $2\pi$。
这一结论由接触点速度约束积分得到，比较的是同一标记半径在固定空间和转臂架中的转角，不需要把两种“转一周”的含义混在一起。

% END DETAILED ANSWER tr19
\exercise{光滑支承上的球板系统哪个量守恒？}\label{test:c3-20}

\QuestionStatement{均匀实心球质量 $m$、半径 $R$，初速 $v>0$，初始角速度为零。
质量 $M$ 的长水平板初始静止，板的支承光滑；
球板间有摩擦且始终接触，板足够长。
求进入无相对滑动状态时的速度和机械能损失。}{tr20}{训练题20题图}
% BEGIN DETAILED ANSWER tr20
\solution \Step{选择动量与冲量方法}
以图~\ref{fig:solution-tr20}~中球相对板的初始滑动与摩擦冲量为依据，先写出图右侧第一行的几何或动力学关系，再代入其下方的计算式。
 水平外力为零，系统水平动量守恒；滑动摩擦会耗散机械能。
记摩擦冲量大小为 $J$，球向右速度为 $v_b$，板速度为 $V$，
并把球的顺时针转速大小记为 $\omega$。
球、板和球自转的冲量方程分别为
\[
mv_b=mv-J,\qquad MV=J,\qquad
\frac25mR^2\omega=JR.
\]
\SolutionFigure{tr20}{训练题20的计算示意}
\Step{用接触速度求过渡条件} 所以
\[
v_b=v-\frac Jm,\qquad V=\frac JM,\qquad
R\omega=\frac{5J}{2m}.
\]
球底部相对板的滑动速度为
\[
u=v_b-R\omega-V
=v-J\left(\frac7{2m}+\frac1M\right).
\]
首次达到 $u=0$ 时，
\[
J_*=\frac{2mMv}{7M+2m},
\]
进而
\[
v_b=\frac{5M+2m}{7M+2m}v,\qquad
V=\frac{2m}{7M+2m}v,\qquad
R\omega=\frac{5M}{7M+2m}v.
\]
代回可验证 $mv_b+MV=mv$ 和 $v_b-R\omega=V$。
\Step{计算完整末动能} 初能为 $E_i=mv^2/2$，末能必须包含三项：
\[
E_f=\frac12mv_b^2+\frac12MV^2+
\frac12\left(\frac25mR^2\right)\omega^2.
\]
化简得到
\[
E_i-E_f=\frac{mMv^2}{7M+2m}.
\]
当 $M\to\infty$ 时，板的速度趋零，球的 $v_b\to5v/7$，
损失趋于 $mv^2/7$，恢复固定平板的结果。
达成无相对滑动后若无其他作用，三个速度均可保持恒定，此时所需静摩擦为零。
\Step{用摩擦总功率再求能量损失}
若滑动摩擦大小为常量 $f$，则 $J=ft$；接触相对速度从 $v$ 线性降至零，过渡时间 $t_*=J_*/f$。
球板系统摩擦总功率为 $-f|u|$，因而
\[
\Delta E=f\int_0^{t_*}|u|\dd t=f\frac{vt_*}{2}=\frac{vJ_*}{2}
=\frac{mMv^2}{7M+2m}.
\]
它与三项末动能的计算相同；摩擦系数控制过渡所需时间，却不改变在板足够长、能达到无滑状态时的末速度与总损失。

% END DETAILED ANSWER tr20
\subsection{自由转动、重陀螺与磁进动}\label{sub:c3-11.3}
\exercise{无外力矩是否保证角速度在空间恒定？}\label{test:c3-21}

\QuestionStatement{刚体绕固定点自由转动，或者以质心为角动量参考点且无外力矩。
说明空间角动量与角速度各满足什么关系，并给出角速度不恒定的例子。}{tr21}{训练题21题图}
% BEGIN DETAILED ANSWER tr21
\solution \Step{明确角动量守恒}
以图~\ref{fig:solution-tr21}~中非主轴角速度与固定角动量的关系为依据，先写出图右侧第一行的几何或动力学关系，再代入其下方的计算式。
 在惯性系中，
\[
\frac{\dd\vect L}{\dd t}=\vect0,\qquad \vect L=\text{空间常矢量}.
\]
\SolutionFigure{tr21}{训练题21的计算示意}
\Step{空间惯量张量的时间变化} 但是 $\vect L=\mathsf I(t)\vect\omega$，非球形刚体的空间惯量张量
$\mathsf I(t)$ 随姿态变化，故不能由 $\vect L$ 恒定直接推出 $\vect\omega$ 恒定。
\Step{构造轴对称反例} 对主惯量 $(A,A,C)$ 的轴对称刚体，设 $\omega_3=s$，
可将角动量写为
\[
\vect L=A\vect\omega+(C-A)s\vect e_3,
\qquad
\vect\omega=\frac{\vect L}{A}-\frac{C-A}{A}s\vect e_3.
\]
无矩时 $s$ 恒定，但倾斜的本体轴 $\vect e_3$ 围绕 $\vect L$ 进动。
因此当 $A\ne C$、$s\ne0$、$\vect e_3$ 不平行于 $\vect L$ 时，
$\vect\omega$ 的空间方向随时间改变。
特殊情形包括惯量各向同性的刚体，此时
$\vect L=I\vect\omega$；也包括沿一根主轴的精确纯转动。
这些特殊解的角速度可在空间恒定，但不能替代一般结论。
\Step{角速度的时间导数}
对轴对称反例的角速度表达式求空间导数，因 $\vect L,s,A,C$ 均恒定，
\[
\dot{\vect\omega}=-\frac{C-A}{A}s\dot{\vect e}_3
=-\frac{C-A}{A^2}s(\vect L\times\vect e_3).
\]
只要 $C\ne A$、$s\ne0$、轴与角动量不平行，这一导数非零，构成明确的反例。
角动量守恒来自零外矩；角速度还受随姿态变化的质量分布约束，二者不能用同一个“守恒”判断替代。

% END DETAILED ANSWER tr21
\exercise{三种主轴中哪一根自由转动不稳定？}\label{test:c3-22}

\QuestionStatement{取 $I_1<I_2<I_3$。围绕中间惯量轴作纯转动
$\omega_2=\Omega\ne0$，写出 $\omega_1,\omega_3$ 的线性扰动方程，
给出带初始扰动的解，并与另两根轴比较。}{tr22}{训练题22题图}
% BEGIN DETAILED ANSWER tr22
\solution \Step{列无矩欧拉方程}
以图~\ref{fig:solution-tr22}~中中间惯量轴上的主转动与横向扰动为依据，先写出图右侧第一行的几何或动力学关系，再代入其下方的计算式。
 无矩欧拉方程为
\[
I_1\dot\omega_1+(I_3-I_2)\omega_2\omega_3=0,\qquad
I_3\dot\omega_3+(I_2-I_1)\omega_1\omega_2=0.
\]
\SolutionFigure{tr22}{训练题22的计算示意}
\Step{保留一阶扰动} 把 $\omega_2=\Omega+O(\delta)$，
$\omega_1=x,\omega_3=z$ 代入并保留一阶项，得
\[
\dot x=-az,\qquad \dot z=-bx,\qquad
a=\frac{(I_3-I_2)\Omega}{I_1},\quad
b=\frac{(I_2-I_1)\Omega}{I_3}.
\]
两系数同号，且
\[
\ddot x=abx,\qquad
\lambda=\sqrt{ab}
=|\Omega|\sqrt{\frac{(I_3-I_2)(I_2-I_1)}{I_1I_3}}>0.
\]
\Step{应用扰动初值} 若初值为 $x_0,z_0$，则
\[
\begin{aligned}
x(t)&=x_0\cosh\lambda t-\frac a\lambda z_0\sinh\lambda t,\\
z(t)&=z_0\cosh\lambda t-\frac b\lambda x_0\sinh\lambda t.
\end{aligned}
\]
一般初值含 $e^{\lambda t}$ 项，扰动增长，故中间轴转动不稳定。
沿最小惯量轴，二阶方程的系数为
$-\Omega^2(I_2-I_1)(I_3-I_1)/(I_2I_3)<0$；
沿最大惯量轴则为
$-\Omega^2(I_3-I_1)(I_3-I_2)/(I_1I_2)<0$。
相应横向扰动为振荡，在线性意义下稳定。
惯量重根时上述增长率可能为零，必须另行考察，不能沿用三根互异的判断。
\Step{初始条件与线性解的有效时段}
在 $t=0$，$\cosh0=1,\sinh0=0$，两解回到 $x_0,z_0$；求导后分别给出 $\dot x(0)=-az_0$、$\dot z(0)=-bx_0$，与一阶方程一致。
即使某一特殊初值恰落在衰减方向上，任意微小偏离仍会激发增长模态，所以不能用这一特殊轨迹证明稳定。
当 $|x|,|z|$ 增长到与 $|\Omega|$ 同量级时，舍去的二阶扰动不再小，线性公式必须停止外推，但不稳定性的判断已经成立。

% END DETAILED ANSWER tr22
\exercise{对称自由陀螺的空间进动率}\label{test:c3-23}

\QuestionStatement{无矩轴对称刚体主惯量为 $(A,A,C)$，
轴向角速度 $s$。以空间 $\zeta$ 轴沿恒定角动量 $\vect L$，
求 $\dot\varphi,\dot\theta,\dot\psi$，并解释空间进动与本体自转的区别。}{tr23}{训练题23题图}
% BEGIN DETAILED ANSWER tr23
\solution \Step{求恒定倾角}
以图~\ref{fig:solution-tr23}~中本体轴绕固定角动量的圆锥运动为依据，先写出图右侧第一行的几何或动力学关系，再代入其下方的计算式。
 欧拉第三式给出 $C\dot s=0$，故 $s$ 恒定。
轴向角动量为 $Cs$，于是
\[
L\cos\theta=Cs
\]
说明 $\theta$ 恒定。\SolutionFigure{tr23}{训练题23的计算示意}
\Step{求空间进动率} 利用前题的矢量式，
\[
\dot{\vect e}_3=\vect\omega\times\vect e_3
=\frac{\vect L}{A}\times\vect e_3.
\]
因此非退化倾斜情形下
\[
\dot\varphi=\frac LA,\qquad \dot\theta=0.
\]
\Step{求欧拉角自转率} 又因 $s=\dot\psi+\dot\varphi\cos\theta$，
\[
\dot\psi=s-\frac LA\frac{Cs}{L}
=-\frac{C-A}{A}s.
\]
所以空间进动率由角动量大小和横向惯量决定，
欧拉自转角率则随 $C-A$ 的符号改变。
如果 $C>A$、$s>0$，$\dot\psi<0$，并不意味着实际轴向角速度 $s$ 也为负；
两者包含的进动投影不同。
在 $\theta=0,\pi$ 时节点轴退化，只能由总角率描述轴向转动，
不宜分别把 $\dot\varphi,\dot\psi$ 当作可独立观测的角率。
\Step{补出本体横向分量并联系两种频率}
若初始 $\omega_1=w,\omega_2=0$，欧拉方程给出
\[
\omega_1=w\cos nt,\quad\omega_2=w\sin nt,\qquad
n=\frac{C-A}{A}s=-\dot\psi.
\]
这里 $n$ 是本体系中的横向分量旋转率，$L/A$ 是空间中轴绕角动量的进动率；两者并非同一参考系中的同一个角速度。
当 $A=C$，$n=0$，球形惯量使空间角速度固定。此时若仍任意标记一根体轴，它可随刚体绕固定角速度转动，不能据 $n=0$ 推出所有体轴空间静止。

% END DETAILED ANSWER tr23
\exercise{重陀螺有哪些守恒量？}\label{test:c3-24}

\QuestionStatement{轴对称刚体由理想支点固定于 $O$，质心为 $l\vect e_3$，
置于均匀恒定重力场中，不受其他外力矩。空间 $\zeta$ 轴竖直向上。
设横向惯量 $A$、轴向惯量 $C$，
用欧拉角写出三项独立积分及其物理来源。}{tr24}{训练题24题图}
% BEGIN DETAILED ANSWER tr24
\solution \Step{写出动能与重力势能}
以图~\ref{fig:solution-tr24}~中竖直方向和本体对称轴上的两个角动量投影为依据，先写出图右侧第一行的几何或动力学关系，再代入其下方的计算式。
 动能和势能分别为
\[
T=\frac A2(\dot\theta^2+\dot\varphi^2\sin^2\theta)
+\frac C2(\dot\psi+\dot\varphi\cos\theta)^2,\qquad
U=mgl\cos\theta.
\]
\SolutionFigure{tr24}{训练题24的计算示意}
\Step{由循环坐标求守恒量} 拉格朗日函数 $\mathcal L=T-U$ 不含 $\psi,\varphi$ 本身，
因此两共轭动量守恒：
\[
P=\frac{\partial\mathcal L}{\partial\dot\psi}
=C(\dot\psi+\dot\varphi\cos\theta)=C\omega_3,
\]
\[
K=\frac{\partial\mathcal L}{\partial\dot\varphi}
=A\dot\varphi\sin^2\theta+P\cos\theta=L_\zeta.
\]
第一项来自轴对称性以及重力矩沿本体轴的分量为零；
第二项来自重力矩沿空间竖直方向的分量为零。
\Step{约化为一维章动} 理想固定支点不做功，且势能无显式时间依赖，所以
\[
E=T+U
\]
也守恒。消去 $\dot\psi,\dot\varphi$ 后，
\[
\frac A2\dot\theta^2+
\frac{(K-P\cos\theta)^2}{2A\sin^2\theta}
+mgl\cos\theta+\frac{P^2}{2C}=E.
\]
但一般情况下
$\dot{\vect L}=\vect r_{OC}\times(-mg\vect e_\zeta)\ne0$，
整个角动量矢量并不守恒。两个固定投影守恒不能替代矢量守恒。
\Step{从守恒量重构角变量的求解次序}
先用约化能量求 $\theta(t)$ 的一维运动；然后依次代入
\[
\dot\varphi=\frac{K-P\cos\theta}{A\sin^2\theta},\qquad
\dot\psi=\frac PC-\dot\varphi\cos\theta.
\]
两次积分还须用初始方位和自转角确定积分常数。
允许的章动区间由约化动能非负确定。三个守恒量减少求解难度，但不足以凭空给出三个角的初始相位，完整运动仍依赖完整初值。

% END DETAILED ANSWER tr24
\exercise{高速重陀螺的慢进动率}\label{test:c3-25}

\QuestionStatement{保持固定倾角 $0<\theta_0<\pi$ 的重陀螺，
轴向角动量为 $P\ne0$。求常进动率的两个分支，
给出慢分支近似的条件，并讨论初始没有进动的倾斜释放。}{tr25}{训练题25题图}
% BEGIN DETAILED ANSWER tr25
\solution \Step{列定常进动条件}
以图~\ref{fig:solution-tr25}~中重陀螺的倾角与进动方向为依据，先写出图右侧第一行的几何或动力学关系，再代入其下方的计算式。
 常进动要求 $\dot\theta=\ddot\theta=0$。
倾角方程化为
\[
A\Omega^2\cos\theta_0-P\Omega+mgl=0.
\]
\SolutionFigure{tr25}{训练题25的计算示意}
\Step{求两支精确根} 当 $\cos\theta_0\ne0$ 且
$D=P^2-4Amgl\cos\theta_0\ge0$ 时，两根为
\[
\Omega_\pm=\frac{P\pm\sqrt D}{2A\cos\theta_0}.
\]
\Step{选择并展开慢分支} 为处理 $P$ 的正负，把固定倾角下随 $|P|$ 增大而趋于零的进动分支写成
\[
\Omega_{\text{慢}}
=\frac{P-\operatorname{sgn}(P)\sqrt D}{2A\cos\theta_0}
=\frac{2mgl}{P+\operatorname{sgn}(P)\sqrt D}.
\]
若
\[
\left|\frac{4Amgl\cos\theta_0}{P^2}\right|\ll1,
\]
则
$\sqrt D=|P|[1-2Amgl\cos\theta_0/P^2+O(P^{-4})]$，
从而
\[
\Omega_{\text{慢}}\simeq\frac{mgl}{P}.
\]
当 $\cos\theta_0=0$，二次项消失，精确关系就是 $\Omega=mgl/P$。
这个近似针对满足常进动条件的慢分支。
若初始 $\dot\varphi=0$、$\dot\theta=0$ 且轴倾斜，
倾角方程起初给出
$\ddot\theta=(mgl/A)\sin\theta_0\ne0$，
物体立即开始章动。因此不能把此类初值直接改成一条常进动轨迹。
\Step{保留下一阶项评估慢进动近似}
把慢分支写成 $\Omega=mgl/P+\delta$ 并代回定常方程，保留最低阶修正，有
\[
\delta\simeq\frac{A(mgl)^2\cos\theta_0}{P^3},\qquad
\Omega_{\text{慢}}\simeq\frac{mgl}{P}
\left[1+\frac{Amgl\cos\theta_0}{P^2}\right].
\]
当 $\cos\theta_0>0$，修正与 $P$ 同号；当 $\cos\theta_0<0$，修正相反。
这一展开仍要求相应小参数远小于1，不能用于两个根接近合并、判别式接近零的位置。

% END DETAILED ANSWER tr25
\exercise{重陀螺回环轨迹的判别}\label{test:c3-26}

\QuestionStatement{已知一段章动允许区间为
$\theta_{\min}\le\theta\le\theta_{\max}$，其中端点均不在 $0,\pi$。
如何利用 $P,K$ 判断方位进动是否反向？端点方位率为零与内部变号有何区别？}{tr26}{训练题26题图}
% BEGIN DETAILED ANSWER tr26
\solution \Step{写方位速度的符号条件}
以图~\ref{fig:solution-tr26}~中章动允许区间内的候选零进动位置 $K/P$为依据，先写出图右侧第一行的几何或动力学关系，再代入其下方的计算式。
 对允许区间内的常规姿态，
\[
\dot\varphi=\frac{K-P\cos\theta}{A\sin^2\theta}.
\]
分母严格为正，故转向只由分子决定。
\SolutionFigure{tr26}{训练题26的计算示意}
\Step{判断内部反向} 若 $P\ne0$，候选反向角满足
\[
\cos\theta_*=\frac KP.
\]
先核查 $|K/P|\le1$，再核查该角是否在章动允许区间内。
若在严格内部且分子跨过零，
$\dot\varphi$ 在经过该角时改变符号，方位运动出现反向；
与章动叠加后可形成回环轨迹。
若分子在整个允许区间同号，则方位单向进动。
\Step{区分端点尖点} 如果分子仅在转向端点为零，则该点同时有
$\dot\theta=0,\dot\varphi=0$。本体轴在单位球面上的速度
\[
\dot{\vect e}_3=\dot\theta\,\vect e_\theta+
\sin\theta\,\dot\varphi\,\vect e_\varphi
\]
此时为零，在非退化初值下可形成尖点。
内部反向处一般仍有非零 $\dot\theta$，
所以轴端轨迹不必在该处停住。
当 $P=0$ 时，分子就是常数 $K$，不应继续使用 $K/P$ 的判据。
\Step{在余弦变量中执行判别}
把章动范围写为 $u\in[u_-,u_+]\subset(-1,1)$，则 $\dot\varphi$ 的符号来自直线 $K-Pu$。
若 $K/P$ 严格位于两端之间，且不是定常双根姿态，运动经过它时分子变号；若仅等于端点，须同时检查 $\dot\theta=0$ 的退化；若位于区间外，方位速度全程同号。
此判别首先需从能量方程取得真实允许区间，单靠 $|K/P|<1$ 只说明有一个候选球面纬度，不能保证轨迹会经过它。

% END DETAILED ANSWER tr26
\exercise{正、负陀螺磁比的拉莫尔转向}\label{test:c3-27}

\QuestionStatement{磁矩满足 $\vect\mu=\gamma\vect L$，
均匀磁场 $\vect B=B\vect k$、$B>0$，忽略耗散。
从力矩方程确定进动角速度的有向矢量，并写出一个横向初值解。}{tr27}{训练题27题图}
% BEGIN DETAILED ANSWER tr27
\solution \Step{由磁力矩求有向进动率}
以图~\ref{fig:solution-tr27}~中磁场方向及角动量的进动圆锥为依据，先写出图右侧第一行的几何或动力学关系，再代入其下方的计算式。
 磁力矩为
\[
\dot{\vect L}=\vect\mu\times\vect B
=\gamma\vect L\times\vect B
=(-\gamma\vect B)\times\vect L.
\]
与 $\dot{\vect L}=\vect\Omega_L\times\vect L$ 比较，
\[
\vect\Omega_L=-\gamma\vect B.
\]
\SolutionFigure{tr27}{训练题27的计算示意}
\Step{求横向初值解} 分量方程为
\[
\dot L_x=\gamma BL_y,\qquad
\dot L_y=-\gamma BL_x,\qquad \dot L_z=0.
\]
若 $L_x(0)=L_\perp,L_y(0)=0$，
\[
L_x=L_\perp\cos(\gamma Bt),\qquad
L_y=-L_\perp\sin(\gamma Bt),\qquad L_z=L_z(0).
\]
\Step{进动方向与守恒量} 所以 $\gamma>0$ 时，横向投影起初从 $+x$ 方向转向 $-y$ 方向；
$\gamma<0$ 时转向相反。频率大小为 $|\gamma|B$，
但只报告正的频率大小会丢掉转向信息。
由方程还可直接验证
$\dd|\vect L|^2/\dd t=0$、$\dd(\vect L\cdot\vect B)/\dd t=0$，
故磁矩端点在垂直于磁场的平面内作圆周运动。
\Step{写出一般横向初值的旋转矩阵}
若初始为 $(L_{x0},L_{y0})$，记 $\chi=\gamma Bt$，则
\[
\begin{pmatrix}L_x\\L_y\end{pmatrix}
=\begin{pmatrix}\cos\chi&\sin\chi\\-\sin\chi&\cos\chi\end{pmatrix}
\begin{pmatrix}L_{x0}\\L_{y0}\end{pmatrix}.
\]
矩阵正交，因而 $L_x^2+L_y^2$ 不变。取 $L_{y0}=0$ 恢复前述特例，取 $\gamma\mapsto-\gamma$ 则仅改变旋转方向。
理想无阻尼模型保持倾角，不会自行让磁矩逐渐对齐磁场。

% END DETAILED ANSWER tr27
\exercise{共振角频率与普通频率的换算}\label{test:c3-28}

\QuestionStatement{某磁矩的陀螺磁比以
$\mathrm{rad\,s^{-1}\,T^{-1}}$ 为单位给出。
写出共振角频率和每秒周数，并以一个指定参数的算例检查 $2\pi$ 因子。并以 $|\gamma|=2\pi\times10^7\,\mathrm{rad\,s^{-1}\,T^{-1}}$、$B_0=2\,\mathrm T$ 为给定数值，求频率和周期。}{tr28}{训练题28题图}
% BEGIN DETAILED ANSWER tr28
\solution \Step{由相位定义换算频率}
以图~\ref{fig:solution-tr28}~中一周相位 $2\pi$ 与普通频率为依据，先写出图右侧第一行的几何或动力学关系，再代入其下方的计算式。
 一周的相位变化为 $2\pi$，因此
\[
\omega_0=|\gamma|B_0,\qquad
f_0=\frac{\omega_0}{2\pi}=\frac{|\gamma|B_0}{2\pi}.
\]
\SolutionFigure{tr28}{训练题28的计算示意}
\Step{完整代入给定单位} 为作纯粹的单位练习，设给定
$|\gamma|=2\pi\times10^7\,\mathrm{rad\,s^{-1}\,T^{-1}}$、
$B_0=2\,\mathrm T$，则
\[
\omega_0=4\pi\times10^7\,\mathrm{rad\,s^{-1}},\qquad
f_0=2\times10^7\,\mathrm{Hz}.
\]
周期 $T_0=1/f_0=5\times10^{-8}\,\mathrm s$；
代入 $\omega_0T_0$ 确实得到 $2\pi$。
\Step{区分参数表约定} 若参数表直接给 $\gamma/(2\pi)$，单位为 $\mathrm{Hz/T}$，
则应直接乘 $B_0$ 求 $f_0$，不能再除一次 $2\pi$。
相位演化、旋转坐标变换和脉冲翻转角通常使用角频率，
频谱横轴则常使用普通频率，计算前需明确单位。
\Step{相位与频率单位}
给定数值的相位为 $\Phi(t)=\omega_0t$（弧度）；也可写为周数 $N(t)=f_0t$，两者满足 $\Phi=2\pi N$。
例如 $t=25\,\mathrm{ns}$ 时，$N=1/2$ 周、$\Phi=\pi$，不是“$\pi$ Hz”。
$\gamma$ 的符号决定旋转方向，频谱中报告的正频率大小使用 $|\gamma|$；若讨论有向相位演化，则应另保留符号，不能混用两种约定。

% END DETAILED ANSWER tr28
\exercise{磁偶极能有额外半因子吗？}\label{test:c3-29}

\QuestionStatement{外磁场沿 $z$ 方向，磁矩满足 $\vect\mu=\gamma\vect S$。
写出磁偶极能和自旋 $1/2$ 的两个能级，
说明能级中的半因子来自何处。}{tr29}{训练题29题图}
% BEGIN DETAILED ANSWER tr29
\solution \Step{写固定外场的相互作用能}
以图~\ref{fig:solution-tr29}~中两能级的完整间距为依据，先写出图右侧第一行的几何或动力学关系，再代入其下方的计算式。
 固定外场中磁偶极相互作用能为
\[
U_Z=-\vect\mu\cdot\vect B=-\gamma B S_z.
\]
\SolutionFigure{tr29}{训练题29的计算示意}
\Step{代入自旋投影} 自旋 $1/2$ 的两个投影为 $S_z=\pm\hbar/2$，所以
\[
E_+=-\frac{\gamma\hbar B}{2},\qquad
E_-=\frac{\gamma\hbar B}{2}.
\]
\Step{计算两能级完整间距} 两能级间距的大小为
\[
|\Delta E|=|\gamma|\hbar B=\hbar\omega_0.
\]
若错把能量本身再乘 $1/2$，会使能级间距和共振频率都少一半。
这里每个能级中的半因子来自自旋投影的取值，
而不是磁偶极能公式里另外存在一个半因子。
Thomas 修正涉及相对论运动中的自旋轨道耦合和参考系转动，
不应移入上述固定外磁场的塞曼能公式。
对 $\gamma<0$，哪个投影处于较低能级会交换，
但能级间距的正值仍为 $|\gamma|\hbar B$。
\Step{从能级差恢复共振条件}
辐射量子能量为 $hf=\hbar\omega$。令其等于能级间距，
\[
hf_0=|\gamma|\hbar B\quad\Longrightarrow\quad
f_0=\frac{|\gamma|B}{2\pi},\qquad \omega_0=|\gamma|B.
\]
它与经典力矩进动的频率大小一致，提供能量与运动学的交叉核验。
每一个能级距零能量的大小是间距的一半，但跃迁跨越的是两个能级的完整差值；这正是不能再把共振频率除以2的原因。

% END DETAILED ANSWER tr29
\exercise{磁场加倍时频率和理想翻转脉冲怎样变化？}\label{test:c3-30}

\QuestionStatement{忽略弛豫，射频磁场的共振旋转分量幅度为 $B_1$。
静磁场由 $B_0$ 改为 $2B_0$，而 $B_1$ 保持不变。
求共振频率及理想 $90^\circ$、$180^\circ$ 脉冲时间的变化。}{tr30}{训练题30题图}
% BEGIN DETAILED ANSWER tr30
\solution \Step{改变静场后的共振频率}
以图~\ref{fig:solution-tr30}~中纵向载频与横向翻转场的独立作用为依据，先写出图右侧第一行的几何或动力学关系，再代入其下方的计算式。
 原共振角频率为 $\omega_0=|\gamma|B_0$，
新频率为
\[
\omega_0'=|\gamma|(2B_0)=2\omega_0,\qquad f_0'=2f_0.
\]
必须把射频调到新的共振频率。
\SolutionFigure{tr30}{训练题30的计算示意}
\Step{共振系内积分翻转角} 共振旋转系内的翻转角率只由横向旋转分量决定：
\[
\Omega_1=|\gamma|B_1,\qquad \beta=\Omega_1\tau.
\]
故
\[
\tau_{90}=\frac{\pi}{2|\gamma|B_1},\qquad
\tau_{180}=\frac{\pi}{|\gamma|B_1}.
\]
\Step{比较脉冲时间与失谐} 两种脉冲时间都不因 $B_0$ 加倍而改变。
若射频仍保持旧频率，则出现失谐
$\Delta=\omega_0'-\omega_{\mathrm{RF}}$，
有效旋转轴含纵向分量，旋转率为
$\sqrt{\Omega_1^2+\Delta^2}$；原脉冲时间不再保证预期翻转角和方向。
另外，$B_1$ 此处是共振旋转分量幅度；
若实验参数给的是线偏振场的峰值，在旋转波近似下其共振分量为峰值的一半，
必须先统一幅度定义再计算时间。
\Step{脉冲时间比例与失谐限制}
由 $\beta=|\gamma|B_1\tau$，$90^\circ$ 与 $180^\circ$ 两时间之比恒为 $1:2$。若要将两时间同时减半，需要把共振旋转幅度 $B_1$ 加倍，而仅增大 $B_0$ 不起这一作用。
保持旧载频时，纵向失谐使有效场倾斜。即使把脉冲时间改成 $\pi/[2\sqrt{\Omega_1^2+\Delta^2}]$，也只是绕倾斜轴转 $90^\circ$，一般不能使原来沿纵向的磁矩完成目标横向翻转；须先重新调到共振。

% END DETAILED ANSWER tr30
\clearpage
\section*{原课件作业题号}\label{appendix:coursework}
\addcontentsline{toc}{section}{原课件作业题号}
以下保留课件给出的教材作业题号。仅列题号的条目须结合原教材使用；
课件中给出完整题干的摆杆转盘、八分之一球等题目，已在正文中逐步求解。
讲义中的例题和训练题另行连续编号。
\begin{center}\begin{tabular}{@{}lll@{}}\toprule
课件&原页码&教材作业题号\\\midrule
c3-2&21&3.1、3.4、3.5\\
c3-3&24&3.6、3.7、3.9\\
c3-4&25&3.10、3.13、3.14\\
c3-5&21&3.20、3.22、3.26\\
c3-6&26&3.30、3.31、3.33\\\bottomrule
\end{tabular}\end{center}
\end{document}
